\documentclass[10pt, ukrainian]{article}
\setlength\textwidth{190mm} \setlength\textheight{277mm}
\setlength\voffset{-38mm} \setlength\hoffset{-38mm}
\setlength\parindent{0mm}
\pagestyle{empty}
\usepackage[T2A]{fontenc}
\usepackage[utf8]{inputenc}
\usepackage[ukrainian]{babel}
\begin{document}
{{\begin {small}\par
{ \center  Київський національний університет імені Тараса Шевченка \par}
{\parbox {5 cm}{Освітня програма \hfill}{\parbox {7 cm}{Інформатика \hfill}}\parbox {6 cm}{\hfill Семестр 6}}\par
{\parbox {5 cm} {Навчальний предмет \hfill}\parbox {7 cm}{Розробка програмного забезпечення \hfill}\parbox {6cm}{\hfill}\par}\vspace{-7pt} \end{small}}{\center Екзаменаційний білет \No \textbf
{ 1 }\par}}
{
\begin {large}
{\begin {center}\textbf{Завдання 1}\end {center}}\par
Що з наведеного найкраще характеризує рефакторінг?\par
\vspace{2mm}



(\,1\,) Докорінна переробка всього коду з нуля,  що не змінює його видиму поведінку%bad



(\,2\,) Довільна переробка коду, що виникає у процесі розробки%bad



(\,3\,) Зміна внутрішньої структури коду, що не змінює його видиму поведінку%good



(\,4\,) Зміна внутрішньої структури коду під час усунення помилок%bad


{\begin {center}\textbf{Завдання 2}\end {center}}\par
Який з наведених фрагментів коду явно потребує поліпшення (форматування рядків та типи до уваги не брати)?\par
\vspace{2mm}



(\,1\,) void load\_data(T \& arg1, double arg2, double arg3, double arg4); %bad



(\,2\,) \{int xtrmprg; …\}%bad



(\,3\,) \{fstream inputFile, outputFile;…\} %good



(\,4\,) void ff(…);%bad


{\begin {center}\textbf{Завдання 3}\end {center}}\par
Що є характерним для структурного програмування?\par
\vspace{2mm}



(\,1\,) Використання оператора безумовного переходу. %bad



(\,2\,) Використання механізму винятків. %bad



(\,3\,) У коді мають використовуватись структурні типи даних. %bad



(\,4\,) Переходи з середини одного блоку в середину іншого неможливі. %good


{\begin {center}\textbf{Завдання 4}\end {center}}\par
Вибрати всі істинні твердження (якщо існують).\par
\vspace{2mm}



(\,1\,) Адаптивні моделі розробки допускають зміну вимог до продукту. %good



(\,2\,) Неперервне проектування здійснює створення та модифікацію проекту системи в міру розробки. %good



(\,3\,) Неперервне проектування має бути ітеративним та інкрементним одночасно. %good



(\,4\,) Неперервність проектування полягає в тім, що процес розробки не можна переривати більше ніж на вихідні. %bad


{\begin {center}\textbf{Завдання 5}\end {center}}\par
Вибрати всі істинні твердження (якщо існують).\par
\vspace{2mm}



(\,1\,) Код програми пишеться для комп'ютера, тому він має бути синтаксично та семантично коректним. Інше значення не має. %bad



(\,2\,) Складні умови виникають в реальних задачах доволі часто, а тому є доволі природніми і не можуть вважатись недоліком коду. %bad



(\,3\,) Гарний код не повинен містити різних фрагментів, що складаються з кількох інструкцій та виконують однакове перетворення даних. %good



(\,4\,) Одним з засобів позбутись повторень у коді є використання функцій. %good

\newpage

{\begin {center}\textbf{Завдання 6}\end {center}}\par
Вибрати всі істинні твердження (якщо існують).\par
\vspace{2mm}



(\,1\,) Водоспадна методологія може мати тенденцію до надлишкового проектування. %good



(\,2\,) За використання гнучких методологій кожен розробник робить, що хоче. %bad



(\,3\,) Адаптивні методології розробки дозволяють уникнути надлишкового проектування. %good



(\,4\,) За використання гнучких методологій розробки припиняти розробку не рекомендується. %good


{\begin {center}\textbf{Завдання 7}\end {center}}\par
Вибрати всі істинні твердження (якщо існують).\par
\vspace{2mm}



(\,1\,) Гнучкі методології можуть передбачати як спільне, так і індивідуальне володіння кодом. %good



(\,2\,) Перша версія коду не завжди може бути якісною -- і це нормально. %good



(\,3\,) Одна з принципових відмінностей RUP від Scrum полягає у використанні численних проміжних технічних артефактів. %good



(\,4\,) Якщо виникають проблеми з модифікацією коду, то слід відразу змінити команду розробників на більш кваліфіковану.%bad


{\begin {center}\textbf{Завдання 8}\end {center}}\par
Що з наступного не має відношення до Scrum або не виконується/не використовується за використання оригінального Scrum?\par
\vspace{2mm}



(\,1\,) порядок виконання запланованої на спрінт роботи визначає власник продукту%good



(\,2\,) журнал невиконаних задач (backlog) продукту%bad



(\,3\,) залучення зовнішніх фахівців для побудови архітектури системи%good



(\,4\,) ретроспектива спрінту%bad


{\begin {center}\textbf{Завдання 9}\end {center}}\par
Що з наведеного нижче не розглядалось та навіть не згадувалось у курсі?\par
\vspace{2mm}



(\,1\,) LAZY%good



(\,2\,) sashimi%bad



(\,3\,) DSDM%bad



(\,4\,) Scrum of scrums%bad


{\begin {center}\textbf{Завдання 10}\end {center}}\par
Вибрати всі істинні твердження (якщо існують).\par
\vspace{2mm}



(\,1\,) Тестовий приклад складається з набору вхідних даних, умов виконання та очікуваних результатів.%good



(\,2\,) І верифікація, і валідація можуть використовувати тестування.%good



(\,3\,) Користувачі бувають непрямими.%good



(\,4\,) Модель якості при використанні є ієрархічною.%good



\end {large}}
{\vspace{5mm}\smallskip \begin {small} Затверджено на засіданні кафедри теоретичної кібернетики, протокол \No 6 від   29.01.2021 р.\par\smallskip\smallskip
{\parbox {6.5 cm}{Зав. кафедри \dotfill Ю.В.~Крак}}\hspace {0.7cm}{\hbox to 6.5 cm{ Екзаменатор \dotfill Т.О.~Карнаух}} \end{small}}
\newpage

{{\begin {small}\par
{ \center  Київський національний університет імені Тараса Шевченка \par}
{\parbox {5 cm}{Освітня програма \hfill}{\parbox {7 cm}{Інформатика \hfill}}\parbox {6 cm}{\hfill Семестр 6}}\par
{\parbox {5 cm} {Навчальний предмет \hfill}\parbox {7 cm}{Розробка програмного забезпечення \hfill}\parbox {6cm}{\hfill}\par}\vspace{-7pt} \end{small}}{\center Екзаменаційний білет \No \textbf
{ 2 }\par}}
{
\begin {large}
{\begin {center}\textbf{Завдання 1}\end {center}}\par
Так звані трикутники балансування зв'язують між собою:\par
\vspace{2mm}



(\,1\,) Завжди рівно 3 параметри%bad



(\,2\,) 4 або 3 параметри, при цьому може бути відсутній обсяг %good



(\,3\,) 4 або 3 параметри, при цьому може бути відсутній час%bad



(\,4\,) 4 або 3 параметри, при цьому може бути відсутня якість%bad


{\begin {center}\textbf{Завдання 2}\end {center}}\par
Який з наведених фрагментів коду явно потребує поліпшення (форматування рядків та типи до уваги не брати)?\par
\vspace{2mm}



(\,1\,) \{int x; int y; …\}%good



(\,2\,) void f13(…);%bad



(\,3\,) \{fstream input\_file, output\_file;…\} %good



(\,4\,) \{double maxPower, minPower;…\} %good


{\begin {center}\textbf{Завдання 3}\end {center}}\par
Що є характерним для структурного програмування?\par
\vspace{2mm}



(\,1\,) У коді мають використовуватись типи даних, що є класами. %bad



(\,2\,) У коді мають використовуватись типи структур, класи використовувати заборонено. %bad



(\,3\,) Програма зображується ієрархічною структурою блоків.%good



(\,4\,) У коді мають використовуватись типи структур. %bad


{\begin {center}\textbf{Завдання 4}\end {center}}\par
Вибрати всі істинні твердження (якщо існують).\par
\vspace{2mm}



(\,1\,) Адаптивні методології має сенс використовувати для проектів, де можлива зміна вимог. %good



(\,2\,) Для адаптивних методологій прогнозувати результат чи терміни його отримання абсолютно неможливо. %bad



(\,3\,) У прогнозній моделі розробки розробка може виконуватись ітеративно. %good



(\,4\,) Адаптивна модель розробки адаптує вимоги до продукту до наявної команди розробників. %bad


{\begin {center}\textbf{Завдання 5}\end {center}}\par
Вибрати всі істинні твердження (якщо існують).\par
\vspace{2mm}



(\,1\,) Зв'язки між окремими модулями системи мають бути якомога більш слабкими. %good



(\,2\,) Одним з прийомів спрощення складної умовної логіки є використання шаблону State. %good



(\,3\,) Чим вища кваліфікація програміста, тим більш складні для розуміння та довгі умовні вирази він використовує. %bad



(\,4\,) Код з довгими функціями має меншу кількість функцій, а тому його швидше та простіше тестувати. %bad

\newpage

{\begin {center}\textbf{Завдання 6}\end {center}}\par
Вибрати всі істинні твердження (якщо існують).\par
\vspace{2mm}



(\,1\,) Гнучкі методології передбачають часті випуски продукту. %good



(\,2\,) На відміну від Scrum, RUP є компонентно-орієнтованим. %bad



(\,3\,) Гнучкі методології базуються на комунікації розробників, а тому передбачають виключно колективне володіння кодом. %bad



(\,4\,) Гнучкі методології повністю нехтують документацією. %bad


{\begin {center}\textbf{Завдання 7}\end {center}}\par
Вибрати всі істинні твердження (якщо існують).\par
\vspace{2mm}



(\,1\,) Методологія розробки Big Design Up Front погана тим, що може призводити до недостатнього проектування.%bad



(\,2\,) Використання гнучких методологій при розробці надвеликих проектів не несе в собі жодних ризиків. %bad



(\,3\,) Використання RUP не заперечує часті випуски продукту, але й не наполягає на них. %good



(\,4\,) Методологія розробки Big Design Up Front погана тим, що може призводити до надлишкового проектування. %good


{\begin {center}\textbf{Завдання 8}\end {center}}\par
Що з наступного не має відношення до Scrum або не виконується/не використовується за використання оригінального Scrum?\par
\vspace{2mm}



(\,1\,) бачення продукту (vision) %bad



(\,2\,) правила визначення готовності (Definition of Done) %bad



(\,3\,) скрам-мастер визначає пріоритети задач%good



(\,4\,) власник продукту остаточно затверджує проект системи%good


{\begin {center}\textbf{Завдання 9}\end {center}}\par
Що з наведеного нижче не розглядалось та навіть не згадувалось у курсі?\par
\vspace{2mm}



(\,1\,) LOAN%good



(\,2\,) BDUF%bad



(\,3\,) водоспадна модель%bad



(\,4\,) LeSS%bad


{\begin {center}\textbf{Завдання 10}\end {center}}\par
Вибрати всі істинні твердження (якщо існують).\par
\vspace{2mm}



(\,1\,) Якість можна визначати через атрибути якості.%good



(\,2\,) Розглядають модель якості даних.%good



(\,3\,) Модель якості даних є ієрархічною.%good



(\,4\,) Якість буває внутрішньою.%good



\end {large}}
{\vspace{5mm}\smallskip \begin {small} Затверджено на засіданні кафедри теоретичної кібернетики, протокол \No 6 від   29.01.2021 р.\par\smallskip\smallskip
{\parbox {6.5 cm}{Зав. кафедри \dotfill Ю.В.~Крак}}\hspace {0.7cm}{\hbox to 6.5 cm{ Екзаменатор \dotfill Т.О.~Карнаух}} \end{small}}
\newpage

{{\begin {small}\par
{ \center  Київський національний університет імені Тараса Шевченка \par}
{\parbox {5 cm}{Освітня програма \hfill}{\parbox {7 cm}{Інформатика \hfill}}\parbox {6 cm}{\hfill Семестр 6}}\par
{\parbox {5 cm} {Навчальний предмет \hfill}\parbox {7 cm}{Розробка програмного забезпечення \hfill}\parbox {6cm}{\hfill}\par}\vspace{-7pt} \end{small}}{\center Екзаменаційний білет \No \textbf
{ 3 }\par}}
{
\begin {large}
{\begin {center}\textbf{Завдання 1}\end {center}}\par
Що є метою рефакторингу?\par
\vspace{2mm}



(\,1\,) Поліпшення внутрішньої структури коду%good



(\,2\,) Усунення помилок у коді%bad



(\,3\,) Зменшення кількості рядків коду%bad



(\,4\,) Додавання нової функціональності%bad


{\begin {center}\textbf{Завдання 2}\end {center}}\par
Який з наведених фрагментів коду явно потребує поліпшення (форматування рядків та типи до уваги не брати)?\par
\vspace{2mm}



(\,1\,) \{double max\_power, min\_power;…\} %good



(\,2\,) \{int a\_, a\_\_; …\}%bad



(\,3\,) void load\_data(…);%good



(\,4\,) \{int input\_file, file\_to\_output; …\}%bad


{\begin {center}\textbf{Завдання 3}\end {center}}\par
Що є характерним для структурного програмування?\par
\vspace{2mm}



(\,1\,) Кожен блок коду повинен мати єдиний вхід та єдиний вихід.%good



(\,2\,) Заборона використання скалярних (простих) типів даних. %bad



(\,3\,) У коді мають використовуватись типи структур. %bad



(\,4\,) У коді мають використовуватись типи структур, класи використовувати заборонено. %bad


{\begin {center}\textbf{Завдання 4}\end {center}}\par
Вибрати всі істинні твердження (якщо існують).\par
\vspace{2mm}



(\,1\,) Неперервне проектування створює специфікацію проекту перед початком розробки/ітерації, а під час ітерації його неперервно модифікує.%bad



(\,2\,) Ітеративні та інкрементні методи розробки є синонімом неперервного проектування. %bad



(\,3\,) Прогнозні моделі розробки з самого початку фіксують вимоги до продукту. %good



(\,4\,) Для адаптивних моделей розробки оцінки часу та бюджету не виконують, бо вимоги є змінними. %bad


{\begin {center}\textbf{Завдання 5}\end {center}}\par
Вибрати всі істинні твердження (якщо існують).\par
\vspace{2mm}



(\,1\,) Залежності між різними модулями системи мають бути якомога більш сильними, бо інакше вона розпадеться. %bad



(\,2\,) Гарний код має читатись як розмовна мова. %good



(\,3\,) Повторюваний код дозволяє програмісту краще зрозуміти його логіку та вибрати з кількох варіантів найкращий. %bad



(\,4\,) Чим менше модулів у проекті, тим краще, бо меншу кількість модулів треба інтегрувати в єдине ціле. %bad

\newpage

{\begin {center}\textbf{Завдання 6}\end {center}}\par
Вибрати всі істинні твердження (якщо існують).\par
\vspace{2mm}



(\,1\,) Гнучкі методології розробки намагаються якомога швидше визначити архітектуру системи, для чого на початку завжди виконується початкове проектування системи в цілому, а тільки потім переходять до ітеративного уточнення деталей та кодування. %bad



(\,2\,) Висока внутрішня якість коду є марною тратою часу на те, що замовник не побачить та/або не зможе зацінити. %bad



(\,3\,) Водоспадна методологія розробки може призводити до недостатнього проектування. %bad



(\,4\,) Гнучкі методи розробки нехтують проміжними технічними артефактами та документацією виключно через їх непотрібність. %bad


{\begin {center}\textbf{Завдання 7}\end {center}}\par
Вибрати всі істинні твердження (якщо існують).\par
\vspace{2mm}



(\,1\,) За відмови від документації одним з методом збереження знань про систему є постійне пересування персоналу між її частинами. %good



(\,2\,) Гнучкі методології розробки не допускають використання UML. %bad



(\,3\,) Програмування в парах (pair programming) спрощує введення в команду нових розробників. %good



(\,4\,) Розробка за гнучкими методологіями цінить співробітництво з замовником та піддатлівість змінам. %good


{\begin {center}\textbf{Завдання 8}\end {center}}\par
Що з наступного не має відношення до Scrum або не виконується/не використовується за використання оригінального Scrum?\par
\vspace{2mm}



(\,1\,) власник продукту призначає відповідальних за конкретні задачі%good



(\,2\,) порядок виконання запланованої на спрінт роботи визначає скрам-мастер%good



(\,3\,) журнал невиконаних задач (backlog) спрінту%bad



(\,4\,) скрам-мастер визначає, які задачі команда буде виконувати на наступному спрінті%good


{\begin {center}\textbf{Завдання 9}\end {center}}\par
Що з наведеного нижче не розглядалось та навіть не згадувалось у курсі?\par
\vspace{2mm}



(\,1\,) SAFe%bad



(\,2\,) poker%bad



(\,3\,) XP%bad



(\,4\,) спіральна модель%bad


{\begin {center}\textbf{Завдання 10}\end {center}}\par
Вибрати всі істинні твердження (якщо існують).\par
\vspace{2mm}



(\,1\,) Коли нема можливості перевірити всі можливі комбінації, тестують їх попарні (pairwise) комбінації. Це може значно зменшити кількість тестів.%good



(\,2\,) Користувачі бувають вторинними.%good



(\,3\,) Для тестування код треба ізолювати. Для того, щоб ізолювати, треба мати тести.%good



(\,4\,) Атрибути якості можна використовувати при уточненні вимог.%good



\end {large}}
{\vspace{5mm}\smallskip \begin {small} Затверджено на засіданні кафедри теоретичної кібернетики, протокол \No 6 від   29.01.2021 р.\par\smallskip\smallskip
{\parbox {6.5 cm}{Зав. кафедри \dotfill Ю.В.~Крак}}\hspace {0.7cm}{\hbox to 6.5 cm{ Екзаменатор \dotfill Т.О.~Карнаух}} \end{small}}
\newpage

{{\begin {small}\par
{ \center  Київський національний університет імені Тараса Шевченка \par}
{\parbox {5 cm}{Освітня програма \hfill}{\parbox {7 cm}{Інформатика \hfill}}\parbox {6 cm}{\hfill Семестр 6}}\par
{\parbox {5 cm} {Навчальний предмет \hfill}\parbox {7 cm}{Розробка програмного забезпечення \hfill}\parbox {6cm}{\hfill}\par}\vspace{-7pt} \end{small}}{\center Екзаменаційний білет \No \textbf
{ 4 }\par}}
{
\begin {large}
{\begin {center}\textbf{Завдання 1}\end {center}}\par
Так звані трикутники балансування зв'язують між собою:\par
\vspace{2mm}



(\,1\,) Завжди рівно 3 параметри%bad



(\,2\,) Завжди рівно 4 параметри%bad



(\,3\,) 4 або 3 параметри, при цьому може бути відсутній час%bad



(\,4\,) 4 або 3 параметри, при цьому може бути відсутній обсяг %good


{\begin {center}\textbf{Завдання 2}\end {center}}\par
Який з наведених фрагментів коду явно потребує поліпшення (форматування рядків та типи до уваги не брати)?\par
\vspace{2mm}



(\,1\,) void help();%good



(\,2\,) \{fstream input\_file, outputFile; …\}%bad



(\,3\,) \{double the\_best\_value\_anyone\_ever\_have;\} %bad



(\,4\,) \{double avg\_temperature;…\} %good


{\begin {center}\textbf{Завдання 3}\end {center}}\par
Що є характерним для структурного програмування?\par
\vspace{2mm}



(\,1\,) Переходи з середини одного блоку в середину іншого неможливі. %good



(\,2\,) Використання оператора безумовного переходу. %bad



(\,3\,) Використання механізму винятків. %bad



(\,4\,) У коді мають використовуватись структурні типи даних. %bad


{\begin {center}\textbf{Завдання 4}\end {center}}\par
Вибрати всі істинні твердження (якщо існують).\par
\vspace{2mm}



(\,1\,) В адаптивних методологіях планування робіт не виконується. %bad



(\,2\,) Прогнозна модель розробки може одночасно бути інкрементною. %good



(\,3\,) Адаптивні моделі розробки не допускають зміну вимог до продукту, але продукт може розроблятись поетапно. %bad



(\,4\,) Адаптивність методології проявляється в тому, що за зміни вимог розробка починається заново. %bad


{\begin {center}\textbf{Завдання 5}\end {center}}\par
Вибрати всі істинні твердження (якщо існують).\par
\vspace{2mm}



(\,1\,) Одним з засобів позбутись повторень у коді є використання успадкування. %good



(\,2\,) Довжина функцій не впливає на внутрішню якість коду. %bad



(\,3\,) Код, що виконує однакову послідовність логічних кроків, у більшості випадків є повторюваним.%good



(\,4\,) Одним з прийомів спрощення умовної логіки є використання поліморфізму. %good

\newpage

{\begin {center}\textbf{Завдання 6}\end {center}}\par
Вибрати всі істинні твердження (якщо існують).\par
\vspace{2mm}



(\,1\,) Якщо виникають проблеми з модифікацією коду, то слід розглянути внесення змін у проект. %good



(\,2\,) RUP є прогнозною методологією. %bad



(\,3\,) Багато хто з тих, що використовує Scrum, складає архітектурну модель розроблюваної системи. %good



(\,4\,) Гнучкі методології розробки є стійкими до припинення розробки на певний час. %bad


{\begin {center}\textbf{Завдання 7}\end {center}}\par
Вибрати всі істинні твердження (якщо існують).\par
\vspace{2mm}



(\,1\,) Розробка за гнучкими методологіями цінить людей та взаємодію між ними. %good



(\,2\,) Основна відмінність V-моделі від водоспадної полягає в тім, що діяльності з тестування виконуються паралельно з усіма іншими діяльностями з розробки. %good



(\,3\,) Програмування в парах (pair programming) підвищує якість коду та дозволяє економити на виправленні помилок. %good



(\,4\,) RUP відрізняється від Scrum тим, що ні в що не ставить персонал.%bad


{\begin {center}\textbf{Завдання 8}\end {center}}\par
Що з наступного не має відношення до Scrum або не виконується/не використовується за використання оригінального Scrum?\par
\vspace{2mm}



(\,1\,) під час спрінту дозволяється внесення змін до вимог, що вже реалізуються на поточному спрінті%good



(\,2\,) виділення в команді одного чи кількох проектувальників%good



(\,3\,) скрам-мастер остаточно затверджує проект системи%good



(\,4\,) під час спрінту зміни, що впливають на досягнення мети спрінту, зазвичай, не допускаються%bad


{\begin {center}\textbf{Завдання 9}\end {center}}\par
Що з наведеного нижче не розглядалось та навіть не згадувалось у курсі?\par
\vspace{2mm}



(\,1\,) CRAZY%good



(\,2\,) TDD%bad



(\,3\,) RUP%bad



(\,4\,) X-модель%good


{\begin {center}\textbf{Завдання 10}\end {center}}\par
Вибрати всі істинні твердження (якщо існують).\par
\vspace{2mm}



(\,1\,) Розглядають модель якості продукту.%good



(\,2\,) Забезпечення якості це не тільки тестування.%good



(\,3\,) Тестування буває позитивне та негативне.%good



(\,4\,) Якість буває зовнішньою.%good



\end {large}}
{\vspace{5mm}\smallskip \begin {small} Затверджено на засіданні кафедри теоретичної кібернетики, протокол \No 6 від   29.01.2021 р.\par\smallskip\smallskip
{\parbox {6.5 cm}{Зав. кафедри \dotfill Ю.В.~Крак}}\hspace {0.7cm}{\hbox to 6.5 cm{ Екзаменатор \dotfill Т.О.~Карнаух}} \end{small}}
\newpage

{{\begin {small}\par
{ \center  Київський національний університет імені Тараса Шевченка \par}
{\parbox {5 cm}{Освітня програма \hfill}{\parbox {7 cm}{Інформатика \hfill}}\parbox {6 cm}{\hfill Семестр 6}}\par
{\parbox {5 cm} {Навчальний предмет \hfill}\parbox {7 cm}{Розробка програмного забезпечення \hfill}\parbox {6cm}{\hfill}\par}\vspace{-7pt} \end{small}}{\center Екзаменаційний білет \No \textbf
{ 5 }\par}}
{
\begin {large}
{\begin {center}\textbf{Завдання 1}\end {center}}\par
Так звані трикутники балансування зв'язують між собою:\par
\vspace{2mm}



(\,1\,) 4 або 3 параметри, при цьому може бути відсутній обсяг %good



(\,2\,) 4 або 3 параметри, при цьому може бути відсутній час%bad



(\,3\,) Завжди рівно 4 параметри%bad



(\,4\,) 4 або 3 параметри, при цьому може бути відсутня якість%bad


{\begin {center}\textbf{Завдання 2}\end {center}}\par
Який з наведених фрагментів коду явно потребує поліпшення (форматування рядків та типи до уваги не брати)?\par
\vspace{2mm}



(\,1\,) \{fstream inputFile, outputFile;…\} %good



(\,2\,) \{int x; int y; …\}%good



(\,3\,) \{double max\_power, min\_power;…\} %good



(\,4\,) void ff(…);%bad


{\begin {center}\textbf{Завдання 3}\end {center}}\par
Що є характерним для структурного програмування?\par
\vspace{2mm}



(\,1\,) Кожен блок коду повинен мати єдиний вхід та єдиний вихід.%good



(\,2\,) Заборона використання скалярних (простих) типів даних. %bad



(\,3\,) У коді мають використовуватись типи даних, що є класами. %bad



(\,4\,) Використання механізму винятків. %bad


{\begin {center}\textbf{Завдання 4}\end {center}}\par
Вибрати всі істинні твердження (якщо існують).\par
\vspace{2mm}



(\,1\,) Прогнозні та адаптивні методології використовують одні ті самі діяльності з розробки ПЗ: визначення та аналіз вимог, проектування, реалізація, тестування, впровадження, супроводження. %good



(\,2\,) І прогнозні, і адаптивні методології передбачають планування виконуваних робіт. %good



(\,3\,) Прогнозна модель розробки не може одночасно бути інкрементною. %bad



(\,4\,) Гнучкі методології не передбачають планування робіт. %bad


{\begin {center}\textbf{Завдання 5}\end {center}}\par
Вибрати всі істинні твердження (якщо існують).\par
\vspace{2mm}



(\,1\,) Від довгих функцій можна позбутися, розклавши їх у послідовність викликів дрібніших функцій. %good



(\,2\,) Залежності всередині модуля можуть бути сильними. %good



(\,3\,) Одним з прийомів спрощення складної умовної логіки є використання шаблону Strategy. %good



(\,4\,) Втілення в коді різних підходів до одного того самого перетворення даних дозволяє програмісту продемонструвати свою професійну підготовку з найкращого боку. %bad

\newpage

{\begin {center}\textbf{Завдання 6}\end {center}}\par
Вибрати всі істинні твердження (якщо існують).\par
\vspace{2mm}



(\,1\,) Розробка за гнучкими методологіями не передбачає узгодження умов контракту. %bad



(\,2\,) Ітеративна та інкрементна розробка передбачає наявність у тому чи іншому вигляді етапу ініціалізації, ітерацій та керуючої таблиці проекту. %good



(\,3\,) RUP є архітектурно-орієнтованим. %good



(\,4\,) Програмування в парах (pair programming) зводиться до того, що кожним фрагментом коду володіють два програмісти. %bad


{\begin {center}\textbf{Завдання 7}\end {center}}\par
Вибрати всі істинні твердження (якщо існують).\par
\vspace{2mm}



(\,1\,) Водоспадна методологія розробки може призводити до надлишкового проектування. %good



(\,2\,) RUP принципово відрізняється від Scrum тим, що є керованим прецедентами. У той же час у Scrum при плануванні розробки прецеденти (чи їх аналоги) до уваги не беруться. %bad



(\,3\,) Гнучкі методології розробки допускають суттєві зміни вимог, проте все одно намагаються якомога швидше зафіксувати архітектуру. %good



(\,4\,) Гнучкі методи розробки нехтують проміжними технічними артефактами та документацією через те, що за зміни вимог та/або проекту треба витрачати додатковий час на підтримання їх у відповідному стані. %good


{\begin {center}\textbf{Завдання 8}\end {center}}\par
Що з наступного не має відношення до Scrum або не виконується/не використовується за використання оригінального Scrum?\par
\vspace{2mm}



(\,1\,) виділення одного з учасників команди як системного архітектора%good



(\,2\,) виділення в команді одного чи кількох тестувальників%good



(\,3\,) власник продукту призначає час, необхідний команді для реалізації конкретної задачі%good



(\,4\,) розповіді користувача%bad


{\begin {center}\textbf{Завдання 9}\end {center}}\par
Що з наведеного нижче не розглядалось та навіть не згадувалось у курсі?\par
\vspace{2mm}



(\,1\,) Lean%bad



(\,2\,) FDD%bad



(\,3\,) PRINCE2%good



(\,4\,) модель хаосу%bad


{\begin {center}\textbf{Завдання 10}\end {center}}\par
Вибрати всі істинні твердження (якщо існують).\par
\vspace{2mm}



(\,1\,) За провалених димових тестів решта тестів не виконується і код повертається на доробку.%good



(\,2\,) Якість можна визначати через дефекти.%good



(\,3\,) Можна розглядати якість під час використання.%good



(\,4\,) Модель якості продукту є ієрархічною.%good



\end {large}}
{\vspace{5mm}\smallskip \begin {small} Затверджено на засіданні кафедри теоретичної кібернетики, протокол \No 6 від   29.01.2021 р.\par\smallskip\smallskip
{\parbox {6.5 cm}{Зав. кафедри \dotfill Ю.В.~Крак}}\hspace {0.7cm}{\hbox to 6.5 cm{ Екзаменатор \dotfill Т.О.~Карнаух}} \end{small}}
\newpage

{{\begin {small}\par
{ \center  Київський національний університет імені Тараса Шевченка \par}
{\parbox {5 cm}{Освітня програма \hfill}{\parbox {7 cm}{Інформатика \hfill}}\parbox {6 cm}{\hfill Семестр 6}}\par
{\parbox {5 cm} {Навчальний предмет \hfill}\parbox {7 cm}{Розробка програмного забезпечення \hfill}\parbox {6cm}{\hfill}\par}\vspace{-7pt} \end{small}}{\center Екзаменаційний білет \No \textbf
{ 6 }\par}}
{
\begin {large}
{\begin {center}\textbf{Завдання 1}\end {center}}\par
Так звані трикутники балансування зв'язують між собою:\par
\vspace{2mm}



(\,1\,) Завжди рівно 4 параметри%bad



(\,2\,) 4 або 3 параметри, при цьому може бути відсутня якість%bad



(\,3\,) Завжди рівно 3 параметри%bad



(\,4\,) 4 або 3 параметри, при цьому може бути відсутній обсяг %good


{\begin {center}\textbf{Завдання 2}\end {center}}\par
Який з наведених фрагментів коду явно потребує поліпшення (форматування рядків та типи до уваги не брати)?\par
\vspace{2mm}



(\,1\,) \{fstream input\_file, output\_file;…\} %good



(\,2\,) void load\_data(T \& arg1, double arg2, double arg3, double arg4); %bad



(\,3\,) \{int a\_, a\_\_; …\}%bad



(\,4\,) \{double the\_best\_value\_anyone\_ever\_have;\} %bad


{\begin {center}\textbf{Завдання 3}\end {center}}\par
Що є характерним для структурного програмування?\par
\vspace{2mm}



(\,1\,) Використання оператора безумовного переходу. %bad



(\,2\,) У коді мають використовуватись типи структур, класи використовувати заборонено. %bad



(\,3\,) Програма зображується ієрархічною структурою блоків.%good



(\,4\,) У коді мають використовуватись типи структур. %bad


{\begin {center}\textbf{Завдання 4}\end {center}}\par
Вибрати всі істинні твердження (якщо існують).\par
\vspace{2mm}



(\,1\,) У прогнозній моделі розробки розробка не може виконуватись ітеративно. %bad



(\,2\,) Адаптивні методології не передбачають визначення вимог. %bad



(\,3\,) У випадку, коли вимоги до продукту відомі заздалегідь та є фіксованими контрактом, адаптивні методології розробки використовувати не можна. %bad



(\,4\,) Для адаптивних методологій розробки, на відміну від прогнозних, попередні оцінки бюджету та/або часу не виконують, бо вони все одно зміняться. %bad


{\begin {center}\textbf{Завдання 5}\end {center}}\par
Вибрати всі істинні твердження (якщо існують).\par
\vspace{2mm}



(\,1\,) Довгі функції доволі часто приховують складну логіку обробки або повторюваний код. %good



(\,2\,) Якщо кожна функція бібліотеки буде здатна в залежності від параметрів виконувати кілька різних дій, то загальна кількість функцій у ній буде менша і її буде простіше використовувати. %bad



(\,3\,) Одним з засобів позбутись повторень у коді є використання шаблону Template Method (але цей шаблон здатен боротися і з іншими недоліками коду). %good



(\,4\,) Повторюваний код має повторюватись у точності до символу, інакше він повторюваним вважатись не може. %bad

\newpage

{\begin {center}\textbf{Завдання 6}\end {center}}\par
Вибрати всі істинні твердження (якщо існують).\par
\vspace{2mm}



(\,1\,) RUP є ітеративним та інкрементним.%good



(\,2\,) За невизначених або піддатливих змінам вимог краще застосовувати гнучкі методології розробки. %good



(\,3\,) Висока внутрішня якість коду здешевшує розробку, бо дозволяє економити на вартості та часі майбутніх змін. %good



(\,4\,) За застосування адаптивних методологій розробки, так само як і за прогнозних, намагаються якомога швидше визначити архітектуру системи. Відмінність полягає в тім, у який спосіб це виконується.%good


{\begin {center}\textbf{Завдання 7}\end {center}}\par
Вибрати всі істинні твердження (якщо існують).\par
\vspace{2mm}



(\,1\,) Використання гнучких методологій вимагає високої кваліфікації персоналу. %good



(\,2\,) За програмування в парах (pair programming) поки один колега набирає код на комп'ютері інший може піти пограти в настільний теніс. %bad



(\,3\,) Без рефакторінгу неперервне проектування неможливе. %good



(\,4\,) Модель хаосу є методологією розробки програмного забезпечення. %bad


{\begin {center}\textbf{Завдання 8}\end {center}}\par
Що з наступного не має відношення до Scrum або не виконується/не використовується за використання оригінального Scrum?\par
\vspace{2mm}



(\,1\,) розповіді користувача%bad



(\,2\,) команда визначає, які задачі вона буде виконувати на наступному спрінті%good



(\,3\,) канбан%bad



(\,4\,) побудова різнопланових моделей системи%good


{\begin {center}\textbf{Завдання 9}\end {center}}\par
Що з наведеного нижче не розглядалось та навіть не згадувалось у курсі?\par
\vspace{2mm}



(\,1\,) X-модель%good



(\,2\,) V-модель%bad



(\,3\,) LeSS%bad



(\,4\,) Scrum%bad


{\begin {center}\textbf{Завдання 10}\end {center}}\par
Вибрати всі істинні твердження (якщо існують).\par
\vspace{2mm}



(\,1\,) Тестування вимог може породжувати їх уточнення.%good



(\,2\,) Відсутність знайдених дефектів на початкових стадіях розробки може бути свідченням поганого покриття коду тестами.%good



(\,3\,) Розглядають модель якості при використанні.%good



(\,4\,) Користувачі бувають основними.%good



\end {large}}
{\vspace{5mm}\smallskip \begin {small} Затверджено на засіданні кафедри теоретичної кібернетики, протокол \No 6 від   29.01.2021 р.\par\smallskip\smallskip
{\parbox {6.5 cm}{Зав. кафедри \dotfill Ю.В.~Крак}}\hspace {0.7cm}{\hbox to 6.5 cm{ Екзаменатор \dotfill Т.О.~Карнаух}} \end{small}}
\newpage

{{\begin {small}\par
{ \center  Київський національний університет імені Тараса Шевченка \par}
{\parbox {5 cm}{Освітня програма \hfill}{\parbox {7 cm}{Інформатика \hfill}}\parbox {6 cm}{\hfill Семестр 6}}\par
{\parbox {5 cm} {Навчальний предмет \hfill}\parbox {7 cm}{Розробка програмного забезпечення \hfill}\parbox {6cm}{\hfill}\par}\vspace{-7pt} \end{small}}{\center Екзаменаційний білет \No \textbf
{ 7 }\par}}
{
\begin {large}
{\begin {center}\textbf{Завдання 1}\end {center}}\par
Шаблон проектування:\par
\vspace{2mm}



(\,1\,) Формулює ідею розв'язання проблеми та може пропонувати одну чи кілька її реалізацій%good



(\,2\,) Задає розв'язання типової проблеми до найдрібніших подробиць%bad



(\,3\,) Тільки формулює ідею розв'язання проблеми%bad



(\,4\,) Є шаблоном, в який треба підставити власні ідентифікатори%bad


{\begin {center}\textbf{Завдання 2}\end {center}}\par
Який з наведених фрагментів коду явно потребує поліпшення (форматування рядків та типи до уваги не брати)?\par
\vspace{2mm}



(\,1\,) \{int xtrmprg; …\}%bad



(\,2\,) \{double maxPower, minPower;…\} %good



(\,3\,) void help();%good



(\,4\,) void load\_data(…);%good


{\begin {center}\textbf{Завдання 3}\end {center}}\par
Що є характерним для структурного програмування?\par
\vspace{2mm}



(\,1\,) У коді мають використовуватись структурні типи даних. %bad



(\,2\,) Заборона використання скалярних (простих) типів даних. %bad



(\,3\,) У коді мають використовуватись типи даних, що є класами. %bad



(\,4\,) Програма зображується ієрархічною структурою блоків.%good


{\begin {center}\textbf{Завдання 4}\end {center}}\par
Вибрати всі істинні твердження (якщо існують).\par
\vspace{2mm}



(\,1\,) Для адаптивних методологій прогнозувати результат чи терміни його отримання абсолютно неможливо. %bad



(\,2\,) Неперервне проектування має бути ітеративним та інкрементним одночасно. %good



(\,3\,) Адаптивні моделі розробки допускають зміну вимог до продукту. %good



(\,4\,) І прогнозні, і адаптивні методології передбачають планування виконуваних робіт. %good


{\begin {center}\textbf{Завдання 5}\end {center}}\par
Вибрати всі істинні твердження (якщо існують).\par
\vspace{2mm}



(\,1\,) Одним з прийомів спрощення складних умов є їх декомпозиція. %good



(\,2\,) Складні комбінації умов починають керувати логікою виконання, тому що такими є вимоги до функціонування ПЗ. Оскільки ПЗ має якнайкраще забезпечувати потреби замовника, то складні комбінації умов у коді вважатись недоліком коду не можуть. %bad



(\,3\,) Чим менше класів у коді, тим менше класів треба налагоджувати, а тому класи за можливості слід об'єднувати. %bad



(\,4\,) Правильно спроектована функція повинна мати рівно один обов'язок. %good

\newpage

{\begin {center}\textbf{Завдання 6}\end {center}}\par
Вибрати всі істинні твердження (якщо існують).\par
\vspace{2mm}



(\,1\,) RUP є керованим прецедентами. %good



(\,2\,) Розробка за гнучкими методологіями ніколи не слідує плану. %bad



(\,3\,) Розповіді користувача (user stories) – це прогресивний інструмент гнучких методологій, який нічого спільного не має з прецедентами (use case). %bad



(\,4\,) RUP є компонентно-орієнтованим. %good


{\begin {center}\textbf{Завдання 7}\end {center}}\par
Вибрати всі істинні твердження (якщо існують).\par
\vspace{2mm}



(\,1\,) Розповіді користувача (user stories) це менш формалізований порівняно з прецедентами спосіб опису вимог до системи. %good



(\,2\,) Scrum передбачає колективне володіння кодом. %good



(\,3\,) Ітеративна методологія розробки у загальному випадку не гарантує, що не буде надлишкового проектування. %good



(\,4\,) Тільки RUP використовує UML.%bad


{\begin {center}\textbf{Завдання 8}\end {center}}\par
Що з наступного не має відношення до Scrum або не виконується/не використовується за використання оригінального Scrum?\par
\vspace{2mm}



(\,1\,) планування спрінту%bad



(\,2\,) покер планування (Planning Poker)%bad



(\,3\,) скрам-мастер визначає час, необхідний команді для реалізації конкретної задачі%good



(\,4\,) приведення у належний вигляд журналу невиконаних задач (backlog grooming) %bad


{\begin {center}\textbf{Завдання 9}\end {center}}\par
Що з наведеного нижче не розглядалось та навіть не згадувалось у курсі?\par
\vspace{2mm}



(\,1\,) RAD%bad



(\,2\,) BDD%bad



(\,3\,) poker%bad



(\,4\,) Lean%bad


{\begin {center}\textbf{Завдання 10}\end {center}}\par
Вибрати всі істинні твердження (якщо існують).\par
\vspace{2mm}



(\,1\,) Тестування буває позитивне та негативне.%good



(\,2\,) Атрибути якості можна використовувати при уточненні вимог.%good



(\,3\,) Користувачі бувають основними.%good



(\,4\,) Тестовий приклад складається з набору вхідних даних, умов виконання та очікуваних результатів.%good



\end {large}}
{\vspace{5mm}\smallskip \begin {small} Затверджено на засіданні кафедри теоретичної кібернетики, протокол \No 6 від   29.01.2021 р.\par\smallskip\smallskip
{\parbox {6.5 cm}{Зав. кафедри \dotfill Ю.В.~Крак}}\hspace {0.7cm}{\hbox to 6.5 cm{ Екзаменатор \dotfill Т.О.~Карнаух}} \end{small}}
\newpage

{{\begin {small}\par
{ \center  Київський національний університет імені Тараса Шевченка \par}
{\parbox {5 cm}{Освітня програма \hfill}{\parbox {7 cm}{Інформатика \hfill}}\parbox {6 cm}{\hfill Семестр 6}}\par
{\parbox {5 cm} {Навчальний предмет \hfill}\parbox {7 cm}{Розробка програмного забезпечення \hfill}\parbox {6cm}{\hfill}\par}\vspace{-7pt} \end{small}}{\center Екзаменаційний білет \No \textbf
{ 8 }\par}}
{
\begin {large}
{\begin {center}\textbf{Завдання 1}\end {center}}\par
Так звані трикутники балансування зв'язують між собою:\par
\vspace{2mm}



(\,1\,) Завжди рівно 4 параметри%bad



(\,2\,) Завжди рівно 3 параметри%bad



(\,3\,) 4 або 3 параметри, при цьому може бути відсутній обсяг %good



(\,4\,) 4 або 3 параметри, при цьому може бути відсутній час%bad


{\begin {center}\textbf{Завдання 2}\end {center}}\par
Який з наведених фрагментів коду явно потребує поліпшення (форматування рядків та типи до уваги не брати)?\par
\vspace{2mm}



(\,1\,) void f13(…);%bad



(\,2\,) \{double avg\_temperature;…\} %good



(\,3\,) \{int input\_file, file\_to\_output; …\}%bad



(\,4\,) \{fstream input\_file, outputFile; …\}%bad


{\begin {center}\textbf{Завдання 3}\end {center}}\par
Що є характерним для структурного програмування?\par
\vspace{2mm}



(\,1\,) Використання механізму винятків. %bad



(\,2\,) Заборона використання скалярних (простих) типів даних. %bad



(\,3\,) Переходи з середини одного блоку в середину іншого неможливі. %good



(\,4\,) У коді мають використовуватись типи структур. %bad


{\begin {center}\textbf{Завдання 4}\end {center}}\par
Вибрати всі істинні твердження (якщо існують).\par
\vspace{2mm}



(\,1\,) В адаптивних методологіях планування робіт не виконується. %bad



(\,2\,) Прогнозні моделі розробки з самого початку фіксують вимоги до продукту. %good



(\,3\,) Адаптивні моделі розробки не допускають зміну вимог до продукту, але продукт може розроблятись поетапно. %bad



(\,4\,) Прогнозна модель розробки може одночасно бути інкрементною. %good


{\begin {center}\textbf{Завдання 5}\end {center}}\par
Вибрати всі істинні твердження (якщо існують).\par
\vspace{2mm}



(\,1\,) Код з довгими функціями має меншу кількість функцій, а тому для нього треба менше тестів. %bad



(\,2\,) Залежності між різними модулями системи мають бути якомога більш сильними, бо інакше вона розпадеться. %bad



(\,3\,) Чим вища кваліфікація програміста, тим більш складні для розуміння та довгі умовні вирази він використовує. %bad



(\,4\,) Залежності всередині модуля можуть бути сильними. %good

\newpage

{\begin {center}\textbf{Завдання 6}\end {center}}\par
Вибрати всі істинні твердження (якщо існують).\par
\vspace{2mm}



(\,1\,) Сьогодні водоспадна методологія не має права на існування. %bad



(\,2\,) Поступове уточнення вимог замовника полегшує адаптацію до змін. %good



(\,3\,) Код, що себе сам документує, -- це код, який виводить користувачу документацію на себе. %bad



(\,4\,) Водоспадна методологія не призначена для внесення змін у вимоги до системи. %good


{\begin {center}\textbf{Завдання 7}\end {center}}\par
Вибрати всі істинні твердження (якщо існують).\par
\vspace{2mm}



(\,1\,) Простота – мистецтво максимізації незробленої роботи. %good



(\,2\,) Програмування в парах (pair programming) зводиться до того, що кожним фрагментом коду володіють два програмісти. %bad



(\,3\,) Ітеративна методологія розробки у загальному випадку не гарантує, що не буде надлишкового проектування. %good



(\,4\,) Висока внутрішня якість коду є марною тратою часу на те, що замовник не побачить та/або не зможе зацінити. %bad


{\begin {center}\textbf{Завдання 8}\end {center}}\par
Що з наступного не має відношення до Scrum або не виконується/не використовується за використання оригінального Scrum?\par
\vspace{2mm}



(\,1\,) щоденний скрам%bad



(\,2\,) обернені вимоги%bad



(\,3\,) блочні тести%bad



(\,4\,) порядок виконання запланованої на спрінт роботи визначають розробники%bad


{\begin {center}\textbf{Завдання 9}\end {center}}\par
Що з наведеного нижче не розглядалось та навіть не згадувалось у курсі?\par
\vspace{2mm}



(\,1\,) CRAZY%good



(\,2\,) Cleanroom%bad



(\,3\,) SAFe%bad



(\,4\,) RAD%bad


{\begin {center}\textbf{Завдання 10}\end {center}}\par
Вибрати всі істинні твердження (якщо існують).\par
\vspace{2mm}



(\,1\,) Якість буває внутрішньою.%good



(\,2\,) Модель якості при використанні є ієрархічною.%good



(\,3\,) Модель якості даних є ієрархічною.%good



(\,4\,) Для тестування код треба ізолювати. Для того, щоб ізолювати, треба мати тести.%good



\end {large}}
{\vspace{5mm}\smallskip \begin {small} Затверджено на засіданні кафедри теоретичної кібернетики, протокол \No 6 від   29.01.2021 р.\par\smallskip\smallskip
{\parbox {6.5 cm}{Зав. кафедри \dotfill Ю.В.~Крак}}\hspace {0.7cm}{\hbox to 6.5 cm{ Екзаменатор \dotfill Т.О.~Карнаух}} \end{small}}
\newpage

{{\begin {small}\par
{ \center  Київський національний університет імені Тараса Шевченка \par}
{\parbox {5 cm}{Освітня програма \hfill}{\parbox {7 cm}{Інформатика \hfill}}\parbox {6 cm}{\hfill Семестр 6}}\par
{\parbox {5 cm} {Навчальний предмет \hfill}\parbox {7 cm}{Розробка програмного забезпечення \hfill}\parbox {6cm}{\hfill}\par}\vspace{-7pt} \end{small}}{\center Екзаменаційний білет \No \textbf
{ 9 }\par}}
{
\begin {large}
{\begin {center}\textbf{Завдання 1}\end {center}}\par
Що є метою рефакторингу?\par
\vspace{2mm}



(\,1\,) Усунення помилок у коді%bad



(\,2\,) Зменшення кількості рядків коду%bad



(\,3\,) Додавання нової функціональності%bad



(\,4\,) Поліпшення внутрішньої структури коду%good


{\begin {center}\textbf{Завдання 2}\end {center}}\par
Який з наведених фрагментів коду явно потребує поліпшення (форматування рядків та типи до уваги не брати)?\par
\vspace{2mm}



(\,1\,) \{int x; int y; …\}%good



(\,2\,) \{int xtrmprg; …\}%bad



(\,3\,) \{double the\_best\_value\_anyone\_ever\_have;\} %bad



(\,4\,) void load\_data(…);%good


{\begin {center}\textbf{Завдання 3}\end {center}}\par
Що є характерним для структурного програмування?\par
\vspace{2mm}



(\,1\,) Використання оператора безумовного переходу. %bad



(\,2\,) Кожен блок коду повинен мати єдиний вхід та єдиний вихід.%good



(\,3\,) У коді мають використовуватись типи структур, класи використовувати заборонено. %bad



(\,4\,) У коді мають використовуватись типи даних, що є класами. %bad


{\begin {center}\textbf{Завдання 4}\end {center}}\par
Вибрати всі істинні твердження (якщо існують).\par
\vspace{2mm}



(\,1\,) Прогнозні та адаптивні методології використовують одні ті самі діяльності з розробки ПЗ: визначення та аналіз вимог, проектування, реалізація, тестування, впровадження, супроводження. %good



(\,2\,) Неперервне проектування створює специфікацію проекту перед початком розробки/ітерації, а під час ітерації його неперервно модифікує.%bad



(\,3\,) У прогнозній моделі розробки розробка може виконуватись ітеративно. %good



(\,4\,) Адаптивні методології має сенс використовувати для проектів, де можлива зміна вимог. %good


{\begin {center}\textbf{Завдання 5}\end {center}}\par
Вибрати всі істинні твердження (якщо існують).\par
\vspace{2mm}



(\,1\,) Код програми пишеться для комп'ютера, тому він має бути синтаксично та семантично коректним. Інше значення не має. %bad



(\,2\,) Одним з засобів позбутись повторень у коді є використання функцій. %good



(\,3\,) Якщо кожна функція бібліотеки буде здатна в залежності від параметрів виконувати кілька різних дій, то загальна кількість функцій у ній буде менша і її буде простіше використовувати. %bad



(\,4\,) Одним з засобів позбутись повторень у коді є використання успадкування. %good

\newpage

{\begin {center}\textbf{Завдання 6}\end {center}}\par
Вибрати всі істинні твердження (якщо існують).\par
\vspace{2mm}



(\,1\,) Без рефакторінгу неперервне проектування неможливе. %good



(\,2\,) Використання гнучких методологій вимагає високої кваліфікації персоналу. %good



(\,3\,) Код, що себе сам документує, -- це код, який виводить користувачу документацію на себе. %bad



(\,4\,) Якщо виникають проблеми з модифікацією коду, то слід відразу змінити команду розробників на більш кваліфіковану.%bad


{\begin {center}\textbf{Завдання 7}\end {center}}\par
Вибрати всі істинні твердження (якщо існують).\par
\vspace{2mm}



(\,1\,) Методологія розробки Big Design Up Front погана тим, що може призводити до недостатнього проектування.%bad



(\,2\,) Використання RUP не заперечує часті випуски продукту, але й не наполягає на них. %good



(\,3\,) Гнучкі методології базуються на комунікації розробників, а тому передбачають виключно колективне володіння кодом. %bad



(\,4\,) На відміну від Scrum, RUP є компонентно-орієнтованим. %bad


{\begin {center}\textbf{Завдання 8}\end {center}}\par
Що з наступного не має відношення до Scrum або не виконується/не використовується за використання оригінального Scrum?\par
\vspace{2mm}



(\,1\,) скрам-мастер призначає відповідальних за конкретні задачі%good



(\,2\,) скрам-мастер визначає пріоритети задач%good



(\,3\,) правила визначення готовності (Definition of Done) %bad



(\,4\,) скрам-мастер визначає час, необхідний команді для реалізації конкретної задачі%good


{\begin {center}\textbf{Завдання 9}\end {center}}\par
Що з наведеного нижче не розглядалось та навіть не згадувалось у курсі?\par
\vspace{2mm}



(\,1\,) sashimi%bad



(\,2\,) FDD%bad



(\,3\,) XP%bad



(\,4\,) PRINCE2%good


{\begin {center}\textbf{Завдання 10}\end {center}}\par
Вибрати всі істинні твердження (якщо існують).\par
\vspace{2mm}



(\,1\,) Якість можна визначати через атрибути якості.%good



(\,2\,) Користувачі бувають вторинними.%good



(\,3\,) Якість буває зовнішньою.%good



(\,4\,) Тестування вимог може породжувати їх уточнення.%good



\end {large}}
{\vspace{5mm}\smallskip \begin {small} Затверджено на засіданні кафедри теоретичної кібернетики, протокол \No 6 від   29.01.2021 р.\par\smallskip\smallskip
{\parbox {6.5 cm}{Зав. кафедри \dotfill Ю.В.~Крак}}\hspace {0.7cm}{\hbox to 6.5 cm{ Екзаменатор \dotfill Т.О.~Карнаух}} \end{small}}
\newpage

{{\begin {small}\par
{ \center  Київський національний університет імені Тараса Шевченка \par}
{\parbox {5 cm}{Освітня програма \hfill}{\parbox {7 cm}{Інформатика \hfill}}\parbox {6 cm}{\hfill Семестр 6}}\par
{\parbox {5 cm} {Навчальний предмет \hfill}\parbox {7 cm}{Розробка програмного забезпечення \hfill}\parbox {6cm}{\hfill}\par}\vspace{-7pt} \end{small}}{\center Екзаменаційний білет \No \textbf
{ 10 }\par}}
{
\begin {large}
{\begin {center}\textbf{Завдання 1}\end {center}}\par
Що з наведеного найкраще характеризує рефакторінг?\par
\vspace{2mm}



(\,1\,) Зміна внутрішньої структури коду під час усунення помилок%bad



(\,2\,) Довільна переробка коду, що виникає у процесі розробки%bad



(\,3\,) Зміна внутрішньої структури коду, що не змінює його видиму поведінку%good



(\,4\,) Докорінна переробка всього коду з нуля,  що не змінює його видиму поведінку%bad


{\begin {center}\textbf{Завдання 2}\end {center}}\par
Який з наведених фрагментів коду явно потребує поліпшення (форматування рядків та типи до уваги не брати)?\par
\vspace{2mm}



(\,1\,) \{int a\_, a\_\_; …\}%bad



(\,2\,) \{double avg\_temperature;…\} %good



(\,3\,) void load\_data(T \& arg1, double arg2, double arg3, double arg4); %bad



(\,4\,) \{fstream inputFile, outputFile;…\} %good


{\begin {center}\textbf{Завдання 3}\end {center}}\par
Що є характерним для структурного програмування?\par
\vspace{2mm}



(\,1\,) Переходи з середини одного блоку в середину іншого неможливі. %good



(\,2\,) У коді мають використовуватись структурні типи даних. %bad



(\,3\,) Використання оператора безумовного переходу. %bad



(\,4\,) У коді мають використовуватись типи даних, що є класами. %bad


{\begin {center}\textbf{Завдання 4}\end {center}}\par
Вибрати всі істинні твердження (якщо існують).\par
\vspace{2mm}



(\,1\,) У прогнозній моделі розробки розробка не може виконуватись ітеративно. %bad



(\,2\,) Неперервне проектування здійснює створення та модифікацію проекту системи в міру розробки. %good



(\,3\,) Ітеративні та інкрементні методи розробки є синонімом неперервного проектування. %bad



(\,4\,) Адаптивні методології не передбачають визначення вимог. %bad


{\begin {center}\textbf{Завдання 5}\end {center}}\par
Вибрати всі істинні твердження (якщо існують).\par
\vspace{2mm}



(\,1\,) Довжина функцій не впливає на внутрішню якість коду. %bad



(\,2\,) Довгі функції доволі часто приховують складну логіку обробки або повторюваний код. %good



(\,3\,) Зв'язки між окремими модулями системи мають бути якомога більш слабкими. %good



(\,4\,) Код, що виконує однакову послідовність логічних кроків, у більшості випадків є повторюваним.%good

\newpage

{\begin {center}\textbf{Завдання 6}\end {center}}\par
Вибрати всі істинні твердження (якщо існують).\par
\vspace{2mm}



(\,1\,) RUP відрізняється від Scrum тим, що ні в що не ставить персонал.%bad



(\,2\,) Гнучкі методології розробки не допускають використання UML. %bad



(\,3\,) Сьогодні водоспадна методологія не має права на існування. %bad



(\,4\,) Водоспадна методологія може мати тенденцію до надлишкового проектування. %good


{\begin {center}\textbf{Завдання 7}\end {center}}\par
Вибрати всі істинні твердження (якщо існують).\par
\vspace{2mm}



(\,1\,) Адаптивні методології розробки дозволяють уникнути надлишкового проектування. %good



(\,2\,) Одна з принципових відмінностей RUP від Scrum полягає у використанні численних проміжних технічних артефактів. %good



(\,3\,) Розробка за гнучкими методологіями цінить співробітництво з замовником та піддатлівість змінам. %good



(\,4\,) Використання гнучких методологій при розробці надвеликих проектів не несе в собі жодних ризиків. %bad


{\begin {center}\textbf{Завдання 8}\end {center}}\par
Що з наступного не має відношення до Scrum або не виконується/не використовується за використання оригінального Scrum?\par
\vspace{2mm}



(\,1\,) ретроспектива спрінту%bad



(\,2\,) журнал невиконаних задач (backlog) спрінту%bad



(\,3\,) порядок виконання запланованої на спрінт роботи визначає скрам-мастер%good



(\,4\,) під час спрінту дозволяється внесення змін до вимог, що вже реалізуються на поточному спрінті%good


{\begin {center}\textbf{Завдання 9}\end {center}}\par
Що з наведеного нижче не розглядалось та навіть не згадувалось у курсі?\par
\vspace{2mm}



(\,1\,) спіральна модель%bad



(\,2\,) LOAN%good



(\,3\,) V-модель%bad



(\,4\,) LAZY%good


{\begin {center}\textbf{Завдання 10}\end {center}}\par
Вибрати всі істинні твердження (якщо існують).\par
\vspace{2mm}



(\,1\,) Користувачі бувають непрямими.%good



(\,2\,) Забезпечення якості це не тільки тестування.%good



(\,3\,) За провалених димових тестів решта тестів не виконується і код повертається на доробку.%good



(\,4\,) І верифікація, і валідація можуть використовувати тестування.%good



\end {large}}
{\vspace{5mm}\smallskip \begin {small} Затверджено на засіданні кафедри теоретичної кібернетики, протокол \No 6 від   29.01.2021 р.\par\smallskip\smallskip
{\parbox {6.5 cm}{Зав. кафедри \dotfill Ю.В.~Крак}}\hspace {0.7cm}{\hbox to 6.5 cm{ Екзаменатор \dotfill Т.О.~Карнаух}} \end{small}}
\newpage

{{\begin {small}\par
{ \center  Київський національний університет імені Тараса Шевченка \par}
{\parbox {5 cm}{Освітня програма \hfill}{\parbox {7 cm}{Інформатика \hfill}}\parbox {6 cm}{\hfill Семестр 6}}\par
{\parbox {5 cm} {Навчальний предмет \hfill}\parbox {7 cm}{Розробка програмного забезпечення \hfill}\parbox {6cm}{\hfill}\par}\vspace{-7pt} \end{small}}{\center Екзаменаційний білет \No \textbf
{ 11 }\par}}
{
\begin {large}
{\begin {center}\textbf{Завдання 1}\end {center}}\par
Шаблон проектування:\par
\vspace{2mm}



(\,1\,) Формулює ідею розв'язання проблеми та може пропонувати одну чи кілька її реалізацій%good



(\,2\,) Є шаблоном, в який треба підставити власні ідентифікатори%bad



(\,3\,) Тільки формулює ідею розв'язання проблеми%bad



(\,4\,) Задає розв'язання типової проблеми до найдрібніших подробиць%bad


{\begin {center}\textbf{Завдання 2}\end {center}}\par
Який з наведених фрагментів коду явно потребує поліпшення (форматування рядків та типи до уваги не брати)?\par
\vspace{2mm}



(\,1\,) void help();%good



(\,2\,) \{fstream input\_file, outputFile; …\}%bad



(\,3\,) void f13(…);%bad



(\,4\,) \{double maxPower, minPower;…\} %good


{\begin {center}\textbf{Завдання 3}\end {center}}\par
Що є характерним для структурного програмування?\par
\vspace{2mm}



(\,1\,) У коді мають використовуватись типи структур, класи використовувати заборонено. %bad



(\,2\,) Використання механізму винятків. %bad



(\,3\,) Програма зображується ієрархічною структурою блоків.%good



(\,4\,) У коді мають використовуватись структурні типи даних. %bad


{\begin {center}\textbf{Завдання 4}\end {center}}\par
Вибрати всі істинні твердження (якщо існують).\par
\vspace{2mm}



(\,1\,) Адаптивна модель розробки адаптує вимоги до продукту до наявної команди розробників. %bad



(\,2\,) Для адаптивних моделей розробки оцінки часу та бюджету не виконують, бо вимоги є змінними. %bad



(\,3\,) Гнучкі методології не передбачають планування робіт. %bad



(\,4\,) Прогнозна модель розробки не може одночасно бути інкрементною. %bad


{\begin {center}\textbf{Завдання 5}\end {center}}\par
Вибрати всі істинні твердження (якщо існують).\par
\vspace{2mm}



(\,1\,) Складні комбінації умов починають керувати логікою виконання, тому що такими є вимоги до функціонування ПЗ. Оскільки ПЗ має якнайкраще забезпечувати потреби замовника, то складні комбінації умов у коді вважатись недоліком коду не можуть. %bad



(\,2\,) Гарний код не повинен містити різних фрагментів, що складаються з кількох інструкцій та виконують однакове перетворення даних. %good



(\,3\,) Гарний код має читатись як розмовна мова. %good



(\,4\,) Повторюваний код має повторюватись у точності до символу, інакше він повторюваним вважатись не може. %bad

\newpage

{\begin {center}\textbf{Завдання 6}\end {center}}\par
Вибрати всі істинні твердження (якщо існують).\par
\vspace{2mm}



(\,1\,) Основна відмінність V-моделі від водоспадної полягає в тім, що діяльності з тестування виконуються паралельно з усіма іншими діяльностями з розробки. %good



(\,2\,) RUP є прогнозною методологією. %bad



(\,3\,) Якщо виникають проблеми з модифікацією коду, то слід розглянути внесення змін у проект. %good



(\,4\,) За застосування адаптивних методологій розробки, так само як і за прогнозних, намагаються якомога швидше визначити архітектуру системи. Відмінність полягає в тім, у який спосіб це виконується.%good


{\begin {center}\textbf{Завдання 7}\end {center}}\par
Вибрати всі істинні твердження (якщо існують).\par
\vspace{2mm}



(\,1\,) Гнучкі методології можуть передбачати як спільне, так і індивідуальне володіння кодом. %good



(\,2\,) За програмування в парах (pair programming) поки один колега набирає код на комп'ютері інший може піти пограти в настільний теніс. %bad



(\,3\,) Простота – мистецтво максимізації незробленої роботи. %good



(\,4\,) Модель хаосу є методологією розробки програмного забезпечення. %bad


{\begin {center}\textbf{Завдання 8}\end {center}}\par
Що з наступного не має відношення до Scrum або не виконується/не використовується за використання оригінального Scrum?\par
\vspace{2mm}



(\,1\,) щоденний скрам%bad



(\,2\,) власник продукту призначає час, необхідний команді для реалізації конкретної задачі%good



(\,3\,) блочні тести%bad



(\,4\,) планування спрінту%bad


{\begin {center}\textbf{Завдання 9}\end {center}}\par
Що з наведеного нижче не розглядалось та навіть не згадувалось у курсі?\par
\vspace{2mm}



(\,1\,) sashimi%bad



(\,2\,) Cleanroom%bad



(\,3\,) X-модель%good



(\,4\,) Scrum of scrums%bad


{\begin {center}\textbf{Завдання 10}\end {center}}\par
Вибрати всі істинні твердження (якщо існують).\par
\vspace{2mm}



(\,1\,) Розглядають модель якості продукту.%good



(\,2\,) Розглядають модель якості при використанні.%good



(\,3\,) Можна розглядати якість під час використання.%good



(\,4\,) Відсутність знайдених дефектів на початкових стадіях розробки може бути свідченням поганого покриття коду тестами.%good



\end {large}}
{\vspace{5mm}\smallskip \begin {small} Затверджено на засіданні кафедри теоретичної кібернетики, протокол \No 6 від   29.01.2021 р.\par\smallskip\smallskip
{\parbox {6.5 cm}{Зав. кафедри \dotfill Ю.В.~Крак}}\hspace {0.7cm}{\hbox to 6.5 cm{ Екзаменатор \dotfill Т.О.~Карнаух}} \end{small}}
\newpage

{{\begin {small}\par
{ \center  Київський національний університет імені Тараса Шевченка \par}
{\parbox {5 cm}{Освітня програма \hfill}{\parbox {7 cm}{Інформатика \hfill}}\parbox {6 cm}{\hfill Семестр 6}}\par
{\parbox {5 cm} {Навчальний предмет \hfill}\parbox {7 cm}{Розробка програмного забезпечення \hfill}\parbox {6cm}{\hfill}\par}\vspace{-7pt} \end{small}}{\center Екзаменаційний білет \No \textbf
{ 12 }\par}}
{
\begin {large}
{\begin {center}\textbf{Завдання 1}\end {center}}\par
Так звані трикутники балансування зв'язують між собою:\par
\vspace{2mm}



(\,1\,) Завжди рівно 3 параметри%bad



(\,2\,) Завжди рівно 4 параметри%bad



(\,3\,) 4 або 3 параметри, при цьому може бути відсутній обсяг %good



(\,4\,) 4 або 3 параметри, при цьому може бути відсутня якість%bad


{\begin {center}\textbf{Завдання 2}\end {center}}\par
Який з наведених фрагментів коду явно потребує поліпшення (форматування рядків та типи до уваги не брати)?\par
\vspace{2mm}



(\,1\,) \{fstream input\_file, output\_file;…\} %good



(\,2\,) \{int input\_file, file\_to\_output; …\}%bad



(\,3\,) void ff(…);%bad



(\,4\,) \{double max\_power, min\_power;…\} %good


{\begin {center}\textbf{Завдання 3}\end {center}}\par
Що є характерним для структурного програмування?\par
\vspace{2mm}



(\,1\,) У коді мають використовуватись типи структур. %bad



(\,2\,) Кожен блок коду повинен мати єдиний вхід та єдиний вихід.%good



(\,3\,) Заборона використання скалярних (простих) типів даних. %bad



(\,4\,) Заборона використання скалярних (простих) типів даних. %bad


{\begin {center}\textbf{Завдання 4}\end {center}}\par
Вибрати всі істинні твердження (якщо існують).\par
\vspace{2mm}



(\,1\,) Адаптивність методології проявляється в тому, що за зміни вимог розробка починається заново. %bad



(\,2\,) Неперервність проектування полягає в тім, що процес розробки не можна переривати більше ніж на вихідні. %bad



(\,3\,) Для адаптивних методологій розробки, на відміну від прогнозних, попередні оцінки бюджету та/або часу не виконують, бо вони все одно зміняться. %bad



(\,4\,) У випадку, коли вимоги до продукту відомі заздалегідь та є фіксованими контрактом, адаптивні методології розробки використовувати не можна. %bad


{\begin {center}\textbf{Завдання 5}\end {center}}\par
Вибрати всі істинні твердження (якщо існують).\par
\vspace{2mm}



(\,1\,) Чим менше модулів у проекті, тим краще, бо меншу кількість модулів треба інтегрувати в єдине ціле. %bad



(\,2\,) Чим менше класів у коді, тим менше класів треба налагоджувати, а тому класи за можливості слід об'єднувати. %bad



(\,3\,) Одним з прийомів спрощення умовної логіки є використання поліморфізму. %good



(\,4\,) Одним з прийомів спрощення складних умов є їх декомпозиція. %good

\newpage

{\begin {center}\textbf{Завдання 6}\end {center}}\par
Вибрати всі істинні твердження (якщо існують).\par
\vspace{2mm}



(\,1\,) Тільки RUP використовує UML.%bad



(\,2\,) Гнучкі методології розробки допускають суттєві зміни вимог, проте все одно намагаються якомога швидше зафіксувати архітектуру. %good



(\,3\,) За невизначених або піддатливих змінам вимог краще застосовувати гнучкі методології розробки. %good



(\,4\,) Гнучкі методології розробки намагаються якомога швидше визначити архітектуру системи, для чого на початку завжди виконується початкове проектування системи в цілому, а тільки потім переходять до ітеративного уточнення деталей та кодування. %bad


{\begin {center}\textbf{Завдання 7}\end {center}}\par
Вибрати всі істинні твердження (якщо існують).\par
\vspace{2mm}



(\,1\,) Багато хто з тих, що використовує Scrum, складає архітектурну модель розроблюваної системи. %good



(\,2\,) Ітеративна та інкрементна розробка передбачає наявність у тому чи іншому вигляді етапу ініціалізації, ітерацій та керуючої таблиці проекту. %good



(\,3\,) Методологія розробки Big Design Up Front погана тим, що може призводити до надлишкового проектування. %good



(\,4\,) Розробка за гнучкими методологіями не передбачає узгодження умов контракту. %bad


{\begin {center}\textbf{Завдання 8}\end {center}}\par
Що з наступного не має відношення до Scrum або не виконується/не використовується за використання оригінального Scrum?\par
\vspace{2mm}



(\,1\,) під час спрінту зміни, що впливають на досягнення мети спрінту, зазвичай, не допускаються%bad



(\,2\,) виділення в команді одного чи кількох тестувальників%good



(\,3\,) скрам-мастер остаточно затверджує проект системи%good



(\,4\,) виділення одного з учасників команди як системного архітектора%good


{\begin {center}\textbf{Завдання 9}\end {center}}\par
Що з наведеного нижче не розглядалось та навіть не згадувалось у курсі?\par
\vspace{2mm}



(\,1\,) LOAN%good



(\,2\,) poker%bad



(\,3\,) TDD%bad



(\,4\,) Scrum%bad


{\begin {center}\textbf{Завдання 10}\end {center}}\par
Вибрати всі істинні твердження (якщо існують).\par
\vspace{2mm}



(\,1\,) Коли нема можливості перевірити всі можливі комбінації, тестують їх попарні (pairwise) комбінації. Це може значно зменшити кількість тестів.%good



(\,2\,) Якість можна визначати через дефекти.%good



(\,3\,) Модель якості продукту є ієрархічною.%good



(\,4\,) Розглядають модель якості даних.%good



\end {large}}
{\vspace{5mm}\smallskip \begin {small} Затверджено на засіданні кафедри теоретичної кібернетики, протокол \No 6 від   29.01.2021 р.\par\smallskip\smallskip
{\parbox {6.5 cm}{Зав. кафедри \dotfill Ю.В.~Крак}}\hspace {0.7cm}{\hbox to 6.5 cm{ Екзаменатор \dotfill Т.О.~Карнаух}} \end{small}}
\newpage

{{\begin {small}\par
{ \center  Київський національний університет імені Тараса Шевченка \par}
{\parbox {5 cm}{Освітня програма \hfill}{\parbox {7 cm}{Інформатика \hfill}}\parbox {6 cm}{\hfill Семестр 6}}\par
{\parbox {5 cm} {Навчальний предмет \hfill}\parbox {7 cm}{Розробка програмного забезпечення \hfill}\parbox {6cm}{\hfill}\par}\vspace{-7pt} \end{small}}{\center Екзаменаційний білет \No \textbf
{ 13 }\par}}
{
\begin {large}
{\begin {center}\textbf{Завдання 1}\end {center}}\par
Так звані трикутники балансування зв'язують між собою:\par
\vspace{2mm}



(\,1\,) 4 або 3 параметри, при цьому може бути відсутній час%bad



(\,2\,) 4 або 3 параметри, при цьому може бути відсутній обсяг %good



(\,3\,) Завжди рівно 4 параметри%bad



(\,4\,) 4 або 3 параметри, при цьому може бути відсутня якість%bad


{\begin {center}\textbf{Завдання 2}\end {center}}\par
Який з наведених фрагментів коду явно потребує поліпшення (форматування рядків та типи до уваги не брати)?\par
\vspace{2mm}



(\,1\,) \{double max\_power, min\_power;…\} %good



(\,2\,) \{int xtrmprg; …\}%bad



(\,3\,) \{fstream input\_file, output\_file;…\} %good



(\,4\,) \{int a\_, a\_\_; …\}%bad


{\begin {center}\textbf{Завдання 3}\end {center}}\par
Що є характерним для структурного програмування?\par
\vspace{2mm}



(\,1\,) У коді мають використовуватись типи структур, класи використовувати заборонено. %bad



(\,2\,) Переходи з середини одного блоку в середину іншого неможливі. %good



(\,3\,) У коді мають використовуватись структурні типи даних. %bad



(\,4\,) У коді мають використовуватись типи даних, що є класами. %bad


{\begin {center}\textbf{Завдання 4}\end {center}}\par
Вибрати всі істинні твердження (якщо існують).\par
\vspace{2mm}



(\,1\,) Прогнозна модель розробки не може одночасно бути інкрементною. %bad



(\,2\,) Адаптивні методології не передбачають визначення вимог. %bad



(\,3\,) Неперервне проектування створює специфікацію проекту перед початком розробки/ітерації, а під час ітерації його неперервно модифікує.%bad



(\,4\,) Ітеративні та інкрементні методи розробки є синонімом неперервного проектування. %bad


{\begin {center}\textbf{Завдання 5}\end {center}}\par
Вибрати всі істинні твердження (якщо існують).\par
\vspace{2mm}



(\,1\,) Код з довгими функціями має меншу кількість функцій, а тому для нього треба менше тестів. %bad



(\,2\,) Від довгих функцій можна позбутися, розклавши їх у послідовність викликів дрібніших функцій. %good



(\,3\,) Код з довгими функціями має меншу кількість функцій, а тому його швидше та простіше тестувати. %bad



(\,4\,) Втілення в коді різних підходів до одного того самого перетворення даних дозволяє програмісту продемонструвати свою професійну підготовку з найкращого боку. %bad

\newpage

{\begin {center}\textbf{Завдання 6}\end {center}}\par
Вибрати всі істинні твердження (якщо існують).\par
\vspace{2mm}



(\,1\,) За використання гнучких методологій кожен розробник робить, що хоче. %bad



(\,2\,) Гнучкі методи розробки нехтують проміжними технічними артефактами та документацією через те, що за зміни вимог та/або проекту треба витрачати додатковий час на підтримання їх у відповідному стані. %good



(\,3\,) Розповіді користувача (user stories) це менш формалізований порівняно з прецедентами спосіб опису вимог до системи. %good



(\,4\,) За відмови від документації одним з методом збереження знань про систему є постійне пересування персоналу між її частинами. %good


{\begin {center}\textbf{Завдання 7}\end {center}}\par
Вибрати всі істинні твердження (якщо існують).\par
\vspace{2mm}



(\,1\,) Гнучкі методології розробки є стійкими до припинення розробки на певний час. %bad



(\,2\,) Розробка за гнучкими методологіями цінить людей та взаємодію між ними. %good



(\,3\,) Перша версія коду не завжди може бути якісною -- і це нормально. %good



(\,4\,) Програмування в парах (pair programming) спрощує введення в команду нових розробників. %good


{\begin {center}\textbf{Завдання 8}\end {center}}\par
Що з наступного не має відношення до Scrum або не виконується/не використовується за використання оригінального Scrum?\par
\vspace{2mm}



(\,1\,) побудова різнопланових моделей системи%good



(\,2\,) виділення в команді одного чи кількох проектувальників%good



(\,3\,) скрам-мастер призначає відповідальних за конкретні задачі%good



(\,4\,) команда визначає, які задачі вона буде виконувати на наступному спрінті%good


{\begin {center}\textbf{Завдання 9}\end {center}}\par
Що з наведеного нижче не розглядалось та навіть не згадувалось у курсі?\par
\vspace{2mm}



(\,1\,) модель хаосу%bad



(\,2\,) DSDM%bad



(\,3\,) LeSS%bad



(\,4\,) Lean%bad


{\begin {center}\textbf{Завдання 10}\end {center}}\par
Вибрати всі істинні твердження (якщо існують).\par
\vspace{2mm}



(\,1\,) Відсутність знайдених дефектів на початкових стадіях розробки може бути свідченням поганого покриття коду тестами.%good



(\,2\,) І верифікація, і валідація можуть використовувати тестування.%good



(\,3\,) Можна розглядати якість під час використання.%good



(\,4\,) Тестування буває позитивне та негативне.%good



\end {large}}
{\vspace{5mm}\smallskip \begin {small} Затверджено на засіданні кафедри теоретичної кібернетики, протокол \No 6 від   29.01.2021 р.\par\smallskip\smallskip
{\parbox {6.5 cm}{Зав. кафедри \dotfill Ю.В.~Крак}}\hspace {0.7cm}{\hbox to 6.5 cm{ Екзаменатор \dotfill Т.О.~Карнаух}} \end{small}}
\newpage

{{\begin {small}\par
{ \center  Київський національний університет імені Тараса Шевченка \par}
{\parbox {5 cm}{Освітня програма \hfill}{\parbox {7 cm}{Інформатика \hfill}}\parbox {6 cm}{\hfill Семестр 6}}\par
{\parbox {5 cm} {Навчальний предмет \hfill}\parbox {7 cm}{Розробка програмного забезпечення \hfill}\parbox {6cm}{\hfill}\par}\vspace{-7pt} \end{small}}{\center Екзаменаційний білет \No \textbf
{ 14 }\par}}
{
\begin {large}
{\begin {center}\textbf{Завдання 1}\end {center}}\par
Так звані трикутники балансування зв'язують між собою:\par
\vspace{2mm}



(\,1\,) 4 або 3 параметри, при цьому може бути відсутній обсяг %good



(\,2\,) 4 або 3 параметри, при цьому може бути відсутній час%bad



(\,3\,) Завжди рівно 3 параметри%bad



(\,4\,) 4 або 3 параметри, при цьому може бути відсутня якість%bad


{\begin {center}\textbf{Завдання 2}\end {center}}\par
Який з наведених фрагментів коду явно потребує поліпшення (форматування рядків та типи до уваги не брати)?\par
\vspace{2mm}



(\,1\,) void load\_data(…);%good



(\,2\,) void load\_data(T \& arg1, double arg2, double arg3, double arg4); %bad



(\,3\,) \{double the\_best\_value\_anyone\_ever\_have;\} %bad



(\,4\,) \{fstream input\_file, outputFile; …\}%bad


{\begin {center}\textbf{Завдання 3}\end {center}}\par
Що є характерним для структурного програмування?\par
\vspace{2mm}



(\,1\,) Використання оператора безумовного переходу. %bad



(\,2\,) Використання механізму винятків. %bad



(\,3\,) Програма зображується ієрархічною структурою блоків.%good



(\,4\,) У коді мають використовуватись типи структур. %bad


{\begin {center}\textbf{Завдання 4}\end {center}}\par
Вибрати всі істинні твердження (якщо існують).\par
\vspace{2mm}



(\,1\,) В адаптивних методологіях планування робіт не виконується. %bad



(\,2\,) У прогнозній моделі розробки розробка може виконуватись ітеративно. %good



(\,3\,) Для адаптивних методологій розробки, на відміну від прогнозних, попередні оцінки бюджету та/або часу не виконують, бо вони все одно зміняться. %bad



(\,4\,) У прогнозній моделі розробки розробка не може виконуватись ітеративно. %bad


{\begin {center}\textbf{Завдання 5}\end {center}}\par
Вибрати всі істинні твердження (якщо існують).\par
\vspace{2mm}



(\,1\,) Повторюваний код дозволяє програмісту краще зрозуміти його логіку та вибрати з кількох варіантів найкращий. %bad



(\,2\,) Одним з прийомів спрощення складної умовної логіки є використання шаблону Strategy. %good



(\,3\,) Одним з засобів позбутись повторень у коді є використання шаблону Template Method (але цей шаблон здатен боротися і з іншими недоліками коду). %good



(\,4\,) Правильно спроектована функція повинна мати рівно один обов'язок. %good

\newpage

{\begin {center}\textbf{Завдання 6}\end {center}}\par
Вибрати всі істинні твердження (якщо існують).\par
\vspace{2mm}



(\,1\,) Водоспадна методологія розробки може призводити до недостатнього проектування. %bad



(\,2\,) Водоспадна методологія не призначена для внесення змін у вимоги до системи. %good



(\,3\,) За використання гнучких методологій розробки припиняти розробку не рекомендується. %good



(\,4\,) RUP є ітеративним та інкрементним.%good


{\begin {center}\textbf{Завдання 7}\end {center}}\par
Вибрати всі істинні твердження (якщо існують).\par
\vspace{2mm}



(\,1\,) RUP є архітектурно-орієнтованим. %good



(\,2\,) Гнучкі методології повністю нехтують документацією. %bad



(\,3\,) Водоспадна методологія розробки може призводити до надлишкового проектування. %good



(\,4\,) RUP є керованим прецедентами. %good


{\begin {center}\textbf{Завдання 8}\end {center}}\par
Що з наступного не має відношення до Scrum або не виконується/не використовується за використання оригінального Scrum?\par
\vspace{2mm}



(\,1\,) розповіді користувача%bad



(\,2\,) приведення у належний вигляд журналу невиконаних задач (backlog grooming) %bad



(\,3\,) залучення зовнішніх фахівців для побудови архітектури системи%good



(\,4\,) порядок виконання запланованої на спрінт роботи визначає власник продукту%good


{\begin {center}\textbf{Завдання 9}\end {center}}\par
Що з наведеного нижче не розглядалось та навіть не згадувалось у курсі?\par
\vspace{2mm}



(\,1\,) LAZY%good



(\,2\,) водоспадна модель%bad



(\,3\,) PRINCE2%good



(\,4\,) BDD%bad


{\begin {center}\textbf{Завдання 10}\end {center}}\par
Вибрати всі істинні твердження (якщо існують).\par
\vspace{2mm}



(\,1\,) Забезпечення якості це не тільки тестування.%good



(\,2\,) Користувачі бувають основними.%good



(\,3\,) Якість можна визначати через атрибути якості.%good



(\,4\,) Атрибути якості можна використовувати при уточненні вимог.%good



\end {large}}
{\vspace{5mm}\smallskip \begin {small} Затверджено на засіданні кафедри теоретичної кібернетики, протокол \No 6 від   29.01.2021 р.\par\smallskip\smallskip
{\parbox {6.5 cm}{Зав. кафедри \dotfill Ю.В.~Крак}}\hspace {0.7cm}{\hbox to 6.5 cm{ Екзаменатор \dotfill Т.О.~Карнаух}} \end{small}}
\newpage

{{\begin {small}\par
{ \center  Київський національний університет імені Тараса Шевченка \par}
{\parbox {5 cm}{Освітня програма \hfill}{\parbox {7 cm}{Інформатика \hfill}}\parbox {6 cm}{\hfill Семестр 6}}\par
{\parbox {5 cm} {Навчальний предмет \hfill}\parbox {7 cm}{Розробка програмного забезпечення \hfill}\parbox {6cm}{\hfill}\par}\vspace{-7pt} \end{small}}{\center Екзаменаційний білет \No \textbf
{ 15 }\par}}
{
\begin {large}
{\begin {center}\textbf{Завдання 1}\end {center}}\par
Так звані трикутники балансування зв'язують між собою:\par
\vspace{2mm}



(\,1\,) 4 або 3 параметри, при цьому може бути відсутня якість%bad



(\,2\,) 4 або 3 параметри, при цьому може бути відсутній час%bad



(\,3\,) 4 або 3 параметри, при цьому може бути відсутній обсяг %good



(\,4\,) Завжди рівно 4 параметри%bad


{\begin {center}\textbf{Завдання 2}\end {center}}\par
Який з наведених фрагментів коду явно потребує поліпшення (форматування рядків та типи до уваги не брати)?\par
\vspace{2mm}



(\,1\,) \{double avg\_temperature;…\} %good



(\,2\,) void f13(…);%bad



(\,3\,) \{int input\_file, file\_to\_output; …\}%bad



(\,4\,) void help();%good


{\begin {center}\textbf{Завдання 3}\end {center}}\par
Що є характерним для структурного програмування?\par
\vspace{2mm}



(\,1\,) Кожен блок коду повинен мати єдиний вхід та єдиний вихід.%good



(\,2\,) Використання оператора безумовного переходу. %bad



(\,3\,) У коді мають використовуватись структурні типи даних. %bad



(\,4\,) Використання механізму винятків. %bad


{\begin {center}\textbf{Завдання 4}\end {center}}\par
Вибрати всі істинні твердження (якщо існують).\par
\vspace{2mm}



(\,1\,) Адаптивність методології проявляється в тому, що за зміни вимог розробка починається заново. %bad



(\,2\,) Для адаптивних моделей розробки оцінки часу та бюджету не виконують, бо вимоги є змінними. %bad



(\,3\,) Для адаптивних методологій прогнозувати результат чи терміни його отримання абсолютно неможливо. %bad



(\,4\,) Адаптивні методології має сенс використовувати для проектів, де можлива зміна вимог. %good


{\begin {center}\textbf{Завдання 5}\end {center}}\par
Вибрати всі істинні твердження (якщо існують).\par
\vspace{2mm}



(\,1\,) Складні умови виникають в реальних задачах доволі часто, а тому є доволі природніми і не можуть вважатись недоліком коду. %bad



(\,2\,) Одним з прийомів спрощення складної умовної логіки є використання шаблону State. %good



(\,3\,) Одним з засобів позбутись повторень у коді є використання шаблону Template Method (але цей шаблон здатен боротися і з іншими недоліками коду). %good



(\,4\,) Якщо кожна функція бібліотеки буде здатна в залежності від параметрів виконувати кілька різних дій, то загальна кількість функцій у ній буде менша і її буде простіше використовувати. %bad

\newpage

{\begin {center}\textbf{Завдання 6}\end {center}}\par
Вибрати всі істинні твердження (якщо існують).\par
\vspace{2mm}



(\,1\,) Гнучкі методології передбачають часті випуски продукту. %good



(\,2\,) Scrum передбачає колективне володіння кодом. %good



(\,3\,) Розповіді користувача (user stories) – це прогресивний інструмент гнучких методологій, який нічого спільного не має з прецедентами (use case). %bad



(\,4\,) Поступове уточнення вимог замовника полегшує адаптацію до змін. %good


{\begin {center}\textbf{Завдання 7}\end {center}}\par
Вибрати всі істинні твердження (якщо існують).\par
\vspace{2mm}



(\,1\,) Висока внутрішня якість коду здешевшує розробку, бо дозволяє економити на вартості та часі майбутніх змін. %good



(\,2\,) RUP принципово відрізняється від Scrum тим, що є керованим прецедентами. У той же час у Scrum при плануванні розробки прецеденти (чи їх аналоги) до уваги не беруться. %bad



(\,3\,) Програмування в парах (pair programming) підвищує якість коду та дозволяє економити на виправленні помилок. %good



(\,4\,) Гнучкі методи розробки нехтують проміжними технічними артефактами та документацією виключно через їх непотрібність. %bad


{\begin {center}\textbf{Завдання 8}\end {center}}\par
Що з наступного не має відношення до Scrum або не виконується/не використовується за використання оригінального Scrum?\par
\vspace{2mm}



(\,1\,) канбан%bad



(\,2\,) розповіді користувача%bad



(\,3\,) власник продукту призначає відповідальних за конкретні задачі%good



(\,4\,) покер планування (Planning Poker)%bad


{\begin {center}\textbf{Завдання 9}\end {center}}\par
Що з наведеного нижче не розглядалось та навіть не згадувалось у курсі?\par
\vspace{2mm}



(\,1\,) BDUF%bad



(\,2\,) CRAZY%good



(\,3\,) RUP%bad



(\,4\,) SAFe%bad


{\begin {center}\textbf{Завдання 10}\end {center}}\par
Вибрати всі істинні твердження (якщо існують).\par
\vspace{2mm}



(\,1\,) Користувачі бувають вторинними.%good



(\,2\,) Коли нема можливості перевірити всі можливі комбінації, тестують їх попарні (pairwise) комбінації. Це може значно зменшити кількість тестів.%good



(\,3\,) Тестування вимог може породжувати їх уточнення.%good



(\,4\,) Модель якості при використанні є ієрархічною.%good



\end {large}}
{\vspace{5mm}\smallskip \begin {small} Затверджено на засіданні кафедри теоретичної кібернетики, протокол \No 6 від   29.01.2021 р.\par\smallskip\smallskip
{\parbox {6.5 cm}{Зав. кафедри \dotfill Ю.В.~Крак}}\hspace {0.7cm}{\hbox to 6.5 cm{ Екзаменатор \dotfill Т.О.~Карнаух}} \end{small}}
\newpage

{{\begin {small}\par
{ \center  Київський національний університет імені Тараса Шевченка \par}
{\parbox {5 cm}{Освітня програма \hfill}{\parbox {7 cm}{Інформатика \hfill}}\parbox {6 cm}{\hfill Семестр 6}}\par
{\parbox {5 cm} {Навчальний предмет \hfill}\parbox {7 cm}{Розробка програмного забезпечення \hfill}\parbox {6cm}{\hfill}\par}\vspace{-7pt} \end{small}}{\center Екзаменаційний білет \No \textbf
{ 16 }\par}}
{
\begin {large}
{\begin {center}\textbf{Завдання 1}\end {center}}\par
Шаблон проектування:\par
\vspace{2mm}



(\,1\,) Тільки формулює ідею розв'язання проблеми%bad



(\,2\,) Задає розв'язання типової проблеми до найдрібніших подробиць%bad



(\,3\,) Є шаблоном, в який треба підставити власні ідентифікатори%bad



(\,4\,) Формулює ідею розв'язання проблеми та може пропонувати одну чи кілька її реалізацій%good


{\begin {center}\textbf{Завдання 2}\end {center}}\par
Який з наведених фрагментів коду явно потребує поліпшення (форматування рядків та типи до уваги не брати)?\par
\vspace{2mm}



(\,1\,) void ff(…);%bad



(\,2\,) \{int x; int y; …\}%good



(\,3\,) \{double maxPower, minPower;…\} %good



(\,4\,) \{fstream inputFile, outputFile;…\} %good


{\begin {center}\textbf{Завдання 3}\end {center}}\par
Що є характерним для структурного програмування?\par
\vspace{2mm}



(\,1\,) Заборона використання скалярних (простих) типів даних. %bad



(\,2\,) Переходи з середини одного блоку в середину іншого неможливі. %good



(\,3\,) У коді мають використовуватись типи структур. %bad



(\,4\,) У коді мають використовуватись типи структур, класи використовувати заборонено. %bad


{\begin {center}\textbf{Завдання 4}\end {center}}\par
Вибрати всі істинні твердження (якщо існують).\par
\vspace{2mm}



(\,1\,) Адаптивна модель розробки адаптує вимоги до продукту до наявної команди розробників. %bad



(\,2\,) Прогнозні моделі розробки з самого початку фіксують вимоги до продукту. %good



(\,3\,) Гнучкі методології не передбачають планування робіт. %bad



(\,4\,) Прогнозна модель розробки може одночасно бути інкрементною. %good


{\begin {center}\textbf{Завдання 5}\end {center}}\par
Вибрати всі істинні твердження (якщо існують).\par
\vspace{2mm}



(\,1\,) Довжина функцій не впливає на внутрішню якість коду. %bad



(\,2\,) Чим вища кваліфікація програміста, тим більш складні для розуміння та довгі умовні вирази він використовує. %bad



(\,3\,) Код програми пишеться для комп'ютера, тому він має бути синтаксично та семантично коректним. Інше значення не має. %bad



(\,4\,) Повторюваний код дозволяє програмісту краще зрозуміти його логіку та вибрати з кількох варіантів найкращий. %bad

\newpage

{\begin {center}\textbf{Завдання 6}\end {center}}\par
Вибрати всі істинні твердження (якщо існують).\par
\vspace{2mm}



(\,1\,) RUP є компонентно-орієнтованим. %good



(\,2\,) Розробка за гнучкими методологіями ніколи не слідує плану. %bad



(\,3\,) Розповіді користувача (user stories) – це прогресивний інструмент гнучких методологій, який нічого спільного не має з прецедентами (use case). %bad



(\,4\,) Основна відмінність V-моделі від водоспадної полягає в тім, що діяльності з тестування виконуються паралельно з усіма іншими діяльностями з розробки. %good


{\begin {center}\textbf{Завдання 7}\end {center}}\par
Вибрати всі істинні твердження (якщо існують).\par
\vspace{2mm}



(\,1\,) Гнучкі методи розробки нехтують проміжними технічними артефактами та документацією виключно через їх непотрібність. %bad



(\,2\,) Тільки RUP використовує UML.%bad



(\,3\,) Водоспадна методологія розробки може призводити до недостатнього проектування. %bad



(\,4\,) Перша версія коду не завжди може бути якісною -- і це нормально. %good


{\begin {center}\textbf{Завдання 8}\end {center}}\par
Що з наступного не має відношення до Scrum або не виконується/не використовується за використання оригінального Scrum?\par
\vspace{2mm}



(\,1\,) скрам-мастер визначає, які задачі команда буде виконувати на наступному спрінті%good



(\,2\,) журнал невиконаних задач (backlog) продукту%bad



(\,3\,) порядок виконання запланованої на спрінт роботи визначають розробники%bad



(\,4\,) обернені вимоги%bad


{\begin {center}\textbf{Завдання 9}\end {center}}\par
Що з наведеного нижче не розглядалось та навіть не згадувалось у курсі?\par
\vspace{2mm}



(\,1\,) RUP%bad



(\,2\,) Lean%bad



(\,3\,) Scrum%bad



(\,4\,) sashimi%bad


{\begin {center}\textbf{Завдання 10}\end {center}}\par
Вибрати всі істинні твердження (якщо існують).\par
\vspace{2mm}



(\,1\,) Для тестування код треба ізолювати. Для того, щоб ізолювати, треба мати тести.%good



(\,2\,) Модель якості продукту є ієрархічною.%good



(\,3\,) Якість буває внутрішньою.%good



(\,4\,) Якість буває зовнішньою.%good



\end {large}}
{\vspace{5mm}\smallskip \begin {small} Затверджено на засіданні кафедри теоретичної кібернетики, протокол \No 6 від   29.01.2021 р.\par\smallskip\smallskip
{\parbox {6.5 cm}{Зав. кафедри \dotfill Ю.В.~Крак}}\hspace {0.7cm}{\hbox to 6.5 cm{ Екзаменатор \dotfill Т.О.~Карнаух}} \end{small}}
\newpage

{{\begin {small}\par
{ \center  Київський національний університет імені Тараса Шевченка \par}
{\parbox {5 cm}{Освітня програма \hfill}{\parbox {7 cm}{Інформатика \hfill}}\parbox {6 cm}{\hfill Семестр 6}}\par
{\parbox {5 cm} {Навчальний предмет \hfill}\parbox {7 cm}{Розробка програмного забезпечення \hfill}\parbox {6cm}{\hfill}\par}\vspace{-7pt} \end{small}}{\center Екзаменаційний білет \No \textbf
{ 17 }\par}}
{
\begin {large}
{\begin {center}\textbf{Завдання 1}\end {center}}\par
Так звані трикутники балансування зв'язують між собою:\par
\vspace{2mm}



(\,1\,) Завжди рівно 3 параметри%bad



(\,2\,) Завжди рівно 4 параметри%bad



(\,3\,) 4 або 3 параметри, при цьому може бути відсутня якість%bad



(\,4\,) 4 або 3 параметри, при цьому може бути відсутній обсяг %good


{\begin {center}\textbf{Завдання 2}\end {center}}\par
Який з наведених фрагментів коду явно потребує поліпшення (форматування рядків та типи до уваги не брати)?\par
\vspace{2mm}



(\,1\,) \{int x; int y; …\}%good



(\,2\,) void help();%good



(\,3\,) \{int input\_file, file\_to\_output; …\}%bad



(\,4\,) void load\_data(…);%good


{\begin {center}\textbf{Завдання 3}\end {center}}\par
Що є характерним для структурного програмування?\par
\vspace{2mm}



(\,1\,) Кожен блок коду повинен мати єдиний вхід та єдиний вихід.%good



(\,2\,) У коді мають використовуватись типи даних, що є класами. %bad



(\,3\,) У коді мають використовуватись типи структур, класи використовувати заборонено. %bad



(\,4\,) Заборона використання скалярних (простих) типів даних. %bad


{\begin {center}\textbf{Завдання 4}\end {center}}\par
Вибрати всі істинні твердження (якщо існують).\par
\vspace{2mm}



(\,1\,) І прогнозні, і адаптивні методології передбачають планування виконуваних робіт. %good



(\,2\,) Неперервність проектування полягає в тім, що процес розробки не можна переривати більше ніж на вихідні. %bad



(\,3\,) У випадку, коли вимоги до продукту відомі заздалегідь та є фіксованими контрактом, адаптивні методології розробки використовувати не можна. %bad



(\,4\,) Прогнозні та адаптивні методології використовують одні ті самі діяльності з розробки ПЗ: визначення та аналіз вимог, проектування, реалізація, тестування, впровадження, супроводження. %good


{\begin {center}\textbf{Завдання 5}\end {center}}\par
Вибрати всі істинні твердження (якщо існують).\par
\vspace{2mm}



(\,1\,) Чим менше класів у коді, тим менше класів треба налагоджувати, а тому класи за можливості слід об'єднувати. %bad



(\,2\,) Одним з прийомів спрощення умовної логіки є використання поліморфізму. %good



(\,3\,) Від довгих функцій можна позбутися, розклавши їх у послідовність викликів дрібніших функцій. %good



(\,4\,) Код, що виконує однакову послідовність логічних кроків, у більшості випадків є повторюваним.%good

\newpage

{\begin {center}\textbf{Завдання 6}\end {center}}\par
Вибрати всі істинні твердження (якщо існують).\par
\vspace{2mm}



(\,1\,) RUP є компонентно-орієнтованим. %good



(\,2\,) Scrum передбачає колективне володіння кодом. %good



(\,3\,) Сьогодні водоспадна методологія не має права на існування. %bad



(\,4\,) RUP є керованим прецедентами. %good


{\begin {center}\textbf{Завдання 7}\end {center}}\par
Вибрати всі істинні твердження (якщо існують).\par
\vspace{2mm}



(\,1\,) Без рефакторінгу неперервне проектування неможливе. %good



(\,2\,) Розробка за гнучкими методологіями цінить людей та взаємодію між ними. %good



(\,3\,) Використання гнучких методологій при розробці надвеликих проектів не несе в собі жодних ризиків. %bad



(\,4\,) Програмування в парах (pair programming) спрощує введення в команду нових розробників. %good


{\begin {center}\textbf{Завдання 8}\end {center}}\par
Що з наступного не має відношення до Scrum або не виконується/не використовується за використання оригінального Scrum?\par
\vspace{2mm}



(\,1\,) бачення продукту (vision) %bad



(\,2\,) власник продукту остаточно затверджує проект системи%good



(\,3\,) розповіді користувача%bad



(\,4\,) скрам-мастер призначає відповідальних за конкретні задачі%good


{\begin {center}\textbf{Завдання 9}\end {center}}\par
Що з наведеного нижче не розглядалось та навіть не згадувалось у курсі?\par
\vspace{2mm}



(\,1\,) SAFe%bad



(\,2\,) LOAN%good



(\,3\,) спіральна модель%bad



(\,4\,) BDD%bad


{\begin {center}\textbf{Завдання 10}\end {center}}\par
Вибрати всі істинні твердження (якщо існують).\par
\vspace{2mm}



(\,1\,) Користувачі бувають непрямими.%good



(\,2\,) Розглядають модель якості при використанні.%good



(\,3\,) За провалених димових тестів решта тестів не виконується і код повертається на доробку.%good



(\,4\,) Якість можна визначати через дефекти.%good



\end {large}}
{\vspace{5mm}\smallskip \begin {small} Затверджено на засіданні кафедри теоретичної кібернетики, протокол \No 6 від   29.01.2021 р.\par\smallskip\smallskip
{\parbox {6.5 cm}{Зав. кафедри \dotfill Ю.В.~Крак}}\hspace {0.7cm}{\hbox to 6.5 cm{ Екзаменатор \dotfill Т.О.~Карнаух}} \end{small}}
\newpage

{{\begin {small}\par
{ \center  Київський національний університет імені Тараса Шевченка \par}
{\parbox {5 cm}{Освітня програма \hfill}{\parbox {7 cm}{Інформатика \hfill}}\parbox {6 cm}{\hfill Семестр 6}}\par
{\parbox {5 cm} {Навчальний предмет \hfill}\parbox {7 cm}{Розробка програмного забезпечення \hfill}\parbox {6cm}{\hfill}\par}\vspace{-7pt} \end{small}}{\center Екзаменаційний білет \No \textbf
{ 18 }\par}}
{
\begin {large}
{\begin {center}\textbf{Завдання 1}\end {center}}\par
Так звані трикутники балансування зв'язують між собою:\par
\vspace{2mm}



(\,1\,) 4 або 3 параметри, при цьому може бути відсутній обсяг %good



(\,2\,) Завжди рівно 3 параметри%bad



(\,3\,) 4 або 3 параметри, при цьому може бути відсутня якість%bad



(\,4\,) 4 або 3 параметри, при цьому може бути відсутній час%bad


{\begin {center}\textbf{Завдання 2}\end {center}}\par
Який з наведених фрагментів коду явно потребує поліпшення (форматування рядків та типи до уваги не брати)?\par
\vspace{2mm}



(\,1\,) void load\_data(T \& arg1, double arg2, double arg3, double arg4); %bad



(\,2\,) \{int xtrmprg; …\}%bad



(\,3\,) void f13(…);%bad



(\,4\,) void ff(…);%bad


{\begin {center}\textbf{Завдання 3}\end {center}}\par
Що є характерним для структурного програмування?\par
\vspace{2mm}



(\,1\,) Програма зображується ієрархічною структурою блоків.%good



(\,2\,) Використання оператора безумовного переходу. %bad



(\,3\,) У коді мають використовуватись типи структур. %bad



(\,4\,) У коді мають використовуватись типи даних, що є класами. %bad


{\begin {center}\textbf{Завдання 4}\end {center}}\par
Вибрати всі істинні твердження (якщо існують).\par
\vspace{2mm}



(\,1\,) Адаптивні моделі розробки не допускають зміну вимог до продукту, але продукт може розроблятись поетапно. %bad



(\,2\,) Неперервне проектування має бути ітеративним та інкрементним одночасно. %good



(\,3\,) Неперервне проектування здійснює створення та модифікацію проекту системи в міру розробки. %good



(\,4\,) Адаптивні моделі розробки допускають зміну вимог до продукту. %good


{\begin {center}\textbf{Завдання 5}\end {center}}\par
Вибрати всі істинні твердження (якщо існують).\par
\vspace{2mm}



(\,1\,) Одним з прийомів спрощення складної умовної логіки є використання шаблону Strategy. %good



(\,2\,) Код з довгими функціями має меншу кількість функцій, а тому для нього треба менше тестів. %bad



(\,3\,) Повторюваний код має повторюватись у точності до символу, інакше він повторюваним вважатись не може. %bad



(\,4\,) Чим менше модулів у проекті, тим краще, бо меншу кількість модулів треба інтегрувати в єдине ціле. %bad

\newpage

{\begin {center}\textbf{Завдання 6}\end {center}}\par
Вибрати всі істинні твердження (якщо існують).\par
\vspace{2mm}



(\,1\,) Висока внутрішня якість коду здешевшує розробку, бо дозволяє економити на вартості та часі майбутніх змін. %good



(\,2\,) RUP відрізняється від Scrum тим, що ні в що не ставить персонал.%bad



(\,3\,) Використання гнучких методологій вимагає високої кваліфікації персоналу. %good



(\,4\,) За використання гнучких методологій розробки припиняти розробку не рекомендується. %good


{\begin {center}\textbf{Завдання 7}\end {center}}\par
Вибрати всі істинні твердження (якщо існують).\par
\vspace{2mm}



(\,1\,) Простота – мистецтво максимізації незробленої роботи. %good



(\,2\,) Водоспадна методологія не призначена для внесення змін у вимоги до системи. %good



(\,3\,) RUP є ітеративним та інкрементним.%good



(\,4\,) Адаптивні методології розробки дозволяють уникнути надлишкового проектування. %good


{\begin {center}\textbf{Завдання 8}\end {center}}\par
Що з наступного не має відношення до Scrum або не виконується/не використовується за використання оригінального Scrum?\par
\vspace{2mm}



(\,1\,) власник продукту призначає відповідальних за конкретні задачі%good



(\,2\,) команда визначає, які задачі вона буде виконувати на наступному спрінті%good



(\,3\,) власник продукту остаточно затверджує проект системи%good



(\,4\,) порядок виконання запланованої на спрінт роботи визначає власник продукту%good


{\begin {center}\textbf{Завдання 9}\end {center}}\par
Що з наведеного нижче не розглядалось та навіть не згадувалось у курсі?\par
\vspace{2mm}



(\,1\,) XP%bad



(\,2\,) X-модель%good



(\,3\,) LAZY%good



(\,4\,) модель хаосу%bad


{\begin {center}\textbf{Завдання 10}\end {center}}\par
Вибрати всі істинні твердження (якщо існують).\par
\vspace{2mm}



(\,1\,) Модель якості даних є ієрархічною.%good



(\,2\,) Розглядають модель якості даних.%good



(\,3\,) Розглядають модель якості продукту.%good



(\,4\,) Тестовий приклад складається з набору вхідних даних, умов виконання та очікуваних результатів.%good



\end {large}}
{\vspace{5mm}\smallskip \begin {small} Затверджено на засіданні кафедри теоретичної кібернетики, протокол \No 6 від   29.01.2021 р.\par\smallskip\smallskip
{\parbox {6.5 cm}{Зав. кафедри \dotfill Ю.В.~Крак}}\hspace {0.7cm}{\hbox to 6.5 cm{ Екзаменатор \dotfill Т.О.~Карнаух}} \end{small}}
\newpage

{{\begin {small}\par
{ \center  Київський національний університет імені Тараса Шевченка \par}
{\parbox {5 cm}{Освітня програма \hfill}{\parbox {7 cm}{Інформатика \hfill}}\parbox {6 cm}{\hfill Семестр 6}}\par
{\parbox {5 cm} {Навчальний предмет \hfill}\parbox {7 cm}{Розробка програмного забезпечення \hfill}\parbox {6cm}{\hfill}\par}\vspace{-7pt} \end{small}}{\center Екзаменаційний білет \No \textbf
{ 19 }\par}}
{
\begin {large}
{\begin {center}\textbf{Завдання 1}\end {center}}\par
Що з наведеного найкраще характеризує рефакторінг?\par
\vspace{2mm}



(\,1\,) Зміна внутрішньої структури коду під час усунення помилок%bad



(\,2\,) Докорінна переробка всього коду з нуля,  що не змінює його видиму поведінку%bad



(\,3\,) Зміна внутрішньої структури коду, що не змінює його видиму поведінку%good



(\,4\,) Довільна переробка коду, що виникає у процесі розробки%bad


{\begin {center}\textbf{Завдання 2}\end {center}}\par
Який з наведених фрагментів коду явно потребує поліпшення (форматування рядків та типи до уваги не брати)?\par
\vspace{2mm}



(\,1\,) \{fstream input\_file, outputFile; …\}%bad



(\,2\,) \{int a\_, a\_\_; …\}%bad



(\,3\,) \{double avg\_temperature;…\} %good



(\,4\,) \{fstream inputFile, outputFile;…\} %good


{\begin {center}\textbf{Завдання 3}\end {center}}\par
Що є характерним для структурного програмування?\par
\vspace{2mm}



(\,1\,) У коді мають використовуватись структурні типи даних. %bad



(\,2\,) Використання механізму винятків. %bad



(\,3\,) Заборона використання скалярних (простих) типів даних. %bad



(\,4\,) Переходи з середини одного блоку в середину іншого неможливі. %good


{\begin {center}\textbf{Завдання 4}\end {center}}\par
Вибрати всі істинні твердження (якщо існують).\par
\vspace{2mm}



(\,1\,) Прогнозна модель розробки не може одночасно бути інкрементною. %bad



(\,2\,) У прогнозній моделі розробки розробка може виконуватись ітеративно. %good



(\,3\,) Для адаптивних методологій розробки, на відміну від прогнозних, попередні оцінки бюджету та/або часу не виконують, бо вони все одно зміняться. %bad



(\,4\,) Адаптивні моделі розробки не допускають зміну вимог до продукту, але продукт може розроблятись поетапно. %bad


{\begin {center}\textbf{Завдання 5}\end {center}}\par
Вибрати всі істинні твердження (якщо існують).\par
\vspace{2mm}



(\,1\,) Одним з засобів позбутись повторень у коді є використання функцій. %good



(\,2\,) Залежності всередині модуля можуть бути сильними. %good



(\,3\,) Залежності між різними модулями системи мають бути якомога більш сильними, бо інакше вона розпадеться. %bad



(\,4\,) Код з довгими функціями має меншу кількість функцій, а тому його швидше та простіше тестувати. %bad

\newpage

{\begin {center}\textbf{Завдання 6}\end {center}}\par
Вибрати всі істинні твердження (якщо існують).\par
\vspace{2mm}



(\,1\,) Модель хаосу є методологією розробки програмного забезпечення. %bad



(\,2\,) Багато хто з тих, що використовує Scrum, складає архітектурну модель розроблюваної системи. %good



(\,3\,) Якщо виникають проблеми з модифікацією коду, то слід відразу змінити команду розробників на більш кваліфіковану.%bad



(\,4\,) Гнучкі методології передбачають часті випуски продукту. %good


{\begin {center}\textbf{Завдання 7}\end {center}}\par
Вибрати всі істинні твердження (якщо існують).\par
\vspace{2mm}



(\,1\,) На відміну від Scrum, RUP є компонентно-орієнтованим. %bad



(\,2\,) Розробка за гнучкими методологіями цінить співробітництво з замовником та піддатлівість змінам. %good



(\,3\,) Гнучкі методології повністю нехтують документацією. %bad



(\,4\,) Розповіді користувача (user stories) це менш формалізований порівняно з прецедентами спосіб опису вимог до системи. %good


{\begin {center}\textbf{Завдання 8}\end {center}}\par
Що з наступного не має відношення до Scrum або не виконується/не використовується за використання оригінального Scrum?\par
\vspace{2mm}



(\,1\,) виділення одного з учасників команди як системного архітектора%good



(\,2\,) під час спрінту зміни, що впливають на досягнення мети спрінту, зазвичай, не допускаються%bad



(\,3\,) виділення в команді одного чи кількох проектувальників%good



(\,4\,) покер планування (Planning Poker)%bad


{\begin {center}\textbf{Завдання 9}\end {center}}\par
Що з наведеного нижче не розглядалось та навіть не згадувалось у курсі?\par
\vspace{2mm}



(\,1\,) poker%bad



(\,2\,) BDUF%bad



(\,3\,) CRAZY%good



(\,4\,) V-модель%bad


{\begin {center}\textbf{Завдання 10}\end {center}}\par
Вибрати всі істинні твердження (якщо існують).\par
\vspace{2mm}



(\,1\,) Розглядають модель якості при використанні.%good



(\,2\,) Якість можна визначати через дефекти.%good



(\,3\,) Коли нема можливості перевірити всі можливі комбінації, тестують їх попарні (pairwise) комбінації. Це може значно зменшити кількість тестів.%good



(\,4\,) Модель якості при використанні є ієрархічною.%good



\end {large}}
{\vspace{5mm}\smallskip \begin {small} Затверджено на засіданні кафедри теоретичної кібернетики, протокол \No 6 від   29.01.2021 р.\par\smallskip\smallskip
{\parbox {6.5 cm}{Зав. кафедри \dotfill Ю.В.~Крак}}\hspace {0.7cm}{\hbox to 6.5 cm{ Екзаменатор \dotfill Т.О.~Карнаух}} \end{small}}
\newpage

{{\begin {small}\par
{ \center  Київський національний університет імені Тараса Шевченка \par}
{\parbox {5 cm}{Освітня програма \hfill}{\parbox {7 cm}{Інформатика \hfill}}\parbox {6 cm}{\hfill Семестр 6}}\par
{\parbox {5 cm} {Навчальний предмет \hfill}\parbox {7 cm}{Розробка програмного забезпечення \hfill}\parbox {6cm}{\hfill}\par}\vspace{-7pt} \end{small}}{\center Екзаменаційний білет \No \textbf
{ 20 }\par}}
{
\begin {large}
{\begin {center}\textbf{Завдання 1}\end {center}}\par
Так звані трикутники балансування зв'язують між собою:\par
\vspace{2mm}



(\,1\,) 4 або 3 параметри, при цьому може бути відсутній час%bad



(\,2\,) Завжди рівно 4 параметри%bad



(\,3\,) Завжди рівно 3 параметри%bad



(\,4\,) 4 або 3 параметри, при цьому може бути відсутній обсяг %good


{\begin {center}\textbf{Завдання 2}\end {center}}\par
Який з наведених фрагментів коду явно потребує поліпшення (форматування рядків та типи до уваги не брати)?\par
\vspace{2mm}



(\,1\,) \{double maxPower, minPower;…\} %good



(\,2\,) \{fstream input\_file, output\_file;…\} %good



(\,3\,) \{double max\_power, min\_power;…\} %good



(\,4\,) \{double the\_best\_value\_anyone\_ever\_have;\} %bad


{\begin {center}\textbf{Завдання 3}\end {center}}\par
Що є характерним для структурного програмування?\par
\vspace{2mm}



(\,1\,) Кожен блок коду повинен мати єдиний вхід та єдиний вихід.%good



(\,2\,) У коді мають використовуватись типи структур, класи використовувати заборонено. %bad



(\,3\,) У коді мають використовуватись типи структур. %bad



(\,4\,) У коді мають використовуватись типи даних, що є класами. %bad


{\begin {center}\textbf{Завдання 4}\end {center}}\par
Вибрати всі істинні твердження (якщо існують).\par
\vspace{2mm}



(\,1\,) Адаптивна модель розробки адаптує вимоги до продукту до наявної команди розробників. %bad



(\,2\,) В адаптивних методологіях планування робіт не виконується. %bad



(\,3\,) Ітеративні та інкрементні методи розробки є синонімом неперервного проектування. %bad



(\,4\,) Прогнозні моделі розробки з самого початку фіксують вимоги до продукту. %good


{\begin {center}\textbf{Завдання 5}\end {center}}\par
Вибрати всі істинні твердження (якщо існують).\par
\vspace{2mm}



(\,1\,) Втілення в коді різних підходів до одного того самого перетворення даних дозволяє програмісту продемонструвати свою професійну підготовку з найкращого боку. %bad



(\,2\,) Довгі функції доволі часто приховують складну логіку обробки або повторюваний код. %good



(\,3\,) Складні умови виникають в реальних задачах доволі часто, а тому є доволі природніми і не можуть вважатись недоліком коду. %bad



(\,4\,) Одним з засобів позбутись повторень у коді є використання успадкування. %good

\newpage

{\begin {center}\textbf{Завдання 6}\end {center}}\par
Вибрати всі істинні твердження (якщо існують).\par
\vspace{2mm}



(\,1\,) Гнучкі методології можуть передбачати як спільне, так і індивідуальне володіння кодом. %good



(\,2\,) RUP принципово відрізняється від Scrum тим, що є керованим прецедентами. У той же час у Scrum при плануванні розробки прецеденти (чи їх аналоги) до уваги не беруться. %bad



(\,3\,) Водоспадна методологія розробки може призводити до надлишкового проектування. %good



(\,4\,) Ітеративна методологія розробки у загальному випадку не гарантує, що не буде надлишкового проектування. %good


{\begin {center}\textbf{Завдання 7}\end {center}}\par
Вибрати всі істинні твердження (якщо існують).\par
\vspace{2mm}



(\,1\,) Методологія розробки Big Design Up Front погана тим, що може призводити до надлишкового проектування. %good



(\,2\,) Ітеративна та інкрементна розробка передбачає наявність у тому чи іншому вигляді етапу ініціалізації, ітерацій та керуючої таблиці проекту. %good



(\,3\,) За використання гнучких методологій кожен розробник робить, що хоче. %bad



(\,4\,) Гнучкі методи розробки нехтують проміжними технічними артефактами та документацією через те, що за зміни вимог та/або проекту треба витрачати додатковий час на підтримання їх у відповідному стані. %good


{\begin {center}\textbf{Завдання 8}\end {center}}\par
Що з наступного не має відношення до Scrum або не виконується/не використовується за використання оригінального Scrum?\par
\vspace{2mm}



(\,1\,) обернені вимоги%bad



(\,2\,) власник продукту призначає час, необхідний команді для реалізації конкретної задачі%good



(\,3\,) скрам-мастер остаточно затверджує проект системи%good



(\,4\,) журнал невиконаних задач (backlog) спрінту%bad


{\begin {center}\textbf{Завдання 9}\end {center}}\par
Що з наведеного нижче не розглядалось та навіть не згадувалось у курсі?\par
\vspace{2mm}



(\,1\,) DSDM%bad



(\,2\,) TDD%bad



(\,3\,) LeSS%bad



(\,4\,) PRINCE2%good


{\begin {center}\textbf{Завдання 10}\end {center}}\par
Вибрати всі істинні твердження (якщо існують).\par
\vspace{2mm}



(\,1\,) Забезпечення якості це не тільки тестування.%good



(\,2\,) Користувачі бувають вторинними.%good



(\,3\,) Тестування вимог може породжувати їх уточнення.%good



(\,4\,) І верифікація, і валідація можуть використовувати тестування.%good



\end {large}}
{\vspace{5mm}\smallskip \begin {small} Затверджено на засіданні кафедри теоретичної кібернетики, протокол \No 6 від   29.01.2021 р.\par\smallskip\smallskip
{\parbox {6.5 cm}{Зав. кафедри \dotfill Ю.В.~Крак}}\hspace {0.7cm}{\hbox to 6.5 cm{ Екзаменатор \dotfill Т.О.~Карнаух}} \end{small}}
\newpage

{{\begin {small}\par
{ \center  Київський національний університет імені Тараса Шевченка \par}
{\parbox {5 cm}{Освітня програма \hfill}{\parbox {7 cm}{Інформатика \hfill}}\parbox {6 cm}{\hfill Семестр 6}}\par
{\parbox {5 cm} {Навчальний предмет \hfill}\parbox {7 cm}{Розробка програмного забезпечення \hfill}\parbox {6cm}{\hfill}\par}\vspace{-7pt} \end{small}}{\center Екзаменаційний білет \No \textbf
{ 21 }\par}}
{
\begin {large}
{\begin {center}\textbf{Завдання 1}\end {center}}\par
Що є метою рефакторингу?\par
\vspace{2mm}



(\,1\,) Додавання нової функціональності%bad



(\,2\,) Поліпшення внутрішньої структури коду%good



(\,3\,) Усунення помилок у коді%bad



(\,4\,) Зменшення кількості рядків коду%bad


{\begin {center}\textbf{Завдання 2}\end {center}}\par
Який з наведених фрагментів коду явно потребує поліпшення (форматування рядків та типи до уваги не брати)?\par
\vspace{2mm}



(\,1\,) \{double maxPower, minPower;…\} %good



(\,2\,) \{int a\_, a\_\_; …\}%bad



(\,3\,) \{int xtrmprg; …\}%bad



(\,4\,) void ff(…);%bad


{\begin {center}\textbf{Завдання 3}\end {center}}\par
Що є характерним для структурного програмування?\par
\vspace{2mm}



(\,1\,) Використання оператора безумовного переходу. %bad



(\,2\,) Використання механізму винятків. %bad



(\,3\,) Програма зображується ієрархічною структурою блоків.%good



(\,4\,) У коді мають використовуватись структурні типи даних. %bad


{\begin {center}\textbf{Завдання 4}\end {center}}\par
Вибрати всі істинні твердження (якщо існують).\par
\vspace{2mm}



(\,1\,) Адаптивні методології має сенс використовувати для проектів, де можлива зміна вимог. %good



(\,2\,) Для адаптивних методологій прогнозувати результат чи терміни його отримання абсолютно неможливо. %bad



(\,3\,) Адаптивні моделі розробки допускають зміну вимог до продукту. %good



(\,4\,) Для адаптивних моделей розробки оцінки часу та бюджету не виконують, бо вимоги є змінними. %bad


{\begin {center}\textbf{Завдання 5}\end {center}}\par
Вибрати всі істинні твердження (якщо існують).\par
\vspace{2mm}



(\,1\,) Одним з прийомів спрощення складних умов є їх декомпозиція. %good



(\,2\,) Зв'язки між окремими модулями системи мають бути якомога більш слабкими. %good



(\,3\,) Складні комбінації умов починають керувати логікою виконання, тому що такими є вимоги до функціонування ПЗ. Оскільки ПЗ має якнайкраще забезпечувати потреби замовника, то складні комбінації умов у коді вважатись недоліком коду не можуть. %bad



(\,4\,) Гарний код не повинен містити різних фрагментів, що складаються з кількох інструкцій та виконують однакове перетворення даних. %good

\newpage

{\begin {center}\textbf{Завдання 6}\end {center}}\par
Вибрати всі істинні твердження (якщо існують).\par
\vspace{2mm}



(\,1\,) RUP є прогнозною методологією. %bad



(\,2\,) Гнучкі методології базуються на комунікації розробників, а тому передбачають виключно колективне володіння кодом. %bad



(\,3\,) Якщо виникають проблеми з модифікацією коду, то слід розглянути внесення змін у проект. %good



(\,4\,) Водоспадна методологія може мати тенденцію до надлишкового проектування. %good


{\begin {center}\textbf{Завдання 7}\end {center}}\par
Вибрати всі істинні твердження (якщо існують).\par
\vspace{2mm}



(\,1\,) За невизначених або піддатливих змінам вимог краще застосовувати гнучкі методології розробки. %good



(\,2\,) Поступове уточнення вимог замовника полегшує адаптацію до змін. %good



(\,3\,) Одна з принципових відмінностей RUP від Scrum полягає у використанні численних проміжних технічних артефактів. %good



(\,4\,) Висока внутрішня якість коду є марною тратою часу на те, що замовник не побачить та/або не зможе зацінити. %bad


{\begin {center}\textbf{Завдання 8}\end {center}}\par
Що з наступного не має відношення до Scrum або не виконується/не використовується за використання оригінального Scrum?\par
\vspace{2mm}



(\,1\,) журнал невиконаних задач (backlog) продукту%bad



(\,2\,) розповіді користувача%bad



(\,3\,) щоденний скрам%bad



(\,4\,) блочні тести%bad


{\begin {center}\textbf{Завдання 9}\end {center}}\par
Що з наведеного нижче не розглядалось та навіть не згадувалось у курсі?\par
\vspace{2mm}



(\,1\,) CRAZY%good



(\,2\,) Scrum of scrums%bad



(\,3\,) FDD%bad



(\,4\,) X-модель%good


{\begin {center}\textbf{Завдання 10}\end {center}}\par
Вибрати всі істинні твердження (якщо існують).\par
\vspace{2mm}



(\,1\,) Якість можна визначати через атрибути якості.%good



(\,2\,) Розглядають модель якості продукту.%good



(\,3\,) Користувачі бувають основними.%good



(\,4\,) Для тестування код треба ізолювати. Для того, щоб ізолювати, треба мати тести.%good



\end {large}}
{\vspace{5mm}\smallskip \begin {small} Затверджено на засіданні кафедри теоретичної кібернетики, протокол \No 6 від   29.01.2021 р.\par\smallskip\smallskip
{\parbox {6.5 cm}{Зав. кафедри \dotfill Ю.В.~Крак}}\hspace {0.7cm}{\hbox to 6.5 cm{ Екзаменатор \dotfill Т.О.~Карнаух}} \end{small}}
\newpage

{{\begin {small}\par
{ \center  Київський національний університет імені Тараса Шевченка \par}
{\parbox {5 cm}{Освітня програма \hfill}{\parbox {7 cm}{Інформатика \hfill}}\parbox {6 cm}{\hfill Семестр 6}}\par
{\parbox {5 cm} {Навчальний предмет \hfill}\parbox {7 cm}{Розробка програмного забезпечення \hfill}\parbox {6cm}{\hfill}\par}\vspace{-7pt} \end{small}}{\center Екзаменаційний білет \No \textbf
{ 22 }\par}}
{
\begin {large}
{\begin {center}\textbf{Завдання 1}\end {center}}\par
Так звані трикутники балансування зв'язують між собою:\par
\vspace{2mm}



(\,1\,) 4 або 3 параметри, при цьому може бути відсутній обсяг %good



(\,2\,) 4 або 3 параметри, при цьому може бути відсутній час%bad



(\,3\,) Завжди рівно 3 параметри%bad



(\,4\,) 4 або 3 параметри, при цьому може бути відсутня якість%bad


{\begin {center}\textbf{Завдання 2}\end {center}}\par
Який з наведених фрагментів коду явно потребує поліпшення (форматування рядків та типи до уваги не брати)?\par
\vspace{2mm}



(\,1\,) void f13(…);%bad



(\,2\,) \{double max\_power, min\_power;…\} %good



(\,3\,) \{fstream input\_file, output\_file;…\} %good



(\,4\,) \{double avg\_temperature;…\} %good


{\begin {center}\textbf{Завдання 3}\end {center}}\par
Що є характерним для структурного програмування?\par
\vspace{2mm}



(\,1\,) Використання оператора безумовного переходу. %bad



(\,2\,) Переходи з середини одного блоку в середину іншого неможливі. %good



(\,3\,) Використання механізму винятків. %bad



(\,4\,) У коді мають використовуватись типи структур. %bad


{\begin {center}\textbf{Завдання 4}\end {center}}\par
Вибрати всі істинні твердження (якщо існують).\par
\vspace{2mm}



(\,1\,) Гнучкі методології не передбачають планування робіт. %bad



(\,2\,) Неперервне проектування має бути ітеративним та інкрементним одночасно. %good



(\,3\,) Прогнозні та адаптивні методології використовують одні ті самі діяльності з розробки ПЗ: визначення та аналіз вимог, проектування, реалізація, тестування, впровадження, супроводження. %good



(\,4\,) Адаптивні методології не передбачають визначення вимог. %bad


{\begin {center}\textbf{Завдання 5}\end {center}}\par
Вибрати всі істинні твердження (якщо існують).\par
\vspace{2mm}



(\,1\,) Гарний код має читатись як розмовна мова. %good



(\,2\,) Одним з прийомів спрощення складної умовної логіки є використання шаблону State. %good



(\,3\,) Правильно спроектована функція повинна мати рівно один обов'язок. %good



(\,4\,) Код програми пишеться для комп'ютера, тому він має бути синтаксично та семантично коректним. Інше значення не має. %bad

\newpage

{\begin {center}\textbf{Завдання 6}\end {center}}\par
Вибрати всі істинні твердження (якщо існують).\par
\vspace{2mm}



(\,1\,) Гнучкі методології розробки є стійкими до припинення розробки на певний час. %bad



(\,2\,) Розробка за гнучкими методологіями ніколи не слідує плану. %bad



(\,3\,) Програмування в парах (pair programming) підвищує якість коду та дозволяє економити на виправленні помилок. %good



(\,4\,) Гнучкі методології розробки намагаються якомога швидше визначити архітектуру системи, для чого на початку завжди виконується початкове проектування системи в цілому, а тільки потім переходять до ітеративного уточнення деталей та кодування. %bad


{\begin {center}\textbf{Завдання 7}\end {center}}\par
Вибрати всі істинні твердження (якщо існують).\par
\vspace{2mm}



(\,1\,) Методологія розробки Big Design Up Front погана тим, що може призводити до недостатнього проектування.%bad



(\,2\,) Програмування в парах (pair programming) зводиться до того, що кожним фрагментом коду володіють два програмісти. %bad



(\,3\,) Гнучкі методології розробки не допускають використання UML. %bad



(\,4\,) За програмування в парах (pair programming) поки один колега набирає код на комп'ютері інший може піти пограти в настільний теніс. %bad


{\begin {center}\textbf{Завдання 8}\end {center}}\par
Що з наступного не має відношення до Scrum або не виконується/не використовується за використання оригінального Scrum?\par
\vspace{2mm}



(\,1\,) бачення продукту (vision) %bad



(\,2\,) порядок виконання запланованої на спрінт роботи визначають розробники%bad



(\,3\,) виділення в команді одного чи кількох тестувальників%good



(\,4\,) залучення зовнішніх фахівців для побудови архітектури системи%good


{\begin {center}\textbf{Завдання 9}\end {center}}\par
Що з наведеного нижче не розглядалось та навіть не згадувалось у курсі?\par
\vspace{2mm}



(\,1\,) LeSS%bad



(\,2\,) poker%bad



(\,3\,) Cleanroom%bad



(\,4\,) RAD%bad


{\begin {center}\textbf{Завдання 10}\end {center}}\par
Вибрати всі істинні твердження (якщо існують).\par
\vspace{2mm}



(\,1\,) Користувачі бувають непрямими.%good



(\,2\,) Тестування буває позитивне та негативне.%good



(\,3\,) За провалених димових тестів решта тестів не виконується і код повертається на доробку.%good



(\,4\,) Можна розглядати якість під час використання.%good



\end {large}}
{\vspace{5mm}\smallskip \begin {small} Затверджено на засіданні кафедри теоретичної кібернетики, протокол \No 6 від   29.01.2021 р.\par\smallskip\smallskip
{\parbox {6.5 cm}{Зав. кафедри \dotfill Ю.В.~Крак}}\hspace {0.7cm}{\hbox to 6.5 cm{ Екзаменатор \dotfill Т.О.~Карнаух}} \end{small}}
\newpage

{{\begin {small}\par
{ \center  Київський національний університет імені Тараса Шевченка \par}
{\parbox {5 cm}{Освітня програма \hfill}{\parbox {7 cm}{Інформатика \hfill}}\parbox {6 cm}{\hfill Семестр 6}}\par
{\parbox {5 cm} {Навчальний предмет \hfill}\parbox {7 cm}{Розробка програмного забезпечення \hfill}\parbox {6cm}{\hfill}\par}\vspace{-7pt} \end{small}}{\center Екзаменаційний білет \No \textbf
{ 23 }\par}}
{
\begin {large}
{\begin {center}\textbf{Завдання 1}\end {center}}\par
Що є метою рефакторингу?\par
\vspace{2mm}



(\,1\,) Усунення помилок у коді%bad



(\,2\,) Зменшення кількості рядків коду%bad



(\,3\,) Поліпшення внутрішньої структури коду%good



(\,4\,) Додавання нової функціональності%bad


{\begin {center}\textbf{Завдання 2}\end {center}}\par
Який з наведених фрагментів коду явно потребує поліпшення (форматування рядків та типи до уваги не брати)?\par
\vspace{2mm}



(\,1\,) \{double the\_best\_value\_anyone\_ever\_have;\} %bad



(\,2\,) void help();%good



(\,3\,) \{fstream input\_file, outputFile; …\}%bad



(\,4\,) void load\_data(T \& arg1, double arg2, double arg3, double arg4); %bad


{\begin {center}\textbf{Завдання 3}\end {center}}\par
Що є характерним для структурного програмування?\par
\vspace{2mm}



(\,1\,) У коді мають використовуватись типи даних, що є класами. %bad



(\,2\,) Заборона використання скалярних (простих) типів даних. %bad



(\,3\,) Кожен блок коду повинен мати єдиний вхід та єдиний вихід.%good



(\,4\,) У коді мають використовуватись типи структур, класи використовувати заборонено. %bad


{\begin {center}\textbf{Завдання 4}\end {center}}\par
Вибрати всі істинні твердження (якщо існують).\par
\vspace{2mm}



(\,1\,) Неперервне проектування створює специфікацію проекту перед початком розробки/ітерації, а під час ітерації його неперервно модифікує.%bad



(\,2\,) Прогнозна модель розробки може одночасно бути інкрементною. %good



(\,3\,) У прогнозній моделі розробки розробка не може виконуватись ітеративно. %bad



(\,4\,) У випадку, коли вимоги до продукту відомі заздалегідь та є фіксованими контрактом, адаптивні методології розробки використовувати не можна. %bad


{\begin {center}\textbf{Завдання 5}\end {center}}\par
Вибрати всі істинні твердження (якщо існують).\par
\vspace{2mm}



(\,1\,) Чим менше класів у коді, тим менше класів треба налагоджувати, а тому класи за можливості слід об'єднувати. %bad



(\,2\,) Гарний код не повинен містити різних фрагментів, що складаються з кількох інструкцій та виконують однакове перетворення даних. %good



(\,3\,) Чим вища кваліфікація програміста, тим більш складні для розуміння та довгі умовні вирази він використовує. %bad



(\,4\,) Код з довгими функціями має меншу кількість функцій, а тому для нього треба менше тестів. %bad

\newpage

{\begin {center}\textbf{Завдання 6}\end {center}}\par
Вибрати всі істинні твердження (якщо існують).\par
\vspace{2mm}



(\,1\,) Код, що себе сам документує, -- це код, який виводить користувачу документацію на себе. %bad



(\,2\,) За відмови від документації одним з методом збереження знань про систему є постійне пересування персоналу між її частинами. %good



(\,3\,) За застосування адаптивних методологій розробки, так само як і за прогнозних, намагаються якомога швидше визначити архітектуру системи. Відмінність полягає в тім, у який спосіб це виконується.%good



(\,4\,) Використання RUP не заперечує часті випуски продукту, але й не наполягає на них. %good


{\begin {center}\textbf{Завдання 7}\end {center}}\par
Вибрати всі істинні твердження (якщо існують).\par
\vspace{2mm}



(\,1\,) RUP є архітектурно-орієнтованим. %good



(\,2\,) Розробка за гнучкими методологіями не передбачає узгодження умов контракту. %bad



(\,3\,) Гнучкі методології розробки допускають суттєві зміни вимог, проте все одно намагаються якомога швидше зафіксувати архітектуру. %good



(\,4\,) Якщо виникають проблеми з модифікацією коду, то слід розглянути внесення змін у проект. %good


{\begin {center}\textbf{Завдання 8}\end {center}}\par
Що з наступного не має відношення до Scrum або не виконується/не використовується за використання оригінального Scrum?\par
\vspace{2mm}



(\,1\,) скрам-мастер визначає, які задачі команда буде виконувати на наступному спрінті%good



(\,2\,) ретроспектива спрінту%bad



(\,3\,) канбан%bad



(\,4\,) під час спрінту дозволяється внесення змін до вимог, що вже реалізуються на поточному спрінті%good


{\begin {center}\textbf{Завдання 9}\end {center}}\par
Що з наведеного нижче не розглядалось та навіть не згадувалось у курсі?\par
\vspace{2mm}



(\,1\,) LAZY%good



(\,2\,) водоспадна модель%bad



(\,3\,) SAFe%bad



(\,4\,) Scrum of scrums%bad


{\begin {center}\textbf{Завдання 10}\end {center}}\par
Вибрати всі істинні твердження (якщо існують).\par
\vspace{2mm}



(\,1\,) Тестовий приклад складається з набору вхідних даних, умов виконання та очікуваних результатів.%good



(\,2\,) Якість буває внутрішньою.%good



(\,3\,) Модель якості продукту є ієрархічною.%good



(\,4\,) Модель якості даних є ієрархічною.%good



\end {large}}
{\vspace{5mm}\smallskip \begin {small} Затверджено на засіданні кафедри теоретичної кібернетики, протокол \No 6 від   29.01.2021 р.\par\smallskip\smallskip
{\parbox {6.5 cm}{Зав. кафедри \dotfill Ю.В.~Крак}}\hspace {0.7cm}{\hbox to 6.5 cm{ Екзаменатор \dotfill Т.О.~Карнаух}} \end{small}}
\newpage

{{\begin {small}\par
{ \center  Київський національний університет імені Тараса Шевченка \par}
{\parbox {5 cm}{Освітня програма \hfill}{\parbox {7 cm}{Інформатика \hfill}}\parbox {6 cm}{\hfill Семестр 6}}\par
{\parbox {5 cm} {Навчальний предмет \hfill}\parbox {7 cm}{Розробка програмного забезпечення \hfill}\parbox {6cm}{\hfill}\par}\vspace{-7pt} \end{small}}{\center Екзаменаційний білет \No \textbf
{ 24 }\par}}
{
\begin {large}
{\begin {center}\textbf{Завдання 1}\end {center}}\par
Так звані трикутники балансування зв'язують між собою:\par
\vspace{2mm}



(\,1\,) Завжди рівно 4 параметри%bad



(\,2\,) Завжди рівно 3 параметри%bad



(\,3\,) 4 або 3 параметри, при цьому може бути відсутня якість%bad



(\,4\,) 4 або 3 параметри, при цьому може бути відсутній обсяг %good


{\begin {center}\textbf{Завдання 2}\end {center}}\par
Який з наведених фрагментів коду явно потребує поліпшення (форматування рядків та типи до уваги не брати)?\par
\vspace{2mm}



(\,1\,) void load\_data(…);%good



(\,2\,) \{int x; int y; …\}%good



(\,3\,) \{int input\_file, file\_to\_output; …\}%bad



(\,4\,) \{fstream inputFile, outputFile;…\} %good


{\begin {center}\textbf{Завдання 3}\end {center}}\par
Що є характерним для структурного програмування?\par
\vspace{2mm}



(\,1\,) У коді мають використовуватись структурні типи даних. %bad



(\,2\,) Використання оператора безумовного переходу. %bad



(\,3\,) У коді мають використовуватись типи структур. %bad



(\,4\,) Програма зображується ієрархічною структурою блоків.%good


{\begin {center}\textbf{Завдання 4}\end {center}}\par
Вибрати всі істинні твердження (якщо існують).\par
\vspace{2mm}



(\,1\,) Неперервне проектування здійснює створення та модифікацію проекту системи в міру розробки. %good



(\,2\,) Неперервність проектування полягає в тім, що процес розробки не можна переривати більше ніж на вихідні. %bad



(\,3\,) Адаптивність методології проявляється в тому, що за зміни вимог розробка починається заново. %bad



(\,4\,) І прогнозні, і адаптивні методології передбачають планування виконуваних робіт. %good


{\begin {center}\textbf{Завдання 5}\end {center}}\par
Вибрати всі істинні твердження (якщо існують).\par
\vspace{2mm}



(\,1\,) Одним з прийомів спрощення складних умов є їх декомпозиція. %good



(\,2\,) Довжина функцій не впливає на внутрішню якість коду. %bad



(\,3\,) Повторюваний код дозволяє програмісту краще зрозуміти його логіку та вибрати з кількох варіантів найкращий. %bad



(\,4\,) Складні умови виникають в реальних задачах доволі часто, а тому є доволі природніми і не можуть вважатись недоліком коду. %bad

\newpage

{\begin {center}\textbf{Завдання 6}\end {center}}\par
Вибрати всі істинні твердження (якщо існують).\par
\vspace{2mm}



(\,1\,) RUP відрізняється від Scrum тим, що ні в що не ставить персонал.%bad



(\,2\,) Використання RUP не заперечує часті випуски продукту, але й не наполягає на них. %good



(\,3\,) На відміну від Scrum, RUP є компонентно-орієнтованим. %bad



(\,4\,) Адаптивні методології розробки дозволяють уникнути надлишкового проектування. %good


{\begin {center}\textbf{Завдання 7}\end {center}}\par
Вибрати всі істинні твердження (якщо існують).\par
\vspace{2mm}



(\,1\,) За програмування в парах (pair programming) поки один колега набирає код на комп'ютері інший може піти пограти в настільний теніс. %bad



(\,2\,) Методологія розробки Big Design Up Front погана тим, що може призводити до надлишкового проектування. %good



(\,3\,) Одна з принципових відмінностей RUP від Scrum полягає у використанні численних проміжних технічних артефактів. %good



(\,4\,) RUP є ітеративним та інкрементним.%good


{\begin {center}\textbf{Завдання 8}\end {center}}\par
Що з наступного не має відношення до Scrum або не виконується/не використовується за використання оригінального Scrum?\par
\vspace{2mm}



(\,1\,) приведення у належний вигляд журналу невиконаних задач (backlog grooming) %bad



(\,2\,) правила визначення готовності (Definition of Done) %bad



(\,3\,) порядок виконання запланованої на спрінт роботи визначає скрам-мастер%good



(\,4\,) скрам-мастер визначає час, необхідний команді для реалізації конкретної задачі%good


{\begin {center}\textbf{Завдання 9}\end {center}}\par
Що з наведеного нижче не розглядалось та навіть не згадувалось у курсі?\par
\vspace{2mm}



(\,1\,) Cleanroom%bad



(\,2\,) LOAN%good



(\,3\,) Lean%bad



(\,4\,) BDD%bad


{\begin {center}\textbf{Завдання 10}\end {center}}\par
Вибрати всі істинні твердження (якщо існують).\par
\vspace{2mm}



(\,1\,) Розглядають модель якості даних.%good



(\,2\,) Якість буває зовнішньою.%good



(\,3\,) Атрибути якості можна використовувати при уточненні вимог.%good



(\,4\,) Відсутність знайдених дефектів на початкових стадіях розробки може бути свідченням поганого покриття коду тестами.%good



\end {large}}
{\vspace{5mm}\smallskip \begin {small} Затверджено на засіданні кафедри теоретичної кібернетики, протокол \No 6 від   29.01.2021 р.\par\smallskip\smallskip
{\parbox {6.5 cm}{Зав. кафедри \dotfill Ю.В.~Крак}}\hspace {0.7cm}{\hbox to 6.5 cm{ Екзаменатор \dotfill Т.О.~Карнаух}} \end{small}}
\newpage

{{\begin {small}\par
{ \center  Київський національний університет імені Тараса Шевченка \par}
{\parbox {5 cm}{Освітня програма \hfill}{\parbox {7 cm}{Інформатика \hfill}}\parbox {6 cm}{\hfill Семестр 6}}\par
{\parbox {5 cm} {Навчальний предмет \hfill}\parbox {7 cm}{Розробка програмного забезпечення \hfill}\parbox {6cm}{\hfill}\par}\vspace{-7pt} \end{small}}{\center Екзаменаційний білет \No \textbf
{ 25 }\par}}
{
\begin {large}
{\begin {center}\textbf{Завдання 1}\end {center}}\par
Так звані трикутники балансування зв'язують між собою:\par
\vspace{2mm}



(\,1\,) 4 або 3 параметри, при цьому може бути відсутній обсяг %good



(\,2\,) Завжди рівно 4 параметри%bad



(\,3\,) 4 або 3 параметри, при цьому може бути відсутній час%bad



(\,4\,) Завжди рівно 3 параметри%bad


{\begin {center}\textbf{Завдання 2}\end {center}}\par
Який з наведених фрагментів коду явно потребує поліпшення (форматування рядків та типи до уваги не брати)?\par
\vspace{2mm}



(\,1\,) \{int a\_, a\_\_; …\}%bad



(\,2\,) void load\_data(…);%good



(\,3\,) \{int input\_file, file\_to\_output; …\}%bad



(\,4\,) \{double avg\_temperature;…\} %good


{\begin {center}\textbf{Завдання 3}\end {center}}\par
Що є характерним для структурного програмування?\par
\vspace{2mm}



(\,1\,) Використання механізму винятків. %bad



(\,2\,) Заборона використання скалярних (простих) типів даних. %bad



(\,3\,) У коді мають використовуватись структурні типи даних. %bad



(\,4\,) Кожен блок коду повинен мати єдиний вхід та єдиний вихід.%good


{\begin {center}\textbf{Завдання 4}\end {center}}\par
Вибрати всі істинні твердження (якщо існують).\par
\vspace{2mm}



(\,1\,) В адаптивних методологіях планування робіт не виконується. %bad



(\,2\,) Для адаптивних методологій прогнозувати результат чи терміни його отримання абсолютно неможливо. %bad



(\,3\,) Прогнозна модель розробки може одночасно бути інкрементною. %good



(\,4\,) Неперервне проектування має бути ітеративним та інкрементним одночасно. %good


{\begin {center}\textbf{Завдання 5}\end {center}}\par
Вибрати всі істинні твердження (якщо існують).\par
\vspace{2mm}



(\,1\,) Одним з засобів позбутись повторень у коді є використання шаблону Template Method (але цей шаблон здатен боротися і з іншими недоліками коду). %good



(\,2\,) Код з довгими функціями має меншу кількість функцій, а тому його швидше та простіше тестувати. %bad



(\,3\,) Втілення в коді різних підходів до одного того самого перетворення даних дозволяє програмісту продемонструвати свою професійну підготовку з найкращого боку. %bad



(\,4\,) Від довгих функцій можна позбутися, розклавши їх у послідовність викликів дрібніших функцій. %good

\newpage

{\begin {center}\textbf{Завдання 6}\end {center}}\par
Вибрати всі істинні твердження (якщо існують).\par
\vspace{2mm}



(\,1\,) Якщо виникають проблеми з модифікацією коду, то слід відразу змінити команду розробників на більш кваліфіковану.%bad



(\,2\,) Без рефакторінгу неперервне проектування неможливе. %good



(\,3\,) Висока внутрішня якість коду здешевшує розробку, бо дозволяє економити на вартості та часі майбутніх змін. %good



(\,4\,) RUP є керованим прецедентами. %good


{\begin {center}\textbf{Завдання 7}\end {center}}\par
Вибрати всі істинні твердження (якщо існують).\par
\vspace{2mm}



(\,1\,) Перша версія коду не завжди може бути якісною -- і це нормально. %good



(\,2\,) Тільки RUP використовує UML.%bad



(\,3\,) Гнучкі методи розробки нехтують проміжними технічними артефактами та документацією виключно через їх непотрібність. %bad



(\,4\,) Scrum передбачає колективне володіння кодом. %good


{\begin {center}\textbf{Завдання 8}\end {center}}\par
Що з наступного не має відношення до Scrum або не виконується/не використовується за використання оригінального Scrum?\par
\vspace{2mm}



(\,1\,) планування спрінту%bad



(\,2\,) побудова різнопланових моделей системи%good



(\,3\,) скрам-мастер визначає пріоритети задач%good



(\,4\,) порядок виконання запланованої на спрінт роботи визначає власник продукту%good


{\begin {center}\textbf{Завдання 9}\end {center}}\par
Що з наведеного нижче не розглядалось та навіть не згадувалось у курсі?\par
\vspace{2mm}



(\,1\,) PRINCE2%good



(\,2\,) TDD%bad



(\,3\,) FDD%bad



(\,4\,) sashimi%bad


{\begin {center}\textbf{Завдання 10}\end {center}}\par
Вибрати всі істинні твердження (якщо існують).\par
\vspace{2mm}



(\,1\,) Коли нема можливості перевірити всі можливі комбінації, тестують їх попарні (pairwise) комбінації. Це може значно зменшити кількість тестів.%good



(\,2\,) Користувачі бувають основними.%good



(\,3\,) За провалених димових тестів решта тестів не виконується і код повертається на доробку.%good



(\,4\,) Модель якості продукту є ієрархічною.%good



\end {large}}
{\vspace{5mm}\smallskip \begin {small} Затверджено на засіданні кафедри теоретичної кібернетики, протокол \No 6 від   29.01.2021 р.\par\smallskip\smallskip
{\parbox {6.5 cm}{Зав. кафедри \dotfill Ю.В.~Крак}}\hspace {0.7cm}{\hbox to 6.5 cm{ Екзаменатор \dotfill Т.О.~Карнаух}} \end{small}}
\newpage

{{\begin {small}\par
{ \center  Київський національний університет імені Тараса Шевченка \par}
{\parbox {5 cm}{Освітня програма \hfill}{\parbox {7 cm}{Інформатика \hfill}}\parbox {6 cm}{\hfill Семестр 6}}\par
{\parbox {5 cm} {Навчальний предмет \hfill}\parbox {7 cm}{Розробка програмного забезпечення \hfill}\parbox {6cm}{\hfill}\par}\vspace{-7pt} \end{small}}{\center Екзаменаційний білет \No \textbf
{ 26 }\par}}
{
\begin {large}
{\begin {center}\textbf{Завдання 1}\end {center}}\par
Так звані трикутники балансування зв'язують між собою:\par
\vspace{2mm}



(\,1\,) Завжди рівно 4 параметри%bad



(\,2\,) 4 або 3 параметри, при цьому може бути відсутній обсяг %good



(\,3\,) 4 або 3 параметри, при цьому може бути відсутня якість%bad



(\,4\,) 4 або 3 параметри, при цьому може бути відсутній час%bad


{\begin {center}\textbf{Завдання 2}\end {center}}\par
Який з наведених фрагментів коду явно потребує поліпшення (форматування рядків та типи до уваги не брати)?\par
\vspace{2mm}



(\,1\,) \{int x; int y; …\}%good



(\,2\,) void help();%good



(\,3\,) \{int xtrmprg; …\}%bad



(\,4\,) \{fstream input\_file, outputFile; …\}%bad


{\begin {center}\textbf{Завдання 3}\end {center}}\par
Що є характерним для структурного програмування?\par
\vspace{2mm}



(\,1\,) Програма зображується ієрархічною структурою блоків.%good



(\,2\,) У коді мають використовуватись типи даних, що є класами. %bad



(\,3\,) У коді мають використовуватись типи структур, класи використовувати заборонено. %bad



(\,4\,) У коді мають використовуватись типи даних, що є класами. %bad


{\begin {center}\textbf{Завдання 4}\end {center}}\par
Вибрати всі істинні твердження (якщо існують).\par
\vspace{2mm}



(\,1\,) Неперервність проектування полягає в тім, що процес розробки не можна переривати більше ніж на вихідні. %bad



(\,2\,) У прогнозній моделі розробки розробка може виконуватись ітеративно. %good



(\,3\,) Неперервне проектування здійснює створення та модифікацію проекту системи в міру розробки. %good



(\,4\,) Адаптивність методології проявляється в тому, що за зміни вимог розробка починається заново. %bad


{\begin {center}\textbf{Завдання 5}\end {center}}\par
Вибрати всі істинні твердження (якщо існують).\par
\vspace{2mm}



(\,1\,) Складні комбінації умов починають керувати логікою виконання, тому що такими є вимоги до функціонування ПЗ. Оскільки ПЗ має якнайкраще забезпечувати потреби замовника, то складні комбінації умов у коді вважатись недоліком коду не можуть. %bad



(\,2\,) Довгі функції доволі часто приховують складну логіку обробки або повторюваний код. %good



(\,3\,) Одним з засобів позбутись повторень у коді є використання успадкування. %good



(\,4\,) Правильно спроектована функція повинна мати рівно один обов'язок. %good

\newpage

{\begin {center}\textbf{Завдання 6}\end {center}}\par
Вибрати всі істинні твердження (якщо існують).\par
\vspace{2mm}



(\,1\,) Розробка за гнучкими методологіями ніколи не слідує плану. %bad



(\,2\,) Гнучкі методології повністю нехтують документацією. %bad



(\,3\,) Гнучкі методології розробки намагаються якомога швидше визначити архітектуру системи, для чого на початку завжди виконується початкове проектування системи в цілому, а тільки потім переходять до ітеративного уточнення деталей та кодування. %bad



(\,4\,) Розробка за гнучкими методологіями цінить людей та взаємодію між ними. %good


{\begin {center}\textbf{Завдання 7}\end {center}}\par
Вибрати всі істинні твердження (якщо існують).\par
\vspace{2mm}



(\,1\,) Гнучкі методології розробки допускають суттєві зміни вимог, проте все одно намагаються якомога швидше зафіксувати архітектуру. %good



(\,2\,) Гнучкі методології можуть передбачати як спільне, так і індивідуальне володіння кодом. %good



(\,3\,) RUP принципово відрізняється від Scrum тим, що є керованим прецедентами. У той же час у Scrum при плануванні розробки прецеденти (чи їх аналоги) до уваги не беруться. %bad



(\,4\,) RUP є прогнозною методологією. %bad


{\begin {center}\textbf{Завдання 8}\end {center}}\par
Що з наступного не має відношення до Scrum або не виконується/не використовується за використання оригінального Scrum?\par
\vspace{2mm}



(\,1\,) щоденний скрам%bad



(\,2\,) розповіді користувача%bad



(\,3\,) розповіді користувача%bad



(\,4\,) власник продукту призначає відповідальних за конкретні задачі%good


{\begin {center}\textbf{Завдання 9}\end {center}}\par
Що з наведеного нижче не розглядалось та навіть не згадувалось у курсі?\par
\vspace{2mm}



(\,1\,) DSDM%bad



(\,2\,) RAD%bad



(\,3\,) PRINCE2%good



(\,4\,) CRAZY%good


{\begin {center}\textbf{Завдання 10}\end {center}}\par
Вибрати всі істинні твердження (якщо існують).\par
\vspace{2mm}



(\,1\,) І верифікація, і валідація можуть використовувати тестування.%good



(\,2\,) Розглядають модель якості при використанні.%good



(\,3\,) Розглядають модель якості даних.%good



(\,4\,) Якість буває внутрішньою.%good



\end {large}}
{\vspace{5mm}\smallskip \begin {small} Затверджено на засіданні кафедри теоретичної кібернетики, протокол \No 6 від   29.01.2021 р.\par\smallskip\smallskip
{\parbox {6.5 cm}{Зав. кафедри \dotfill Ю.В.~Крак}}\hspace {0.7cm}{\hbox to 6.5 cm{ Екзаменатор \dotfill Т.О.~Карнаух}} \end{small}}
\newpage

{{\begin {small}\par
{ \center  Київський національний університет імені Тараса Шевченка \par}
{\parbox {5 cm}{Освітня програма \hfill}{\parbox {7 cm}{Інформатика \hfill}}\parbox {6 cm}{\hfill Семестр 6}}\par
{\parbox {5 cm} {Навчальний предмет \hfill}\parbox {7 cm}{Розробка програмного забезпечення \hfill}\parbox {6cm}{\hfill}\par}\vspace{-7pt} \end{small}}{\center Екзаменаційний білет \No \textbf
{ 27 }\par}}
{
\begin {large}
{\begin {center}\textbf{Завдання 1}\end {center}}\par
Що з наведеного найкраще характеризує рефакторінг?\par
\vspace{2mm}



(\,1\,) Зміна внутрішньої структури коду під час усунення помилок%bad



(\,2\,) Зміна внутрішньої структури коду, що не змінює його видиму поведінку%good



(\,3\,) Довільна переробка коду, що виникає у процесі розробки%bad



(\,4\,) Докорінна переробка всього коду з нуля,  що не змінює його видиму поведінку%bad


{\begin {center}\textbf{Завдання 2}\end {center}}\par
Який з наведених фрагментів коду явно потребує поліпшення (форматування рядків та типи до уваги не брати)?\par
\vspace{2mm}



(\,1\,) \{fstream inputFile, outputFile;…\} %good



(\,2\,) void f13(…);%bad



(\,3\,) \{double the\_best\_value\_anyone\_ever\_have;\} %bad



(\,4\,) void ff(…);%bad


{\begin {center}\textbf{Завдання 3}\end {center}}\par
Що є характерним для структурного програмування?\par
\vspace{2mm}



(\,1\,) У коді мають використовуватись типи структур. %bad



(\,2\,) Переходи з середини одного блоку в середину іншого неможливі. %good



(\,3\,) У коді мають використовуватись структурні типи даних. %bad



(\,4\,) У коді мають використовуватись типи структур, класи використовувати заборонено. %bad


{\begin {center}\textbf{Завдання 4}\end {center}}\par
Вибрати всі істинні твердження (якщо існують).\par
\vspace{2mm}



(\,1\,) І прогнозні, і адаптивні методології передбачають планування виконуваних робіт. %good



(\,2\,) У випадку, коли вимоги до продукту відомі заздалегідь та є фіксованими контрактом, адаптивні методології розробки використовувати не можна. %bad



(\,3\,) Адаптивні методології має сенс використовувати для проектів, де можлива зміна вимог. %good



(\,4\,) Для адаптивних методологій розробки, на відміну від прогнозних, попередні оцінки бюджету та/або часу не виконують, бо вони все одно зміняться. %bad


{\begin {center}\textbf{Завдання 5}\end {center}}\par
Вибрати всі істинні твердження (якщо існують).\par
\vspace{2mm}



(\,1\,) Зв'язки між окремими модулями системи мають бути якомога більш слабкими. %good



(\,2\,) Одним з прийомів спрощення складної умовної логіки є використання шаблону State. %good



(\,3\,) Залежності між різними модулями системи мають бути якомога більш сильними, бо інакше вона розпадеться. %bad



(\,4\,) Гарний код має читатись як розмовна мова. %good

\newpage

{\begin {center}\textbf{Завдання 6}\end {center}}\par
Вибрати всі істинні твердження (якщо існують).\par
\vspace{2mm}



(\,1\,) Використання гнучких методологій при розробці надвеликих проектів не несе в собі жодних ризиків. %bad



(\,2\,) Гнучкі методології розробки не допускають використання UML. %bad



(\,3\,) Розробка за гнучкими методологіями цінить співробітництво з замовником та піддатлівість змінам. %good



(\,4\,) Розробка за гнучкими методологіями не передбачає узгодження умов контракту. %bad


{\begin {center}\textbf{Завдання 7}\end {center}}\par
Вибрати всі істинні твердження (якщо існують).\par
\vspace{2mm}



(\,1\,) Гнучкі методології базуються на комунікації розробників, а тому передбачають виключно колективне володіння кодом. %bad



(\,2\,) Код, що себе сам документує, -- це код, який виводить користувачу документацію на себе. %bad



(\,3\,) Водоспадна методологія не призначена для внесення змін у вимоги до системи. %good



(\,4\,) Розповіді користувача (user stories) це менш формалізований порівняно з прецедентами спосіб опису вимог до системи. %good


{\begin {center}\textbf{Завдання 8}\end {center}}\par
Що з наступного не має відношення до Scrum або не виконується/не використовується за використання оригінального Scrum?\par
\vspace{2mm}



(\,1\,) побудова різнопланових моделей системи%good



(\,2\,) під час спрінту дозволяється внесення змін до вимог, що вже реалізуються на поточному спрінті%good



(\,3\,) ретроспектива спрінту%bad



(\,4\,) канбан%bad


{\begin {center}\textbf{Завдання 9}\end {center}}\par
Що з наведеного нижче не розглядалось та навіть не згадувалось у курсі?\par
\vspace{2mm}



(\,1\,) модель хаосу%bad



(\,2\,) BDUF%bad



(\,3\,) LAZY%good



(\,4\,) LOAN%good


{\begin {center}\textbf{Завдання 10}\end {center}}\par
Вибрати всі істинні твердження (якщо існують).\par
\vspace{2mm}



(\,1\,) Для тестування код треба ізолювати. Для того, щоб ізолювати, треба мати тести.%good



(\,2\,) Можна розглядати якість під час використання.%good



(\,3\,) Відсутність знайдених дефектів на початкових стадіях розробки може бути свідченням поганого покриття коду тестами.%good



(\,4\,) Модель якості при використанні є ієрархічною.%good



\end {large}}
{\vspace{5mm}\smallskip \begin {small} Затверджено на засіданні кафедри теоретичної кібернетики, протокол \No 6 від   29.01.2021 р.\par\smallskip\smallskip
{\parbox {6.5 cm}{Зав. кафедри \dotfill Ю.В.~Крак}}\hspace {0.7cm}{\hbox to 6.5 cm{ Екзаменатор \dotfill Т.О.~Карнаух}} \end{small}}
\newpage

{{\begin {small}\par
{ \center  Київський національний університет імені Тараса Шевченка \par}
{\parbox {5 cm}{Освітня програма \hfill}{\parbox {7 cm}{Інформатика \hfill}}\parbox {6 cm}{\hfill Семестр 6}}\par
{\parbox {5 cm} {Навчальний предмет \hfill}\parbox {7 cm}{Розробка програмного забезпечення \hfill}\parbox {6cm}{\hfill}\par}\vspace{-7pt} \end{small}}{\center Екзаменаційний білет \No \textbf
{ 28 }\par}}
{
\begin {large}
{\begin {center}\textbf{Завдання 1}\end {center}}\par
Шаблон проектування:\par
\vspace{2mm}



(\,1\,) Є шаблоном, в який треба підставити власні ідентифікатори%bad



(\,2\,) Тільки формулює ідею розв'язання проблеми%bad



(\,3\,) Задає розв'язання типової проблеми до найдрібніших подробиць%bad



(\,4\,) Формулює ідею розв'язання проблеми та може пропонувати одну чи кілька її реалізацій%good


{\begin {center}\textbf{Завдання 2}\end {center}}\par
Який з наведених фрагментів коду явно потребує поліпшення (форматування рядків та типи до уваги не брати)?\par
\vspace{2mm}



(\,1\,) void load\_data(T \& arg1, double arg2, double arg3, double arg4); %bad



(\,2\,) \{double maxPower, minPower;…\} %good



(\,3\,) \{fstream input\_file, output\_file;…\} %good



(\,4\,) \{double max\_power, min\_power;…\} %good


{\begin {center}\textbf{Завдання 3}\end {center}}\par
Що є характерним для структурного програмування?\par
\vspace{2mm}



(\,1\,) Заборона використання скалярних (простих) типів даних. %bad



(\,2\,) Використання оператора безумовного переходу. %bad



(\,3\,) Використання механізму винятків. %bad



(\,4\,) Програма зображується ієрархічною структурою блоків.%good


{\begin {center}\textbf{Завдання 4}\end {center}}\par
Вибрати всі істинні твердження (якщо існують).\par
\vspace{2mm}



(\,1\,) Адаптивна модель розробки адаптує вимоги до продукту до наявної команди розробників. %bad



(\,2\,) Ітеративні та інкрементні методи розробки є синонімом неперервного проектування. %bad



(\,3\,) Адаптивні моделі розробки допускають зміну вимог до продукту. %good



(\,4\,) Для адаптивних моделей розробки оцінки часу та бюджету не виконують, бо вимоги є змінними. %bad


{\begin {center}\textbf{Завдання 5}\end {center}}\par
Вибрати всі істинні твердження (якщо існують).\par
\vspace{2mm}



(\,1\,) Чим менше модулів у проекті, тим краще, бо меншу кількість модулів треба інтегрувати в єдине ціле. %bad



(\,2\,) Якщо кожна функція бібліотеки буде здатна в залежності від параметрів виконувати кілька різних дій, то загальна кількість функцій у ній буде менша і її буде простіше використовувати. %bad



(\,3\,) Одним з засобів позбутись повторень у коді є використання функцій. %good



(\,4\,) Одним з прийомів спрощення умовної логіки є використання поліморфізму. %good

\newpage

{\begin {center}\textbf{Завдання 6}\end {center}}\par
Вибрати всі істинні твердження (якщо існують).\par
\vspace{2mm}



(\,1\,) RUP є архітектурно-орієнтованим. %good



(\,2\,) Основна відмінність V-моделі від водоспадної полягає в тім, що діяльності з тестування виконуються паралельно з усіма іншими діяльностями з розробки. %good



(\,3\,) Гнучкі методології передбачають часті випуски продукту. %good



(\,4\,) Водоспадна методологія розробки може призводити до надлишкового проектування. %good


{\begin {center}\textbf{Завдання 7}\end {center}}\par
Вибрати всі істинні твердження (якщо існують).\par
\vspace{2mm}



(\,1\,) Програмування в парах (pair programming) підвищує якість коду та дозволяє економити на виправленні помилок. %good



(\,2\,) Ітеративна та інкрементна розробка передбачає наявність у тому чи іншому вигляді етапу ініціалізації, ітерацій та керуючої таблиці проекту. %good



(\,3\,) Використання гнучких методологій вимагає високої кваліфікації персоналу. %good



(\,4\,) Багато хто з тих, що використовує Scrum, складає архітектурну модель розроблюваної системи. %good


{\begin {center}\textbf{Завдання 8}\end {center}}\par
Що з наступного не має відношення до Scrum або не виконується/не використовується за використання оригінального Scrum?\par
\vspace{2mm}



(\,1\,) власник продукту призначає час, необхідний команді для реалізації конкретної задачі%good



(\,2\,) порядок виконання запланованої на спрінт роботи визначають розробники%bad



(\,3\,) виділення одного з учасників команди як системного архітектора%good



(\,4\,) власник продукту остаточно затверджує проект системи%good


{\begin {center}\textbf{Завдання 9}\end {center}}\par
Що з наведеного нижче не розглядалось та навіть не згадувалось у курсі?\par
\vspace{2mm}



(\,1\,) SAFe%bad



(\,2\,) водоспадна модель%bad



(\,3\,) sashimi%bad



(\,4\,) Scrum%bad


{\begin {center}\textbf{Завдання 10}\end {center}}\par
Вибрати всі істинні твердження (якщо існують).\par
\vspace{2mm}



(\,1\,) Тестування вимог може породжувати їх уточнення.%good



(\,2\,) Модель якості даних є ієрархічною.%good



(\,3\,) Якість можна визначати через атрибути якості.%good



(\,4\,) Користувачі бувають непрямими.%good



\end {large}}
{\vspace{5mm}\smallskip \begin {small} Затверджено на засіданні кафедри теоретичної кібернетики, протокол \No 6 від   29.01.2021 р.\par\smallskip\smallskip
{\parbox {6.5 cm}{Зав. кафедри \dotfill Ю.В.~Крак}}\hspace {0.7cm}{\hbox to 6.5 cm{ Екзаменатор \dotfill Т.О.~Карнаух}} \end{small}}
\newpage

{{\begin {small}\par
{ \center  Київський національний університет імені Тараса Шевченка \par}
{\parbox {5 cm}{Освітня програма \hfill}{\parbox {7 cm}{Інформатика \hfill}}\parbox {6 cm}{\hfill Семестр 6}}\par
{\parbox {5 cm} {Навчальний предмет \hfill}\parbox {7 cm}{Розробка програмного забезпечення \hfill}\parbox {6cm}{\hfill}\par}\vspace{-7pt} \end{small}}{\center Екзаменаційний білет \No \textbf
{ 29 }\par}}
{
\begin {large}
{\begin {center}\textbf{Завдання 1}\end {center}}\par
Так звані трикутники балансування зв'язують між собою:\par
\vspace{2mm}



(\,1\,) 4 або 3 параметри, при цьому може бути відсутній обсяг %good



(\,2\,) Завжди рівно 4 параметри%bad



(\,3\,) Завжди рівно 3 параметри%bad



(\,4\,) 4 або 3 параметри, при цьому може бути відсутній час%bad


{\begin {center}\textbf{Завдання 2}\end {center}}\par
Який з наведених фрагментів коду явно потребує поліпшення (форматування рядків та типи до уваги не брати)?\par
\vspace{2mm}



(\,1\,) \{double max\_power, min\_power;…\} %good



(\,2\,) void load\_data(T \& arg1, double arg2, double arg3, double arg4); %bad



(\,3\,) \{int input\_file, file\_to\_output; …\}%bad



(\,4\,) \{int xtrmprg; …\}%bad


{\begin {center}\textbf{Завдання 3}\end {center}}\par
Що є характерним для структурного програмування?\par
\vspace{2mm}



(\,1\,) У коді мають використовуватись структурні типи даних. %bad



(\,2\,) У коді мають використовуватись типи структур, класи використовувати заборонено. %bad



(\,3\,) Переходи з середини одного блоку в середину іншого неможливі. %good



(\,4\,) Використання механізму винятків. %bad


{\begin {center}\textbf{Завдання 4}\end {center}}\par
Вибрати всі істинні твердження (якщо існують).\par
\vspace{2mm}



(\,1\,) Прогнозна модель розробки не може одночасно бути інкрементною. %bad



(\,2\,) Неперервне проектування створює специфікацію проекту перед початком розробки/ітерації, а під час ітерації його неперервно модифікує.%bad



(\,3\,) Адаптивні методології не передбачають визначення вимог. %bad



(\,4\,) Гнучкі методології не передбачають планування робіт. %bad


{\begin {center}\textbf{Завдання 5}\end {center}}\par
Вибрати всі істинні твердження (якщо існують).\par
\vspace{2mm}



(\,1\,) Код, що виконує однакову послідовність логічних кроків, у більшості випадків є повторюваним.%good



(\,2\,) Одним з прийомів спрощення складної умовної логіки є використання шаблону Strategy. %good



(\,3\,) Повторюваний код має повторюватись у точності до символу, інакше він повторюваним вважатись не може. %bad



(\,4\,) Залежності всередині модуля можуть бути сильними. %good

\newpage

{\begin {center}\textbf{Завдання 6}\end {center}}\par
Вибрати всі істинні твердження (якщо існують).\par
\vspace{2mm}



(\,1\,) Ітеративна методологія розробки у загальному випадку не гарантує, що не буде надлишкового проектування. %good



(\,2\,) Програмування в парах (pair programming) зводиться до того, що кожним фрагментом коду володіють два програмісти. %bad



(\,3\,) За використання гнучких методологій розробки припиняти розробку не рекомендується. %good



(\,4\,) RUP є компонентно-орієнтованим. %good


{\begin {center}\textbf{Завдання 7}\end {center}}\par
Вибрати всі істинні твердження (якщо існують).\par
\vspace{2mm}



(\,1\,) За використання гнучких методологій кожен розробник робить, що хоче. %bad



(\,2\,) Водоспадна методологія може мати тенденцію до надлишкового проектування. %good



(\,3\,) За відмови від документації одним з методом збереження знань про систему є постійне пересування персоналу між її частинами. %good



(\,4\,) Сьогодні водоспадна методологія не має права на існування. %bad


{\begin {center}\textbf{Завдання 8}\end {center}}\par
Що з наступного не має відношення до Scrum або не виконується/не використовується за використання оригінального Scrum?\par
\vspace{2mm}



(\,1\,) скрам-мастер остаточно затверджує проект системи%good



(\,2\,) команда визначає, які задачі вона буде виконувати на наступному спрінті%good



(\,3\,) скрам-мастер призначає відповідальних за конкретні задачі%good



(\,4\,) покер планування (Planning Poker)%bad


{\begin {center}\textbf{Завдання 9}\end {center}}\par
Що з наведеного нижче не розглядалось та навіть не згадувалось у курсі?\par
\vspace{2mm}



(\,1\,) Lean%bad



(\,2\,) poker%bad



(\,3\,) спіральна модель%bad



(\,4\,) RUP%bad


{\begin {center}\textbf{Завдання 10}\end {center}}\par
Вибрати всі істинні твердження (якщо існують).\par
\vspace{2mm}



(\,1\,) Тестовий приклад складається з набору вхідних даних, умов виконання та очікуваних результатів.%good



(\,2\,) Якість буває зовнішньою.%good



(\,3\,) Атрибути якості можна використовувати при уточненні вимог.%good



(\,4\,) Забезпечення якості це не тільки тестування.%good



\end {large}}
{\vspace{5mm}\smallskip \begin {small} Затверджено на засіданні кафедри теоретичної кібернетики, протокол \No 6 від   29.01.2021 р.\par\smallskip\smallskip
{\parbox {6.5 cm}{Зав. кафедри \dotfill Ю.В.~Крак}}\hspace {0.7cm}{\hbox to 6.5 cm{ Екзаменатор \dotfill Т.О.~Карнаух}} \end{small}}
\newpage

{{\begin {small}\par
{ \center  Київський національний університет імені Тараса Шевченка \par}
{\parbox {5 cm}{Освітня програма \hfill}{\parbox {7 cm}{Інформатика \hfill}}\parbox {6 cm}{\hfill Семестр 6}}\par
{\parbox {5 cm} {Навчальний предмет \hfill}\parbox {7 cm}{Розробка програмного забезпечення \hfill}\parbox {6cm}{\hfill}\par}\vspace{-7pt} \end{small}}{\center Екзаменаційний білет \No \textbf
{ 30 }\par}}
{
\begin {large}
{\begin {center}\textbf{Завдання 1}\end {center}}\par
Так звані трикутники балансування зв'язують між собою:\par
\vspace{2mm}



(\,1\,) Завжди рівно 4 параметри%bad



(\,2\,) 4 або 3 параметри, при цьому може бути відсутній обсяг %good



(\,3\,) 4 або 3 параметри, при цьому може бути відсутній час%bad



(\,4\,) 4 або 3 параметри, при цьому може бути відсутня якість%bad


{\begin {center}\textbf{Завдання 2}\end {center}}\par
Який з наведених фрагментів коду явно потребує поліпшення (форматування рядків та типи до уваги не брати)?\par
\vspace{2mm}



(\,1\,) void help();%good



(\,2\,) \{double maxPower, minPower;…\} %good



(\,3\,) void f13(…);%bad



(\,4\,) \{int a\_, a\_\_; …\}%bad


{\begin {center}\textbf{Завдання 3}\end {center}}\par
Що є характерним для структурного програмування?\par
\vspace{2mm}



(\,1\,) Кожен блок коду повинен мати єдиний вхід та єдиний вихід.%good



(\,2\,) Використання оператора безумовного переходу. %bad



(\,3\,) У коді мають використовуватись типи структур. %bad



(\,4\,) Заборона використання скалярних (простих) типів даних. %bad


{\begin {center}\textbf{Завдання 4}\end {center}}\par
Вибрати всі істинні твердження (якщо існують).\par
\vspace{2mm}



(\,1\,) У прогнозній моделі розробки розробка не може виконуватись ітеративно. %bad



(\,2\,) Прогнозні та адаптивні методології використовують одні ті самі діяльності з розробки ПЗ: визначення та аналіз вимог, проектування, реалізація, тестування, впровадження, супроводження. %good



(\,3\,) Прогнозні моделі розробки з самого початку фіксують вимоги до продукту. %good



(\,4\,) Адаптивні моделі розробки не допускають зміну вимог до продукту, але продукт може розроблятись поетапно. %bad


{\begin {center}\textbf{Завдання 5}\end {center}}\par
Вибрати всі істинні твердження (якщо існують).\par
\vspace{2mm}



(\,1\,) Складні комбінації умов починають керувати логікою виконання, тому що такими є вимоги до функціонування ПЗ. Оскільки ПЗ має якнайкраще забезпечувати потреби замовника, то складні комбінації умов у коді вважатись недоліком коду не можуть. %bad



(\,2\,) Гарний код має читатись як розмовна мова. %good



(\,3\,) Код, що виконує однакову послідовність логічних кроків, у більшості випадків є повторюваним.%good



(\,4\,) Чим менше класів у коді, тим менше класів треба налагоджувати, а тому класи за можливості слід об'єднувати. %bad

\newpage

{\begin {center}\textbf{Завдання 6}\end {center}}\par
Вибрати всі істинні твердження (якщо існують).\par
\vspace{2mm}



(\,1\,) Програмування в парах (pair programming) спрощує введення в команду нових розробників. %good



(\,2\,) Висока внутрішня якість коду є марною тратою часу на те, що замовник не побачить та/або не зможе зацінити. %bad



(\,3\,) Модель хаосу є методологією розробки програмного забезпечення. %bad



(\,4\,) За застосування адаптивних методологій розробки, так само як і за прогнозних, намагаються якомога швидше визначити архітектуру системи. Відмінність полягає в тім, у який спосіб це виконується.%good


{\begin {center}\textbf{Завдання 7}\end {center}}\par
Вибрати всі істинні твердження (якщо існують).\par
\vspace{2mm}



(\,1\,) За невизначених або піддатливих змінам вимог краще застосовувати гнучкі методології розробки. %good



(\,2\,) Поступове уточнення вимог замовника полегшує адаптацію до змін. %good



(\,3\,) Гнучкі методи розробки нехтують проміжними технічними артефактами та документацією через те, що за зміни вимог та/або проекту треба витрачати додатковий час на підтримання їх у відповідному стані. %good



(\,4\,) Водоспадна методологія розробки може призводити до недостатнього проектування. %bad


{\begin {center}\textbf{Завдання 8}\end {center}}\par
Що з наступного не має відношення до Scrum або не виконується/не використовується за використання оригінального Scrum?\par
\vspace{2mm}



(\,1\,) правила визначення готовності (Definition of Done) %bad



(\,2\,) виділення в команді одного чи кількох тестувальників%good



(\,3\,) планування спрінту%bad



(\,4\,) скрам-мастер визначає пріоритети задач%good


{\begin {center}\textbf{Завдання 9}\end {center}}\par
Що з наведеного нижче не розглядалось та навіть не згадувалось у курсі?\par
\vspace{2mm}



(\,1\,) XP%bad



(\,2\,) V-модель%bad



(\,3\,) X-модель%good



(\,4\,) LeSS%bad


{\begin {center}\textbf{Завдання 10}\end {center}}\par
Вибрати всі істинні твердження (якщо існують).\par
\vspace{2mm}



(\,1\,) Користувачі бувають вторинними.%good



(\,2\,) Якість можна визначати через дефекти.%good



(\,3\,) Тестування буває позитивне та негативне.%good



(\,4\,) Розглядають модель якості продукту.%good



\end {large}}
{\vspace{5mm}\smallskip \begin {small} Затверджено на засіданні кафедри теоретичної кібернетики, протокол \No 6 від   29.01.2021 р.\par\smallskip\smallskip
{\parbox {6.5 cm}{Зав. кафедри \dotfill Ю.В.~Крак}}\hspace {0.7cm}{\hbox to 6.5 cm{ Екзаменатор \dotfill Т.О.~Карнаух}} \end{small}}
\newpage

{{\begin {small}\par
{ \center  Київський національний університет імені Тараса Шевченка \par}
{\parbox {5 cm}{Освітня програма \hfill}{\parbox {7 cm}{Інформатика \hfill}}\parbox {6 cm}{\hfill Семестр 6}}\par
{\parbox {5 cm} {Навчальний предмет \hfill}\parbox {7 cm}{Розробка програмного забезпечення \hfill}\parbox {6cm}{\hfill}\par}\vspace{-7pt} \end{small}}{\center Екзаменаційний білет \No \textbf
{ 31 }\par}}
{
\begin {large}
{\begin {center}\textbf{Завдання 1}\end {center}}\par
Так звані трикутники балансування зв'язують між собою:\par
\vspace{2mm}



(\,1\,) 4 або 3 параметри, при цьому може бути відсутній час%bad



(\,2\,) Завжди рівно 3 параметри%bad



(\,3\,) 4 або 3 параметри, при цьому може бути відсутня якість%bad



(\,4\,) 4 або 3 параметри, при цьому може бути відсутній обсяг %good


{\begin {center}\textbf{Завдання 2}\end {center}}\par
Який з наведених фрагментів коду явно потребує поліпшення (форматування рядків та типи до уваги не брати)?\par
\vspace{2mm}



(\,1\,) void load\_data(…);%good



(\,2\,) \{double avg\_temperature;…\} %good



(\,3\,) \{int x; int y; …\}%good



(\,4\,) \{fstream input\_file, output\_file;…\} %good


{\begin {center}\textbf{Завдання 3}\end {center}}\par
Що є характерним для структурного програмування?\par
\vspace{2mm}



(\,1\,) У коді мають використовуватись типи даних, що є класами. %bad



(\,2\,) Заборона використання скалярних (простих) типів даних. %bad



(\,3\,) У коді мають використовуватись типи структур. %bad



(\,4\,) Переходи з середини одного блоку в середину іншого неможливі. %good


{\begin {center}\textbf{Завдання 4}\end {center}}\par
Вибрати всі істинні твердження (якщо існують).\par
\vspace{2mm}



(\,1\,) Для адаптивних методологій розробки, на відміну від прогнозних, попередні оцінки бюджету та/або часу не виконують, бо вони все одно зміняться. %bad



(\,2\,) У випадку, коли вимоги до продукту відомі заздалегідь та є фіксованими контрактом, адаптивні методології розробки використовувати не можна. %bad



(\,3\,) І прогнозні, і адаптивні методології передбачають планування виконуваних робіт. %good



(\,4\,) В адаптивних методологіях планування робіт не виконується. %bad


{\begin {center}\textbf{Завдання 5}\end {center}}\par
Вибрати всі істинні твердження (якщо існують).\par
\vspace{2mm}



(\,1\,) Одним з прийомів спрощення умовної логіки є використання поліморфізму. %good



(\,2\,) Залежності всередині модуля можуть бути сильними. %good



(\,3\,) Залежності між різними модулями системи мають бути якомога більш сильними, бо інакше вона розпадеться. %bad



(\,4\,) Довжина функцій не впливає на внутрішню якість коду. %bad

\newpage

{\begin {center}\textbf{Завдання 6}\end {center}}\par
Вибрати всі істинні твердження (якщо існують).\par
\vspace{2mm}



(\,1\,) Простота – мистецтво максимізації незробленої роботи. %good



(\,2\,) Гнучкі методології розробки є стійкими до припинення розробки на певний час. %bad



(\,3\,) Розповіді користувача (user stories) – це прогресивний інструмент гнучких методологій, який нічого спільного не має з прецедентами (use case). %bad



(\,4\,) Методологія розробки Big Design Up Front погана тим, що може призводити до недостатнього проектування.%bad


{\begin {center}\textbf{Завдання 7}\end {center}}\par
Вибрати всі істинні твердження (якщо існують).\par
\vspace{2mm}



(\,1\,) Перша версія коду не завжди може бути якісною -- і це нормально. %good



(\,2\,) Розповіді користувача (user stories) це менш формалізований порівняно з прецедентами спосіб опису вимог до системи. %good



(\,3\,) Розробка за гнучкими методологіями цінить співробітництво з замовником та піддатлівість змінам. %good



(\,4\,) Гнучкі методології розробки намагаються якомога швидше визначити архітектуру системи, для чого на початку завжди виконується початкове проектування системи в цілому, а тільки потім переходять до ітеративного уточнення деталей та кодування. %bad


{\begin {center}\textbf{Завдання 8}\end {center}}\par
Що з наступного не має відношення до Scrum або не виконується/не використовується за використання оригінального Scrum?\par
\vspace{2mm}



(\,1\,) під час спрінту зміни, що впливають на досягнення мети спрінту, зазвичай, не допускаються%bad



(\,2\,) скрам-мастер визначає, які задачі команда буде виконувати на наступному спрінті%good



(\,3\,) виділення в команді одного чи кількох проектувальників%good



(\,4\,) бачення продукту (vision) %bad


{\begin {center}\textbf{Завдання 9}\end {center}}\par
Що з наведеного нижче не розглядалось та навіть не згадувалось у курсі?\par
\vspace{2mm}



(\,1\,) спіральна модель%bad



(\,2\,) V-модель%bad



(\,3\,) PRINCE2%good



(\,4\,) LeSS%bad


{\begin {center}\textbf{Завдання 10}\end {center}}\par
Вибрати всі істинні твердження (якщо існують).\par
\vspace{2mm}



(\,1\,) Тестовий приклад складається з набору вхідних даних, умов виконання та очікуваних результатів.%good



(\,2\,) Відсутність знайдених дефектів на початкових стадіях розробки може бути свідченням поганого покриття коду тестами.%good



(\,3\,) Користувачі бувають основними.%good



(\,4\,) Модель якості даних є ієрархічною.%good



\end {large}}
{\vspace{5mm}\smallskip \begin {small} Затверджено на засіданні кафедри теоретичної кібернетики, протокол \No 6 від   29.01.2021 р.\par\smallskip\smallskip
{\parbox {6.5 cm}{Зав. кафедри \dotfill Ю.В.~Крак}}\hspace {0.7cm}{\hbox to 6.5 cm{ Екзаменатор \dotfill Т.О.~Карнаух}} \end{small}}
\newpage

{{\begin {small}\par
{ \center  Київський національний університет імені Тараса Шевченка \par}
{\parbox {5 cm}{Освітня програма \hfill}{\parbox {7 cm}{Інформатика \hfill}}\parbox {6 cm}{\hfill Семестр 6}}\par
{\parbox {5 cm} {Навчальний предмет \hfill}\parbox {7 cm}{Розробка програмного забезпечення \hfill}\parbox {6cm}{\hfill}\par}\vspace{-7pt} \end{small}}{\center Екзаменаційний білет \No \textbf
{ 32 }\par}}
{
\begin {large}
{\begin {center}\textbf{Завдання 1}\end {center}}\par
Що з наведеного найкраще характеризує рефакторінг?\par
\vspace{2mm}



(\,1\,) Докорінна переробка всього коду з нуля,  що не змінює його видиму поведінку%bad



(\,2\,) Зміна внутрішньої структури коду під час усунення помилок%bad



(\,3\,) Довільна переробка коду, що виникає у процесі розробки%bad



(\,4\,) Зміна внутрішньої структури коду, що не змінює його видиму поведінку%good


{\begin {center}\textbf{Завдання 2}\end {center}}\par
Який з наведених фрагментів коду явно потребує поліпшення (форматування рядків та типи до уваги не брати)?\par
\vspace{2mm}



(\,1\,) \{double the\_best\_value\_anyone\_ever\_have;\} %bad



(\,2\,) void ff(…);%bad



(\,3\,) \{fstream inputFile, outputFile;…\} %good



(\,4\,) \{fstream input\_file, outputFile; …\}%bad


{\begin {center}\textbf{Завдання 3}\end {center}}\par
Що є характерним для структурного програмування?\par
\vspace{2mm}



(\,1\,) У коді мають використовуватись структурні типи даних. %bad



(\,2\,) У коді мають використовуватись типи структур, класи використовувати заборонено. %bad



(\,3\,) Використання оператора безумовного переходу. %bad



(\,4\,) Програма зображується ієрархічною структурою блоків.%good


{\begin {center}\textbf{Завдання 4}\end {center}}\par
Вибрати всі істинні твердження (якщо існують).\par
\vspace{2mm}



(\,1\,) У прогнозній моделі розробки розробка не може виконуватись ітеративно. %bad



(\,2\,) Неперервне проектування створює специфікацію проекту перед початком розробки/ітерації, а під час ітерації його неперервно модифікує.%bad



(\,3\,) Неперервне проектування здійснює створення та модифікацію проекту системи в міру розробки. %good



(\,4\,) Адаптивні моделі розробки допускають зміну вимог до продукту. %good


{\begin {center}\textbf{Завдання 5}\end {center}}\par
Вибрати всі істинні твердження (якщо існують).\par
\vspace{2mm}



(\,1\,) Правильно спроектована функція повинна мати рівно один обов'язок. %good



(\,2\,) Код програми пишеться для комп'ютера, тому він має бути синтаксично та семантично коректним. Інше значення не має. %bad



(\,3\,) Чим менше модулів у проекті, тим краще, бо меншу кількість модулів треба інтегрувати в єдине ціле. %bad



(\,4\,) Зв'язки між окремими модулями системи мають бути якомога більш слабкими. %good

\newpage

{\begin {center}\textbf{Завдання 6}\end {center}}\par
Вибрати всі істинні твердження (якщо існують).\par
\vspace{2mm}



(\,1\,) Якщо виникають проблеми з модифікацією коду, то слід розглянути внесення змін у проект. %good



(\,2\,) За відмови від документації одним з методом збереження знань про систему є постійне пересування персоналу між її частинами. %good



(\,3\,) Водоспадна методологія може мати тенденцію до надлишкового проектування. %good



(\,4\,) Гнучкі методології повністю нехтують документацією. %bad


{\begin {center}\textbf{Завдання 7}\end {center}}\par
Вибрати всі істинні твердження (якщо існують).\par
\vspace{2mm}



(\,1\,) Методологія розробки Big Design Up Front погана тим, що може призводити до надлишкового проектування. %good



(\,2\,) Код, що себе сам документує, -- це код, який виводить користувачу документацію на себе. %bad



(\,3\,) Гнучкі методології розробки допускають суттєві зміни вимог, проте все одно намагаються якомога швидше зафіксувати архітектуру. %good



(\,4\,) Простота – мистецтво максимізації незробленої роботи. %good


{\begin {center}\textbf{Завдання 8}\end {center}}\par
Що з наступного не має відношення до Scrum або не виконується/не використовується за використання оригінального Scrum?\par
\vspace{2mm}



(\,1\,) приведення у належний вигляд журналу невиконаних задач (backlog grooming) %bad



(\,2\,) журнал невиконаних задач (backlog) спрінту%bad



(\,3\,) порядок виконання запланованої на спрінт роботи визначає скрам-мастер%good



(\,4\,) скрам-мастер визначає час, необхідний команді для реалізації конкретної задачі%good


{\begin {center}\textbf{Завдання 9}\end {center}}\par
Що з наведеного нижче не розглядалось та навіть не згадувалось у курсі?\par
\vspace{2mm}



(\,1\,) X-модель%good



(\,2\,) Scrum of scrums%bad



(\,3\,) sashimi%bad



(\,4\,) RUP%bad


{\begin {center}\textbf{Завдання 10}\end {center}}\par
Вибрати всі істинні твердження (якщо існують).\par
\vspace{2mm}



(\,1\,) Розглядають модель якості продукту.%good



(\,2\,) Тестування буває позитивне та негативне.%good



(\,3\,) Якість буває зовнішньою.%good



(\,4\,) Якість буває внутрішньою.%good



\end {large}}
{\vspace{5mm}\smallskip \begin {small} Затверджено на засіданні кафедри теоретичної кібернетики, протокол \No 6 від   29.01.2021 р.\par\smallskip\smallskip
{\parbox {6.5 cm}{Зав. кафедри \dotfill Ю.В.~Крак}}\hspace {0.7cm}{\hbox to 6.5 cm{ Екзаменатор \dotfill Т.О.~Карнаух}} \end{small}}
\newpage

{{\begin {small}\par
{ \center  Київський національний університет імені Тараса Шевченка \par}
{\parbox {5 cm}{Освітня програма \hfill}{\parbox {7 cm}{Інформатика \hfill}}\parbox {6 cm}{\hfill Семестр 6}}\par
{\parbox {5 cm} {Навчальний предмет \hfill}\parbox {7 cm}{Розробка програмного забезпечення \hfill}\parbox {6cm}{\hfill}\par}\vspace{-7pt} \end{small}}{\center Екзаменаційний білет \No \textbf
{ 33 }\par}}
{
\begin {large}
{\begin {center}\textbf{Завдання 1}\end {center}}\par
Шаблон проектування:\par
\vspace{2mm}



(\,1\,) Задає розв'язання типової проблеми до найдрібніших подробиць%bad



(\,2\,) Формулює ідею розв'язання проблеми та може пропонувати одну чи кілька її реалізацій%good



(\,3\,) Тільки формулює ідею розв'язання проблеми%bad



(\,4\,) Є шаблоном, в який треба підставити власні ідентифікатори%bad


{\begin {center}\textbf{Завдання 2}\end {center}}\par
Який з наведених фрагментів коду явно потребує поліпшення (форматування рядків та типи до уваги не брати)?\par
\vspace{2mm}



(\,1\,) \{double avg\_temperature;…\} %good



(\,2\,) \{double max\_power, min\_power;…\} %good



(\,3\,) \{int a\_, a\_\_; …\}%bad



(\,4\,) void load\_data(…);%good


{\begin {center}\textbf{Завдання 3}\end {center}}\par
Що є характерним для структурного програмування?\par
\vspace{2mm}



(\,1\,) У коді мають використовуватись типи даних, що є класами. %bad



(\,2\,) Кожен блок коду повинен мати єдиний вхід та єдиний вихід.%good



(\,3\,) Використання механізму винятків. %bad



(\,4\,) У коді мають використовуватись типи структур, класи використовувати заборонено. %bad


{\begin {center}\textbf{Завдання 4}\end {center}}\par
Вибрати всі істинні твердження (якщо існують).\par
\vspace{2mm}



(\,1\,) Гнучкі методології не передбачають планування робіт. %bad



(\,2\,) Прогнозні моделі розробки з самого початку фіксують вимоги до продукту. %good



(\,3\,) У прогнозній моделі розробки розробка може виконуватись ітеративно. %good



(\,4\,) Ітеративні та інкрементні методи розробки є синонімом неперервного проектування. %bad


{\begin {center}\textbf{Завдання 5}\end {center}}\par
Вибрати всі істинні твердження (якщо існують).\par
\vspace{2mm}



(\,1\,) Одним з прийомів спрощення складної умовної логіки є використання шаблону Strategy. %good



(\,2\,) Гарний код не повинен містити різних фрагментів, що складаються з кількох інструкцій та виконують однакове перетворення даних. %good



(\,3\,) Чим вища кваліфікація програміста, тим більш складні для розуміння та довгі умовні вирази він використовує. %bad



(\,4\,) Повторюваний код має повторюватись у точності до символу, інакше він повторюваним вважатись не може. %bad

\newpage

{\begin {center}\textbf{Завдання 6}\end {center}}\par
Вибрати всі істинні твердження (якщо існують).\par
\vspace{2mm}



(\,1\,) Розробка за гнучкими методологіями ніколи не слідує плану. %bad



(\,2\,) Адаптивні методології розробки дозволяють уникнути надлишкового проектування. %good



(\,3\,) RUP є ітеративним та інкрементним.%good



(\,4\,) RUP є архітектурно-орієнтованим. %good


{\begin {center}\textbf{Завдання 7}\end {center}}\par
Вибрати всі істинні твердження (якщо існують).\par
\vspace{2mm}



(\,1\,) RUP принципово відрізняється від Scrum тим, що є керованим прецедентами. У той же час у Scrum при плануванні розробки прецеденти (чи їх аналоги) до уваги не беруться. %bad



(\,2\,) Якщо виникають проблеми з модифікацією коду, то слід відразу змінити команду розробників на більш кваліфіковану.%bad



(\,3\,) За застосування адаптивних методологій розробки, так само як і за прогнозних, намагаються якомога швидше визначити архітектуру системи. Відмінність полягає в тім, у який спосіб це виконується.%good



(\,4\,) Гнучкі методи розробки нехтують проміжними технічними артефактами та документацією виключно через їх непотрібність. %bad


{\begin {center}\textbf{Завдання 8}\end {center}}\par
Що з наступного не має відношення до Scrum або не виконується/не використовується за використання оригінального Scrum?\par
\vspace{2mm}



(\,1\,) блочні тести%bad



(\,2\,) обернені вимоги%bad



(\,3\,) залучення зовнішніх фахівців для побудови архітектури системи%good



(\,4\,) журнал невиконаних задач (backlog) продукту%bad


{\begin {center}\textbf{Завдання 9}\end {center}}\par
Що з наведеного нижче не розглядалось та навіть не згадувалось у курсі?\par
\vspace{2mm}



(\,1\,) FDD%bad



(\,2\,) LOAN%good



(\,3\,) LAZY%good



(\,4\,) RAD%bad


{\begin {center}\textbf{Завдання 10}\end {center}}\par
Вибрати всі істинні твердження (якщо існують).\par
\vspace{2mm}



(\,1\,) За провалених димових тестів решта тестів не виконується і код повертається на доробку.%good



(\,2\,) Можна розглядати якість під час використання.%good



(\,3\,) Розглядають модель якості даних.%good



(\,4\,) Тестування вимог може породжувати їх уточнення.%good



\end {large}}
{\vspace{5mm}\smallskip \begin {small} Затверджено на засіданні кафедри теоретичної кібернетики, протокол \No 6 від   29.01.2021 р.\par\smallskip\smallskip
{\parbox {6.5 cm}{Зав. кафедри \dotfill Ю.В.~Крак}}\hspace {0.7cm}{\hbox to 6.5 cm{ Екзаменатор \dotfill Т.О.~Карнаух}} \end{small}}
\newpage

{{\begin {small}\par
{ \center  Київський національний університет імені Тараса Шевченка \par}
{\parbox {5 cm}{Освітня програма \hfill}{\parbox {7 cm}{Інформатика \hfill}}\parbox {6 cm}{\hfill Семестр 6}}\par
{\parbox {5 cm} {Навчальний предмет \hfill}\parbox {7 cm}{Розробка програмного забезпечення \hfill}\parbox {6cm}{\hfill}\par}\vspace{-7pt} \end{small}}{\center Екзаменаційний білет \No \textbf
{ 34 }\par}}
{
\begin {large}
{\begin {center}\textbf{Завдання 1}\end {center}}\par
Так звані трикутники балансування зв'язують між собою:\par
\vspace{2mm}



(\,1\,) 4 або 3 параметри, при цьому може бути відсутня якість%bad



(\,2\,) Завжди рівно 4 параметри%bad



(\,3\,) Завжди рівно 3 параметри%bad



(\,4\,) 4 або 3 параметри, при цьому може бути відсутній обсяг %good


{\begin {center}\textbf{Завдання 2}\end {center}}\par
Який з наведених фрагментів коду явно потребує поліпшення (форматування рядків та типи до уваги не брати)?\par
\vspace{2mm}



(\,1\,) \{fstream inputFile, outputFile;…\} %good



(\,2\,) \{double maxPower, minPower;…\} %good



(\,3\,) \{int xtrmprg; …\}%bad



(\,4\,) void ff(…);%bad


{\begin {center}\textbf{Завдання 3}\end {center}}\par
Що є характерним для структурного програмування?\par
\vspace{2mm}



(\,1\,) Кожен блок коду повинен мати єдиний вхід та єдиний вихід.%good



(\,2\,) У коді мають використовуватись типи даних, що є класами. %bad



(\,3\,) У коді мають використовуватись структурні типи даних. %bad



(\,4\,) Використання механізму винятків. %bad


{\begin {center}\textbf{Завдання 4}\end {center}}\par
Вибрати всі істинні твердження (якщо існують).\par
\vspace{2mm}



(\,1\,) Адаптивні моделі розробки не допускають зміну вимог до продукту, але продукт може розроблятись поетапно. %bad



(\,2\,) Адаптивність методології проявляється в тому, що за зміни вимог розробка починається заново. %bad



(\,3\,) Адаптивні методології має сенс використовувати для проектів, де можлива зміна вимог. %good



(\,4\,) Прогнозна модель розробки не може одночасно бути інкрементною. %bad


{\begin {center}\textbf{Завдання 5}\end {center}}\par
Вибрати всі істинні твердження (якщо існують).\par
\vspace{2mm}



(\,1\,) Якщо кожна функція бібліотеки буде здатна в залежності від параметрів виконувати кілька різних дій, то загальна кількість функцій у ній буде менша і її буде простіше використовувати. %bad



(\,2\,) Від довгих функцій можна позбутися, розклавши їх у послідовність викликів дрібніших функцій. %good



(\,3\,) Втілення в коді різних підходів до одного того самого перетворення даних дозволяє програмісту продемонструвати свою професійну підготовку з найкращого боку. %bad



(\,4\,) Одним з прийомів спрощення складних умов є їх декомпозиція. %good

\newpage

{\begin {center}\textbf{Завдання 6}\end {center}}\par
Вибрати всі істинні твердження (якщо існують).\par
\vspace{2mm}



(\,1\,) Програмування в парах (pair programming) спрощує введення в команду нових розробників. %good



(\,2\,) Гнучкі методології розробки не допускають використання UML. %bad



(\,3\,) Використання RUP не заперечує часті випуски продукту, але й не наполягає на них. %good



(\,4\,) Гнучкі методології можуть передбачати як спільне, так і індивідуальне володіння кодом. %good


{\begin {center}\textbf{Завдання 7}\end {center}}\par
Вибрати всі істинні твердження (якщо існують).\par
\vspace{2mm}



(\,1\,) Багато хто з тих, що використовує Scrum, складає архітектурну модель розроблюваної системи. %good



(\,2\,) Основна відмінність V-моделі від водоспадної полягає в тім, що діяльності з тестування виконуються паралельно з усіма іншими діяльностями з розробки. %good



(\,3\,) Гнучкі методології базуються на комунікації розробників, а тому передбачають виключно колективне володіння кодом. %bad



(\,4\,) Сьогодні водоспадна методологія не має права на існування. %bad


{\begin {center}\textbf{Завдання 8}\end {center}}\par
Що з наступного не має відношення до Scrum або не виконується/не використовується за використання оригінального Scrum?\par
\vspace{2mm}



(\,1\,) щоденний скрам%bad



(\,2\,) журнал невиконаних задач (backlog) продукту%bad



(\,3\,) розповіді користувача%bad



(\,4\,) побудова різнопланових моделей системи%good


{\begin {center}\textbf{Завдання 9}\end {center}}\par
Що з наведеного нижче не розглядалось та навіть не згадувалось у курсі?\par
\vspace{2mm}



(\,1\,) TDD%bad



(\,2\,) SAFe%bad



(\,3\,) Scrum%bad



(\,4\,) CRAZY%good


{\begin {center}\textbf{Завдання 10}\end {center}}\par
Вибрати всі істинні твердження (якщо існують).\par
\vspace{2mm}



(\,1\,) Модель якості продукту є ієрархічною.%good



(\,2\,) Користувачі бувають непрямими.%good



(\,3\,) Модель якості при використанні є ієрархічною.%good



(\,4\,) Користувачі бувають вторинними.%good



\end {large}}
{\vspace{5mm}\smallskip \begin {small} Затверджено на засіданні кафедри теоретичної кібернетики, протокол \No 6 від   29.01.2021 р.\par\smallskip\smallskip
{\parbox {6.5 cm}{Зав. кафедри \dotfill Ю.В.~Крак}}\hspace {0.7cm}{\hbox to 6.5 cm{ Екзаменатор \dotfill Т.О.~Карнаух}} \end{small}}
\newpage

{{\begin {small}\par
{ \center  Київський національний університет імені Тараса Шевченка \par}
{\parbox {5 cm}{Освітня програма \hfill}{\parbox {7 cm}{Інформатика \hfill}}\parbox {6 cm}{\hfill Семестр 6}}\par
{\parbox {5 cm} {Навчальний предмет \hfill}\parbox {7 cm}{Розробка програмного забезпечення \hfill}\parbox {6cm}{\hfill}\par}\vspace{-7pt} \end{small}}{\center Екзаменаційний білет \No \textbf
{ 35 }\par}}
{
\begin {large}
{\begin {center}\textbf{Завдання 1}\end {center}}\par
Що є метою рефакторингу?\par
\vspace{2mm}



(\,1\,) Додавання нової функціональності%bad



(\,2\,) Усунення помилок у коді%bad



(\,3\,) Поліпшення внутрішньої структури коду%good



(\,4\,) Зменшення кількості рядків коду%bad


{\begin {center}\textbf{Завдання 2}\end {center}}\par
Який з наведених фрагментів коду явно потребує поліпшення (форматування рядків та типи до уваги не брати)?\par
\vspace{2mm}



(\,1\,) void help();%good



(\,2\,) void f13(…);%bad



(\,3\,) \{fstream input\_file, output\_file;…\} %good



(\,4\,) \{int input\_file, file\_to\_output; …\}%bad


{\begin {center}\textbf{Завдання 3}\end {center}}\par
Що є характерним для структурного програмування?\par
\vspace{2mm}



(\,1\,) У коді мають використовуватись типи структур. %bad



(\,2\,) Переходи з середини одного блоку в середину іншого неможливі. %good



(\,3\,) Заборона використання скалярних (простих) типів даних. %bad



(\,4\,) Використання оператора безумовного переходу. %bad


{\begin {center}\textbf{Завдання 4}\end {center}}\par
Вибрати всі істинні твердження (якщо існують).\par
\vspace{2mm}



(\,1\,) Для адаптивних методологій прогнозувати результат чи терміни його отримання абсолютно неможливо. %bad



(\,2\,) Прогнозні та адаптивні методології використовують одні ті самі діяльності з розробки ПЗ: визначення та аналіз вимог, проектування, реалізація, тестування, впровадження, супроводження. %good



(\,3\,) Неперервність проектування полягає в тім, що процес розробки не можна переривати більше ніж на вихідні. %bad



(\,4\,) Неперервне проектування має бути ітеративним та інкрементним одночасно. %good


{\begin {center}\textbf{Завдання 5}\end {center}}\par
Вибрати всі істинні твердження (якщо існують).\par
\vspace{2mm}



(\,1\,) Одним з засобів позбутись повторень у коді є використання функцій. %good



(\,2\,) Одним з засобів позбутись повторень у коді є використання успадкування. %good



(\,3\,) Повторюваний код дозволяє програмісту краще зрозуміти його логіку та вибрати з кількох варіантів найкращий. %bad



(\,4\,) Одним з прийомів спрощення складної умовної логіки є використання шаблону State. %good

\newpage

{\begin {center}\textbf{Завдання 6}\end {center}}\par
Вибрати всі істинні твердження (якщо існують).\par
\vspace{2mm}



(\,1\,) Водоспадна методологія розробки може призводити до надлишкового проектування. %good



(\,2\,) Розробка за гнучкими методологіями не передбачає узгодження умов контракту. %bad



(\,3\,) На відміну від Scrum, RUP є компонентно-орієнтованим. %bad



(\,4\,) Водоспадна методологія розробки може призводити до недостатнього проектування. %bad


{\begin {center}\textbf{Завдання 7}\end {center}}\par
Вибрати всі істинні твердження (якщо існують).\par
\vspace{2mm}



(\,1\,) Без рефакторінгу неперервне проектування неможливе. %good



(\,2\,) Ітеративна та інкрементна розробка передбачає наявність у тому чи іншому вигляді етапу ініціалізації, ітерацій та керуючої таблиці проекту. %good



(\,3\,) За використання гнучких методологій розробки припиняти розробку не рекомендується. %good



(\,4\,) Тільки RUP використовує UML.%bad


{\begin {center}\textbf{Завдання 8}\end {center}}\par
Що з наступного не має відношення до Scrum або не виконується/не використовується за використання оригінального Scrum?\par
\vspace{2mm}



(\,1\,) виділення в команді одного чи кількох проектувальників%good



(\,2\,) залучення зовнішніх фахівців для побудови архітектури системи%good



(\,3\,) власник продукту призначає час, необхідний команді для реалізації конкретної задачі%good



(\,4\,) обернені вимоги%bad


{\begin {center}\textbf{Завдання 9}\end {center}}\par
Що з наведеного нижче не розглядалось та навіть не згадувалось у курсі?\par
\vspace{2mm}



(\,1\,) Lean%bad



(\,2\,) poker%bad



(\,3\,) XP%bad



(\,4\,) BDUF%bad


{\begin {center}\textbf{Завдання 10}\end {center}}\par
Вибрати всі істинні твердження (якщо існують).\par
\vspace{2mm}



(\,1\,) Атрибути якості можна використовувати при уточненні вимог.%good



(\,2\,) Для тестування код треба ізолювати. Для того, щоб ізолювати, треба мати тести.%good



(\,3\,) Якість можна визначати через атрибути якості.%good



(\,4\,) Розглядають модель якості при використанні.%good



\end {large}}
{\vspace{5mm}\smallskip \begin {small} Затверджено на засіданні кафедри теоретичної кібернетики, протокол \No 6 від   29.01.2021 р.\par\smallskip\smallskip
{\parbox {6.5 cm}{Зав. кафедри \dotfill Ю.В.~Крак}}\hspace {0.7cm}{\hbox to 6.5 cm{ Екзаменатор \dotfill Т.О.~Карнаух}} \end{small}}
\newpage

{{\begin {small}\par
{ \center  Київський національний університет імені Тараса Шевченка \par}
{\parbox {5 cm}{Освітня програма \hfill}{\parbox {7 cm}{Інформатика \hfill}}\parbox {6 cm}{\hfill Семестр 6}}\par
{\parbox {5 cm} {Навчальний предмет \hfill}\parbox {7 cm}{Розробка програмного забезпечення \hfill}\parbox {6cm}{\hfill}\par}\vspace{-7pt} \end{small}}{\center Екзаменаційний білет \No \textbf
{ 36 }\par}}
{
\begin {large}
{\begin {center}\textbf{Завдання 1}\end {center}}\par
Що з наведеного найкраще характеризує рефакторінг?\par
\vspace{2mm}



(\,1\,) Докорінна переробка всього коду з нуля,  що не змінює його видиму поведінку%bad



(\,2\,) Довільна переробка коду, що виникає у процесі розробки%bad



(\,3\,) Зміна внутрішньої структури коду під час усунення помилок%bad



(\,4\,) Зміна внутрішньої структури коду, що не змінює його видиму поведінку%good


{\begin {center}\textbf{Завдання 2}\end {center}}\par
Який з наведених фрагментів коду явно потребує поліпшення (форматування рядків та типи до уваги не брати)?\par
\vspace{2mm}



(\,1\,) \{int x; int y; …\}%good



(\,2\,) \{fstream input\_file, outputFile; …\}%bad



(\,3\,) void load\_data(T \& arg1, double arg2, double arg3, double arg4); %bad



(\,4\,) \{double the\_best\_value\_anyone\_ever\_have;\} %bad


{\begin {center}\textbf{Завдання 3}\end {center}}\par
Що є характерним для структурного програмування?\par
\vspace{2mm}



(\,1\,) У коді мають використовуватись типи структур, класи використовувати заборонено. %bad



(\,2\,) Заборона використання скалярних (простих) типів даних. %bad



(\,3\,) У коді мають використовуватись типи структур. %bad



(\,4\,) Програма зображується ієрархічною структурою блоків.%good


{\begin {center}\textbf{Завдання 4}\end {center}}\par
Вибрати всі істинні твердження (якщо існують).\par
\vspace{2mm}



(\,1\,) Адаптивні методології не передбачають визначення вимог. %bad



(\,2\,) Для адаптивних моделей розробки оцінки часу та бюджету не виконують, бо вимоги є змінними. %bad



(\,3\,) Адаптивна модель розробки адаптує вимоги до продукту до наявної команди розробників. %bad



(\,4\,) Прогнозна модель розробки може одночасно бути інкрементною. %good


{\begin {center}\textbf{Завдання 5}\end {center}}\par
Вибрати всі істинні твердження (якщо існують).\par
\vspace{2mm}



(\,1\,) Складні умови виникають в реальних задачах доволі часто, а тому є доволі природніми і не можуть вважатись недоліком коду. %bad



(\,2\,) Код з довгими функціями має меншу кількість функцій, а тому для нього треба менше тестів. %bad



(\,3\,) Код з довгими функціями має меншу кількість функцій, а тому його швидше та простіше тестувати. %bad



(\,4\,) Одним з засобів позбутись повторень у коді є використання шаблону Template Method (але цей шаблон здатен боротися і з іншими недоліками коду). %good

\newpage

{\begin {center}\textbf{Завдання 6}\end {center}}\par
Вибрати всі істинні твердження (якщо існують).\par
\vspace{2mm}



(\,1\,) Використання гнучких методологій вимагає високої кваліфікації персоналу. %good



(\,2\,) За використання гнучких методологій кожен розробник робить, що хоче. %bad



(\,3\,) Гнучкі методи розробки нехтують проміжними технічними артефактами та документацією через те, що за зміни вимог та/або проекту треба витрачати додатковий час на підтримання їх у відповідному стані. %good



(\,4\,) Висока внутрішня якість коду здешевшує розробку, бо дозволяє економити на вартості та часі майбутніх змін. %good


{\begin {center}\textbf{Завдання 7}\end {center}}\par
Вибрати всі істинні твердження (якщо існують).\par
\vspace{2mm}



(\,1\,) Ітеративна методологія розробки у загальному випадку не гарантує, що не буде надлишкового проектування. %good



(\,2\,) Розробка за гнучкими методологіями цінить людей та взаємодію між ними. %good



(\,3\,) Висока внутрішня якість коду є марною тратою часу на те, що замовник не побачить та/або не зможе зацінити. %bad



(\,4\,) Водоспадна методологія не призначена для внесення змін у вимоги до системи. %good


{\begin {center}\textbf{Завдання 8}\end {center}}\par
Що з наступного не має відношення до Scrum або не виконується/не використовується за використання оригінального Scrum?\par
\vspace{2mm}



(\,1\,) власник продукту остаточно затверджує проект системи%good



(\,2\,) правила визначення готовності (Definition of Done) %bad



(\,3\,) скрам-мастер призначає відповідальних за конкретні задачі%good



(\,4\,) скрам-мастер визначає, які задачі команда буде виконувати на наступному спрінті%good


{\begin {center}\textbf{Завдання 9}\end {center}}\par
Що з наведеного нижче не розглядалось та навіть не згадувалось у курсі?\par
\vspace{2mm}



(\,1\,) водоспадна модель%bad



(\,2\,) LAZY%good



(\,3\,) модель хаосу%bad



(\,4\,) LOAN%good


{\begin {center}\textbf{Завдання 10}\end {center}}\par
Вибрати всі істинні твердження (якщо існують).\par
\vspace{2mm}



(\,1\,) І верифікація, і валідація можуть використовувати тестування.%good



(\,2\,) Забезпечення якості це не тільки тестування.%good



(\,3\,) Якість можна визначати через дефекти.%good



(\,4\,) Коли нема можливості перевірити всі можливі комбінації, тестують їх попарні (pairwise) комбінації. Це може значно зменшити кількість тестів.%good



\end {large}}
{\vspace{5mm}\smallskip \begin {small} Затверджено на засіданні кафедри теоретичної кібернетики, протокол \No 6 від   29.01.2021 р.\par\smallskip\smallskip
{\parbox {6.5 cm}{Зав. кафедри \dotfill Ю.В.~Крак}}\hspace {0.7cm}{\hbox to 6.5 cm{ Екзаменатор \dotfill Т.О.~Карнаух}} \end{small}}
\newpage

{{\begin {small}\par
{ \center  Київський національний університет імені Тараса Шевченка \par}
{\parbox {5 cm}{Освітня програма \hfill}{\parbox {7 cm}{Інформатика \hfill}}\parbox {6 cm}{\hfill Семестр 6}}\par
{\parbox {5 cm} {Навчальний предмет \hfill}\parbox {7 cm}{Розробка програмного забезпечення \hfill}\parbox {6cm}{\hfill}\par}\vspace{-7pt} \end{small}}{\center Екзаменаційний білет \No \textbf
{ 37 }\par}}
{
\begin {large}
{\begin {center}\textbf{Завдання 1}\end {center}}\par
Що є метою рефакторингу?\par
\vspace{2mm}



(\,1\,) Додавання нової функціональності%bad



(\,2\,) Зменшення кількості рядків коду%bad



(\,3\,) Усунення помилок у коді%bad



(\,4\,) Поліпшення внутрішньої структури коду%good


{\begin {center}\textbf{Завдання 2}\end {center}}\par
Який з наведених фрагментів коду явно потребує поліпшення (форматування рядків та типи до уваги не брати)?\par
\vspace{2mm}



(\,1\,) \{int xtrmprg; …\}%bad



(\,2\,) \{int a\_, a\_\_; …\}%bad



(\,3\,) \{int input\_file, file\_to\_output; …\}%bad



(\,4\,) \{double maxPower, minPower;…\} %good


{\begin {center}\textbf{Завдання 3}\end {center}}\par
Що є характерним для структурного програмування?\par
\vspace{2mm}



(\,1\,) Використання оператора безумовного переходу. %bad



(\,2\,) У коді мають використовуватись типи даних, що є класами. %bad



(\,3\,) Програма зображується ієрархічною структурою блоків.%good



(\,4\,) У коді мають використовуватись структурні типи даних. %bad


{\begin {center}\textbf{Завдання 4}\end {center}}\par
Вибрати всі істинні твердження (якщо існують).\par
\vspace{2mm}



(\,1\,) Неперервне проектування має бути ітеративним та інкрементним одночасно. %good



(\,2\,) Неперервне проектування здійснює створення та модифікацію проекту системи в міру розробки. %good



(\,3\,) У прогнозній моделі розробки розробка може виконуватись ітеративно. %good



(\,4\,) Прогнозна модель розробки не може одночасно бути інкрементною. %bad


{\begin {center}\textbf{Завдання 5}\end {center}}\par
Вибрати всі істинні твердження (якщо існують).\par
\vspace{2mm}



(\,1\,) Довгі функції доволі часто приховують складну логіку обробки або повторюваний код. %good



(\,2\,) Одним з засобів позбутись повторень у коді є використання функцій. %good



(\,3\,) Повторюваний код дозволяє програмісту краще зрозуміти його логіку та вибрати з кількох варіантів найкращий. %bad



(\,4\,) Код з довгими функціями має меншу кількість функцій, а тому для нього треба менше тестів. %bad

\newpage

{\begin {center}\textbf{Завдання 6}\end {center}}\par
Вибрати всі істинні твердження (якщо існують).\par
\vspace{2mm}



(\,1\,) Використання гнучких методологій при розробці надвеликих проектів не несе в собі жодних ризиків. %bad



(\,2\,) Методологія розробки Big Design Up Front погана тим, що може призводити до недостатнього проектування.%bad



(\,3\,) RUP є керованим прецедентами. %good



(\,4\,) Scrum передбачає колективне володіння кодом. %good


{\begin {center}\textbf{Завдання 7}\end {center}}\par
Вибрати всі істинні твердження (якщо існують).\par
\vspace{2mm}



(\,1\,) Програмування в парах (pair programming) зводиться до того, що кожним фрагментом коду володіють два програмісти. %bad



(\,2\,) Поступове уточнення вимог замовника полегшує адаптацію до змін. %good



(\,3\,) Гнучкі методології передбачають часті випуски продукту. %good



(\,4\,) RUP є компонентно-орієнтованим. %good


{\begin {center}\textbf{Завдання 8}\end {center}}\par
Що з наступного не має відношення до Scrum або не виконується/не використовується за використання оригінального Scrum?\par
\vspace{2mm}



(\,1\,) виділення одного з учасників команди як системного архітектора%good



(\,2\,) журнал невиконаних задач (backlog) спрінту%bad



(\,3\,) скрам-мастер остаточно затверджує проект системи%good



(\,4\,) виділення в команді одного чи кількох тестувальників%good


{\begin {center}\textbf{Завдання 9}\end {center}}\par
Що з наведеного нижче не розглядалось та навіть не згадувалось у курсі?\par
\vspace{2mm}



(\,1\,) poker%bad



(\,2\,) Cleanroom%bad



(\,3\,) X-модель%good



(\,4\,) DSDM%bad


{\begin {center}\textbf{Завдання 10}\end {center}}\par
Вибрати всі істинні твердження (якщо існують).\par
\vspace{2mm}



(\,1\,) Тестування буває позитивне та негативне.%good



(\,2\,) Якість можна визначати через атрибути якості.%good



(\,3\,) Якість можна визначати через дефекти.%good



(\,4\,) Модель якості даних є ієрархічною.%good



\end {large}}
{\vspace{5mm}\smallskip \begin {small} Затверджено на засіданні кафедри теоретичної кібернетики, протокол \No 6 від   29.01.2021 р.\par\smallskip\smallskip
{\parbox {6.5 cm}{Зав. кафедри \dotfill Ю.В.~Крак}}\hspace {0.7cm}{\hbox to 6.5 cm{ Екзаменатор \dotfill Т.О.~Карнаух}} \end{small}}
\newpage

{{\begin {small}\par
{ \center  Київський національний університет імені Тараса Шевченка \par}
{\parbox {5 cm}{Освітня програма \hfill}{\parbox {7 cm}{Інформатика \hfill}}\parbox {6 cm}{\hfill Семестр 6}}\par
{\parbox {5 cm} {Навчальний предмет \hfill}\parbox {7 cm}{Розробка програмного забезпечення \hfill}\parbox {6cm}{\hfill}\par}\vspace{-7pt} \end{small}}{\center Екзаменаційний білет \No \textbf
{ 38 }\par}}
{
\begin {large}
{\begin {center}\textbf{Завдання 1}\end {center}}\par
Так звані трикутники балансування зв'язують між собою:\par
\vspace{2mm}



(\,1\,) 4 або 3 параметри, при цьому може бути відсутня якість%bad



(\,2\,) Завжди рівно 4 параметри%bad



(\,3\,) Завжди рівно 3 параметри%bad



(\,4\,) 4 або 3 параметри, при цьому може бути відсутній обсяг %good


{\begin {center}\textbf{Завдання 2}\end {center}}\par
Який з наведених фрагментів коду явно потребує поліпшення (форматування рядків та типи до уваги не брати)?\par
\vspace{2mm}



(\,1\,) \{fstream input\_file, outputFile; …\}%bad



(\,2\,) void load\_data(T \& arg1, double arg2, double arg3, double arg4); %bad



(\,3\,) void ff(…);%bad



(\,4\,) void load\_data(…);%good


{\begin {center}\textbf{Завдання 3}\end {center}}\par
Що є характерним для структурного програмування?\par
\vspace{2mm}



(\,1\,) Використання механізму винятків. %bad



(\,2\,) Використання механізму винятків. %bad



(\,3\,) Кожен блок коду повинен мати єдиний вхід та єдиний вихід.%good



(\,4\,) У коді мають використовуватись типи даних, що є класами. %bad


{\begin {center}\textbf{Завдання 4}\end {center}}\par
Вибрати всі істинні твердження (якщо існують).\par
\vspace{2mm}



(\,1\,) В адаптивних методологіях планування робіт не виконується. %bad



(\,2\,) Ітеративні та інкрементні методи розробки є синонімом неперервного проектування. %bad



(\,3\,) Неперервність проектування полягає в тім, що процес розробки не можна переривати більше ніж на вихідні. %bad



(\,4\,) Неперервне проектування створює специфікацію проекту перед початком розробки/ітерації, а під час ітерації його неперервно модифікує.%bad


{\begin {center}\textbf{Завдання 5}\end {center}}\par
Вибрати всі істинні твердження (якщо існують).\par
\vspace{2mm}



(\,1\,) Одним з прийомів спрощення складної умовної логіки є використання шаблону State. %good



(\,2\,) Довжина функцій не впливає на внутрішню якість коду. %bad



(\,3\,) Одним з прийомів спрощення умовної логіки є використання поліморфізму. %good



(\,4\,) Гарний код не повинен містити різних фрагментів, що складаються з кількох інструкцій та виконують однакове перетворення даних. %good

\newpage

{\begin {center}\textbf{Завдання 6}\end {center}}\par
Вибрати всі істинні твердження (якщо існують).\par
\vspace{2mm}



(\,1\,) Модель хаосу є методологією розробки програмного забезпечення. %bad



(\,2\,) RUP є прогнозною методологією. %bad



(\,3\,) За програмування в парах (pair programming) поки один колега набирає код на комп'ютері інший може піти пограти в настільний теніс. %bad



(\,4\,) Одна з принципових відмінностей RUP від Scrum полягає у використанні численних проміжних технічних артефактів. %good


{\begin {center}\textbf{Завдання 7}\end {center}}\par
Вибрати всі істинні твердження (якщо існують).\par
\vspace{2mm}



(\,1\,) Гнучкі методології розробки є стійкими до припинення розробки на певний час. %bad



(\,2\,) Розповіді користувача (user stories) – це прогресивний інструмент гнучких методологій, який нічого спільного не має з прецедентами (use case). %bad



(\,3\,) Програмування в парах (pair programming) підвищує якість коду та дозволяє економити на виправленні помилок. %good



(\,4\,) За невизначених або піддатливих змінам вимог краще застосовувати гнучкі методології розробки. %good


{\begin {center}\textbf{Завдання 8}\end {center}}\par
Що з наступного не має відношення до Scrum або не виконується/не використовується за використання оригінального Scrum?\par
\vspace{2mm}



(\,1\,) порядок виконання запланованої на спрінт роботи визначає власник продукту%good



(\,2\,) скрам-мастер визначає пріоритети задач%good



(\,3\,) під час спрінту зміни, що впливають на досягнення мети спрінту, зазвичай, не допускаються%bad



(\,4\,) планування спрінту%bad


{\begin {center}\textbf{Завдання 9}\end {center}}\par
Що з наведеного нижче не розглядалось та навіть не згадувалось у курсі?\par
\vspace{2mm}



(\,1\,) BDD%bad



(\,2\,) CRAZY%good



(\,3\,) PRINCE2%good



(\,4\,) RUP%bad


{\begin {center}\textbf{Завдання 10}\end {center}}\par
Вибрати всі істинні твердження (якщо існують).\par
\vspace{2mm}



(\,1\,) Відсутність знайдених дефектів на початкових стадіях розробки може бути свідченням поганого покриття коду тестами.%good



(\,2\,) Користувачі бувають непрямими.%good



(\,3\,) Тестовий приклад складається з набору вхідних даних, умов виконання та очікуваних результатів.%good



(\,4\,) За провалених димових тестів решта тестів не виконується і код повертається на доробку.%good



\end {large}}
{\vspace{5mm}\smallskip \begin {small} Затверджено на засіданні кафедри теоретичної кібернетики, протокол \No 6 від   29.01.2021 р.\par\smallskip\smallskip
{\parbox {6.5 cm}{Зав. кафедри \dotfill Ю.В.~Крак}}\hspace {0.7cm}{\hbox to 6.5 cm{ Екзаменатор \dotfill Т.О.~Карнаух}} \end{small}}
\newpage

{{\begin {small}\par
{ \center  Київський національний університет імені Тараса Шевченка \par}
{\parbox {5 cm}{Освітня програма \hfill}{\parbox {7 cm}{Інформатика \hfill}}\parbox {6 cm}{\hfill Семестр 6}}\par
{\parbox {5 cm} {Навчальний предмет \hfill}\parbox {7 cm}{Розробка програмного забезпечення \hfill}\parbox {6cm}{\hfill}\par}\vspace{-7pt} \end{small}}{\center Екзаменаційний білет \No \textbf
{ 39 }\par}}
{
\begin {large}
{\begin {center}\textbf{Завдання 1}\end {center}}\par
Шаблон проектування:\par
\vspace{2mm}



(\,1\,) Є шаблоном, в який треба підставити власні ідентифікатори%bad



(\,2\,) Задає розв'язання типової проблеми до найдрібніших подробиць%bad



(\,3\,) Тільки формулює ідею розв'язання проблеми%bad



(\,4\,) Формулює ідею розв'язання проблеми та може пропонувати одну чи кілька її реалізацій%good


{\begin {center}\textbf{Завдання 2}\end {center}}\par
Який з наведених фрагментів коду явно потребує поліпшення (форматування рядків та типи до уваги не брати)?\par
\vspace{2mm}



(\,1\,) \{double max\_power, min\_power;…\} %good



(\,2\,) \{fstream input\_file, output\_file;…\} %good



(\,3\,) void help();%good



(\,4\,) \{double avg\_temperature;…\} %good


{\begin {center}\textbf{Завдання 3}\end {center}}\par
Що є характерним для структурного програмування?\par
\vspace{2mm}



(\,1\,) Переходи з середини одного блоку в середину іншого неможливі. %good



(\,2\,) Заборона використання скалярних (простих) типів даних. %bad



(\,3\,) Використання оператора безумовного переходу. %bad



(\,4\,) У коді мають використовуватись типи структур. %bad


{\begin {center}\textbf{Завдання 4}\end {center}}\par
Вибрати всі істинні твердження (якщо існують).\par
\vspace{2mm}



(\,1\,) Прогнозні моделі розробки з самого початку фіксують вимоги до продукту. %good



(\,2\,) У прогнозній моделі розробки розробка не може виконуватись ітеративно. %bad



(\,3\,) У випадку, коли вимоги до продукту відомі заздалегідь та є фіксованими контрактом, адаптивні методології розробки використовувати не можна. %bad



(\,4\,) Адаптивна модель розробки адаптує вимоги до продукту до наявної команди розробників. %bad


{\begin {center}\textbf{Завдання 5}\end {center}}\par
Вибрати всі істинні твердження (якщо існують).\par
\vspace{2mm}



(\,1\,) Складні комбінації умов починають керувати логікою виконання, тому що такими є вимоги до функціонування ПЗ. Оскільки ПЗ має якнайкраще забезпечувати потреби замовника, то складні комбінації умов у коді вважатись недоліком коду не можуть. %bad



(\,2\,) Одним з прийомів спрощення складної умовної логіки є використання шаблону Strategy. %good



(\,3\,) Залежності всередині модуля можуть бути сильними. %good



(\,4\,) Одним з засобів позбутись повторень у коді є використання шаблону Template Method (але цей шаблон здатен боротися і з іншими недоліками коду). %good

\newpage

{\begin {center}\textbf{Завдання 6}\end {center}}\par
Вибрати всі істинні твердження (якщо існують).\par
\vspace{2mm}



(\,1\,) RUP відрізняється від Scrum тим, що ні в що не ставить персонал.%bad



(\,2\,) Методологія розробки Big Design Up Front погана тим, що може призводити до недостатнього проектування.%bad



(\,3\,) Простота – мистецтво максимізації незробленої роботи. %good



(\,4\,) Гнучкі методології повністю нехтують документацією. %bad


{\begin {center}\textbf{Завдання 7}\end {center}}\par
Вибрати всі істинні твердження (якщо існують).\par
\vspace{2mm}



(\,1\,) Використання гнучких методологій вимагає високої кваліфікації персоналу. %good



(\,2\,) Ітеративна та інкрементна розробка передбачає наявність у тому чи іншому вигляді етапу ініціалізації, ітерацій та керуючої таблиці проекту. %good



(\,3\,) RUP є керованим прецедентами. %good



(\,4\,) Основна відмінність V-моделі від водоспадної полягає в тім, що діяльності з тестування виконуються паралельно з усіма іншими діяльностями з розробки. %good


{\begin {center}\textbf{Завдання 8}\end {center}}\par
Що з наступного не має відношення до Scrum або не виконується/не використовується за використання оригінального Scrum?\par
\vspace{2mm}



(\,1\,) власник продукту призначає відповідальних за конкретні задачі%good



(\,2\,) розповіді користувача%bad



(\,3\,) ретроспектива спрінту%bad



(\,4\,) скрам-мастер визначає час, необхідний команді для реалізації конкретної задачі%good


{\begin {center}\textbf{Завдання 9}\end {center}}\par
Що з наведеного нижче не розглядалось та навіть не згадувалось у курсі?\par
\vspace{2mm}



(\,1\,) V-модель%bad



(\,2\,) TDD%bad



(\,3\,) Lean%bad



(\,4\,) sashimi%bad


{\begin {center}\textbf{Завдання 10}\end {center}}\par
Вибрати всі істинні твердження (якщо існують).\par
\vspace{2mm}



(\,1\,) Модель якості продукту є ієрархічною.%good



(\,2\,) Якість буває зовнішньою.%good



(\,3\,) Користувачі бувають вторинними.%good



(\,4\,) Розглядають модель якості продукту.%good



\end {large}}
{\vspace{5mm}\smallskip \begin {small} Затверджено на засіданні кафедри теоретичної кібернетики, протокол \No 6 від   29.01.2021 р.\par\smallskip\smallskip
{\parbox {6.5 cm}{Зав. кафедри \dotfill Ю.В.~Крак}}\hspace {0.7cm}{\hbox to 6.5 cm{ Екзаменатор \dotfill Т.О.~Карнаух}} \end{small}}
\newpage

{{\begin {small}\par
{ \center  Київський національний університет імені Тараса Шевченка \par}
{\parbox {5 cm}{Освітня програма \hfill}{\parbox {7 cm}{Інформатика \hfill}}\parbox {6 cm}{\hfill Семестр 6}}\par
{\parbox {5 cm} {Навчальний предмет \hfill}\parbox {7 cm}{Розробка програмного забезпечення \hfill}\parbox {6cm}{\hfill}\par}\vspace{-7pt} \end{small}}{\center Екзаменаційний білет \No \textbf
{ 40 }\par}}
{
\begin {large}
{\begin {center}\textbf{Завдання 1}\end {center}}\par
Так звані трикутники балансування зв'язують між собою:\par
\vspace{2mm}



(\,1\,) 4 або 3 параметри, при цьому може бути відсутній час%bad



(\,2\,) 4 або 3 параметри, при цьому може бути відсутній обсяг %good



(\,3\,) Завжди рівно 4 параметри%bad



(\,4\,) Завжди рівно 3 параметри%bad


{\begin {center}\textbf{Завдання 2}\end {center}}\par
Який з наведених фрагментів коду явно потребує поліпшення (форматування рядків та типи до уваги не брати)?\par
\vspace{2mm}



(\,1\,) \{double the\_best\_value\_anyone\_ever\_have;\} %bad



(\,2\,) \{fstream inputFile, outputFile;…\} %good



(\,3\,) \{int x; int y; …\}%good



(\,4\,) void f13(…);%bad


{\begin {center}\textbf{Завдання 3}\end {center}}\par
Що є характерним для структурного програмування?\par
\vspace{2mm}



(\,1\,) У коді мають використовуватись типи структур, класи використовувати заборонено. %bad



(\,2\,) Кожен блок коду повинен мати єдиний вхід та єдиний вихід.%good



(\,3\,) У коді мають використовуватись структурні типи даних. %bad



(\,4\,) Використання механізму винятків. %bad


{\begin {center}\textbf{Завдання 4}\end {center}}\par
Вибрати всі істинні твердження (якщо існують).\par
\vspace{2mm}



(\,1\,) Для адаптивних методологій прогнозувати результат чи терміни його отримання абсолютно неможливо. %bad



(\,2\,) Адаптивні моделі розробки допускають зміну вимог до продукту. %good



(\,3\,) Адаптивні моделі розробки не допускають зміну вимог до продукту, але продукт може розроблятись поетапно. %bad



(\,4\,) І прогнозні, і адаптивні методології передбачають планування виконуваних робіт. %good


{\begin {center}\textbf{Завдання 5}\end {center}}\par
Вибрати всі істинні твердження (якщо існують).\par
\vspace{2mm}



(\,1\,) Чим менше класів у коді, тим менше класів треба налагоджувати, а тому класи за можливості слід об'єднувати. %bad



(\,2\,) Одним з засобів позбутись повторень у коді є використання успадкування. %good



(\,3\,) Зв'язки між окремими модулями системи мають бути якомога більш слабкими. %good



(\,4\,) Якщо кожна функція бібліотеки буде здатна в залежності від параметрів виконувати кілька різних дій, то загальна кількість функцій у ній буде менша і її буде простіше використовувати. %bad

\newpage

{\begin {center}\textbf{Завдання 6}\end {center}}\par
Вибрати всі істинні твердження (якщо існують).\par
\vspace{2mm}



(\,1\,) Розробка за гнучкими методологіями ніколи не слідує плану. %bad



(\,2\,) Тільки RUP використовує UML.%bad



(\,3\,) Гнучкі методології можуть передбачати як спільне, так і індивідуальне володіння кодом. %good



(\,4\,) Багато хто з тих, що використовує Scrum, складає архітектурну модель розроблюваної системи. %good


{\begin {center}\textbf{Завдання 7}\end {center}}\par
Вибрати всі істинні твердження (якщо існують).\par
\vspace{2mm}



(\,1\,) Розробка за гнучкими методологіями не передбачає узгодження умов контракту. %bad



(\,2\,) RUP принципово відрізняється від Scrum тим, що є керованим прецедентами. У той же час у Scrum при плануванні розробки прецеденти (чи їх аналоги) до уваги не беруться. %bad



(\,3\,) RUP відрізняється від Scrum тим, що ні в що не ставить персонал.%bad



(\,4\,) Водоспадна методологія розробки може призводити до надлишкового проектування. %good


{\begin {center}\textbf{Завдання 8}\end {center}}\par
Що з наступного не має відношення до Scrum або не виконується/не використовується за використання оригінального Scrum?\par
\vspace{2mm}



(\,1\,) канбан%bad



(\,2\,) команда визначає, які задачі вона буде виконувати на наступному спрінті%good



(\,3\,) бачення продукту (vision) %bad



(\,4\,) приведення у належний вигляд журналу невиконаних задач (backlog grooming) %bad


{\begin {center}\textbf{Завдання 9}\end {center}}\par
Що з наведеного нижче не розглядалось та навіть не згадувалось у курсі?\par
\vspace{2mm}



(\,1\,) водоспадна модель%bad



(\,2\,) SAFe%bad



(\,3\,) LeSS%bad



(\,4\,) RAD%bad


{\begin {center}\textbf{Завдання 10}\end {center}}\par
Вибрати всі істинні твердження (якщо існують).\par
\vspace{2mm}



(\,1\,) Модель якості при використанні є ієрархічною.%good



(\,2\,) Розглядають модель якості даних.%good



(\,3\,) Тестування вимог може породжувати їх уточнення.%good



(\,4\,) Якість буває внутрішньою.%good



\end {large}}
{\vspace{5mm}\smallskip \begin {small} Затверджено на засіданні кафедри теоретичної кібернетики, протокол \No 6 від   29.01.2021 р.\par\smallskip\smallskip
{\parbox {6.5 cm}{Зав. кафедри \dotfill Ю.В.~Крак}}\hspace {0.7cm}{\hbox to 6.5 cm{ Екзаменатор \dotfill Т.О.~Карнаух}} \end{small}}
\newpage

{{\begin {small}\par
{ \center  Київський національний університет імені Тараса Шевченка \par}
{\parbox {5 cm}{Освітня програма \hfill}{\parbox {7 cm}{Інформатика \hfill}}\parbox {6 cm}{\hfill Семестр 6}}\par
{\parbox {5 cm} {Навчальний предмет \hfill}\parbox {7 cm}{Розробка програмного забезпечення \hfill}\parbox {6cm}{\hfill}\par}\vspace{-7pt} \end{small}}{\center Екзаменаційний білет \No \textbf
{ 41 }\par}}
{
\begin {large}
{\begin {center}\textbf{Завдання 1}\end {center}}\par
Так звані трикутники балансування зв'язують між собою:\par
\vspace{2mm}



(\,1\,) 4 або 3 параметри, при цьому може бути відсутній час%bad



(\,2\,) 4 або 3 параметри, при цьому може бути відсутня якість%bad



(\,3\,) Завжди рівно 3 параметри%bad



(\,4\,) 4 або 3 параметри, при цьому може бути відсутній обсяг %good


{\begin {center}\textbf{Завдання 2}\end {center}}\par
Який з наведених фрагментів коду явно потребує поліпшення (форматування рядків та типи до уваги не брати)?\par
\vspace{2mm}



(\,1\,) \{int xtrmprg; …\}%bad



(\,2\,) \{fstream inputFile, outputFile;…\} %good



(\,3\,) \{double maxPower, minPower;…\} %good



(\,4\,) \{fstream input\_file, output\_file;…\} %good


{\begin {center}\textbf{Завдання 3}\end {center}}\par
Що є характерним для структурного програмування?\par
\vspace{2mm}



(\,1\,) У коді мають використовуватись структурні типи даних. %bad



(\,2\,) Програма зображується ієрархічною структурою блоків.%good



(\,3\,) У коді мають використовуватись типи даних, що є класами. %bad



(\,4\,) У коді мають використовуватись типи структур. %bad


{\begin {center}\textbf{Завдання 4}\end {center}}\par
Вибрати всі істинні твердження (якщо існують).\par
\vspace{2mm}



(\,1\,) Для адаптивних моделей розробки оцінки часу та бюджету не виконують, бо вимоги є змінними. %bad



(\,2\,) Адаптивні методології має сенс використовувати для проектів, де можлива зміна вимог. %good



(\,3\,) Для адаптивних методологій розробки, на відміну від прогнозних, попередні оцінки бюджету та/або часу не виконують, бо вони все одно зміняться. %bad



(\,4\,) Адаптивність методології проявляється в тому, що за зміни вимог розробка починається заново. %bad


{\begin {center}\textbf{Завдання 5}\end {center}}\par
Вибрати всі істинні твердження (якщо існують).\par
\vspace{2mm}



(\,1\,) Втілення в коді різних підходів до одного того самого перетворення даних дозволяє програмісту продемонструвати свою професійну підготовку з найкращого боку. %bad



(\,2\,) Чим менше модулів у проекті, тим краще, бо меншу кількість модулів треба інтегрувати в єдине ціле. %bad



(\,3\,) Гарний код має читатись як розмовна мова. %good



(\,4\,) Повторюваний код має повторюватись у точності до символу, інакше він повторюваним вважатись не може. %bad

\newpage

{\begin {center}\textbf{Завдання 6}\end {center}}\par
Вибрати всі істинні твердження (якщо існують).\par
\vspace{2mm}



(\,1\,) Якщо виникають проблеми з модифікацією коду, то слід розглянути внесення змін у проект. %good



(\,2\,) Водоспадна методологія розробки може призводити до недостатнього проектування. %bad



(\,3\,) Використання гнучких методологій при розробці надвеликих проектів не несе в собі жодних ризиків. %bad



(\,4\,) Гнучкі методи розробки нехтують проміжними технічними артефактами та документацією виключно через їх непотрібність. %bad


{\begin {center}\textbf{Завдання 7}\end {center}}\par
Вибрати всі істинні твердження (якщо існують).\par
\vspace{2mm}



(\,1\,) За відмови від документації одним з методом збереження знань про систему є постійне пересування персоналу між її частинами. %good



(\,2\,) Ітеративна методологія розробки у загальному випадку не гарантує, що не буде надлишкового проектування. %good



(\,3\,) Гнучкі методології передбачають часті випуски продукту. %good



(\,4\,) Перша версія коду не завжди може бути якісною -- і це нормально. %good


{\begin {center}\textbf{Завдання 8}\end {center}}\par
Що з наступного не має відношення до Scrum або не виконується/не використовується за використання оригінального Scrum?\par
\vspace{2mm}



(\,1\,) під час спрінту дозволяється внесення змін до вимог, що вже реалізуються на поточному спрінті%good



(\,2\,) порядок виконання запланованої на спрінт роботи визначають розробники%bad



(\,3\,) порядок виконання запланованої на спрінт роботи визначає скрам-мастер%good



(\,4\,) покер планування (Planning Poker)%bad


{\begin {center}\textbf{Завдання 9}\end {center}}\par
Що з наведеного нижче не розглядалось та навіть не згадувалось у курсі?\par
\vspace{2mm}



(\,1\,) PRINCE2%good



(\,2\,) LAZY%good



(\,3\,) DSDM%bad



(\,4\,) Scrum%bad


{\begin {center}\textbf{Завдання 10}\end {center}}\par
Вибрати всі істинні твердження (якщо існують).\par
\vspace{2mm}



(\,1\,) Користувачі бувають основними.%good



(\,2\,) Для тестування код треба ізолювати. Для того, щоб ізолювати, треба мати тести.%good



(\,3\,) Коли нема можливості перевірити всі можливі комбінації, тестують їх попарні (pairwise) комбінації. Це може значно зменшити кількість тестів.%good



(\,4\,) І верифікація, і валідація можуть використовувати тестування.%good



\end {large}}
{\vspace{5mm}\smallskip \begin {small} Затверджено на засіданні кафедри теоретичної кібернетики, протокол \No 6 від   29.01.2021 р.\par\smallskip\smallskip
{\parbox {6.5 cm}{Зав. кафедри \dotfill Ю.В.~Крак}}\hspace {0.7cm}{\hbox to 6.5 cm{ Екзаменатор \dotfill Т.О.~Карнаух}} \end{small}}
\newpage

{{\begin {small}\par
{ \center  Київський національний університет імені Тараса Шевченка \par}
{\parbox {5 cm}{Освітня програма \hfill}{\parbox {7 cm}{Інформатика \hfill}}\parbox {6 cm}{\hfill Семестр 6}}\par
{\parbox {5 cm} {Навчальний предмет \hfill}\parbox {7 cm}{Розробка програмного забезпечення \hfill}\parbox {6cm}{\hfill}\par}\vspace{-7pt} \end{small}}{\center Екзаменаційний білет \No \textbf
{ 42 }\par}}
{
\begin {large}
{\begin {center}\textbf{Завдання 1}\end {center}}\par
Так звані трикутники балансування зв'язують між собою:\par
\vspace{2mm}



(\,1\,) Завжди рівно 4 параметри%bad



(\,2\,) 4 або 3 параметри, при цьому може бути відсутня якість%bad



(\,3\,) 4 або 3 параметри, при цьому може бути відсутній обсяг %good



(\,4\,) 4 або 3 параметри, при цьому може бути відсутній час%bad


{\begin {center}\textbf{Завдання 2}\end {center}}\par
Який з наведених фрагментів коду явно потребує поліпшення (форматування рядків та типи до уваги не брати)?\par
\vspace{2mm}



(\,1\,) void load\_data(T \& arg1, double arg2, double arg3, double arg4); %bad



(\,2\,) \{fstream input\_file, outputFile; …\}%bad



(\,3\,) void help();%good



(\,4\,) void load\_data(…);%good


{\begin {center}\textbf{Завдання 3}\end {center}}\par
Що є характерним для структурного програмування?\par
\vspace{2mm}



(\,1\,) Заборона використання скалярних (простих) типів даних. %bad



(\,2\,) Переходи з середини одного блоку в середину іншого неможливі. %good



(\,3\,) У коді мають використовуватись типи структур, класи використовувати заборонено. %bad



(\,4\,) Використання оператора безумовного переходу. %bad


{\begin {center}\textbf{Завдання 4}\end {center}}\par
Вибрати всі істинні твердження (якщо існують).\par
\vspace{2mm}



(\,1\,) Прогнозна модель розробки може одночасно бути інкрементною. %good



(\,2\,) Прогнозні та адаптивні методології використовують одні ті самі діяльності з розробки ПЗ: визначення та аналіз вимог, проектування, реалізація, тестування, впровадження, супроводження. %good



(\,3\,) Адаптивні методології не передбачають визначення вимог. %bad



(\,4\,) Гнучкі методології не передбачають планування робіт. %bad


{\begin {center}\textbf{Завдання 5}\end {center}}\par
Вибрати всі істинні твердження (якщо існують).\par
\vspace{2mm}



(\,1\,) Правильно спроектована функція повинна мати рівно один обов'язок. %good



(\,2\,) Від довгих функцій можна позбутися, розклавши їх у послідовність викликів дрібніших функцій. %good



(\,3\,) Довгі функції доволі часто приховують складну логіку обробки або повторюваний код. %good



(\,4\,) Код програми пишеться для комп'ютера, тому він має бути синтаксично та семантично коректним. Інше значення не має. %bad

\newpage

{\begin {center}\textbf{Завдання 6}\end {center}}\par
Вибрати всі істинні твердження (якщо існують).\par
\vspace{2mm}



(\,1\,) Код, що себе сам документує, -- це код, який виводить користувачу документацію на себе. %bad



(\,2\,) Без рефакторінгу неперервне проектування неможливе. %good



(\,3\,) Програмування в парах (pair programming) підвищує якість коду та дозволяє економити на виправленні помилок. %good



(\,4\,) За використання гнучких методологій кожен розробник робить, що хоче. %bad


{\begin {center}\textbf{Завдання 7}\end {center}}\par
Вибрати всі істинні твердження (якщо існують).\par
\vspace{2mm}



(\,1\,) Гнучкі методи розробки нехтують проміжними технічними артефактами та документацією через те, що за зміни вимог та/або проекту треба витрачати додатковий час на підтримання їх у відповідному стані. %good



(\,2\,) Гнучкі методології розробки є стійкими до припинення розробки на певний час. %bad



(\,3\,) Гнучкі методології розробки не допускають використання UML. %bad



(\,4\,) За використання гнучких методологій розробки припиняти розробку не рекомендується. %good


{\begin {center}\textbf{Завдання 8}\end {center}}\par
Що з наступного не має відношення до Scrum або не виконується/не використовується за використання оригінального Scrum?\par
\vspace{2mm}



(\,1\,) блочні тести%bad



(\,2\,) виділення в команді одного чи кількох тестувальників%good



(\,3\,) під час спрінту зміни, що впливають на досягнення мети спрінту, зазвичай, не допускаються%bad



(\,4\,) бачення продукту (vision) %bad


{\begin {center}\textbf{Завдання 9}\end {center}}\par
Що з наведеного нижче не розглядалось та навіть не згадувалось у курсі?\par
\vspace{2mm}



(\,1\,) CRAZY%good



(\,2\,) Cleanroom%bad



(\,3\,) LeSS%bad



(\,4\,) FDD%bad


{\begin {center}\textbf{Завдання 10}\end {center}}\par
Вибрати всі істинні твердження (якщо існують).\par
\vspace{2mm}



(\,1\,) Атрибути якості можна використовувати при уточненні вимог.%good



(\,2\,) Можна розглядати якість під час використання.%good



(\,3\,) Забезпечення якості це не тільки тестування.%good



(\,4\,) Розглядають модель якості при використанні.%good



\end {large}}
{\vspace{5mm}\smallskip \begin {small} Затверджено на засіданні кафедри теоретичної кібернетики, протокол \No 6 від   29.01.2021 р.\par\smallskip\smallskip
{\parbox {6.5 cm}{Зав. кафедри \dotfill Ю.В.~Крак}}\hspace {0.7cm}{\hbox to 6.5 cm{ Екзаменатор \dotfill Т.О.~Карнаух}} \end{small}}
\newpage

{{\begin {small}\par
{ \center  Київський національний університет імені Тараса Шевченка \par}
{\parbox {5 cm}{Освітня програма \hfill}{\parbox {7 cm}{Інформатика \hfill}}\parbox {6 cm}{\hfill Семестр 6}}\par
{\parbox {5 cm} {Навчальний предмет \hfill}\parbox {7 cm}{Розробка програмного забезпечення \hfill}\parbox {6cm}{\hfill}\par}\vspace{-7pt} \end{small}}{\center Екзаменаційний білет \No \textbf
{ 43 }\par}}
{
\begin {large}
{\begin {center}\textbf{Завдання 1}\end {center}}\par
Що є метою рефакторингу?\par
\vspace{2mm}



(\,1\,) Зменшення кількості рядків коду%bad



(\,2\,) Усунення помилок у коді%bad



(\,3\,) Додавання нової функціональності%bad



(\,4\,) Поліпшення внутрішньої структури коду%good


{\begin {center}\textbf{Завдання 2}\end {center}}\par
Який з наведених фрагментів коду явно потребує поліпшення (форматування рядків та типи до уваги не брати)?\par
\vspace{2mm}



(\,1\,) \{double max\_power, min\_power;…\} %good



(\,2\,) \{int a\_, a\_\_; …\}%bad



(\,3\,) \{double avg\_temperature;…\} %good



(\,4\,) \{int x; int y; …\}%good


{\begin {center}\textbf{Завдання 3}\end {center}}\par
Що є характерним для структурного програмування?\par
\vspace{2mm}



(\,1\,) Переходи з середини одного блоку в середину іншого неможливі. %good



(\,2\,) Використання оператора безумовного переходу. %bad



(\,3\,) У коді мають використовуватись типи структур, класи використовувати заборонено. %bad



(\,4\,) Заборона використання скалярних (простих) типів даних. %bad


{\begin {center}\textbf{Завдання 4}\end {center}}\par
Вибрати всі істинні твердження (якщо існують).\par
\vspace{2mm}



(\,1\,) Неперервне проектування має бути ітеративним та інкрементним одночасно. %good



(\,2\,) Адаптивність методології проявляється в тому, що за зміни вимог розробка починається заново. %bad



(\,3\,) Для адаптивних методологій прогнозувати результат чи терміни його отримання абсолютно неможливо. %bad



(\,4\,) Гнучкі методології не передбачають планування робіт. %bad


{\begin {center}\textbf{Завдання 5}\end {center}}\par
Вибрати всі істинні твердження (якщо існують).\par
\vspace{2mm}



(\,1\,) Код, що виконує однакову послідовність логічних кроків, у більшості випадків є повторюваним.%good



(\,2\,) Чим вища кваліфікація програміста, тим більш складні для розуміння та довгі умовні вирази він використовує. %bad



(\,3\,) Код з довгими функціями має меншу кількість функцій, а тому його швидше та простіше тестувати. %bad



(\,4\,) Залежності між різними модулями системи мають бути якомога більш сильними, бо інакше вона розпадеться. %bad

\newpage

{\begin {center}\textbf{Завдання 6}\end {center}}\par
Вибрати всі істинні твердження (якщо існують).\par
\vspace{2mm}



(\,1\,) Розробка за гнучкими методологіями цінить людей та взаємодію між ними. %good



(\,2\,) Гнучкі методології базуються на комунікації розробників, а тому передбачають виключно колективне володіння кодом. %bad



(\,3\,) На відміну від Scrum, RUP є компонентно-орієнтованим. %bad



(\,4\,) Водоспадна методологія може мати тенденцію до надлишкового проектування. %good


{\begin {center}\textbf{Завдання 7}\end {center}}\par
Вибрати всі істинні твердження (якщо існують).\par
\vspace{2mm}



(\,1\,) Гнучкі методології розробки намагаються якомога швидше визначити архітектуру системи, для чого на початку завжди виконується початкове проектування системи в цілому, а тільки потім переходять до ітеративного уточнення деталей та кодування. %bad



(\,2\,) Адаптивні методології розробки дозволяють уникнути надлишкового проектування. %good



(\,3\,) Водоспадна методологія не призначена для внесення змін у вимоги до системи. %good



(\,4\,) Scrum передбачає колективне володіння кодом. %good


{\begin {center}\textbf{Завдання 8}\end {center}}\par
Що з наступного не має відношення до Scrum або не виконується/не використовується за використання оригінального Scrum?\par
\vspace{2mm}



(\,1\,) команда визначає, які задачі вона буде виконувати на наступному спрінті%good



(\,2\,) журнал невиконаних задач (backlog) продукту%bad



(\,3\,) скрам-мастер визначає час, необхідний команді для реалізації конкретної задачі%good



(\,4\,) правила визначення готовності (Definition of Done) %bad


{\begin {center}\textbf{Завдання 9}\end {center}}\par
Що з наведеного нижче не розглядалось та навіть не згадувалось у курсі?\par
\vspace{2mm}



(\,1\,) модель хаосу%bad



(\,2\,) poker%bad



(\,3\,) BDD%bad



(\,4\,) Lean%bad


{\begin {center}\textbf{Завдання 10}\end {center}}\par
Вибрати всі істинні твердження (якщо існують).\par
\vspace{2mm}



(\,1\,) Атрибути якості можна використовувати при уточненні вимог.%good



(\,2\,) Розглядають модель якості при використанні.%good



(\,3\,) Можна розглядати якість під час використання.%good



(\,4\,) Якість можна визначати через дефекти.%good



\end {large}}
{\vspace{5mm}\smallskip \begin {small} Затверджено на засіданні кафедри теоретичної кібернетики, протокол \No 6 від   29.01.2021 р.\par\smallskip\smallskip
{\parbox {6.5 cm}{Зав. кафедри \dotfill Ю.В.~Крак}}\hspace {0.7cm}{\hbox to 6.5 cm{ Екзаменатор \dotfill Т.О.~Карнаух}} \end{small}}
\newpage

{{\begin {small}\par
{ \center  Київський національний університет імені Тараса Шевченка \par}
{\parbox {5 cm}{Освітня програма \hfill}{\parbox {7 cm}{Інформатика \hfill}}\parbox {6 cm}{\hfill Семестр 6}}\par
{\parbox {5 cm} {Навчальний предмет \hfill}\parbox {7 cm}{Розробка програмного забезпечення \hfill}\parbox {6cm}{\hfill}\par}\vspace{-7pt} \end{small}}{\center Екзаменаційний білет \No \textbf
{ 44 }\par}}
{
\begin {large}
{\begin {center}\textbf{Завдання 1}\end {center}}\par
Так звані трикутники балансування зв'язують між собою:\par
\vspace{2mm}



(\,1\,) 4 або 3 параметри, при цьому може бути відсутня якість%bad



(\,2\,) Завжди рівно 3 параметри%bad



(\,3\,) 4 або 3 параметри, при цьому може бути відсутній час%bad



(\,4\,) 4 або 3 параметри, при цьому може бути відсутній обсяг %good


{\begin {center}\textbf{Завдання 2}\end {center}}\par
Який з наведених фрагментів коду явно потребує поліпшення (форматування рядків та типи до уваги не брати)?\par
\vspace{2mm}



(\,1\,) void ff(…);%bad



(\,2\,) \{int input\_file, file\_to\_output; …\}%bad



(\,3\,) \{double the\_best\_value\_anyone\_ever\_have;\} %bad



(\,4\,) void f13(…);%bad


{\begin {center}\textbf{Завдання 3}\end {center}}\par
Що є характерним для структурного програмування?\par
\vspace{2mm}



(\,1\,) Використання механізму винятків. %bad



(\,2\,) У коді мають використовуватись типи даних, що є класами. %bad



(\,3\,) Кожен блок коду повинен мати єдиний вхід та єдиний вихід.%good



(\,4\,) У коді мають використовуватись структурні типи даних. %bad


{\begin {center}\textbf{Завдання 4}\end {center}}\par
Вибрати всі істинні твердження (якщо існують).\par
\vspace{2mm}



(\,1\,) Для адаптивних методологій розробки, на відміну від прогнозних, попередні оцінки бюджету та/або часу не виконують, бо вони все одно зміняться. %bad



(\,2\,) Адаптивна модель розробки адаптує вимоги до продукту до наявної команди розробників. %bad



(\,3\,) Прогнозна модель розробки може одночасно бути інкрементною. %good



(\,4\,) Прогнозна модель розробки не може одночасно бути інкрементною. %bad


{\begin {center}\textbf{Завдання 5}\end {center}}\par
Вибрати всі істинні твердження (якщо існують).\par
\vspace{2mm}



(\,1\,) Складні умови виникають в реальних задачах доволі часто, а тому є доволі природніми і не можуть вважатись недоліком коду. %bad



(\,2\,) Одним з прийомів спрощення складних умов є їх декомпозиція. %good



(\,3\,) Гарний код не повинен містити різних фрагментів, що складаються з кількох інструкцій та виконують однакове перетворення даних. %good



(\,4\,) Одним з прийомів спрощення умовної логіки є використання поліморфізму. %good

\newpage

{\begin {center}\textbf{Завдання 6}\end {center}}\par
Вибрати всі істинні твердження (якщо існують).\par
\vspace{2mm}



(\,1\,) Розповіді користувача (user stories) – це прогресивний інструмент гнучких методологій, який нічого спільного не має з прецедентами (use case). %bad



(\,2\,) Поступове уточнення вимог замовника полегшує адаптацію до змін. %good



(\,3\,) За програмування в парах (pair programming) поки один колега набирає код на комп'ютері інший може піти пограти в настільний теніс. %bad



(\,4\,) За невизначених або піддатливих змінам вимог краще застосовувати гнучкі методології розробки. %good


{\begin {center}\textbf{Завдання 7}\end {center}}\par
Вибрати всі істинні твердження (якщо існують).\par
\vspace{2mm}



(\,1\,) Використання RUP не заперечує часті випуски продукту, але й не наполягає на них. %good



(\,2\,) Програмування в парах (pair programming) зводиться до того, що кожним фрагментом коду володіють два програмісти. %bad



(\,3\,) RUP є компонентно-орієнтованим. %good



(\,4\,) Висока внутрішня якість коду є марною тратою часу на те, що замовник не побачить та/або не зможе зацінити. %bad


{\begin {center}\textbf{Завдання 8}\end {center}}\par
Що з наступного не має відношення до Scrum або не виконується/не використовується за використання оригінального Scrum?\par
\vspace{2mm}



(\,1\,) власник продукту призначає час, необхідний команді для реалізації конкретної задачі%good



(\,2\,) власник продукту остаточно затверджує проект системи%good



(\,3\,) скрам-мастер призначає відповідальних за конкретні задачі%good



(\,4\,) журнал невиконаних задач (backlog) спрінту%bad


{\begin {center}\textbf{Завдання 9}\end {center}}\par
Що з наведеного нижче не розглядалось та навіть не згадувалось у курсі?\par
\vspace{2mm}



(\,1\,) XP%bad



(\,2\,) SAFe%bad



(\,3\,) спіральна модель%bad



(\,4\,) X-модель%good


{\begin {center}\textbf{Завдання 10}\end {center}}\par
Вибрати всі істинні твердження (якщо існують).\par
\vspace{2mm}



(\,1\,) Розглядають модель якості даних.%good



(\,2\,) Коли нема можливості перевірити всі можливі комбінації, тестують їх попарні (pairwise) комбінації. Це може значно зменшити кількість тестів.%good



(\,3\,) Тестовий приклад складається з набору вхідних даних, умов виконання та очікуваних результатів.%good



(\,4\,) І верифікація, і валідація можуть використовувати тестування.%good



\end {large}}
{\vspace{5mm}\smallskip \begin {small} Затверджено на засіданні кафедри теоретичної кібернетики, протокол \No 6 від   29.01.2021 р.\par\smallskip\smallskip
{\parbox {6.5 cm}{Зав. кафедри \dotfill Ю.В.~Крак}}\hspace {0.7cm}{\hbox to 6.5 cm{ Екзаменатор \dotfill Т.О.~Карнаух}} \end{small}}
\newpage

{{\begin {small}\par
{ \center  Київський національний університет імені Тараса Шевченка \par}
{\parbox {5 cm}{Освітня програма \hfill}{\parbox {7 cm}{Інформатика \hfill}}\parbox {6 cm}{\hfill Семестр 6}}\par
{\parbox {5 cm} {Навчальний предмет \hfill}\parbox {7 cm}{Розробка програмного забезпечення \hfill}\parbox {6cm}{\hfill}\par}\vspace{-7pt} \end{small}}{\center Екзаменаційний білет \No \textbf
{ 45 }\par}}
{
\begin {large}
{\begin {center}\textbf{Завдання 1}\end {center}}\par
Так звані трикутники балансування зв'язують між собою:\par
\vspace{2mm}



(\,1\,) 4 або 3 параметри, при цьому може бути відсутній обсяг %good



(\,2\,) Завжди рівно 4 параметри%bad



(\,3\,) Завжди рівно 3 параметри%bad



(\,4\,) 4 або 3 параметри, при цьому може бути відсутня якість%bad


{\begin {center}\textbf{Завдання 2}\end {center}}\par
Який з наведених фрагментів коду явно потребує поліпшення (форматування рядків та типи до уваги не брати)?\par
\vspace{2mm}



(\,1\,) \{int a\_, a\_\_; …\}%bad



(\,2\,) \{double max\_power, min\_power;…\} %good



(\,3\,) \{fstream input\_file, outputFile; …\}%bad



(\,4\,) \{double the\_best\_value\_anyone\_ever\_have;\} %bad


{\begin {center}\textbf{Завдання 3}\end {center}}\par
Що є характерним для структурного програмування?\par
\vspace{2mm}



(\,1\,) У коді мають використовуватись типи структур. %bad



(\,2\,) Заборона використання скалярних (простих) типів даних. %bad



(\,3\,) Програма зображується ієрархічною структурою блоків.%good



(\,4\,) Використання оператора безумовного переходу. %bad


{\begin {center}\textbf{Завдання 4}\end {center}}\par
Вибрати всі істинні твердження (якщо існують).\par
\vspace{2mm}



(\,1\,) Неперервне проектування здійснює створення та модифікацію проекту системи в міру розробки. %good



(\,2\,) Адаптивні методології не передбачають визначення вимог. %bad



(\,3\,) Для адаптивних моделей розробки оцінки часу та бюджету не виконують, бо вимоги є змінними. %bad



(\,4\,) Неперервність проектування полягає в тім, що процес розробки не можна переривати більше ніж на вихідні. %bad


{\begin {center}\textbf{Завдання 5}\end {center}}\par
Вибрати всі істинні твердження (якщо існують).\par
\vspace{2mm}



(\,1\,) Чим менше класів у коді, тим менше класів треба налагоджувати, а тому класи за можливості слід об'єднувати. %bad



(\,2\,) Чим вища кваліфікація програміста, тим більш складні для розуміння та довгі умовні вирази він використовує. %bad



(\,3\,) Одним з засобів позбутись повторень у коді є використання успадкування. %good



(\,4\,) Одним з засобів позбутись повторень у коді є використання шаблону Template Method (але цей шаблон здатен боротися і з іншими недоліками коду). %good

\newpage

{\begin {center}\textbf{Завдання 6}\end {center}}\par
Вибрати всі істинні твердження (якщо існують).\par
\vspace{2mm}



(\,1\,) Якщо виникають проблеми з модифікацією коду, то слід відразу змінити команду розробників на більш кваліфіковану.%bad



(\,2\,) Одна з принципових відмінностей RUP від Scrum полягає у використанні численних проміжних технічних артефактів. %good



(\,3\,) Висока внутрішня якість коду здешевшує розробку, бо дозволяє економити на вартості та часі майбутніх змін. %good



(\,4\,) Розробка за гнучкими методологіями цінить співробітництво з замовником та піддатлівість змінам. %good


{\begin {center}\textbf{Завдання 7}\end {center}}\par
Вибрати всі істинні твердження (якщо існують).\par
\vspace{2mm}



(\,1\,) Сьогодні водоспадна методологія не має права на існування. %bad



(\,2\,) RUP є ітеративним та інкрементним.%good



(\,3\,) Методологія розробки Big Design Up Front погана тим, що може призводити до надлишкового проектування. %good



(\,4\,) Гнучкі методології розробки допускають суттєві зміни вимог, проте все одно намагаються якомога швидше зафіксувати архітектуру. %good


{\begin {center}\textbf{Завдання 8}\end {center}}\par
Що з наступного не має відношення до Scrum або не виконується/не використовується за використання оригінального Scrum?\par
\vspace{2mm}



(\,1\,) порядок виконання запланованої на спрінт роботи визначає власник продукту%good



(\,2\,) розповіді користувача%bad



(\,3\,) розповіді користувача%bad



(\,4\,) виділення в команді одного чи кількох проектувальників%good


{\begin {center}\textbf{Завдання 9}\end {center}}\par
Що з наведеного нижче не розглядалось та навіть не згадувалось у курсі?\par
\vspace{2mm}



(\,1\,) LOAN%good



(\,2\,) sashimi%bad



(\,3\,) Scrum of scrums%bad



(\,4\,) BDUF%bad


{\begin {center}\textbf{Завдання 10}\end {center}}\par
Вибрати всі істинні твердження (якщо існують).\par
\vspace{2mm}



(\,1\,) Тестування вимог може породжувати їх уточнення.%good



(\,2\,) Якість буває зовнішньою.%good



(\,3\,) Якість буває внутрішньою.%good



(\,4\,) Забезпечення якості це не тільки тестування.%good



\end {large}}
{\vspace{5mm}\smallskip \begin {small} Затверджено на засіданні кафедри теоретичної кібернетики, протокол \No 6 від   29.01.2021 р.\par\smallskip\smallskip
{\parbox {6.5 cm}{Зав. кафедри \dotfill Ю.В.~Крак}}\hspace {0.7cm}{\hbox to 6.5 cm{ Екзаменатор \dotfill Т.О.~Карнаух}} \end{small}}
\newpage

{{\begin {small}\par
{ \center  Київський національний університет імені Тараса Шевченка \par}
{\parbox {5 cm}{Освітня програма \hfill}{\parbox {7 cm}{Інформатика \hfill}}\parbox {6 cm}{\hfill Семестр 6}}\par
{\parbox {5 cm} {Навчальний предмет \hfill}\parbox {7 cm}{Розробка програмного забезпечення \hfill}\parbox {6cm}{\hfill}\par}\vspace{-7pt} \end{small}}{\center Екзаменаційний білет \No \textbf
{ 46 }\par}}
{
\begin {large}
{\begin {center}\textbf{Завдання 1}\end {center}}\par
Так звані трикутники балансування зв'язують між собою:\par
\vspace{2mm}



(\,1\,) Завжди рівно 4 параметри%bad



(\,2\,) 4 або 3 параметри, при цьому може бути відсутній час%bad



(\,3\,) 4 або 3 параметри, при цьому може бути відсутній обсяг %good



(\,4\,) Завжди рівно 3 параметри%bad


{\begin {center}\textbf{Завдання 2}\end {center}}\par
Який з наведених фрагментів коду явно потребує поліпшення (форматування рядків та типи до уваги не брати)?\par
\vspace{2mm}



(\,1\,) void load\_data(T \& arg1, double arg2, double arg3, double arg4); %bad



(\,2\,) void ff(…);%bad



(\,3\,) \{fstream inputFile, outputFile;…\} %good



(\,4\,) \{int xtrmprg; …\}%bad


{\begin {center}\textbf{Завдання 3}\end {center}}\par
Що є характерним для структурного програмування?\par
\vspace{2mm}



(\,1\,) У коді мають використовуватись типи структур. %bad



(\,2\,) У коді мають використовуватись структурні типи даних. %bad



(\,3\,) Кожен блок коду повинен мати єдиний вхід та єдиний вихід.%good



(\,4\,) У коді мають використовуватись типи даних, що є класами. %bad


{\begin {center}\textbf{Завдання 4}\end {center}}\par
Вибрати всі істинні твердження (якщо існують).\par
\vspace{2mm}



(\,1\,) Адаптивні методології має сенс використовувати для проектів, де можлива зміна вимог. %good



(\,2\,) У прогнозній моделі розробки розробка може виконуватись ітеративно. %good



(\,3\,) Адаптивні моделі розробки не допускають зміну вимог до продукту, але продукт може розроблятись поетапно. %bad



(\,4\,) І прогнозні, і адаптивні методології передбачають планування виконуваних робіт. %good


{\begin {center}\textbf{Завдання 5}\end {center}}\par
Вибрати всі істинні твердження (якщо існують).\par
\vspace{2mm}



(\,1\,) Код з довгими функціями має меншу кількість функцій, а тому його швидше та простіше тестувати. %bad



(\,2\,) Довгі функції доволі часто приховують складну логіку обробки або повторюваний код. %good



(\,3\,) Довжина функцій не впливає на внутрішню якість коду. %bad



(\,4\,) Повторюваний код дозволяє програмісту краще зрозуміти його логіку та вибрати з кількох варіантів найкращий. %bad

\newpage

{\begin {center}\textbf{Завдання 6}\end {center}}\par
Вибрати всі істинні твердження (якщо існують).\par
\vspace{2mm}



(\,1\,) RUP є архітектурно-орієнтованим. %good



(\,2\,) Модель хаосу є методологією розробки програмного забезпечення. %bad



(\,3\,) RUP є прогнозною методологією. %bad



(\,4\,) За застосування адаптивних методологій розробки, так само як і за прогнозних, намагаються якомога швидше визначити архітектуру системи. Відмінність полягає в тім, у який спосіб це виконується.%good


{\begin {center}\textbf{Завдання 7}\end {center}}\par
Вибрати всі істинні твердження (якщо існують).\par
\vspace{2mm}



(\,1\,) Розповіді користувача (user stories) це менш формалізований порівняно з прецедентами спосіб опису вимог до системи. %good



(\,2\,) Програмування в парах (pair programming) спрощує введення в команду нових розробників. %good



(\,3\,) Ітеративна та інкрементна розробка передбачає наявність у тому чи іншому вигляді етапу ініціалізації, ітерацій та керуючої таблиці проекту. %good



(\,4\,) RUP є прогнозною методологією. %bad


{\begin {center}\textbf{Завдання 8}\end {center}}\par
Що з наступного не має відношення до Scrum або не виконується/не використовується за використання оригінального Scrum?\par
\vspace{2mm}



(\,1\,) залучення зовнішніх фахівців для побудови архітектури системи%good



(\,2\,) покер планування (Planning Poker)%bad



(\,3\,) скрам-мастер остаточно затверджує проект системи%good



(\,4\,) скрам-мастер визначає, які задачі команда буде виконувати на наступному спрінті%good


{\begin {center}\textbf{Завдання 9}\end {center}}\par
Що з наведеного нижче не розглядалось та навіть не згадувалось у курсі?\par
\vspace{2mm}



(\,1\,) Scrum of scrums%bad



(\,2\,) Cleanroom%bad



(\,3\,) LOAN%good



(\,4\,) poker%bad


{\begin {center}\textbf{Завдання 10}\end {center}}\par
Вибрати всі істинні твердження (якщо існують).\par
\vspace{2mm}



(\,1\,) Модель якості продукту є ієрархічною.%good



(\,2\,) Для тестування код треба ізолювати. Для того, щоб ізолювати, треба мати тести.%good



(\,3\,) За провалених димових тестів решта тестів не виконується і код повертається на доробку.%good



(\,4\,) Відсутність знайдених дефектів на початкових стадіях розробки може бути свідченням поганого покриття коду тестами.%good



\end {large}}
{\vspace{5mm}\smallskip \begin {small} Затверджено на засіданні кафедри теоретичної кібернетики, протокол \No 6 від   29.01.2021 р.\par\smallskip\smallskip
{\parbox {6.5 cm}{Зав. кафедри \dotfill Ю.В.~Крак}}\hspace {0.7cm}{\hbox to 6.5 cm{ Екзаменатор \dotfill Т.О.~Карнаух}} \end{small}}
\newpage

{{\begin {small}\par
{ \center  Київський національний університет імені Тараса Шевченка \par}
{\parbox {5 cm}{Освітня програма \hfill}{\parbox {7 cm}{Інформатика \hfill}}\parbox {6 cm}{\hfill Семестр 6}}\par
{\parbox {5 cm} {Навчальний предмет \hfill}\parbox {7 cm}{Розробка програмного забезпечення \hfill}\parbox {6cm}{\hfill}\par}\vspace{-7pt} \end{small}}{\center Екзаменаційний білет \No \textbf
{ 47 }\par}}
{
\begin {large}
{\begin {center}\textbf{Завдання 1}\end {center}}\par
Що з наведеного найкраще характеризує рефакторінг?\par
\vspace{2mm}



(\,1\,) Зміна внутрішньої структури коду під час усунення помилок%bad



(\,2\,) Зміна внутрішньої структури коду, що не змінює його видиму поведінку%good



(\,3\,) Довільна переробка коду, що виникає у процесі розробки%bad



(\,4\,) Докорінна переробка всього коду з нуля,  що не змінює його видиму поведінку%bad


{\begin {center}\textbf{Завдання 2}\end {center}}\par
Який з наведених фрагментів коду явно потребує поліпшення (форматування рядків та типи до уваги не брати)?\par
\vspace{2mm}



(\,1\,) \{double maxPower, minPower;…\} %good



(\,2\,) void load\_data(…);%good



(\,3\,) void help();%good



(\,4\,) \{double avg\_temperature;…\} %good


{\begin {center}\textbf{Завдання 3}\end {center}}\par
Що є характерним для структурного програмування?\par
\vspace{2mm}



(\,1\,) Програма зображується ієрархічною структурою блоків.%good



(\,2\,) У коді мають використовуватись типи структур, класи використовувати заборонено. %bad



(\,3\,) Використання механізму винятків. %bad



(\,4\,) У коді мають використовуватись типи даних, що є класами. %bad


{\begin {center}\textbf{Завдання 4}\end {center}}\par
Вибрати всі істинні твердження (якщо існують).\par
\vspace{2mm}



(\,1\,) У випадку, коли вимоги до продукту відомі заздалегідь та є фіксованими контрактом, адаптивні методології розробки використовувати не можна. %bad



(\,2\,) Прогнозні моделі розробки з самого початку фіксують вимоги до продукту. %good



(\,3\,) У прогнозній моделі розробки розробка не може виконуватись ітеративно. %bad



(\,4\,) Прогнозні та адаптивні методології використовують одні ті самі діяльності з розробки ПЗ: визначення та аналіз вимог, проектування, реалізація, тестування, впровадження, супроводження. %good


{\begin {center}\textbf{Завдання 5}\end {center}}\par
Вибрати всі істинні твердження (якщо існують).\par
\vspace{2mm}



(\,1\,) Гарний код має читатись як розмовна мова. %good



(\,2\,) Від довгих функцій можна позбутися, розклавши їх у послідовність викликів дрібніших функцій. %good



(\,3\,) Зв'язки між окремими модулями системи мають бути якомога більш слабкими. %good



(\,4\,) Одним з прийомів спрощення складних умов є їх декомпозиція. %good

\newpage

{\begin {center}\textbf{Завдання 6}\end {center}}\par
Вибрати всі істинні твердження (якщо існують).\par
\vspace{2mm}



(\,1\,) Розробка за гнучкими методологіями ніколи не слідує плану. %bad



(\,2\,) Код, що себе сам документує, -- це код, який виводить користувачу документацію на себе. %bad



(\,3\,) Висока внутрішня якість коду здешевшує розробку, бо дозволяє економити на вартості та часі майбутніх змін. %good



(\,4\,) Гнучкі методи розробки нехтують проміжними технічними артефактами та документацією виключно через їх непотрібність. %bad


{\begin {center}\textbf{Завдання 7}\end {center}}\par
Вибрати всі істинні твердження (якщо існують).\par
\vspace{2mm}



(\,1\,) Водоспадна методологія розробки може призводити до недостатнього проектування. %bad



(\,2\,) Методологія розробки Big Design Up Front погана тим, що може призводити до надлишкового проектування. %good



(\,3\,) Висока внутрішня якість коду є марною тратою часу на те, що замовник не побачить та/або не зможе зацінити. %bad



(\,4\,) Гнучкі методології передбачають часті випуски продукту. %good


{\begin {center}\textbf{Завдання 8}\end {center}}\par
Що з наступного не має відношення до Scrum або не виконується/не використовується за використання оригінального Scrum?\par
\vspace{2mm}



(\,1\,) скрам-мастер визначає пріоритети задач%good



(\,2\,) власник продукту призначає відповідальних за конкретні задачі%good



(\,3\,) побудова різнопланових моделей системи%good



(\,4\,) порядок виконання запланованої на спрінт роботи визначає скрам-мастер%good


{\begin {center}\textbf{Завдання 9}\end {center}}\par
Що з наведеного нижче не розглядалось та навіть не згадувалось у курсі?\par
\vspace{2mm}



(\,1\,) FDD%bad



(\,2\,) Lean%bad



(\,3\,) DSDM%bad



(\,4\,) SAFe%bad


{\begin {center}\textbf{Завдання 10}\end {center}}\par
Вибрати всі істинні твердження (якщо існують).\par
\vspace{2mm}



(\,1\,) Модель якості даних є ієрархічною.%good



(\,2\,) Якість можна визначати через атрибути якості.%good



(\,3\,) Користувачі бувають вторинними.%good



(\,4\,) Користувачі бувають непрямими.%good



\end {large}}
{\vspace{5mm}\smallskip \begin {small} Затверджено на засіданні кафедри теоретичної кібернетики, протокол \No 6 від   29.01.2021 р.\par\smallskip\smallskip
{\parbox {6.5 cm}{Зав. кафедри \dotfill Ю.В.~Крак}}\hspace {0.7cm}{\hbox to 6.5 cm{ Екзаменатор \dotfill Т.О.~Карнаух}} \end{small}}
\newpage

{{\begin {small}\par
{ \center  Київський національний університет імені Тараса Шевченка \par}
{\parbox {5 cm}{Освітня програма \hfill}{\parbox {7 cm}{Інформатика \hfill}}\parbox {6 cm}{\hfill Семестр 6}}\par
{\parbox {5 cm} {Навчальний предмет \hfill}\parbox {7 cm}{Розробка програмного забезпечення \hfill}\parbox {6cm}{\hfill}\par}\vspace{-7pt} \end{small}}{\center Екзаменаційний білет \No \textbf
{ 48 }\par}}
{
\begin {large}
{\begin {center}\textbf{Завдання 1}\end {center}}\par
Так звані трикутники балансування зв'язують між собою:\par
\vspace{2mm}



(\,1\,) 4 або 3 параметри, при цьому може бути відсутня якість%bad



(\,2\,) Завжди рівно 4 параметри%bad



(\,3\,) 4 або 3 параметри, при цьому може бути відсутній обсяг %good



(\,4\,) 4 або 3 параметри, при цьому може бути відсутній час%bad


{\begin {center}\textbf{Завдання 2}\end {center}}\par
Який з наведених фрагментів коду явно потребує поліпшення (форматування рядків та типи до уваги не брати)?\par
\vspace{2mm}



(\,1\,) \{fstream input\_file, output\_file;…\} %good



(\,2\,) void f13(…);%bad



(\,3\,) \{int input\_file, file\_to\_output; …\}%bad



(\,4\,) \{int x; int y; …\}%good


{\begin {center}\textbf{Завдання 3}\end {center}}\par
Що є характерним для структурного програмування?\par
\vspace{2mm}



(\,1\,) Використання механізму винятків. %bad



(\,2\,) У коді мають використовуватись типи структур. %bad



(\,3\,) Використання оператора безумовного переходу. %bad



(\,4\,) Переходи з середини одного блоку в середину іншого неможливі. %good


{\begin {center}\textbf{Завдання 4}\end {center}}\par
Вибрати всі істинні твердження (якщо існують).\par
\vspace{2mm}



(\,1\,) В адаптивних методологіях планування робіт не виконується. %bad



(\,2\,) Ітеративні та інкрементні методи розробки є синонімом неперервного проектування. %bad



(\,3\,) Адаптивні моделі розробки допускають зміну вимог до продукту. %good



(\,4\,) Неперервне проектування створює специфікацію проекту перед початком розробки/ітерації, а під час ітерації його неперервно модифікує.%bad


{\begin {center}\textbf{Завдання 5}\end {center}}\par
Вибрати всі істинні твердження (якщо існують).\par
\vspace{2mm}



(\,1\,) Втілення в коді різних підходів до одного того самого перетворення даних дозволяє програмісту продемонструвати свою професійну підготовку з найкращого боку. %bad



(\,2\,) Код з довгими функціями має меншу кількість функцій, а тому для нього треба менше тестів. %bad



(\,3\,) Правильно спроектована функція повинна мати рівно один обов'язок. %good



(\,4\,) Одним з прийомів спрощення складної умовної логіки є використання шаблону Strategy. %good

\newpage

{\begin {center}\textbf{Завдання 6}\end {center}}\par
Вибрати всі істинні твердження (якщо існують).\par
\vspace{2mm}



(\,1\,) Ітеративна методологія розробки у загальному випадку не гарантує, що не буде надлишкового проектування. %good



(\,2\,) Гнучкі методології розробки не допускають використання UML. %bad



(\,3\,) Простота – мистецтво максимізації незробленої роботи. %good



(\,4\,) Розповіді користувача (user stories) – це прогресивний інструмент гнучких методологій, який нічого спільного не має з прецедентами (use case). %bad


{\begin {center}\textbf{Завдання 7}\end {center}}\par
Вибрати всі істинні твердження (якщо існують).\par
\vspace{2mm}



(\,1\,) Без рефакторінгу неперервне проектування неможливе. %good



(\,2\,) Водоспадна методологія може мати тенденцію до надлишкового проектування. %good



(\,3\,) Якщо виникають проблеми з модифікацією коду, то слід відразу змінити команду розробників на більш кваліфіковану.%bad



(\,4\,) Водоспадна методологія розробки може призводити до надлишкового проектування. %good


{\begin {center}\textbf{Завдання 8}\end {center}}\par
Що з наступного не має відношення до Scrum або не виконується/не використовується за використання оригінального Scrum?\par
\vspace{2mm}



(\,1\,) ретроспектива спрінту%bad



(\,2\,) під час спрінту дозволяється внесення змін до вимог, що вже реалізуються на поточному спрінті%good



(\,3\,) порядок виконання запланованої на спрінт роботи визначають розробники%bad



(\,4\,) обернені вимоги%bad


{\begin {center}\textbf{Завдання 9}\end {center}}\par
Що з наведеного нижче не розглядалось та навіть не згадувалось у курсі?\par
\vspace{2mm}



(\,1\,) sashimi%bad



(\,2\,) BDD%bad



(\,3\,) CRAZY%good



(\,4\,) водоспадна модель%bad


{\begin {center}\textbf{Завдання 10}\end {center}}\par
Вибрати всі істинні твердження (якщо існують).\par
\vspace{2mm}



(\,1\,) Модель якості при використанні є ієрархічною.%good



(\,2\,) Розглядають модель якості продукту.%good



(\,3\,) Тестування буває позитивне та негативне.%good



(\,4\,) Користувачі бувають основними.%good



\end {large}}
{\vspace{5mm}\smallskip \begin {small} Затверджено на засіданні кафедри теоретичної кібернетики, протокол \No 6 від   29.01.2021 р.\par\smallskip\smallskip
{\parbox {6.5 cm}{Зав. кафедри \dotfill Ю.В.~Крак}}\hspace {0.7cm}{\hbox to 6.5 cm{ Екзаменатор \dotfill Т.О.~Карнаух}} \end{small}}
\newpage

{{\begin {small}\par
{ \center  Київський національний університет імені Тараса Шевченка \par}
{\parbox {5 cm}{Освітня програма \hfill}{\parbox {7 cm}{Інформатика \hfill}}\parbox {6 cm}{\hfill Семестр 6}}\par
{\parbox {5 cm} {Навчальний предмет \hfill}\parbox {7 cm}{Розробка програмного забезпечення \hfill}\parbox {6cm}{\hfill}\par}\vspace{-7pt} \end{small}}{\center Екзаменаційний білет \No \textbf
{ 49 }\par}}
{
\begin {large}
{\begin {center}\textbf{Завдання 1}\end {center}}\par
Шаблон проектування:\par
\vspace{2mm}



(\,1\,) Формулює ідею розв'язання проблеми та може пропонувати одну чи кілька її реалізацій%good



(\,2\,) Тільки формулює ідею розв'язання проблеми%bad



(\,3\,) Задає розв'язання типової проблеми до найдрібніших подробиць%bad



(\,4\,) Є шаблоном, в який треба підставити власні ідентифікатори%bad


{\begin {center}\textbf{Завдання 2}\end {center}}\par
Який з наведених фрагментів коду явно потребує поліпшення (форматування рядків та типи до уваги не брати)?\par
\vspace{2mm}



(\,1\,) \{fstream input\_file, output\_file;…\} %good



(\,2\,) void f13(…);%bad



(\,3\,) \{fstream input\_file, outputFile; …\}%bad



(\,4\,) \{fstream inputFile, outputFile;…\} %good


{\begin {center}\textbf{Завдання 3}\end {center}}\par
Що є характерним для структурного програмування?\par
\vspace{2mm}



(\,1\,) Заборона використання скалярних (простих) типів даних. %bad



(\,2\,) Програма зображується ієрархічною структурою блоків.%good



(\,3\,) У коді мають використовуватись структурні типи даних. %bad



(\,4\,) У коді мають використовуватись типи структур, класи використовувати заборонено. %bad


{\begin {center}\textbf{Завдання 4}\end {center}}\par
Вибрати всі істинні твердження (якщо існують).\par
\vspace{2mm}



(\,1\,) Прогнозні та адаптивні методології використовують одні ті самі діяльності з розробки ПЗ: визначення та аналіз вимог, проектування, реалізація, тестування, впровадження, супроводження. %good



(\,2\,) У прогнозній моделі розробки розробка може виконуватись ітеративно. %good



(\,3\,) Для адаптивних методологій розробки, на відміну від прогнозних, попередні оцінки бюджету та/або часу не виконують, бо вони все одно зміняться. %bad



(\,4\,) Гнучкі методології не передбачають планування робіт. %bad


{\begin {center}\textbf{Завдання 5}\end {center}}\par
Вибрати всі істинні твердження (якщо існують).\par
\vspace{2mm}



(\,1\,) Складні умови виникають в реальних задачах доволі часто, а тому є доволі природніми і не можуть вважатись недоліком коду. %bad



(\,2\,) Код програми пишеться для комп'ютера, тому він має бути синтаксично та семантично коректним. Інше значення не має. %bad



(\,3\,) Залежності між різними модулями системи мають бути якомога більш сильними, бо інакше вона розпадеться. %bad



(\,4\,) Якщо кожна функція бібліотеки буде здатна в залежності від параметрів виконувати кілька різних дій, то загальна кількість функцій у ній буде менша і її буде простіше використовувати. %bad

\newpage

{\begin {center}\textbf{Завдання 6}\end {center}}\par
Вибрати всі істинні твердження (якщо існують).\par
\vspace{2mm}



(\,1\,) Розробка за гнучкими методологіями цінить співробітництво з замовником та піддатлівість змінам. %good



(\,2\,) Адаптивні методології розробки дозволяють уникнути надлишкового проектування. %good



(\,3\,) Сьогодні водоспадна методологія не має права на існування. %bad



(\,4\,) Модель хаосу є методологією розробки програмного забезпечення. %bad


{\begin {center}\textbf{Завдання 7}\end {center}}\par
Вибрати всі істинні твердження (якщо існують).\par
\vspace{2mm}



(\,1\,) Багато хто з тих, що використовує Scrum, складає архітектурну модель розроблюваної системи. %good



(\,2\,) Програмування в парах (pair programming) спрощує введення в команду нових розробників. %good



(\,3\,) Перша версія коду не завжди може бути якісною -- і це нормально. %good



(\,4\,) Водоспадна методологія не призначена для внесення змін у вимоги до системи. %good


{\begin {center}\textbf{Завдання 8}\end {center}}\par
Що з наступного не має відношення до Scrum або не виконується/не використовується за використання оригінального Scrum?\par
\vspace{2mm}



(\,1\,) блочні тести%bad



(\,2\,) канбан%bad



(\,3\,) планування спрінту%bad



(\,4\,) щоденний скрам%bad


{\begin {center}\textbf{Завдання 9}\end {center}}\par
Що з наведеного нижче не розглядалось та навіть не згадувалось у курсі?\par
\vspace{2mm}



(\,1\,) TDD%bad



(\,2\,) спіральна модель%bad



(\,3\,) LAZY%good



(\,4\,) X-модель%good


{\begin {center}\textbf{Завдання 10}\end {center}}\par
Вибрати всі істинні твердження (якщо існують).\par
\vspace{2mm}



(\,1\,) Якість можна визначати через дефекти.%good



(\,2\,) Якість буває внутрішньою.%good



(\,3\,) Користувачі бувають вторинними.%good



(\,4\,) Розглядають модель якості даних.%good



\end {large}}
{\vspace{5mm}\smallskip \begin {small} Затверджено на засіданні кафедри теоретичної кібернетики, протокол \No 6 від   29.01.2021 р.\par\smallskip\smallskip
{\parbox {6.5 cm}{Зав. кафедри \dotfill Ю.В.~Крак}}\hspace {0.7cm}{\hbox to 6.5 cm{ Екзаменатор \dotfill Т.О.~Карнаух}} \end{small}}
\newpage

{{\begin {small}\par
{ \center  Київський національний університет імені Тараса Шевченка \par}
{\parbox {5 cm}{Освітня програма \hfill}{\parbox {7 cm}{Інформатика \hfill}}\parbox {6 cm}{\hfill Семестр 6}}\par
{\parbox {5 cm} {Навчальний предмет \hfill}\parbox {7 cm}{Розробка програмного забезпечення \hfill}\parbox {6cm}{\hfill}\par}\vspace{-7pt} \end{small}}{\center Екзаменаційний білет \No \textbf
{ 50 }\par}}
{
\begin {large}
{\begin {center}\textbf{Завдання 1}\end {center}}\par
Що є метою рефакторингу?\par
\vspace{2mm}



(\,1\,) Усунення помилок у коді%bad



(\,2\,) Додавання нової функціональності%bad



(\,3\,) Зменшення кількості рядків коду%bad



(\,4\,) Поліпшення внутрішньої структури коду%good


{\begin {center}\textbf{Завдання 2}\end {center}}\par
Який з наведених фрагментів коду явно потребує поліпшення (форматування рядків та типи до уваги не брати)?\par
\vspace{2mm}



(\,1\,) \{double the\_best\_value\_anyone\_ever\_have;\} %bad



(\,2\,) \{int xtrmprg; …\}%bad



(\,3\,) void load\_data(…);%good



(\,4\,) \{double maxPower, minPower;…\} %good


{\begin {center}\textbf{Завдання 3}\end {center}}\par
Що є характерним для структурного програмування?\par
\vspace{2mm}



(\,1\,) У коді мають використовуватись типи структур, класи використовувати заборонено. %bad



(\,2\,) Заборона використання скалярних (простих) типів даних. %bad



(\,3\,) Переходи з середини одного блоку в середину іншого неможливі. %good



(\,4\,) У коді мають використовуватись типи даних, що є класами. %bad


{\begin {center}\textbf{Завдання 4}\end {center}}\par
Вибрати всі істинні твердження (якщо існують).\par
\vspace{2mm}



(\,1\,) Для адаптивних моделей розробки оцінки часу та бюджету не виконують, бо вимоги є змінними. %bad



(\,2\,) Неперервне проектування створює специфікацію проекту перед початком розробки/ітерації, а під час ітерації його неперервно модифікує.%bad



(\,3\,) Для адаптивних методологій прогнозувати результат чи терміни його отримання абсолютно неможливо. %bad



(\,4\,) Адаптивні методології має сенс використовувати для проектів, де можлива зміна вимог. %good


{\begin {center}\textbf{Завдання 5}\end {center}}\par
Вибрати всі істинні твердження (якщо існують).\par
\vspace{2mm}



(\,1\,) Залежності всередині модуля можуть бути сильними. %good



(\,2\,) Код, що виконує однакову послідовність логічних кроків, у більшості випадків є повторюваним.%good



(\,3\,) Одним з засобів позбутись повторень у коді є використання функцій. %good



(\,4\,) Складні комбінації умов починають керувати логікою виконання, тому що такими є вимоги до функціонування ПЗ. Оскільки ПЗ має якнайкраще забезпечувати потреби замовника, то складні комбінації умов у коді вважатись недоліком коду не можуть. %bad

\newpage

{\begin {center}\textbf{Завдання 6}\end {center}}\par
Вибрати всі істинні твердження (якщо існують).\par
\vspace{2mm}



(\,1\,) За відмови від документації одним з методом збереження знань про систему є постійне пересування персоналу між її частинами. %good



(\,2\,) RUP є компонентно-орієнтованим. %good



(\,3\,) Використання гнучких методологій при розробці надвеликих проектів не несе в собі жодних ризиків. %bad



(\,4\,) Використання гнучких методологій вимагає високої кваліфікації персоналу. %good


{\begin {center}\textbf{Завдання 7}\end {center}}\par
Вибрати всі істинні твердження (якщо існують).\par
\vspace{2mm}



(\,1\,) Розробка за гнучкими методологіями цінить людей та взаємодію між ними. %good



(\,2\,) Розповіді користувача (user stories) це менш формалізований порівняно з прецедентами спосіб опису вимог до системи. %good



(\,3\,) RUP є керованим прецедентами. %good



(\,4\,) На відміну від Scrum, RUP є компонентно-орієнтованим. %bad


{\begin {center}\textbf{Завдання 8}\end {center}}\par
Що з наступного не має відношення до Scrum або не виконується/не використовується за використання оригінального Scrum?\par
\vspace{2mm}



(\,1\,) приведення у належний вигляд журналу невиконаних задач (backlog grooming) %bad



(\,2\,) виділення одного з учасників команди як системного архітектора%good



(\,3\,) щоденний скрам%bad



(\,4\,) приведення у належний вигляд журналу невиконаних задач (backlog grooming) %bad


{\begin {center}\textbf{Завдання 9}\end {center}}\par
Що з наведеного нижче не розглядалось та навіть не згадувалось у курсі?\par
\vspace{2mm}



(\,1\,) LeSS%bad



(\,2\,) Scrum%bad



(\,3\,) V-модель%bad



(\,4\,) PRINCE2%good


{\begin {center}\textbf{Завдання 10}\end {center}}\par
Вибрати всі істинні твердження (якщо існують).\par
\vspace{2mm}



(\,1\,) Користувачі бувають основними.%good



(\,2\,) І верифікація, і валідація можуть використовувати тестування.%good



(\,3\,) За провалених димових тестів решта тестів не виконується і код повертається на доробку.%good



(\,4\,) Модель якості даних є ієрархічною.%good



\end {large}}
{\vspace{5mm}\smallskip \begin {small} Затверджено на засіданні кафедри теоретичної кібернетики, протокол \No 6 від   29.01.2021 р.\par\smallskip\smallskip
{\parbox {6.5 cm}{Зав. кафедри \dotfill Ю.В.~Крак}}\hspace {0.7cm}{\hbox to 6.5 cm{ Екзаменатор \dotfill Т.О.~Карнаух}} \end{small}}
\newpage

{{\begin {small}\par
{ \center  Київський національний університет імені Тараса Шевченка \par}
{\parbox {5 cm}{Освітня програма \hfill}{\parbox {7 cm}{Інформатика \hfill}}\parbox {6 cm}{\hfill Семестр 6}}\par
{\parbox {5 cm} {Навчальний предмет \hfill}\parbox {7 cm}{Розробка програмного забезпечення \hfill}\parbox {6cm}{\hfill}\par}\vspace{-7pt} \end{small}}{\center Екзаменаційний білет \No \textbf
{ 51 }\par}}
{
\begin {large}
{\begin {center}\textbf{Завдання 1}\end {center}}\par
Так звані трикутники балансування зв'язують між собою:\par
\vspace{2mm}



(\,1\,) 4 або 3 параметри, при цьому може бути відсутня якість%bad



(\,2\,) 4 або 3 параметри, при цьому може бути відсутній час%bad



(\,3\,) 4 або 3 параметри, при цьому може бути відсутній обсяг %good



(\,4\,) Завжди рівно 3 параметри%bad


{\begin {center}\textbf{Завдання 2}\end {center}}\par
Який з наведених фрагментів коду явно потребує поліпшення (форматування рядків та типи до уваги не брати)?\par
\vspace{2mm}



(\,1\,) \{double avg\_temperature;…\} %good



(\,2\,) void ff(…);%bad



(\,3\,) \{int x; int y; …\}%good



(\,4\,) \{int a\_, a\_\_; …\}%bad


{\begin {center}\textbf{Завдання 3}\end {center}}\par
Що є характерним для структурного програмування?\par
\vspace{2mm}



(\,1\,) Кожен блок коду повинен мати єдиний вхід та єдиний вихід.%good



(\,2\,) У коді мають використовуватись структурні типи даних. %bad



(\,3\,) У коді мають використовуватись типи структур. %bad



(\,4\,) Використання механізму винятків. %bad


{\begin {center}\textbf{Завдання 4}\end {center}}\par
Вибрати всі істинні твердження (якщо існують).\par
\vspace{2mm}



(\,1\,) Неперервне проектування здійснює створення та модифікацію проекту системи в міру розробки. %good



(\,2\,) Прогнозна модель розробки не може одночасно бути інкрементною. %bad



(\,3\,) Адаптивні моделі розробки не допускають зміну вимог до продукту, але продукт може розроблятись поетапно. %bad



(\,4\,) Адаптивні моделі розробки допускають зміну вимог до продукту. %good


{\begin {center}\textbf{Завдання 5}\end {center}}\par
Вибрати всі істинні твердження (якщо існують).\par
\vspace{2mm}



(\,1\,) Чим менше модулів у проекті, тим краще, бо меншу кількість модулів треба інтегрувати в єдине ціле. %bad



(\,2\,) Одним з прийомів спрощення складної умовної логіки є використання шаблону State. %good



(\,3\,) Повторюваний код має повторюватись у точності до символу, інакше він повторюваним вважатись не може. %bad



(\,4\,) Довгі функції доволі часто приховують складну логіку обробки або повторюваний код. %good

\newpage

{\begin {center}\textbf{Завдання 6}\end {center}}\par
Вибрати всі істинні твердження (якщо існують).\par
\vspace{2mm}



(\,1\,) Гнучкі методології розробки допускають суттєві зміни вимог, проте все одно намагаються якомога швидше зафіксувати архітектуру. %good



(\,2\,) За застосування адаптивних методологій розробки, так само як і за прогнозних, намагаються якомога швидше визначити архітектуру системи. Відмінність полягає в тім, у який спосіб це виконується.%good



(\,3\,) За використання гнучких методологій кожен розробник робить, що хоче. %bad



(\,4\,) Програмування в парах (pair programming) зводиться до того, що кожним фрагментом коду володіють два програмісти. %bad


{\begin {center}\textbf{Завдання 7}\end {center}}\par
Вибрати всі істинні твердження (якщо існують).\par
\vspace{2mm}



(\,1\,) Гнучкі методології можуть передбачати як спільне, так і індивідуальне володіння кодом. %good



(\,2\,) Методологія розробки Big Design Up Front погана тим, що може призводити до недостатнього проектування.%bad



(\,3\,) RUP є ітеративним та інкрементним.%good



(\,4\,) Тільки RUP використовує UML.%bad


{\begin {center}\textbf{Завдання 8}\end {center}}\par
Що з наступного не має відношення до Scrum або не виконується/не використовується за використання оригінального Scrum?\par
\vspace{2mm}



(\,1\,) скрам-мастер визначає пріоритети задач%good



(\,2\,) покер планування (Planning Poker)%bad



(\,3\,) журнал невиконаних задач (backlog) спрінту%bad



(\,4\,) обернені вимоги%bad


{\begin {center}\textbf{Завдання 9}\end {center}}\par
Що з наведеного нижче не розглядалось та навіть не згадувалось у курсі?\par
\vspace{2mm}



(\,1\,) модель хаосу%bad



(\,2\,) XP%bad



(\,3\,) SAFe%bad



(\,4\,) LOAN%good


{\begin {center}\textbf{Завдання 10}\end {center}}\par
Вибрати всі істинні твердження (якщо існують).\par
\vspace{2mm}



(\,1\,) Для тестування код треба ізолювати. Для того, щоб ізолювати, треба мати тести.%good



(\,2\,) Можна розглядати якість під час використання.%good



(\,3\,) Тестування вимог може породжувати їх уточнення.%good



(\,4\,) Якість буває зовнішньою.%good



\end {large}}
{\vspace{5mm}\smallskip \begin {small} Затверджено на засіданні кафедри теоретичної кібернетики, протокол \No 6 від   29.01.2021 р.\par\smallskip\smallskip
{\parbox {6.5 cm}{Зав. кафедри \dotfill Ю.В.~Крак}}\hspace {0.7cm}{\hbox to 6.5 cm{ Екзаменатор \dotfill Т.О.~Карнаух}} \end{small}}
\newpage

{{\begin {small}\par
{ \center  Київський національний університет імені Тараса Шевченка \par}
{\parbox {5 cm}{Освітня програма \hfill}{\parbox {7 cm}{Інформатика \hfill}}\parbox {6 cm}{\hfill Семестр 6}}\par
{\parbox {5 cm} {Навчальний предмет \hfill}\parbox {7 cm}{Розробка програмного забезпечення \hfill}\parbox {6cm}{\hfill}\par}\vspace{-7pt} \end{small}}{\center Екзаменаційний білет \No \textbf
{ 52 }\par}}
{
\begin {large}
{\begin {center}\textbf{Завдання 1}\end {center}}\par
Так звані трикутники балансування зв'язують між собою:\par
\vspace{2mm}



(\,1\,) 4 або 3 параметри, при цьому може бути відсутня якість%bad



(\,2\,) Завжди рівно 4 параметри%bad



(\,3\,) 4 або 3 параметри, при цьому може бути відсутній обсяг %good



(\,4\,) 4 або 3 параметри, при цьому може бути відсутній час%bad


{\begin {center}\textbf{Завдання 2}\end {center}}\par
Який з наведених фрагментів коду явно потребує поліпшення (форматування рядків та типи до уваги не брати)?\par
\vspace{2mm}



(\,1\,) \{double max\_power, min\_power;…\} %good



(\,2\,) \{int input\_file, file\_to\_output; …\}%bad



(\,3\,) void load\_data(T \& arg1, double arg2, double arg3, double arg4); %bad



(\,4\,) void help();%good


{\begin {center}\textbf{Завдання 3}\end {center}}\par
Що є характерним для структурного програмування?\par
\vspace{2mm}



(\,1\,) Використання оператора безумовного переходу. %bad



(\,2\,) Заборона використання скалярних (простих) типів даних. %bad



(\,3\,) У коді мають використовуватись типи даних, що є класами. %bad



(\,4\,) Кожен блок коду повинен мати єдиний вхід та єдиний вихід.%good


{\begin {center}\textbf{Завдання 4}\end {center}}\par
Вибрати всі істинні твердження (якщо існують).\par
\vspace{2mm}



(\,1\,) Адаптивна модель розробки адаптує вимоги до продукту до наявної команди розробників. %bad



(\,2\,) Неперервність проектування полягає в тім, що процес розробки не можна переривати більше ніж на вихідні. %bad



(\,3\,) Адаптивність методології проявляється в тому, що за зміни вимог розробка починається заново. %bad



(\,4\,) Ітеративні та інкрементні методи розробки є синонімом неперервного проектування. %bad


{\begin {center}\textbf{Завдання 5}\end {center}}\par
Вибрати всі істинні твердження (якщо існують).\par
\vspace{2mm}



(\,1\,) Код з довгими функціями має меншу кількість функцій, а тому його швидше та простіше тестувати. %bad



(\,2\,) Одним з прийомів спрощення складної умовної логіки є використання шаблону State. %good



(\,3\,) Повторюваний код має повторюватись у точності до символу, інакше він повторюваним вважатись не може. %bad



(\,4\,) Чим вища кваліфікація програміста, тим більш складні для розуміння та довгі умовні вирази він використовує. %bad

\newpage

{\begin {center}\textbf{Завдання 6}\end {center}}\par
Вибрати всі істинні твердження (якщо існують).\par
\vspace{2mm}



(\,1\,) Scrum передбачає колективне володіння кодом. %good



(\,2\,) Одна з принципових відмінностей RUP від Scrum полягає у використанні численних проміжних технічних артефактів. %good



(\,3\,) Гнучкі методології базуються на комунікації розробників, а тому передбачають виключно колективне володіння кодом. %bad



(\,4\,) Поступове уточнення вимог замовника полегшує адаптацію до змін. %good


{\begin {center}\textbf{Завдання 7}\end {center}}\par
Вибрати всі істинні твердження (якщо існують).\par
\vspace{2mm}



(\,1\,) За програмування в парах (pair programming) поки один колега набирає код на комп'ютері інший може піти пограти в настільний теніс. %bad



(\,2\,) RUP є архітектурно-орієнтованим. %good



(\,3\,) RUP відрізняється від Scrum тим, що ні в що не ставить персонал.%bad



(\,4\,) Гнучкі методології повністю нехтують документацією. %bad


{\begin {center}\textbf{Завдання 8}\end {center}}\par
Що з наступного не має відношення до Scrum або не виконується/не використовується за використання оригінального Scrum?\par
\vspace{2mm}



(\,1\,) виділення одного з учасників команди як системного архітектора%good



(\,2\,) власник продукту призначає відповідальних за конкретні задачі%good



(\,3\,) побудова різнопланових моделей системи%good



(\,4\,) порядок виконання запланованої на спрінт роботи визначають розробники%bad


{\begin {center}\textbf{Завдання 9}\end {center}}\par
Що з наведеного нижче не розглядалось та навіть не згадувалось у курсі?\par
\vspace{2mm}



(\,1\,) Lean%bad



(\,2\,) RUP%bad



(\,3\,) BDUF%bad



(\,4\,) PRINCE2%good


{\begin {center}\textbf{Завдання 10}\end {center}}\par
Вибрати всі істинні твердження (якщо існують).\par
\vspace{2mm}



(\,1\,) Коли нема можливості перевірити всі можливі комбінації, тестують їх попарні (pairwise) комбінації. Це може значно зменшити кількість тестів.%good



(\,2\,) Якість можна визначати через атрибути якості.%good



(\,3\,) Модель якості при використанні є ієрархічною.%good



(\,4\,) Відсутність знайдених дефектів на початкових стадіях розробки може бути свідченням поганого покриття коду тестами.%good



\end {large}}
{\vspace{5mm}\smallskip \begin {small} Затверджено на засіданні кафедри теоретичної кібернетики, протокол \No 6 від   29.01.2021 р.\par\smallskip\smallskip
{\parbox {6.5 cm}{Зав. кафедри \dotfill Ю.В.~Крак}}\hspace {0.7cm}{\hbox to 6.5 cm{ Екзаменатор \dotfill Т.О.~Карнаух}} \end{small}}
\newpage

{{\begin {small}\par
{ \center  Київський національний університет імені Тараса Шевченка \par}
{\parbox {5 cm}{Освітня програма \hfill}{\parbox {7 cm}{Інформатика \hfill}}\parbox {6 cm}{\hfill Семестр 6}}\par
{\parbox {5 cm} {Навчальний предмет \hfill}\parbox {7 cm}{Розробка програмного забезпечення \hfill}\parbox {6cm}{\hfill}\par}\vspace{-7pt} \end{small}}{\center Екзаменаційний білет \No \textbf
{ 53 }\par}}
{
\begin {large}
{\begin {center}\textbf{Завдання 1}\end {center}}\par
Що з наведеного найкраще характеризує рефакторінг?\par
\vspace{2mm}



(\,1\,) Довільна переробка коду, що виникає у процесі розробки%bad



(\,2\,) Зміна внутрішньої структури коду під час усунення помилок%bad



(\,3\,) Докорінна переробка всього коду з нуля,  що не змінює його видиму поведінку%bad



(\,4\,) Зміна внутрішньої структури коду, що не змінює його видиму поведінку%good


{\begin {center}\textbf{Завдання 2}\end {center}}\par
Який з наведених фрагментів коду явно потребує поліпшення (форматування рядків та типи до уваги не брати)?\par
\vspace{2mm}



(\,1\,) void help();%good



(\,2\,) \{fstream inputFile, outputFile;…\} %good



(\,3\,) \{double avg\_temperature;…\} %good



(\,4\,) void f13(…);%bad


{\begin {center}\textbf{Завдання 3}\end {center}}\par
Що є характерним для структурного програмування?\par
\vspace{2mm}



(\,1\,) Використання оператора безумовного переходу. %bad



(\,2\,) У коді мають використовуватись структурні типи даних. %bad



(\,3\,) Програма зображується ієрархічною структурою блоків.%good



(\,4\,) Використання механізму винятків. %bad


{\begin {center}\textbf{Завдання 4}\end {center}}\par
Вибрати всі істинні твердження (якщо існують).\par
\vspace{2mm}



(\,1\,) Неперервне проектування має бути ітеративним та інкрементним одночасно. %good



(\,2\,) В адаптивних методологіях планування робіт не виконується. %bad



(\,3\,) У випадку, коли вимоги до продукту відомі заздалегідь та є фіксованими контрактом, адаптивні методології розробки використовувати не можна. %bad



(\,4\,) І прогнозні, і адаптивні методології передбачають планування виконуваних робіт. %good


{\begin {center}\textbf{Завдання 5}\end {center}}\par
Вибрати всі істинні твердження (якщо існують).\par
\vspace{2mm}



(\,1\,) Зв'язки між окремими модулями системи мають бути якомога більш слабкими. %good



(\,2\,) Одним з прийомів спрощення складних умов є їх декомпозиція. %good



(\,3\,) Повторюваний код дозволяє програмісту краще зрозуміти його логіку та вибрати з кількох варіантів найкращий. %bad



(\,4\,) Правильно спроектована функція повинна мати рівно один обов'язок. %good

\newpage

{\begin {center}\textbf{Завдання 6}\end {center}}\par
Вибрати всі істинні твердження (якщо існують).\par
\vspace{2mm}



(\,1\,) Програмування в парах (pair programming) підвищує якість коду та дозволяє економити на виправленні помилок. %good



(\,2\,) За невизначених або піддатливих змінам вимог краще застосовувати гнучкі методології розробки. %good



(\,3\,) Якщо виникають проблеми з модифікацією коду, то слід розглянути внесення змін у проект. %good



(\,4\,) Розробка за гнучкими методологіями не передбачає узгодження умов контракту. %bad


{\begin {center}\textbf{Завдання 7}\end {center}}\par
Вибрати всі істинні твердження (якщо існують).\par
\vspace{2mm}



(\,1\,) RUP принципово відрізняється від Scrum тим, що є керованим прецедентами. У той же час у Scrum при плануванні розробки прецеденти (чи їх аналоги) до уваги не беруться. %bad



(\,2\,) Використання RUP не заперечує часті випуски продукту, але й не наполягає на них. %good



(\,3\,) За використання гнучких методологій розробки припиняти розробку не рекомендується. %good



(\,4\,) Гнучкі методології розробки намагаються якомога швидше визначити архітектуру системи, для чого на початку завжди виконується початкове проектування системи в цілому, а тільки потім переходять до ітеративного уточнення деталей та кодування. %bad


{\begin {center}\textbf{Завдання 8}\end {center}}\par
Що з наступного не має відношення до Scrum або не виконується/не використовується за використання оригінального Scrum?\par
\vspace{2mm}



(\,1\,) скрам-мастер визначає, які задачі команда буде виконувати на наступному спрінті%good



(\,2\,) власник продукту остаточно затверджує проект системи%good



(\,3\,) порядок виконання запланованої на спрінт роботи визначає скрам-мастер%good



(\,4\,) під час спрінту зміни, що впливають на досягнення мети спрінту, зазвичай, не допускаються%bad


{\begin {center}\textbf{Завдання 9}\end {center}}\par
Що з наведеного нижче не розглядалось та навіть не згадувалось у курсі?\par
\vspace{2mm}



(\,1\,) X-модель%good



(\,2\,) RAD%bad



(\,3\,) DSDM%bad



(\,4\,) CRAZY%good


{\begin {center}\textbf{Завдання 10}\end {center}}\par
Вибрати всі істинні твердження (якщо існують).\par
\vspace{2mm}



(\,1\,) Розглядають модель якості при використанні.%good



(\,2\,) Тестовий приклад складається з набору вхідних даних, умов виконання та очікуваних результатів.%good



(\,3\,) Розглядають модель якості продукту.%good



(\,4\,) Забезпечення якості це не тільки тестування.%good



\end {large}}
{\vspace{5mm}\smallskip \begin {small} Затверджено на засіданні кафедри теоретичної кібернетики, протокол \No 6 від   29.01.2021 р.\par\smallskip\smallskip
{\parbox {6.5 cm}{Зав. кафедри \dotfill Ю.В.~Крак}}\hspace {0.7cm}{\hbox to 6.5 cm{ Екзаменатор \dotfill Т.О.~Карнаух}} \end{small}}
\newpage

{{\begin {small}\par
{ \center  Київський національний університет імені Тараса Шевченка \par}
{\parbox {5 cm}{Освітня програма \hfill}{\parbox {7 cm}{Інформатика \hfill}}\parbox {6 cm}{\hfill Семестр 6}}\par
{\parbox {5 cm} {Навчальний предмет \hfill}\parbox {7 cm}{Розробка програмного забезпечення \hfill}\parbox {6cm}{\hfill}\par}\vspace{-7pt} \end{small}}{\center Екзаменаційний білет \No \textbf
{ 54 }\par}}
{
\begin {large}
{\begin {center}\textbf{Завдання 1}\end {center}}\par
Так звані трикутники балансування зв'язують між собою:\par
\vspace{2mm}



(\,1\,) 4 або 3 параметри, при цьому може бути відсутній час%bad



(\,2\,) Завжди рівно 3 параметри%bad



(\,3\,) Завжди рівно 4 параметри%bad



(\,4\,) 4 або 3 параметри, при цьому може бути відсутній обсяг %good


{\begin {center}\textbf{Завдання 2}\end {center}}\par
Який з наведених фрагментів коду явно потребує поліпшення (форматування рядків та типи до уваги не брати)?\par
\vspace{2mm}



(\,1\,) void load\_data(T \& arg1, double arg2, double arg3, double arg4); %bad



(\,2\,) \{fstream input\_file, output\_file;…\} %good



(\,3\,) void load\_data(…);%good



(\,4\,) \{int x; int y; …\}%good


{\begin {center}\textbf{Завдання 3}\end {center}}\par
Що є характерним для структурного програмування?\par
\vspace{2mm}



(\,1\,) У коді мають використовуватись типи структур. %bad



(\,2\,) У коді мають використовуватись типи структур, класи використовувати заборонено. %bad



(\,3\,) Використання оператора безумовного переходу. %bad



(\,4\,) Переходи з середини одного блоку в середину іншого неможливі. %good


{\begin {center}\textbf{Завдання 4}\end {center}}\par
Вибрати всі істинні твердження (якщо існують).\par
\vspace{2mm}



(\,1\,) У прогнозній моделі розробки розробка не може виконуватись ітеративно. %bad



(\,2\,) Адаптивні методології не передбачають визначення вимог. %bad



(\,3\,) Прогнозна модель розробки може одночасно бути інкрементною. %good



(\,4\,) Прогнозні моделі розробки з самого початку фіксують вимоги до продукту. %good


{\begin {center}\textbf{Завдання 5}\end {center}}\par
Вибрати всі істинні твердження (якщо існують).\par
\vspace{2mm}



(\,1\,) Код з довгими функціями має меншу кількість функцій, а тому для нього треба менше тестів. %bad



(\,2\,) Одним з прийомів спрощення умовної логіки є використання поліморфізму. %good



(\,3\,) Залежності всередині модуля можуть бути сильними. %good



(\,4\,) Складні комбінації умов починають керувати логікою виконання, тому що такими є вимоги до функціонування ПЗ. Оскільки ПЗ має якнайкраще забезпечувати потреби замовника, то складні комбінації умов у коді вважатись недоліком коду не можуть. %bad

\newpage

{\begin {center}\textbf{Завдання 6}\end {center}}\par
Вибрати всі істинні твердження (якщо існують).\par
\vspace{2mm}



(\,1\,) Гнучкі методології розробки є стійкими до припинення розробки на певний час. %bad



(\,2\,) Основна відмінність V-моделі від водоспадної полягає в тім, що діяльності з тестування виконуються паралельно з усіма іншими діяльностями з розробки. %good



(\,3\,) Гнучкі методи розробки нехтують проміжними технічними артефактами та документацією через те, що за зміни вимог та/або проекту треба витрачати додатковий час на підтримання їх у відповідному стані. %good



(\,4\,) Програмування в парах (pair programming) підвищує якість коду та дозволяє економити на виправленні помилок. %good


{\begin {center}\textbf{Завдання 7}\end {center}}\par
Вибрати всі істинні твердження (якщо існують).\par
\vspace{2mm}



(\,1\,) Без рефакторінгу неперервне проектування неможливе. %good



(\,2\,) Поступове уточнення вимог замовника полегшує адаптацію до змін. %good



(\,3\,) RUP принципово відрізняється від Scrum тим, що є керованим прецедентами. У той же час у Scrum при плануванні розробки прецеденти (чи їх аналоги) до уваги не беруться. %bad



(\,4\,) Розповіді користувача (user stories) це менш формалізований порівняно з прецедентами спосіб опису вимог до системи. %good


{\begin {center}\textbf{Завдання 8}\end {center}}\par
Що з наступного не має відношення до Scrum або не виконується/не використовується за використання оригінального Scrum?\par
\vspace{2mm}



(\,1\,) команда визначає, які задачі вона буде виконувати на наступному спрінті%good



(\,2\,) розповіді користувача%bad



(\,3\,) власник продукту призначає час, необхідний команді для реалізації конкретної задачі%good



(\,4\,) бачення продукту (vision) %bad


{\begin {center}\textbf{Завдання 9}\end {center}}\par
Що з наведеного нижче не розглядалось та навіть не згадувалось у курсі?\par
\vspace{2mm}



(\,1\,) TDD%bad



(\,2\,) LeSS%bad



(\,3\,) модель хаосу%bad



(\,4\,) sashimi%bad


{\begin {center}\textbf{Завдання 10}\end {center}}\par
Вибрати всі істинні твердження (якщо існують).\par
\vspace{2mm}



(\,1\,) Атрибути якості можна використовувати при уточненні вимог.%good



(\,2\,) Модель якості продукту є ієрархічною.%good



(\,3\,) Користувачі бувають непрямими.%good



(\,4\,) Тестування буває позитивне та негативне.%good



\end {large}}
{\vspace{5mm}\smallskip \begin {small} Затверджено на засіданні кафедри теоретичної кібернетики, протокол \No 6 від   29.01.2021 р.\par\smallskip\smallskip
{\parbox {6.5 cm}{Зав. кафедри \dotfill Ю.В.~Крак}}\hspace {0.7cm}{\hbox to 6.5 cm{ Екзаменатор \dotfill Т.О.~Карнаух}} \end{small}}
\newpage

{{\begin {small}\par
{ \center  Київський національний університет імені Тараса Шевченка \par}
{\parbox {5 cm}{Освітня програма \hfill}{\parbox {7 cm}{Інформатика \hfill}}\parbox {6 cm}{\hfill Семестр 6}}\par
{\parbox {5 cm} {Навчальний предмет \hfill}\parbox {7 cm}{Розробка програмного забезпечення \hfill}\parbox {6cm}{\hfill}\par}\vspace{-7pt} \end{small}}{\center Екзаменаційний білет \No \textbf
{ 55 }\par}}
{
\begin {large}
{\begin {center}\textbf{Завдання 1}\end {center}}\par
Шаблон проектування:\par
\vspace{2mm}



(\,1\,) Формулює ідею розв'язання проблеми та може пропонувати одну чи кілька її реалізацій%good



(\,2\,) Є шаблоном, в який треба підставити власні ідентифікатори%bad



(\,3\,) Задає розв'язання типової проблеми до найдрібніших подробиць%bad



(\,4\,) Тільки формулює ідею розв'язання проблеми%bad


{\begin {center}\textbf{Завдання 2}\end {center}}\par
Який з наведених фрагментів коду явно потребує поліпшення (форматування рядків та типи до уваги не брати)?\par
\vspace{2mm}



(\,1\,) \{double maxPower, minPower;…\} %good



(\,2\,) \{fstream input\_file, outputFile; …\}%bad



(\,3\,) \{double max\_power, min\_power;…\} %good



(\,4\,) \{int a\_, a\_\_; …\}%bad


{\begin {center}\textbf{Завдання 3}\end {center}}\par
Що є характерним для структурного програмування?\par
\vspace{2mm}



(\,1\,) Переходи з середини одного блоку в середину іншого неможливі. %good



(\,2\,) У коді мають використовуватись типи структур, класи використовувати заборонено. %bad



(\,3\,) У коді мають використовуватись типи даних, що є класами. %bad



(\,4\,) У коді мають використовуватись типи структур. %bad


{\begin {center}\textbf{Завдання 4}\end {center}}\par
Вибрати всі істинні твердження (якщо існують).\par
\vspace{2mm}



(\,1\,) Адаптивні методології не передбачають визначення вимог. %bad



(\,2\,) Неперервність проектування полягає в тім, що процес розробки не можна переривати більше ніж на вихідні. %bad



(\,3\,) Для адаптивних методологій розробки, на відміну від прогнозних, попередні оцінки бюджету та/або часу не виконують, бо вони все одно зміняться. %bad



(\,4\,) Для адаптивних моделей розробки оцінки часу та бюджету не виконують, бо вимоги є змінними. %bad


{\begin {center}\textbf{Завдання 5}\end {center}}\par
Вибрати всі істинні твердження (якщо існують).\par
\vspace{2mm}



(\,1\,) Залежності між різними модулями системи мають бути якомога більш сильними, бо інакше вона розпадеться. %bad



(\,2\,) Гарний код має читатись як розмовна мова. %good



(\,3\,) Складні умови виникають в реальних задачах доволі часто, а тому є доволі природніми і не можуть вважатись недоліком коду. %bad



(\,4\,) Втілення в коді різних підходів до одного того самого перетворення даних дозволяє програмісту продемонструвати свою професійну підготовку з найкращого боку. %bad

\newpage

{\begin {center}\textbf{Завдання 6}\end {center}}\par
Вибрати всі істинні твердження (якщо існують).\par
\vspace{2mm}



(\,1\,) Адаптивні методології розробки дозволяють уникнути надлишкового проектування. %good



(\,2\,) Модель хаосу є методологією розробки програмного забезпечення. %bad



(\,3\,) Розробка за гнучкими методологіями цінить людей та взаємодію між ними. %good



(\,4\,) Використання гнучких методологій при розробці надвеликих проектів не несе в собі жодних ризиків. %bad


{\begin {center}\textbf{Завдання 7}\end {center}}\par
Вибрати всі істинні твердження (якщо існують).\par
\vspace{2mm}



(\,1\,) Використання RUP не заперечує часті випуски продукту, але й не наполягає на них. %good



(\,2\,) RUP є архітектурно-орієнтованим. %good



(\,3\,) Розробка за гнучкими методологіями ніколи не слідує плану. %bad



(\,4\,) За використання гнучких методологій розробки припиняти розробку не рекомендується. %good


{\begin {center}\textbf{Завдання 8}\end {center}}\par
Що з наступного не має відношення до Scrum або не виконується/не використовується за використання оригінального Scrum?\par
\vspace{2mm}



(\,1\,) ретроспектива спрінту%bad



(\,2\,) канбан%bad



(\,3\,) розповіді користувача%bad



(\,4\,) журнал невиконаних задач (backlog) продукту%bad


{\begin {center}\textbf{Завдання 9}\end {center}}\par
Що з наведеного нижче не розглядалось та навіть не згадувалось у курсі?\par
\vspace{2mm}



(\,1\,) poker%bad



(\,2\,) V-модель%bad



(\,3\,) LAZY%good



(\,4\,) XP%bad


{\begin {center}\textbf{Завдання 10}\end {center}}\par
Вибрати всі істинні твердження (якщо існують).\par
\vspace{2mm}



(\,1\,) Розглядають модель якості даних.%good



(\,2\,) За провалених димових тестів решта тестів не виконується і код повертається на доробку.%good



(\,3\,) Тестування буває позитивне та негативне.%good



(\,4\,) Якість буває внутрішньою.%good



\end {large}}
{\vspace{5mm}\smallskip \begin {small} Затверджено на засіданні кафедри теоретичної кібернетики, протокол \No 6 від   29.01.2021 р.\par\smallskip\smallskip
{\parbox {6.5 cm}{Зав. кафедри \dotfill Ю.В.~Крак}}\hspace {0.7cm}{\hbox to 6.5 cm{ Екзаменатор \dotfill Т.О.~Карнаух}} \end{small}}
\newpage

{{\begin {small}\par
{ \center  Київський національний університет імені Тараса Шевченка \par}
{\parbox {5 cm}{Освітня програма \hfill}{\parbox {7 cm}{Інформатика \hfill}}\parbox {6 cm}{\hfill Семестр 6}}\par
{\parbox {5 cm} {Навчальний предмет \hfill}\parbox {7 cm}{Розробка програмного забезпечення \hfill}\parbox {6cm}{\hfill}\par}\vspace{-7pt} \end{small}}{\center Екзаменаційний білет \No \textbf
{ 56 }\par}}
{
\begin {large}
{\begin {center}\textbf{Завдання 1}\end {center}}\par
Так звані трикутники балансування зв'язують між собою:\par
\vspace{2mm}



(\,1\,) Завжди рівно 3 параметри%bad



(\,2\,) Завжди рівно 4 параметри%bad



(\,3\,) 4 або 3 параметри, при цьому може бути відсутній обсяг %good



(\,4\,) 4 або 3 параметри, при цьому може бути відсутня якість%bad


{\begin {center}\textbf{Завдання 2}\end {center}}\par
Який з наведених фрагментів коду явно потребує поліпшення (форматування рядків та типи до уваги не брати)?\par
\vspace{2mm}



(\,1\,) \{int xtrmprg; …\}%bad



(\,2\,) void ff(…);%bad



(\,3\,) \{double the\_best\_value\_anyone\_ever\_have;\} %bad



(\,4\,) \{int input\_file, file\_to\_output; …\}%bad


{\begin {center}\textbf{Завдання 3}\end {center}}\par
Що є характерним для структурного програмування?\par
\vspace{2mm}



(\,1\,) Використання механізму винятків. %bad



(\,2\,) Заборона використання скалярних (простих) типів даних. %bad



(\,3\,) Кожен блок коду повинен мати єдиний вхід та єдиний вихід.%good



(\,4\,) У коді мають використовуватись структурні типи даних. %bad


{\begin {center}\textbf{Завдання 4}\end {center}}\par
Вибрати всі істинні твердження (якщо існують).\par
\vspace{2mm}



(\,1\,) Неперервне проектування має бути ітеративним та інкрементним одночасно. %good



(\,2\,) Ітеративні та інкрементні методи розробки є синонімом неперервного проектування. %bad



(\,3\,) Для адаптивних методологій прогнозувати результат чи терміни його отримання абсолютно неможливо. %bad



(\,4\,) Адаптивні методології має сенс використовувати для проектів, де можлива зміна вимог. %good


{\begin {center}\textbf{Завдання 5}\end {center}}\par
Вибрати всі істинні твердження (якщо існують).\par
\vspace{2mm}



(\,1\,) Якщо кожна функція бібліотеки буде здатна в залежності від параметрів виконувати кілька різних дій, то загальна кількість функцій у ній буде менша і її буде простіше використовувати. %bad



(\,2\,) Одним з засобів позбутись повторень у коді є використання успадкування. %good



(\,3\,) Одним з засобів позбутись повторень у коді є використання функцій. %good



(\,4\,) Від довгих функцій можна позбутися, розклавши їх у послідовність викликів дрібніших функцій. %good

\newpage

{\begin {center}\textbf{Завдання 6}\end {center}}\par
Вибрати всі істинні твердження (якщо існують).\par
\vspace{2mm}



(\,1\,) Гнучкі методології розробки намагаються якомога швидше визначити архітектуру системи, для чого на початку завжди виконується початкове проектування системи в цілому, а тільки потім переходять до ітеративного уточнення деталей та кодування. %bad



(\,2\,) Гнучкі методології можуть передбачати як спільне, так і індивідуальне володіння кодом. %good



(\,3\,) Гнучкі методології розробки не допускають використання UML. %bad



(\,4\,) Перша версія коду не завжди може бути якісною -- і це нормально. %good


{\begin {center}\textbf{Завдання 7}\end {center}}\par
Вибрати всі істинні твердження (якщо існують).\par
\vspace{2mm}



(\,1\,) Висока внутрішня якість коду здешевшує розробку, бо дозволяє економити на вартості та часі майбутніх змін. %good



(\,2\,) Розробка за гнучкими методологіями цінить співробітництво з замовником та піддатлівість змінам. %good



(\,3\,) Одна з принципових відмінностей RUP від Scrum полягає у використанні численних проміжних технічних артефактів. %good



(\,4\,) Scrum передбачає колективне володіння кодом. %good


{\begin {center}\textbf{Завдання 8}\end {center}}\par
Що з наступного не має відношення до Scrum або не виконується/не використовується за використання оригінального Scrum?\par
\vspace{2mm}



(\,1\,) правила визначення готовності (Definition of Done) %bad



(\,2\,) виділення в команді одного чи кількох проектувальників%good



(\,3\,) планування спрінту%bad



(\,4\,) залучення зовнішніх фахівців для побудови архітектури системи%good


{\begin {center}\textbf{Завдання 9}\end {center}}\par
Що з наведеного нижче не розглядалось та навіть не згадувалось у курсі?\par
\vspace{2mm}



(\,1\,) FDD%bad



(\,2\,) Cleanroom%bad



(\,3\,) LAZY%good



(\,4\,) sashimi%bad


{\begin {center}\textbf{Завдання 10}\end {center}}\par
Вибрати всі істинні твердження (якщо існують).\par
\vspace{2mm}



(\,1\,) Тестування вимог може породжувати їх уточнення.%good



(\,2\,) Користувачі бувають вторинними.%good



(\,3\,) Модель якості даних є ієрархічною.%good



(\,4\,) Розглядають модель якості при використанні.%good



\end {large}}
{\vspace{5mm}\smallskip \begin {small} Затверджено на засіданні кафедри теоретичної кібернетики, протокол \No 6 від   29.01.2021 р.\par\smallskip\smallskip
{\parbox {6.5 cm}{Зав. кафедри \dotfill Ю.В.~Крак}}\hspace {0.7cm}{\hbox to 6.5 cm{ Екзаменатор \dotfill Т.О.~Карнаух}} \end{small}}
\newpage

{{\begin {small}\par
{ \center  Київський національний університет імені Тараса Шевченка \par}
{\parbox {5 cm}{Освітня програма \hfill}{\parbox {7 cm}{Інформатика \hfill}}\parbox {6 cm}{\hfill Семестр 6}}\par
{\parbox {5 cm} {Навчальний предмет \hfill}\parbox {7 cm}{Розробка програмного забезпечення \hfill}\parbox {6cm}{\hfill}\par}\vspace{-7pt} \end{small}}{\center Екзаменаційний білет \No \textbf
{ 57 }\par}}
{
\begin {large}
{\begin {center}\textbf{Завдання 1}\end {center}}\par
Так звані трикутники балансування зв'язують між собою:\par
\vspace{2mm}



(\,1\,) 4 або 3 параметри, при цьому може бути відсутній обсяг %good



(\,2\,) 4 або 3 параметри, при цьому може бути відсутня якість%bad



(\,3\,) 4 або 3 параметри, при цьому може бути відсутній час%bad



(\,4\,) Завжди рівно 3 параметри%bad


{\begin {center}\textbf{Завдання 2}\end {center}}\par
Який з наведених фрагментів коду явно потребує поліпшення (форматування рядків та типи до уваги не брати)?\par
\vspace{2mm}



(\,1\,) \{double the\_best\_value\_anyone\_ever\_have;\} %bad



(\,2\,) void f13(…);%bad



(\,3\,) void load\_data(…);%good



(\,4\,) \{fstream input\_file, outputFile; …\}%bad


{\begin {center}\textbf{Завдання 3}\end {center}}\par
Що є характерним для структурного програмування?\par
\vspace{2mm}



(\,1\,) Заборона використання скалярних (простих) типів даних. %bad



(\,2\,) Використання механізму винятків. %bad



(\,3\,) У коді мають використовуватись типи структур. %bad



(\,4\,) Програма зображується ієрархічною структурою блоків.%good


{\begin {center}\textbf{Завдання 4}\end {center}}\par
Вибрати всі істинні твердження (якщо існують).\par
\vspace{2mm}



(\,1\,) Неперервне проектування здійснює створення та модифікацію проекту системи в міру розробки. %good



(\,2\,) У прогнозній моделі розробки розробка не може виконуватись ітеративно. %bad



(\,3\,) Прогнозні моделі розробки з самого початку фіксують вимоги до продукту. %good



(\,4\,) Адаптивна модель розробки адаптує вимоги до продукту до наявної команди розробників. %bad


{\begin {center}\textbf{Завдання 5}\end {center}}\par
Вибрати всі істинні твердження (якщо існують).\par
\vspace{2mm}



(\,1\,) Гарний код не повинен містити різних фрагментів, що складаються з кількох інструкцій та виконують однакове перетворення даних. %good



(\,2\,) Довжина функцій не впливає на внутрішню якість коду. %bad



(\,3\,) Чим менше класів у коді, тим менше класів треба налагоджувати, а тому класи за можливості слід об'єднувати. %bad



(\,4\,) Одним з прийомів спрощення складної умовної логіки є використання шаблону Strategy. %good

\newpage

{\begin {center}\textbf{Завдання 6}\end {center}}\par
Вибрати всі істинні твердження (якщо існують).\par
\vspace{2mm}



(\,1\,) Ітеративна методологія розробки у загальному випадку не гарантує, що не буде надлишкового проектування. %good



(\,2\,) Гнучкі методи розробки нехтують проміжними технічними артефактами та документацією через те, що за зміни вимог та/або проекту треба витрачати додатковий час на підтримання їх у відповідному стані. %good



(\,3\,) Якщо виникають проблеми з модифікацією коду, то слід відразу змінити команду розробників на більш кваліфіковану.%bad



(\,4\,) RUP є компонентно-орієнтованим. %good


{\begin {center}\textbf{Завдання 7}\end {center}}\par
Вибрати всі істинні твердження (якщо існують).\par
\vspace{2mm}



(\,1\,) За відмови від документації одним з методом збереження знань про систему є постійне пересування персоналу між її частинами. %good



(\,2\,) Водоспадна методологія розробки може призводити до недостатнього проектування. %bad



(\,3\,) Програмування в парах (pair programming) зводиться до того, що кожним фрагментом коду володіють два програмісти. %bad



(\,4\,) Гнучкі методології повністю нехтують документацією. %bad


{\begin {center}\textbf{Завдання 8}\end {center}}\par
Що з наступного не має відношення до Scrum або не виконується/не використовується за використання оригінального Scrum?\par
\vspace{2mm}



(\,1\,) під час спрінту дозволяється внесення змін до вимог, що вже реалізуються на поточному спрінті%good



(\,2\,) скрам-мастер визначає час, необхідний команді для реалізації конкретної задачі%good



(\,3\,) виділення в команді одного чи кількох тестувальників%good



(\,4\,) скрам-мастер призначає відповідальних за конкретні задачі%good


{\begin {center}\textbf{Завдання 9}\end {center}}\par
Що з наведеного нижче не розглядалось та навіть не згадувалось у курсі?\par
\vspace{2mm}



(\,1\,) Lean%bad



(\,2\,) CRAZY%good



(\,3\,) Scrum%bad



(\,4\,) BDUF%bad


{\begin {center}\textbf{Завдання 10}\end {center}}\par
Вибрати всі істинні твердження (якщо існують).\par
\vspace{2mm}



(\,1\,) Якість буває зовнішньою.%good



(\,2\,) Можна розглядати якість під час використання.%good



(\,3\,) Атрибути якості можна використовувати при уточненні вимог.%good



(\,4\,) Користувачі бувають непрямими.%good



\end {large}}
{\vspace{5mm}\smallskip \begin {small} Затверджено на засіданні кафедри теоретичної кібернетики, протокол \No 6 від   29.01.2021 р.\par\smallskip\smallskip
{\parbox {6.5 cm}{Зав. кафедри \dotfill Ю.В.~Крак}}\hspace {0.7cm}{\hbox to 6.5 cm{ Екзаменатор \dotfill Т.О.~Карнаух}} \end{small}}
\newpage

{{\begin {small}\par
{ \center  Київський національний університет імені Тараса Шевченка \par}
{\parbox {5 cm}{Освітня програма \hfill}{\parbox {7 cm}{Інформатика \hfill}}\parbox {6 cm}{\hfill Семестр 6}}\par
{\parbox {5 cm} {Навчальний предмет \hfill}\parbox {7 cm}{Розробка програмного забезпечення \hfill}\parbox {6cm}{\hfill}\par}\vspace{-7pt} \end{small}}{\center Екзаменаційний білет \No \textbf
{ 58 }\par}}
{
\begin {large}
{\begin {center}\textbf{Завдання 1}\end {center}}\par
Так звані трикутники балансування зв'язують між собою:\par
\vspace{2mm}



(\,1\,) 4 або 3 параметри, при цьому може бути відсутній час%bad



(\,2\,) Завжди рівно 4 параметри%bad



(\,3\,) 4 або 3 параметри, при цьому може бути відсутній обсяг %good



(\,4\,) Завжди рівно 3 параметри%bad


{\begin {center}\textbf{Завдання 2}\end {center}}\par
Який з наведених фрагментів коду явно потребує поліпшення (форматування рядків та типи до уваги не брати)?\par
\vspace{2mm}



(\,1\,) void help();%good



(\,2\,) \{fstream input\_file, output\_file;…\} %good



(\,3\,) void ff(…);%bad



(\,4\,) \{int xtrmprg; …\}%bad


{\begin {center}\textbf{Завдання 3}\end {center}}\par
Що є характерним для структурного програмування?\par
\vspace{2mm}



(\,1\,) У коді мають використовуватись структурні типи даних. %bad



(\,2\,) Переходи з середини одного блоку в середину іншого неможливі. %good



(\,3\,) Використання оператора безумовного переходу. %bad



(\,4\,) У коді мають використовуватись типи даних, що є класами. %bad


{\begin {center}\textbf{Завдання 4}\end {center}}\par
Вибрати всі істинні твердження (якщо існують).\par
\vspace{2mm}



(\,1\,) Прогнозні та адаптивні методології використовують одні ті самі діяльності з розробки ПЗ: визначення та аналіз вимог, проектування, реалізація, тестування, впровадження, супроводження. %good



(\,2\,) Адаптивні моделі розробки не допускають зміну вимог до продукту, але продукт може розроблятись поетапно. %bad



(\,3\,) У випадку, коли вимоги до продукту відомі заздалегідь та є фіксованими контрактом, адаптивні методології розробки використовувати не можна. %bad



(\,4\,) Неперервне проектування створює специфікацію проекту перед початком розробки/ітерації, а під час ітерації його неперервно модифікує.%bad


{\begin {center}\textbf{Завдання 5}\end {center}}\par
Вибрати всі істинні твердження (якщо існують).\par
\vspace{2mm}



(\,1\,) Одним з засобів позбутись повторень у коді є використання шаблону Template Method (але цей шаблон здатен боротися і з іншими недоліками коду). %good



(\,2\,) Чим менше модулів у проекті, тим краще, бо меншу кількість модулів треба інтегрувати в єдине ціле. %bad



(\,3\,) Код, що виконує однакову послідовність логічних кроків, у більшості випадків є повторюваним.%good



(\,4\,) Код програми пишеться для комп'ютера, тому він має бути синтаксично та семантично коректним. Інше значення не має. %bad

\newpage

{\begin {center}\textbf{Завдання 6}\end {center}}\par
Вибрати всі істинні твердження (якщо існують).\par
\vspace{2mm}



(\,1\,) RUP є ітеративним та інкрементним.%good



(\,2\,) Використання гнучких методологій вимагає високої кваліфікації персоналу. %good



(\,3\,) Гнучкі методології розробки допускають суттєві зміни вимог, проте все одно намагаються якомога швидше зафіксувати архітектуру. %good



(\,4\,) За невизначених або піддатливих змінам вимог краще застосовувати гнучкі методології розробки. %good


{\begin {center}\textbf{Завдання 7}\end {center}}\par
Вибрати всі істинні твердження (якщо існують).\par
\vspace{2mm}



(\,1\,) Водоспадна методологія розробки може призводити до надлишкового проектування. %good



(\,2\,) Гнучкі методи розробки нехтують проміжними технічними артефактами та документацією виключно через їх непотрібність. %bad



(\,3\,) Методологія розробки Big Design Up Front погана тим, що може призводити до надлишкового проектування. %good



(\,4\,) Висока внутрішня якість коду є марною тратою часу на те, що замовник не побачить та/або не зможе зацінити. %bad


{\begin {center}\textbf{Завдання 8}\end {center}}\par
Що з наступного не має відношення до Scrum або не виконується/не використовується за використання оригінального Scrum?\par
\vspace{2mm}



(\,1\,) скрам-мастер остаточно затверджує проект системи%good



(\,2\,) порядок виконання запланованої на спрінт роботи визначає власник продукту%good



(\,3\,) блочні тести%bad



(\,4\,) скрам-мастер призначає відповідальних за конкретні задачі%good


{\begin {center}\textbf{Завдання 9}\end {center}}\par
Що з наведеного нижче не розглядалось та навіть не згадувалось у курсі?\par
\vspace{2mm}



(\,1\,) RAD%bad



(\,2\,) водоспадна модель%bad



(\,3\,) poker%bad



(\,4\,) SAFe%bad


{\begin {center}\textbf{Завдання 10}\end {center}}\par
Вибрати всі істинні твердження (якщо існують).\par
\vspace{2mm}



(\,1\,) Користувачі бувають основними.%good



(\,2\,) Модель якості продукту є ієрархічною.%good



(\,3\,) Тестовий приклад складається з набору вхідних даних, умов виконання та очікуваних результатів.%good



(\,4\,) Модель якості при використанні є ієрархічною.%good



\end {large}}
{\vspace{5mm}\smallskip \begin {small} Затверджено на засіданні кафедри теоретичної кібернетики, протокол \No 6 від   29.01.2021 р.\par\smallskip\smallskip
{\parbox {6.5 cm}{Зав. кафедри \dotfill Ю.В.~Крак}}\hspace {0.7cm}{\hbox to 6.5 cm{ Екзаменатор \dotfill Т.О.~Карнаух}} \end{small}}
\newpage

{{\begin {small}\par
{ \center  Київський національний університет імені Тараса Шевченка \par}
{\parbox {5 cm}{Освітня програма \hfill}{\parbox {7 cm}{Інформатика \hfill}}\parbox {6 cm}{\hfill Семестр 6}}\par
{\parbox {5 cm} {Навчальний предмет \hfill}\parbox {7 cm}{Розробка програмного забезпечення \hfill}\parbox {6cm}{\hfill}\par}\vspace{-7pt} \end{small}}{\center Екзаменаційний білет \No \textbf
{ 59 }\par}}
{
\begin {large}
{\begin {center}\textbf{Завдання 1}\end {center}}\par
Що є метою рефакторингу?\par
\vspace{2mm}



(\,1\,) Зменшення кількості рядків коду%bad



(\,2\,) Додавання нової функціональності%bad



(\,3\,) Усунення помилок у коді%bad



(\,4\,) Поліпшення внутрішньої структури коду%good


{\begin {center}\textbf{Завдання 2}\end {center}}\par
Який з наведених фрагментів коду явно потребує поліпшення (форматування рядків та типи до уваги не брати)?\par
\vspace{2mm}



(\,1\,) \{double max\_power, min\_power;…\} %good



(\,2\,) \{int input\_file, file\_to\_output; …\}%bad



(\,3\,) \{fstream inputFile, outputFile;…\} %good



(\,4\,) \{int a\_, a\_\_; …\}%bad


{\begin {center}\textbf{Завдання 3}\end {center}}\par
Що є характерним для структурного програмування?\par
\vspace{2mm}



(\,1\,) Кожен блок коду повинен мати єдиний вхід та єдиний вихід.%good



(\,2\,) У коді мають використовуватись типи структур, класи використовувати заборонено. %bad



(\,3\,) У коді мають використовуватись типи структур. %bad



(\,4\,) У коді мають використовуватись типи даних, що є класами. %bad


{\begin {center}\textbf{Завдання 4}\end {center}}\par
Вибрати всі істинні твердження (якщо існують).\par
\vspace{2mm}



(\,1\,) В адаптивних методологіях планування робіт не виконується. %bad



(\,2\,) Адаптивні моделі розробки допускають зміну вимог до продукту. %good



(\,3\,) Прогнозна модель розробки не може одночасно бути інкрементною. %bad



(\,4\,) Гнучкі методології не передбачають планування робіт. %bad


{\begin {center}\textbf{Завдання 5}\end {center}}\par
Вибрати всі істинні твердження (якщо існують).\par
\vspace{2mm}



(\,1\,) Втілення в коді різних підходів до одного того самого перетворення даних дозволяє програмісту продемонструвати свою професійну підготовку з найкращого боку. %bad



(\,2\,) Одним з засобів позбутись повторень у коді є використання функцій. %good



(\,3\,) Код з довгими функціями має меншу кількість функцій, а тому для нього треба менше тестів. %bad



(\,4\,) Чим вища кваліфікація програміста, тим більш складні для розуміння та довгі умовні вирази він використовує. %bad

\newpage

{\begin {center}\textbf{Завдання 6}\end {center}}\par
Вибрати всі істинні твердження (якщо існують).\par
\vspace{2mm}



(\,1\,) RUP є керованим прецедентами. %good



(\,2\,) На відміну від Scrum, RUP є компонентно-орієнтованим. %bad



(\,3\,) Програмування в парах (pair programming) спрощує введення в команду нових розробників. %good



(\,4\,) Розробка за гнучкими методологіями не передбачає узгодження умов контракту. %bad


{\begin {center}\textbf{Завдання 7}\end {center}}\par
Вибрати всі істинні твердження (якщо існують).\par
\vspace{2mm}



(\,1\,) За застосування адаптивних методологій розробки, так само як і за прогнозних, намагаються якомога швидше визначити архітектуру системи. Відмінність полягає в тім, у який спосіб це виконується.%good



(\,2\,) Водоспадна методологія не призначена для внесення змін у вимоги до системи. %good



(\,3\,) Код, що себе сам документує, -- це код, який виводить користувачу документацію на себе. %bad



(\,4\,) Гнучкі методології розробки є стійкими до припинення розробки на певний час. %bad


{\begin {center}\textbf{Завдання 8}\end {center}}\par
Що з наступного не має відношення до Scrum або не виконується/не використовується за використання оригінального Scrum?\par
\vspace{2mm}



(\,1\,) щоденний скрам%bad



(\,2\,) скрам-мастер визначає пріоритети задач%good



(\,3\,) власник продукту остаточно затверджує проект системи%good



(\,4\,) виділення одного з учасників команди як системного архітектора%good


{\begin {center}\textbf{Завдання 9}\end {center}}\par
Що з наведеного нижче не розглядалось та навіть не згадувалось у курсі?\par
\vspace{2mm}



(\,1\,) BDD%bad



(\,2\,) Scrum of scrums%bad



(\,3\,) LOAN%good



(\,4\,) X-модель%good


{\begin {center}\textbf{Завдання 10}\end {center}}\par
Вибрати всі істинні твердження (якщо існують).\par
\vspace{2mm}



(\,1\,) І верифікація, і валідація можуть використовувати тестування.%good



(\,2\,) Відсутність знайдених дефектів на початкових стадіях розробки може бути свідченням поганого покриття коду тестами.%good



(\,3\,) Коли нема можливості перевірити всі можливі комбінації, тестують їх попарні (pairwise) комбінації. Це може значно зменшити кількість тестів.%good



(\,4\,) Забезпечення якості це не тільки тестування.%good



\end {large}}
{\vspace{5mm}\smallskip \begin {small} Затверджено на засіданні кафедри теоретичної кібернетики, протокол \No 6 від   29.01.2021 р.\par\smallskip\smallskip
{\parbox {6.5 cm}{Зав. кафедри \dotfill Ю.В.~Крак}}\hspace {0.7cm}{\hbox to 6.5 cm{ Екзаменатор \dotfill Т.О.~Карнаух}} \end{small}}
\newpage

{{\begin {small}\par
{ \center  Київський національний університет імені Тараса Шевченка \par}
{\parbox {5 cm}{Освітня програма \hfill}{\parbox {7 cm}{Інформатика \hfill}}\parbox {6 cm}{\hfill Семестр 6}}\par
{\parbox {5 cm} {Навчальний предмет \hfill}\parbox {7 cm}{Розробка програмного забезпечення \hfill}\parbox {6cm}{\hfill}\par}\vspace{-7pt} \end{small}}{\center Екзаменаційний білет \No \textbf
{ 60 }\par}}
{
\begin {large}
{\begin {center}\textbf{Завдання 1}\end {center}}\par
Так звані трикутники балансування зв'язують між собою:\par
\vspace{2mm}



(\,1\,) Завжди рівно 4 параметри%bad



(\,2\,) 4 або 3 параметри, при цьому може бути відсутній обсяг %good



(\,3\,) Завжди рівно 3 параметри%bad



(\,4\,) 4 або 3 параметри, при цьому може бути відсутня якість%bad


{\begin {center}\textbf{Завдання 2}\end {center}}\par
Який з наведених фрагментів коду явно потребує поліпшення (форматування рядків та типи до уваги не брати)?\par
\vspace{2mm}



(\,1\,) void load\_data(T \& arg1, double arg2, double arg3, double arg4); %bad



(\,2\,) \{double avg\_temperature;…\} %good



(\,3\,) \{int x; int y; …\}%good



(\,4\,) \{double maxPower, minPower;…\} %good


{\begin {center}\textbf{Завдання 3}\end {center}}\par
Що є характерним для структурного програмування?\par
\vspace{2mm}



(\,1\,) У коді мають використовуватись структурні типи даних. %bad



(\,2\,) Використання механізму винятків. %bad



(\,3\,) Використання оператора безумовного переходу. %bad



(\,4\,) Програма зображується ієрархічною структурою блоків.%good


{\begin {center}\textbf{Завдання 4}\end {center}}\par
Вибрати всі істинні твердження (якщо існують).\par
\vspace{2mm}



(\,1\,) У прогнозній моделі розробки розробка може виконуватись ітеративно. %good



(\,2\,) І прогнозні, і адаптивні методології передбачають планування виконуваних робіт. %good



(\,3\,) Прогнозна модель розробки може одночасно бути інкрементною. %good



(\,4\,) Адаптивність методології проявляється в тому, що за зміни вимог розробка починається заново. %bad


{\begin {center}\textbf{Завдання 5}\end {center}}\par
Вибрати всі істинні твердження (якщо існують).\par
\vspace{2mm}



(\,1\,) Одним з прийомів спрощення складної умовної логіки є використання шаблону Strategy. %good



(\,2\,) Чим менше класів у коді, тим менше класів треба налагоджувати, а тому класи за можливості слід об'єднувати. %bad



(\,3\,) Складні комбінації умов починають керувати логікою виконання, тому що такими є вимоги до функціонування ПЗ. Оскільки ПЗ має якнайкраще забезпечувати потреби замовника, то складні комбінації умов у коді вважатись недоліком коду не можуть. %bad



(\,4\,) Чим менше модулів у проекті, тим краще, бо меншу кількість модулів треба інтегрувати в єдине ціле. %bad

\newpage

{\begin {center}\textbf{Завдання 6}\end {center}}\par
Вибрати всі істинні твердження (якщо існують).\par
\vspace{2mm}



(\,1\,) Основна відмінність V-моделі від водоспадної полягає в тім, що діяльності з тестування виконуються паралельно з усіма іншими діяльностями з розробки. %good



(\,2\,) RUP є прогнозною методологією. %bad



(\,3\,) Якщо виникають проблеми з модифікацією коду, то слід розглянути внесення змін у проект. %good



(\,4\,) Багато хто з тих, що використовує Scrum, складає архітектурну модель розроблюваної системи. %good


{\begin {center}\textbf{Завдання 7}\end {center}}\par
Вибрати всі істинні твердження (якщо існують).\par
\vspace{2mm}



(\,1\,) Сьогодні водоспадна методологія не має права на існування. %bad



(\,2\,) Методологія розробки Big Design Up Front погана тим, що може призводити до недостатнього проектування.%bad



(\,3\,) За використання гнучких методологій кожен розробник робить, що хоче. %bad



(\,4\,) Тільки RUP використовує UML.%bad


{\begin {center}\textbf{Завдання 8}\end {center}}\par
Що з наступного не має відношення до Scrum або не виконується/не використовується за використання оригінального Scrum?\par
\vspace{2mm}



(\,1\,) виділення в команді одного чи кількох проектувальників%good



(\,2\,) команда визначає, які задачі вона буде виконувати на наступному спрінті%good



(\,3\,) залучення зовнішніх фахівців для побудови архітектури системи%good



(\,4\,) порядок виконання запланованої на спрінт роботи визначає скрам-мастер%good


{\begin {center}\textbf{Завдання 9}\end {center}}\par
Що з наведеного нижче не розглядалось та навіть не згадувалось у курсі?\par
\vspace{2mm}



(\,1\,) спіральна модель%bad



(\,2\,) RUP%bad



(\,3\,) LeSS%bad



(\,4\,) PRINCE2%good


{\begin {center}\textbf{Завдання 10}\end {center}}\par
Вибрати всі істинні твердження (якщо існують).\par
\vspace{2mm}



(\,1\,) Для тестування код треба ізолювати. Для того, щоб ізолювати, треба мати тести.%good



(\,2\,) Якість можна визначати через дефекти.%good



(\,3\,) Якість можна визначати через атрибути якості.%good



(\,4\,) Розглядають модель якості продукту.%good



\end {large}}
{\vspace{5mm}\smallskip \begin {small} Затверджено на засіданні кафедри теоретичної кібернетики, протокол \No 6 від   29.01.2021 р.\par\smallskip\smallskip
{\parbox {6.5 cm}{Зав. кафедри \dotfill Ю.В.~Крак}}\hspace {0.7cm}{\hbox to 6.5 cm{ Екзаменатор \dotfill Т.О.~Карнаух}} \end{small}}
\newpage

{{\begin {small}\par
{ \center  Київський національний університет імені Тараса Шевченка \par}
{\parbox {5 cm}{Освітня програма \hfill}{\parbox {7 cm}{Інформатика \hfill}}\parbox {6 cm}{\hfill Семестр 6}}\par
{\parbox {5 cm} {Навчальний предмет \hfill}\parbox {7 cm}{Розробка програмного забезпечення \hfill}\parbox {6cm}{\hfill}\par}\vspace{-7pt} \end{small}}{\center Екзаменаційний білет \No \textbf
{ 61 }\par}}
{
\begin {large}
{\begin {center}\textbf{Завдання 1}\end {center}}\par
Що з наведеного найкраще характеризує рефакторінг?\par
\vspace{2mm}



(\,1\,) Зміна внутрішньої структури коду під час усунення помилок%bad



(\,2\,) Довільна переробка коду, що виникає у процесі розробки%bad



(\,3\,) Зміна внутрішньої структури коду, що не змінює його видиму поведінку%good



(\,4\,) Докорінна переробка всього коду з нуля,  що не змінює його видиму поведінку%bad


{\begin {center}\textbf{Завдання 2}\end {center}}\par
Який з наведених фрагментів коду явно потребує поліпшення (форматування рядків та типи до уваги не брати)?\par
\vspace{2mm}



(\,1\,) void ff(…);%bad



(\,2\,) \{int xtrmprg; …\}%bad



(\,3\,) \{double the\_best\_value\_anyone\_ever\_have;\} %bad



(\,4\,) void f13(…);%bad


{\begin {center}\textbf{Завдання 3}\end {center}}\par
Що є характерним для структурного програмування?\par
\vspace{2mm}



(\,1\,) Заборона використання скалярних (простих) типів даних. %bad



(\,2\,) Програма зображується ієрархічною структурою блоків.%good



(\,3\,) У коді мають використовуватись типи структур, класи використовувати заборонено. %bad



(\,4\,) У коді мають використовуватись типи даних, що є класами. %bad


{\begin {center}\textbf{Завдання 4}\end {center}}\par
Вибрати всі істинні твердження (якщо існують).\par
\vspace{2mm}



(\,1\,) Прогнозні моделі розробки з самого початку фіксують вимоги до продукту. %good



(\,2\,) Прогнозні та адаптивні методології використовують одні ті самі діяльності з розробки ПЗ: визначення та аналіз вимог, проектування, реалізація, тестування, впровадження, супроводження. %good



(\,3\,) Прогнозна модель розробки може одночасно бути інкрементною. %good



(\,4\,) Адаптивні методології має сенс використовувати для проектів, де можлива зміна вимог. %good


{\begin {center}\textbf{Завдання 5}\end {center}}\par
Вибрати всі істинні твердження (якщо існують).\par
\vspace{2mm}



(\,1\,) Залежності всередині модуля можуть бути сильними. %good



(\,2\,) Повторюваний код має повторюватись у точності до символу, інакше він повторюваним вважатись не може. %bad



(\,3\,) Повторюваний код дозволяє програмісту краще зрозуміти його логіку та вибрати з кількох варіантів найкращий. %bad



(\,4\,) Гарний код має читатись як розмовна мова. %good

\newpage

{\begin {center}\textbf{Завдання 6}\end {center}}\par
Вибрати всі істинні твердження (якщо існують).\par
\vspace{2mm}



(\,1\,) Розповіді користувача (user stories) – це прогресивний інструмент гнучких методологій, який нічого спільного не має з прецедентами (use case). %bad



(\,2\,) Простота – мистецтво максимізації незробленої роботи. %good



(\,3\,) Ітеративна та інкрементна розробка передбачає наявність у тому чи іншому вигляді етапу ініціалізації, ітерацій та керуючої таблиці проекту. %good



(\,4\,) Гнучкі методології базуються на комунікації розробників, а тому передбачають виключно колективне володіння кодом. %bad


{\begin {center}\textbf{Завдання 7}\end {center}}\par
Вибрати всі істинні твердження (якщо існують).\par
\vspace{2mm}



(\,1\,) Водоспадна методологія може мати тенденцію до надлишкового проектування. %good



(\,2\,) RUP відрізняється від Scrum тим, що ні в що не ставить персонал.%bad



(\,3\,) Гнучкі методології передбачають часті випуски продукту. %good



(\,4\,) За програмування в парах (pair programming) поки один колега набирає код на комп'ютері інший може піти пограти в настільний теніс. %bad


{\begin {center}\textbf{Завдання 8}\end {center}}\par
Що з наступного не має відношення до Scrum або не виконується/не використовується за використання оригінального Scrum?\par
\vspace{2mm}



(\,1\,) побудова різнопланових моделей системи%good



(\,2\,) планування спрінту%bad



(\,3\,) виділення в команді одного чи кількох тестувальників%good



(\,4\,) бачення продукту (vision) %bad


{\begin {center}\textbf{Завдання 9}\end {center}}\par
Що з наведеного нижче не розглядалось та навіть не згадувалось у курсі?\par
\vspace{2mm}



(\,1\,) водоспадна модель%bad



(\,2\,) TDD%bad



(\,3\,) X-модель%good



(\,4\,) LeSS%bad


{\begin {center}\textbf{Завдання 10}\end {center}}\par
Вибрати всі істинні твердження (якщо існують).\par
\vspace{2mm}



(\,1\,) Для тестування код треба ізолювати. Для того, щоб ізолювати, треба мати тести.%good



(\,2\,) І верифікація, і валідація можуть використовувати тестування.%good



(\,3\,) Тестування вимог може породжувати їх уточнення.%good



(\,4\,) Відсутність знайдених дефектів на початкових стадіях розробки може бути свідченням поганого покриття коду тестами.%good



\end {large}}
{\vspace{5mm}\smallskip \begin {small} Затверджено на засіданні кафедри теоретичної кібернетики, протокол \No 6 від   29.01.2021 р.\par\smallskip\smallskip
{\parbox {6.5 cm}{Зав. кафедри \dotfill Ю.В.~Крак}}\hspace {0.7cm}{\hbox to 6.5 cm{ Екзаменатор \dotfill Т.О.~Карнаух}} \end{small}}
\newpage

{{\begin {small}\par
{ \center  Київський національний університет імені Тараса Шевченка \par}
{\parbox {5 cm}{Освітня програма \hfill}{\parbox {7 cm}{Інформатика \hfill}}\parbox {6 cm}{\hfill Семестр 6}}\par
{\parbox {5 cm} {Навчальний предмет \hfill}\parbox {7 cm}{Розробка програмного забезпечення \hfill}\parbox {6cm}{\hfill}\par}\vspace{-7pt} \end{small}}{\center Екзаменаційний білет \No \textbf
{ 62 }\par}}
{
\begin {large}
{\begin {center}\textbf{Завдання 1}\end {center}}\par
Так звані трикутники балансування зв'язують між собою:\par
\vspace{2mm}



(\,1\,) 4 або 3 параметри, при цьому може бути відсутній час%bad



(\,2\,) 4 або 3 параметри, при цьому може бути відсутній обсяг %good



(\,3\,) Завжди рівно 4 параметри%bad



(\,4\,) 4 або 3 параметри, при цьому може бути відсутня якість%bad


{\begin {center}\textbf{Завдання 2}\end {center}}\par
Який з наведених фрагментів коду явно потребує поліпшення (форматування рядків та типи до уваги не брати)?\par
\vspace{2mm}



(\,1\,) \{fstream input\_file, output\_file;…\} %good



(\,2\,) \{double avg\_temperature;…\} %good



(\,3\,) \{int input\_file, file\_to\_output; …\}%bad



(\,4\,) void load\_data(…);%good


{\begin {center}\textbf{Завдання 3}\end {center}}\par
Що є характерним для структурного програмування?\par
\vspace{2mm}



(\,1\,) Використання механізму винятків. %bad



(\,2\,) Кожен блок коду повинен мати єдиний вхід та єдиний вихід.%good



(\,3\,) У коді мають використовуватись типи структур, класи використовувати заборонено. %bad



(\,4\,) У коді мають використовуватись типи структур. %bad


{\begin {center}\textbf{Завдання 4}\end {center}}\par
Вибрати всі істинні твердження (якщо існують).\par
\vspace{2mm}



(\,1\,) Прогнозна модель розробки не може одночасно бути інкрементною. %bad



(\,2\,) В адаптивних методологіях планування робіт не виконується. %bad



(\,3\,) Неперервне проектування створює специфікацію проекту перед початком розробки/ітерації, а під час ітерації його неперервно модифікує.%bad



(\,4\,) Неперервне проектування здійснює створення та модифікацію проекту системи в міру розробки. %good


{\begin {center}\textbf{Завдання 5}\end {center}}\par
Вибрати всі істинні твердження (якщо існують).\par
\vspace{2mm}



(\,1\,) Довжина функцій не впливає на внутрішню якість коду. %bad



(\,2\,) Одним з прийомів спрощення складних умов є їх декомпозиція. %good



(\,3\,) Довгі функції доволі часто приховують складну логіку обробки або повторюваний код. %good



(\,4\,) Від довгих функцій можна позбутися, розклавши їх у послідовність викликів дрібніших функцій. %good

\newpage

{\begin {center}\textbf{Завдання 6}\end {center}}\par
Вибрати всі істинні твердження (якщо існують).\par
\vspace{2mm}



(\,1\,) Водоспадна методологія може мати тенденцію до надлишкового проектування. %good



(\,2\,) За використання гнучких методологій кожен розробник робить, що хоче. %bad



(\,3\,) Гнучкі методології розробки не допускають використання UML. %bad



(\,4\,) RUP є компонентно-орієнтованим. %good


{\begin {center}\textbf{Завдання 7}\end {center}}\par
Вибрати всі істинні твердження (якщо існують).\par
\vspace{2mm}



(\,1\,) Програмування в парах (pair programming) підвищує якість коду та дозволяє економити на виправленні помилок. %good



(\,2\,) Використання RUP не заперечує часті випуски продукту, але й не наполягає на них. %good



(\,3\,) Багато хто з тих, що використовує Scrum, складає архітектурну модель розроблюваної системи. %good



(\,4\,) Адаптивні методології розробки дозволяють уникнути надлишкового проектування. %good


{\begin {center}\textbf{Завдання 8}\end {center}}\par
Що з наступного не має відношення до Scrum або не виконується/не використовується за використання оригінального Scrum?\par
\vspace{2mm}



(\,1\,) під час спрінту дозволяється внесення змін до вимог, що вже реалізуються на поточному спрінті%good



(\,2\,) покер планування (Planning Poker)%bad



(\,3\,) розповіді користувача%bad



(\,4\,) ретроспектива спрінту%bad


{\begin {center}\textbf{Завдання 9}\end {center}}\par
Що з наведеного нижче не розглядалось та навіть не згадувалось у курсі?\par
\vspace{2mm}



(\,1\,) LOAN%good



(\,2\,) BDD%bad



(\,3\,) Cleanroom%bad



(\,4\,) CRAZY%good


{\begin {center}\textbf{Завдання 10}\end {center}}\par
Вибрати всі істинні твердження (якщо існують).\par
\vspace{2mm}



(\,1\,) Користувачі бувають основними.%good



(\,2\,) Якість буває внутрішньою.%good



(\,3\,) Атрибути якості можна використовувати при уточненні вимог.%good



(\,4\,) Якість можна визначати через атрибути якості.%good



\end {large}}
{\vspace{5mm}\smallskip \begin {small} Затверджено на засіданні кафедри теоретичної кібернетики, протокол \No 6 від   29.01.2021 р.\par\smallskip\smallskip
{\parbox {6.5 cm}{Зав. кафедри \dotfill Ю.В.~Крак}}\hspace {0.7cm}{\hbox to 6.5 cm{ Екзаменатор \dotfill Т.О.~Карнаух}} \end{small}}
\newpage

{{\begin {small}\par
{ \center  Київський національний університет імені Тараса Шевченка \par}
{\parbox {5 cm}{Освітня програма \hfill}{\parbox {7 cm}{Інформатика \hfill}}\parbox {6 cm}{\hfill Семестр 6}}\par
{\parbox {5 cm} {Навчальний предмет \hfill}\parbox {7 cm}{Розробка програмного забезпечення \hfill}\parbox {6cm}{\hfill}\par}\vspace{-7pt} \end{small}}{\center Екзаменаційний білет \No \textbf
{ 63 }\par}}
{
\begin {large}
{\begin {center}\textbf{Завдання 1}\end {center}}\par
Шаблон проектування:\par
\vspace{2mm}



(\,1\,) Є шаблоном, в який треба підставити власні ідентифікатори%bad



(\,2\,) Тільки формулює ідею розв'язання проблеми%bad



(\,3\,) Задає розв'язання типової проблеми до найдрібніших подробиць%bad



(\,4\,) Формулює ідею розв'язання проблеми та може пропонувати одну чи кілька її реалізацій%good


{\begin {center}\textbf{Завдання 2}\end {center}}\par
Який з наведених фрагментів коду явно потребує поліпшення (форматування рядків та типи до уваги не брати)?\par
\vspace{2mm}



(\,1\,) \{fstream inputFile, outputFile;…\} %good



(\,2\,) \{double max\_power, min\_power;…\} %good



(\,3\,) void help();%good



(\,4\,) \{int x; int y; …\}%good


{\begin {center}\textbf{Завдання 3}\end {center}}\par
Що є характерним для структурного програмування?\par
\vspace{2mm}



(\,1\,) Заборона використання скалярних (простих) типів даних. %bad



(\,2\,) Використання оператора безумовного переходу. %bad



(\,3\,) Переходи з середини одного блоку в середину іншого неможливі. %good



(\,4\,) У коді мають використовуватись структурні типи даних. %bad


{\begin {center}\textbf{Завдання 4}\end {center}}\par
Вибрати всі істинні твердження (якщо існують).\par
\vspace{2mm}



(\,1\,) І прогнозні, і адаптивні методології передбачають планування виконуваних робіт. %good



(\,2\,) Адаптивність методології проявляється в тому, що за зміни вимог розробка починається заново. %bad



(\,3\,) Гнучкі методології не передбачають планування робіт. %bad



(\,4\,) Для адаптивних моделей розробки оцінки часу та бюджету не виконують, бо вимоги є змінними. %bad


{\begin {center}\textbf{Завдання 5}\end {center}}\par
Вибрати всі істинні твердження (якщо існують).\par
\vspace{2mm}



(\,1\,) Одним з прийомів спрощення умовної логіки є використання поліморфізму. %good



(\,2\,) Одним з засобів позбутись повторень у коді є використання успадкування. %good



(\,3\,) Зв'язки між окремими модулями системи мають бути якомога більш слабкими. %good



(\,4\,) Залежності між різними модулями системи мають бути якомога більш сильними, бо інакше вона розпадеться. %bad

\newpage

{\begin {center}\textbf{Завдання 6}\end {center}}\par
Вибрати всі істинні твердження (якщо існують).\par
\vspace{2mm}



(\,1\,) Водоспадна методологія розробки може призводити до надлишкового проектування. %good



(\,2\,) Ітеративна та інкрементна розробка передбачає наявність у тому чи іншому вигляді етапу ініціалізації, ітерацій та керуючої таблиці проекту. %good



(\,3\,) Без рефакторінгу неперервне проектування неможливе. %good



(\,4\,) Поступове уточнення вимог замовника полегшує адаптацію до змін. %good


{\begin {center}\textbf{Завдання 7}\end {center}}\par
Вибрати всі істинні твердження (якщо існують).\par
\vspace{2mm}



(\,1\,) Програмування в парах (pair programming) спрощує введення в команду нових розробників. %good



(\,2\,) Одна з принципових відмінностей RUP від Scrum полягає у використанні численних проміжних технічних артефактів. %good



(\,3\,) За використання гнучких методологій розробки припиняти розробку не рекомендується. %good



(\,4\,) Розповіді користувача (user stories) це менш формалізований порівняно з прецедентами спосіб опису вимог до системи. %good


{\begin {center}\textbf{Завдання 8}\end {center}}\par
Що з наступного не має відношення до Scrum або не виконується/не використовується за використання оригінального Scrum?\par
\vspace{2mm}



(\,1\,) обернені вимоги%bad



(\,2\,) приведення у належний вигляд журналу невиконаних задач (backlog grooming) %bad



(\,3\,) скрам-мастер визначає, які задачі команда буде виконувати на наступному спрінті%good



(\,4\,) блочні тести%bad


{\begin {center}\textbf{Завдання 9}\end {center}}\par
Що з наведеного нижче не розглядалось та навіть не згадувалось у курсі?\par
\vspace{2mm}



(\,1\,) RAD%bad



(\,2\,) poker%bad



(\,3\,) SAFe%bad



(\,4\,) Scrum%bad


{\begin {center}\textbf{Завдання 10}\end {center}}\par
Вибрати всі істинні твердження (якщо існують).\par
\vspace{2mm}



(\,1\,) Коли нема можливості перевірити всі можливі комбінації, тестують їх попарні (pairwise) комбінації. Це може значно зменшити кількість тестів.%good



(\,2\,) За провалених димових тестів решта тестів не виконується і код повертається на доробку.%good



(\,3\,) Модель якості даних є ієрархічною.%good



(\,4\,) Тестування буває позитивне та негативне.%good



\end {large}}
{\vspace{5mm}\smallskip \begin {small} Затверджено на засіданні кафедри теоретичної кібернетики, протокол \No 6 від   29.01.2021 р.\par\smallskip\smallskip
{\parbox {6.5 cm}{Зав. кафедри \dotfill Ю.В.~Крак}}\hspace {0.7cm}{\hbox to 6.5 cm{ Екзаменатор \dotfill Т.О.~Карнаух}} \end{small}}
\newpage

{{\begin {small}\par
{ \center  Київський національний університет імені Тараса Шевченка \par}
{\parbox {5 cm}{Освітня програма \hfill}{\parbox {7 cm}{Інформатика \hfill}}\parbox {6 cm}{\hfill Семестр 6}}\par
{\parbox {5 cm} {Навчальний предмет \hfill}\parbox {7 cm}{Розробка програмного забезпечення \hfill}\parbox {6cm}{\hfill}\par}\vspace{-7pt} \end{small}}{\center Екзаменаційний білет \No \textbf
{ 64 }\par}}
{
\begin {large}
{\begin {center}\textbf{Завдання 1}\end {center}}\par
Що є метою рефакторингу?\par
\vspace{2mm}



(\,1\,) Поліпшення внутрішньої структури коду%good



(\,2\,) Зменшення кількості рядків коду%bad



(\,3\,) Додавання нової функціональності%bad



(\,4\,) Усунення помилок у коді%bad


{\begin {center}\textbf{Завдання 2}\end {center}}\par
Який з наведених фрагментів коду явно потребує поліпшення (форматування рядків та типи до уваги не брати)?\par
\vspace{2mm}



(\,1\,) void load\_data(T \& arg1, double arg2, double arg3, double arg4); %bad



(\,2\,) \{double maxPower, minPower;…\} %good



(\,3\,) \{int a\_, a\_\_; …\}%bad



(\,4\,) \{fstream input\_file, outputFile; …\}%bad


{\begin {center}\textbf{Завдання 3}\end {center}}\par
Що є характерним для структурного програмування?\par
\vspace{2mm}



(\,1\,) Кожен блок коду повинен мати єдиний вхід та єдиний вихід.%good



(\,2\,) У коді мають використовуватись структурні типи даних. %bad



(\,3\,) Використання оператора безумовного переходу. %bad



(\,4\,) У коді мають використовуватись типи структур, класи використовувати заборонено. %bad


{\begin {center}\textbf{Завдання 4}\end {center}}\par
Вибрати всі істинні твердження (якщо існують).\par
\vspace{2mm}



(\,1\,) Адаптивна модель розробки адаптує вимоги до продукту до наявної команди розробників. %bad



(\,2\,) Адаптивні моделі розробки не допускають зміну вимог до продукту, але продукт може розроблятись поетапно. %bad



(\,3\,) Ітеративні та інкрементні методи розробки є синонімом неперервного проектування. %bad



(\,4\,) У прогнозній моделі розробки розробка може виконуватись ітеративно. %good


{\begin {center}\textbf{Завдання 5}\end {center}}\par
Вибрати всі істинні твердження (якщо існують).\par
\vspace{2mm}



(\,1\,) Одним з прийомів спрощення складної умовної логіки є використання шаблону State. %good



(\,2\,) Код програми пишеться для комп'ютера, тому він має бути синтаксично та семантично коректним. Інше значення не має. %bad



(\,3\,) Код, що виконує однакову послідовність логічних кроків, у більшості випадків є повторюваним.%good



(\,4\,) Код з довгими функціями має меншу кількість функцій, а тому його швидше та простіше тестувати. %bad

\newpage

{\begin {center}\textbf{Завдання 6}\end {center}}\par
Вибрати всі істинні твердження (якщо існують).\par
\vspace{2mm}



(\,1\,) Використання гнучких методологій при розробці надвеликих проектів не несе в собі жодних ризиків. %bad



(\,2\,) Розробка за гнучкими методологіями ніколи не слідує плану. %bad



(\,3\,) Тільки RUP використовує UML.%bad



(\,4\,) Розробка за гнучкими методологіями цінить людей та взаємодію між ними. %good


{\begin {center}\textbf{Завдання 7}\end {center}}\par
Вибрати всі істинні твердження (якщо існують).\par
\vspace{2mm}



(\,1\,) Водоспадна методологія не призначена для внесення змін у вимоги до системи. %good



(\,2\,) Розповіді користувача (user stories) – це прогресивний інструмент гнучких методологій, який нічого спільного не має з прецедентами (use case). %bad



(\,3\,) Перша версія коду не завжди може бути якісною -- і це нормально. %good



(\,4\,) Scrum передбачає колективне володіння кодом. %good


{\begin {center}\textbf{Завдання 8}\end {center}}\par
Що з наступного не має відношення до Scrum або не виконується/не використовується за використання оригінального Scrum?\par
\vspace{2mm}



(\,1\,) порядок виконання запланованої на спрінт роботи визначають розробники%bad



(\,2\,) порядок виконання запланованої на спрінт роботи визначає власник продукту%good



(\,3\,) журнал невиконаних задач (backlog) продукту%bad



(\,4\,) канбан%bad


{\begin {center}\textbf{Завдання 9}\end {center}}\par
Що з наведеного нижче не розглядалось та навіть не згадувалось у курсі?\par
\vspace{2mm}



(\,1\,) sashimi%bad



(\,2\,) PRINCE2%good



(\,3\,) FDD%bad



(\,4\,) DSDM%bad


{\begin {center}\textbf{Завдання 10}\end {center}}\par
Вибрати всі істинні твердження (якщо існують).\par
\vspace{2mm}



(\,1\,) Тестовий приклад складається з набору вхідних даних, умов виконання та очікуваних результатів.%good



(\,2\,) Користувачі бувають непрямими.%good



(\,3\,) Модель якості при використанні є ієрархічною.%good



(\,4\,) Модель якості продукту є ієрархічною.%good



\end {large}}
{\vspace{5mm}\smallskip \begin {small} Затверджено на засіданні кафедри теоретичної кібернетики, протокол \No 6 від   29.01.2021 р.\par\smallskip\smallskip
{\parbox {6.5 cm}{Зав. кафедри \dotfill Ю.В.~Крак}}\hspace {0.7cm}{\hbox to 6.5 cm{ Екзаменатор \dotfill Т.О.~Карнаух}} \end{small}}
\newpage

{{\begin {small}\par
{ \center  Київський національний університет імені Тараса Шевченка \par}
{\parbox {5 cm}{Освітня програма \hfill}{\parbox {7 cm}{Інформатика \hfill}}\parbox {6 cm}{\hfill Семестр 6}}\par
{\parbox {5 cm} {Навчальний предмет \hfill}\parbox {7 cm}{Розробка програмного забезпечення \hfill}\parbox {6cm}{\hfill}\par}\vspace{-7pt} \end{small}}{\center Екзаменаційний білет \No \textbf
{ 65 }\par}}
{
\begin {large}
{\begin {center}\textbf{Завдання 1}\end {center}}\par
Так звані трикутники балансування зв'язують між собою:\par
\vspace{2mm}



(\,1\,) 4 або 3 параметри, при цьому може бути відсутня якість%bad



(\,2\,) Завжди рівно 4 параметри%bad



(\,3\,) 4 або 3 параметри, при цьому може бути відсутній обсяг %good



(\,4\,) 4 або 3 параметри, при цьому може бути відсутній час%bad


{\begin {center}\textbf{Завдання 2}\end {center}}\par
Який з наведених фрагментів коду явно потребує поліпшення (форматування рядків та типи до уваги не брати)?\par
\vspace{2mm}



(\,1\,) void load\_data(…);%good



(\,2\,) \{fstream input\_file, output\_file;…\} %good



(\,3\,) \{double avg\_temperature;…\} %good



(\,4\,) void load\_data(T \& arg1, double arg2, double arg3, double arg4); %bad


{\begin {center}\textbf{Завдання 3}\end {center}}\par
Що є характерним для структурного програмування?\par
\vspace{2mm}



(\,1\,) У коді мають використовуватись типи структур. %bad



(\,2\,) У коді мають використовуватись типи даних, що є класами. %bad



(\,3\,) Заборона використання скалярних (простих) типів даних. %bad



(\,4\,) Програма зображується ієрархічною структурою блоків.%good


{\begin {center}\textbf{Завдання 4}\end {center}}\par
Вибрати всі істинні твердження (якщо існують).\par
\vspace{2mm}



(\,1\,) Адаптивні методології не передбачають визначення вимог. %bad



(\,2\,) У випадку, коли вимоги до продукту відомі заздалегідь та є фіксованими контрактом, адаптивні методології розробки використовувати не можна. %bad



(\,3\,) Неперервне проектування має бути ітеративним та інкрементним одночасно. %good



(\,4\,) Адаптивні моделі розробки допускають зміну вимог до продукту. %good


{\begin {center}\textbf{Завдання 5}\end {center}}\par
Вибрати всі істинні твердження (якщо існують).\par
\vspace{2mm}



(\,1\,) Гарний код не повинен містити різних фрагментів, що складаються з кількох інструкцій та виконують однакове перетворення даних. %good



(\,2\,) Правильно спроектована функція повинна мати рівно один обов'язок. %good



(\,3\,) Якщо кожна функція бібліотеки буде здатна в залежності від параметрів виконувати кілька різних дій, то загальна кількість функцій у ній буде менша і її буде простіше використовувати. %bad



(\,4\,) Складні умови виникають в реальних задачах доволі часто, а тому є доволі природніми і не можуть вважатись недоліком коду. %bad

\newpage

{\begin {center}\textbf{Завдання 6}\end {center}}\par
Вибрати всі істинні твердження (якщо існують).\par
\vspace{2mm}



(\,1\,) Сьогодні водоспадна методологія не має права на існування. %bad



(\,2\,) RUP є керованим прецедентами. %good



(\,3\,) Гнучкі методології повністю нехтують документацією. %bad



(\,4\,) Якщо виникають проблеми з модифікацією коду, то слід відразу змінити команду розробників на більш кваліфіковану.%bad


{\begin {center}\textbf{Завдання 7}\end {center}}\par
Вибрати всі істинні твердження (якщо існують).\par
\vspace{2mm}



(\,1\,) Висока внутрішня якість коду здешевшує розробку, бо дозволяє економити на вартості та часі майбутніх змін. %good



(\,2\,) Гнучкі методології базуються на комунікації розробників, а тому передбачають виключно колективне володіння кодом. %bad



(\,3\,) За відмови від документації одним з методом збереження знань про систему є постійне пересування персоналу між її частинами. %good



(\,4\,) Гнучкі методи розробки нехтують проміжними технічними артефактами та документацією через те, що за зміни вимог та/або проекту треба витрачати додатковий час на підтримання їх у відповідному стані. %good


{\begin {center}\textbf{Завдання 8}\end {center}}\par
Що з наступного не має відношення до Scrum або не виконується/не використовується за використання оригінального Scrum?\par
\vspace{2mm}



(\,1\,) розповіді користувача%bad



(\,2\,) власник продукту призначає час, необхідний команді для реалізації конкретної задачі%good



(\,3\,) правила визначення готовності (Definition of Done) %bad



(\,4\,) скрам-мастер визначає час, необхідний команді для реалізації конкретної задачі%good


{\begin {center}\textbf{Завдання 9}\end {center}}\par
Що з наведеного нижче не розглядалось та навіть не згадувалось у курсі?\par
\vspace{2mm}



(\,1\,) LAZY%good



(\,2\,) Lean%bad



(\,3\,) BDUF%bad



(\,4\,) XP%bad


{\begin {center}\textbf{Завдання 10}\end {center}}\par
Вибрати всі істинні твердження (якщо існують).\par
\vspace{2mm}



(\,1\,) Можна розглядати якість під час використання.%good



(\,2\,) Розглядають модель якості продукту.%good



(\,3\,) Розглядають модель якості при використанні.%good



(\,4\,) Якість буває зовнішньою.%good



\end {large}}
{\vspace{5mm}\smallskip \begin {small} Затверджено на засіданні кафедри теоретичної кібернетики, протокол \No 6 від   29.01.2021 р.\par\smallskip\smallskip
{\parbox {6.5 cm}{Зав. кафедри \dotfill Ю.В.~Крак}}\hspace {0.7cm}{\hbox to 6.5 cm{ Екзаменатор \dotfill Т.О.~Карнаух}} \end{small}}
\newpage

{{\begin {small}\par
{ \center  Київський національний університет імені Тараса Шевченка \par}
{\parbox {5 cm}{Освітня програма \hfill}{\parbox {7 cm}{Інформатика \hfill}}\parbox {6 cm}{\hfill Семестр 6}}\par
{\parbox {5 cm} {Навчальний предмет \hfill}\parbox {7 cm}{Розробка програмного забезпечення \hfill}\parbox {6cm}{\hfill}\par}\vspace{-7pt} \end{small}}{\center Екзаменаційний білет \No \textbf
{ 66 }\par}}
{
\begin {large}
{\begin {center}\textbf{Завдання 1}\end {center}}\par
Що з наведеного найкраще характеризує рефакторінг?\par
\vspace{2mm}



(\,1\,) Зміна внутрішньої структури коду, що не змінює його видиму поведінку%good



(\,2\,) Зміна внутрішньої структури коду під час усунення помилок%bad



(\,3\,) Докорінна переробка всього коду з нуля,  що не змінює його видиму поведінку%bad



(\,4\,) Довільна переробка коду, що виникає у процесі розробки%bad


{\begin {center}\textbf{Завдання 2}\end {center}}\par
Який з наведених фрагментів коду явно потребує поліпшення (форматування рядків та типи до уваги не брати)?\par
\vspace{2mm}



(\,1\,) \{double maxPower, minPower;…\} %good



(\,2\,) void f13(…);%bad



(\,3\,) \{fstream input\_file, outputFile; …\}%bad



(\,4\,) \{double the\_best\_value\_anyone\_ever\_have;\} %bad


{\begin {center}\textbf{Завдання 3}\end {center}}\par
Що є характерним для структурного програмування?\par
\vspace{2mm}



(\,1\,) Переходи з середини одного блоку в середину іншого неможливі. %good



(\,2\,) Використання механізму винятків. %bad



(\,3\,) У коді мають використовуватись структурні типи даних. %bad



(\,4\,) У коді мають використовуватись типи структур. %bad


{\begin {center}\textbf{Завдання 4}\end {center}}\par
Вибрати всі істинні твердження (якщо існують).\par
\vspace{2mm}



(\,1\,) Неперервність проектування полягає в тім, що процес розробки не можна переривати більше ніж на вихідні. %bad



(\,2\,) Для адаптивних методологій прогнозувати результат чи терміни його отримання абсолютно неможливо. %bad



(\,3\,) У прогнозній моделі розробки розробка не може виконуватись ітеративно. %bad



(\,4\,) Для адаптивних методологій розробки, на відміну від прогнозних, попередні оцінки бюджету та/або часу не виконують, бо вони все одно зміняться. %bad


{\begin {center}\textbf{Завдання 5}\end {center}}\par
Вибрати всі істинні твердження (якщо існують).\par
\vspace{2mm}



(\,1\,) Одним з засобів позбутись повторень у коді є використання шаблону Template Method (але цей шаблон здатен боротися і з іншими недоліками коду). %good



(\,2\,) Довгі функції доволі часто приховують складну логіку обробки або повторюваний код. %good



(\,3\,) Повторюваний код має повторюватись у точності до символу, інакше він повторюваним вважатись не може. %bad



(\,4\,) Одним з прийомів спрощення складних умов є їх декомпозиція. %good

\newpage

{\begin {center}\textbf{Завдання 6}\end {center}}\par
Вибрати всі істинні твердження (якщо існують).\par
\vspace{2mm}



(\,1\,) RUP принципово відрізняється від Scrum тим, що є керованим прецедентами. У той же час у Scrum при плануванні розробки прецеденти (чи їх аналоги) до уваги не беруться. %bad



(\,2\,) Методологія розробки Big Design Up Front погана тим, що може призводити до недостатнього проектування.%bad



(\,3\,) На відміну від Scrum, RUP є компонентно-орієнтованим. %bad



(\,4\,) RUP є архітектурно-орієнтованим. %good


{\begin {center}\textbf{Завдання 7}\end {center}}\par
Вибрати всі істинні твердження (якщо існують).\par
\vspace{2mm}



(\,1\,) Якщо виникають проблеми з модифікацією коду, то слід розглянути внесення змін у проект. %good



(\,2\,) Модель хаосу є методологією розробки програмного забезпечення. %bad



(\,3\,) Ітеративна методологія розробки у загальному випадку не гарантує, що не буде надлишкового проектування. %good



(\,4\,) За програмування в парах (pair programming) поки один колега набирає код на комп'ютері інший може піти пограти в настільний теніс. %bad


{\begin {center}\textbf{Завдання 8}\end {center}}\par
Що з наступного не має відношення до Scrum або не виконується/не використовується за використання оригінального Scrum?\par
\vspace{2mm}



(\,1\,) скрам-мастер остаточно затверджує проект системи%good



(\,2\,) власник продукту призначає відповідальних за конкретні задачі%good



(\,3\,) під час спрінту зміни, що впливають на досягнення мети спрінту, зазвичай, не допускаються%bad



(\,4\,) журнал невиконаних задач (backlog) спрінту%bad


{\begin {center}\textbf{Завдання 9}\end {center}}\par
Що з наведеного нижче не розглядалось та навіть не згадувалось у курсі?\par
\vspace{2mm}



(\,1\,) SAFe%bad



(\,2\,) RUP%bad



(\,3\,) Lean%bad



(\,4\,) V-модель%bad


{\begin {center}\textbf{Завдання 10}\end {center}}\par
Вибрати всі істинні твердження (якщо існують).\par
\vspace{2mm}



(\,1\,) Розглядають модель якості даних.%good



(\,2\,) Користувачі бувають вторинними.%good



(\,3\,) Якість можна визначати через дефекти.%good



(\,4\,) Забезпечення якості це не тільки тестування.%good



\end {large}}
{\vspace{5mm}\smallskip \begin {small} Затверджено на засіданні кафедри теоретичної кібернетики, протокол \No 6 від   29.01.2021 р.\par\smallskip\smallskip
{\parbox {6.5 cm}{Зав. кафедри \dotfill Ю.В.~Крак}}\hspace {0.7cm}{\hbox to 6.5 cm{ Екзаменатор \dotfill Т.О.~Карнаух}} \end{small}}
\newpage

{{\begin {small}\par
{ \center  Київський національний університет імені Тараса Шевченка \par}
{\parbox {5 cm}{Освітня програма \hfill}{\parbox {7 cm}{Інформатика \hfill}}\parbox {6 cm}{\hfill Семестр 6}}\par
{\parbox {5 cm} {Навчальний предмет \hfill}\parbox {7 cm}{Розробка програмного забезпечення \hfill}\parbox {6cm}{\hfill}\par}\vspace{-7pt} \end{small}}{\center Екзаменаційний білет \No \textbf
{ 67 }\par}}
{
\begin {large}
{\begin {center}\textbf{Завдання 1}\end {center}}\par
Так звані трикутники балансування зв'язують між собою:\par
\vspace{2mm}



(\,1\,) Завжди рівно 3 параметри%bad



(\,2\,) 4 або 3 параметри, при цьому може бути відсутній час%bad



(\,3\,) Завжди рівно 4 параметри%bad



(\,4\,) 4 або 3 параметри, при цьому може бути відсутній обсяг %good


{\begin {center}\textbf{Завдання 2}\end {center}}\par
Який з наведених фрагментів коду явно потребує поліпшення (форматування рядків та типи до уваги не брати)?\par
\vspace{2mm}



(\,1\,) \{int xtrmprg; …\}%bad



(\,2\,) void help();%good



(\,3\,) \{int x; int y; …\}%good



(\,4\,) \{fstream inputFile, outputFile;…\} %good


{\begin {center}\textbf{Завдання 3}\end {center}}\par
Що є характерним для структурного програмування?\par
\vspace{2mm}



(\,1\,) У коді мають використовуватись типи структур, класи використовувати заборонено. %bad



(\,2\,) Використання механізму винятків. %bad



(\,3\,) Кожен блок коду повинен мати єдиний вхід та єдиний вихід.%good



(\,4\,) Заборона використання скалярних (простих) типів даних. %bad


{\begin {center}\textbf{Завдання 4}\end {center}}\par
Вибрати всі істинні твердження (якщо існують).\par
\vspace{2mm}



(\,1\,) Адаптивна модель розробки адаптує вимоги до продукту до наявної команди розробників. %bad



(\,2\,) Адаптивність методології проявляється в тому, що за зміни вимог розробка починається заново. %bad



(\,3\,) Адаптивні моделі розробки допускають зміну вимог до продукту. %good



(\,4\,) І прогнозні, і адаптивні методології передбачають планування виконуваних робіт. %good


{\begin {center}\textbf{Завдання 5}\end {center}}\par
Вибрати всі істинні твердження (якщо існують).\par
\vspace{2mm}



(\,1\,) Складні комбінації умов починають керувати логікою виконання, тому що такими є вимоги до функціонування ПЗ. Оскільки ПЗ має якнайкраще забезпечувати потреби замовника, то складні комбінації умов у коді вважатись недоліком коду не можуть. %bad



(\,2\,) Втілення в коді різних підходів до одного того самого перетворення даних дозволяє програмісту продемонструвати свою професійну підготовку з найкращого боку. %bad



(\,3\,) Залежності всередині модуля можуть бути сильними. %good



(\,4\,) Гарний код не повинен містити різних фрагментів, що складаються з кількох інструкцій та виконують однакове перетворення даних. %good

\newpage

{\begin {center}\textbf{Завдання 6}\end {center}}\par
Вибрати всі істинні твердження (якщо існують).\par
\vspace{2mm}



(\,1\,) Програмування в парах (pair programming) зводиться до того, що кожним фрагментом коду володіють два програмісти. %bad



(\,2\,) Гнучкі методології розробки допускають суттєві зміни вимог, проте все одно намагаються якомога швидше зафіксувати архітектуру. %good



(\,3\,) Простота – мистецтво максимізації незробленої роботи. %good



(\,4\,) За застосування адаптивних методологій розробки, так само як і за прогнозних, намагаються якомога швидше визначити архітектуру системи. Відмінність полягає в тім, у який спосіб це виконується.%good


{\begin {center}\textbf{Завдання 7}\end {center}}\par
Вибрати всі істинні твердження (якщо існують).\par
\vspace{2mm}



(\,1\,) Методологія розробки Big Design Up Front погана тим, що може призводити до надлишкового проектування. %good



(\,2\,) Код, що себе сам документує, -- це код, який виводить користувачу документацію на себе. %bad



(\,3\,) Гнучкі методи розробки нехтують проміжними технічними артефактами та документацією виключно через їх непотрібність. %bad



(\,4\,) Розробка за гнучкими методологіями цінить співробітництво з замовником та піддатлівість змінам. %good


{\begin {center}\textbf{Завдання 8}\end {center}}\par
Що з наступного не має відношення до Scrum або не виконується/не використовується за використання оригінального Scrum?\par
\vspace{2mm}



(\,1\,) порядок виконання запланованої на спрінт роботи визначає скрам-мастер%good



(\,2\,) скрам-мастер призначає відповідальних за конкретні задачі%good



(\,3\,) власник продукту призначає час, необхідний команді для реалізації конкретної задачі%good



(\,4\,) власник продукту призначає відповідальних за конкретні задачі%good


{\begin {center}\textbf{Завдання 9}\end {center}}\par
Що з наведеного нижче не розглядалось та навіть не згадувалось у курсі?\par
\vspace{2mm}



(\,1\,) спіральна модель%bad



(\,2\,) LOAN%good



(\,3\,) модель хаосу%bad



(\,4\,) X-модель%good


{\begin {center}\textbf{Завдання 10}\end {center}}\par
Вибрати всі істинні твердження (якщо існують).\par
\vspace{2mm}



(\,1\,) Якість можна визначати через дефекти.%good



(\,2\,) Розглядають модель якості даних.%good



(\,3\,) Коли нема можливості перевірити всі можливі комбінації, тестують їх попарні (pairwise) комбінації. Це може значно зменшити кількість тестів.%good



(\,4\,) Модель якості даних є ієрархічною.%good



\end {large}}
{\vspace{5mm}\smallskip \begin {small} Затверджено на засіданні кафедри теоретичної кібернетики, протокол \No 6 від   29.01.2021 р.\par\smallskip\smallskip
{\parbox {6.5 cm}{Зав. кафедри \dotfill Ю.В.~Крак}}\hspace {0.7cm}{\hbox to 6.5 cm{ Екзаменатор \dotfill Т.О.~Карнаух}} \end{small}}
\newpage

{{\begin {small}\par
{ \center  Київський національний університет імені Тараса Шевченка \par}
{\parbox {5 cm}{Освітня програма \hfill}{\parbox {7 cm}{Інформатика \hfill}}\parbox {6 cm}{\hfill Семестр 6}}\par
{\parbox {5 cm} {Навчальний предмет \hfill}\parbox {7 cm}{Розробка програмного забезпечення \hfill}\parbox {6cm}{\hfill}\par}\vspace{-7pt} \end{small}}{\center Екзаменаційний білет \No \textbf
{ 68 }\par}}
{
\begin {large}
{\begin {center}\textbf{Завдання 1}\end {center}}\par
Шаблон проектування:\par
\vspace{2mm}



(\,1\,) Формулює ідею розв'язання проблеми та може пропонувати одну чи кілька її реалізацій%good



(\,2\,) Задає розв'язання типової проблеми до найдрібніших подробиць%bad



(\,3\,) Тільки формулює ідею розв'язання проблеми%bad



(\,4\,) Є шаблоном, в який треба підставити власні ідентифікатори%bad


{\begin {center}\textbf{Завдання 2}\end {center}}\par
Який з наведених фрагментів коду явно потребує поліпшення (форматування рядків та типи до уваги не брати)?\par
\vspace{2mm}



(\,1\,) void ff(…);%bad



(\,2\,) \{double max\_power, min\_power;…\} %good



(\,3\,) \{int a\_, a\_\_; …\}%bad



(\,4\,) \{int input\_file, file\_to\_output; …\}%bad


{\begin {center}\textbf{Завдання 3}\end {center}}\par
Що є характерним для структурного програмування?\par
\vspace{2mm}



(\,1\,) У коді мають використовуватись типи даних, що є класами. %bad



(\,2\,) Використання оператора безумовного переходу. %bad



(\,3\,) У коді мають використовуватись типи структур, класи використовувати заборонено. %bad



(\,4\,) Програма зображується ієрархічною структурою блоків.%good


{\begin {center}\textbf{Завдання 4}\end {center}}\par
Вибрати всі істинні твердження (якщо існують).\par
\vspace{2mm}



(\,1\,) У випадку, коли вимоги до продукту відомі заздалегідь та є фіксованими контрактом, адаптивні методології розробки використовувати не можна. %bad



(\,2\,) В адаптивних методологіях планування робіт не виконується. %bad



(\,3\,) Прогнозна модель розробки може одночасно бути інкрементною. %good



(\,4\,) У прогнозній моделі розробки розробка не може виконуватись ітеративно. %bad


{\begin {center}\textbf{Завдання 5}\end {center}}\par
Вибрати всі істинні твердження (якщо існують).\par
\vspace{2mm}



(\,1\,) Одним з засобів позбутись повторень у коді є використання успадкування. %good



(\,2\,) Одним з прийомів спрощення умовної логіки є використання поліморфізму. %good



(\,3\,) Одним з прийомів спрощення складної умовної логіки є використання шаблону Strategy. %good



(\,4\,) Одним з засобів позбутись повторень у коді є використання функцій. %good

\newpage

{\begin {center}\textbf{Завдання 6}\end {center}}\par
Вибрати всі істинні твердження (якщо існують).\par
\vspace{2mm}



(\,1\,) Водоспадна методологія розробки може призводити до недостатнього проектування. %bad



(\,2\,) RUP відрізняється від Scrum тим, що ні в що не ставить персонал.%bad



(\,3\,) Основна відмінність V-моделі від водоспадної полягає в тім, що діяльності з тестування виконуються паралельно з усіма іншими діяльностями з розробки. %good



(\,4\,) Розробка за гнучкими методологіями не передбачає узгодження умов контракту. %bad


{\begin {center}\textbf{Завдання 7}\end {center}}\par
Вибрати всі істинні твердження (якщо існують).\par
\vspace{2mm}



(\,1\,) Використання гнучких методологій вимагає високої кваліфікації персоналу. %good



(\,2\,) Висока внутрішня якість коду є марною тратою часу на те, що замовник не побачить та/або не зможе зацінити. %bad



(\,3\,) За невизначених або піддатливих змінам вимог краще застосовувати гнучкі методології розробки. %good



(\,4\,) Гнучкі методології передбачають часті випуски продукту. %good


{\begin {center}\textbf{Завдання 8}\end {center}}\par
Що з наступного не має відношення до Scrum або не виконується/не використовується за використання оригінального Scrum?\par
\vspace{2mm}



(\,1\,) бачення продукту (vision) %bad



(\,2\,) виділення в команді одного чи кількох тестувальників%good



(\,3\,) ретроспектива спрінту%bad



(\,4\,) журнал невиконаних задач (backlog) продукту%bad


{\begin {center}\textbf{Завдання 9}\end {center}}\par
Що з наведеного нижче не розглядалось та навіть не згадувалось у курсі?\par
\vspace{2mm}



(\,1\,) Scrum of scrums%bad



(\,2\,) LAZY%good



(\,3\,) poker%bad



(\,4\,) FDD%bad


{\begin {center}\textbf{Завдання 10}\end {center}}\par
Вибрати всі істинні твердження (якщо існують).\par
\vspace{2mm}



(\,1\,) Можна розглядати якість під час використання.%good



(\,2\,) І верифікація, і валідація можуть використовувати тестування.%good



(\,3\,) Розглядають модель якості при використанні.%good



(\,4\,) Розглядають модель якості продукту.%good



\end {large}}
{\vspace{5mm}\smallskip \begin {small} Затверджено на засіданні кафедри теоретичної кібернетики, протокол \No 6 від   29.01.2021 р.\par\smallskip\smallskip
{\parbox {6.5 cm}{Зав. кафедри \dotfill Ю.В.~Крак}}\hspace {0.7cm}{\hbox to 6.5 cm{ Екзаменатор \dotfill Т.О.~Карнаух}} \end{small}}
\newpage

{{\begin {small}\par
{ \center  Київський національний університет імені Тараса Шевченка \par}
{\parbox {5 cm}{Освітня програма \hfill}{\parbox {7 cm}{Інформатика \hfill}}\parbox {6 cm}{\hfill Семестр 6}}\par
{\parbox {5 cm} {Навчальний предмет \hfill}\parbox {7 cm}{Розробка програмного забезпечення \hfill}\parbox {6cm}{\hfill}\par}\vspace{-7pt} \end{small}}{\center Екзаменаційний білет \No \textbf
{ 69 }\par}}
{
\begin {large}
{\begin {center}\textbf{Завдання 1}\end {center}}\par
Так звані трикутники балансування зв'язують між собою:\par
\vspace{2mm}



(\,1\,) 4 або 3 параметри, при цьому може бути відсутній обсяг %good



(\,2\,) 4 або 3 параметри, при цьому може бути відсутній час%bad



(\,3\,) 4 або 3 параметри, при цьому може бути відсутня якість%bad



(\,4\,) Завжди рівно 3 параметри%bad


{\begin {center}\textbf{Завдання 2}\end {center}}\par
Який з наведених фрагментів коду явно потребує поліпшення (форматування рядків та типи до уваги не брати)?\par
\vspace{2mm}



(\,1\,) void load\_data(T \& arg1, double arg2, double arg3, double arg4); %bad



(\,2\,) \{double maxPower, minPower;…\} %good



(\,3\,) \{double the\_best\_value\_anyone\_ever\_have;\} %bad



(\,4\,) \{fstream input\_file, outputFile; …\}%bad


{\begin {center}\textbf{Завдання 3}\end {center}}\par
Що є характерним для структурного програмування?\par
\vspace{2mm}



(\,1\,) У коді мають використовуватись структурні типи даних. %bad



(\,2\,) У коді мають використовуватись типи структур. %bad



(\,3\,) Заборона використання скалярних (простих) типів даних. %bad



(\,4\,) Переходи з середини одного блоку в середину іншого неможливі. %good


{\begin {center}\textbf{Завдання 4}\end {center}}\par
Вибрати всі істинні твердження (якщо існують).\par
\vspace{2mm}



(\,1\,) Прогнозні моделі розробки з самого початку фіксують вимоги до продукту. %good



(\,2\,) Гнучкі методології не передбачають планування робіт. %bad



(\,3\,) Неперервне проектування здійснює створення та модифікацію проекту системи в міру розробки. %good



(\,4\,) Ітеративні та інкрементні методи розробки є синонімом неперервного проектування. %bad


{\begin {center}\textbf{Завдання 5}\end {center}}\par
Вибрати всі істинні твердження (якщо існують).\par
\vspace{2mm}



(\,1\,) Одним з засобів позбутись повторень у коді є використання шаблону Template Method (але цей шаблон здатен боротися і з іншими недоліками коду). %good



(\,2\,) Якщо кожна функція бібліотеки буде здатна в залежності від параметрів виконувати кілька різних дій, то загальна кількість функцій у ній буде менша і її буде простіше використовувати. %bad



(\,3\,) Від довгих функцій можна позбутися, розклавши їх у послідовність викликів дрібніших функцій. %good



(\,4\,) Гарний код має читатись як розмовна мова. %good

\newpage

{\begin {center}\textbf{Завдання 6}\end {center}}\par
Вибрати всі істинні твердження (якщо існують).\par
\vspace{2mm}



(\,1\,) Гнучкі методології можуть передбачати як спільне, так і індивідуальне володіння кодом. %good



(\,2\,) RUP є ітеративним та інкрементним.%good



(\,3\,) Гнучкі методології розробки є стійкими до припинення розробки на певний час. %bad



(\,4\,) Гнучкі методології розробки намагаються якомога швидше визначити архітектуру системи, для чого на початку завжди виконується початкове проектування системи в цілому, а тільки потім переходять до ітеративного уточнення деталей та кодування. %bad


{\begin {center}\textbf{Завдання 7}\end {center}}\par
Вибрати всі істинні твердження (якщо існують).\par
\vspace{2mm}



(\,1\,) RUP є прогнозною методологією. %bad



(\,2\,) Методологія розробки Big Design Up Front погана тим, що може призводити до недостатнього проектування.%bad



(\,3\,) Гнучкі методології можуть передбачати як спільне, так і індивідуальне володіння кодом. %good



(\,4\,) За невизначених або піддатливих змінам вимог краще застосовувати гнучкі методології розробки. %good


{\begin {center}\textbf{Завдання 8}\end {center}}\par
Що з наступного не має відношення до Scrum або не виконується/не використовується за використання оригінального Scrum?\par
\vspace{2mm}



(\,1\,) побудова різнопланових моделей системи%good



(\,2\,) залучення зовнішніх фахівців для побудови архітектури системи%good



(\,3\,) обернені вимоги%bad



(\,4\,) команда визначає, які задачі вона буде виконувати на наступному спрінті%good


{\begin {center}\textbf{Завдання 9}\end {center}}\par
Що з наведеного нижче не розглядалось та навіть не згадувалось у курсі?\par
\vspace{2mm}



(\,1\,) CRAZY%good



(\,2\,) PRINCE2%good



(\,3\,) спіральна модель%bad



(\,4\,) BDD%bad


{\begin {center}\textbf{Завдання 10}\end {center}}\par
Вибрати всі істинні твердження (якщо існують).\par
\vspace{2mm}



(\,1\,) Якість буває внутрішньою.%good



(\,2\,) Якість можна визначати через атрибути якості.%good



(\,3\,) Користувачі бувають непрямими.%good



(\,4\,) Користувачі бувають основними.%good



\end {large}}
{\vspace{5mm}\smallskip \begin {small} Затверджено на засіданні кафедри теоретичної кібернетики, протокол \No 6 від   29.01.2021 р.\par\smallskip\smallskip
{\parbox {6.5 cm}{Зав. кафедри \dotfill Ю.В.~Крак}}\hspace {0.7cm}{\hbox to 6.5 cm{ Екзаменатор \dotfill Т.О.~Карнаух}} \end{small}}
\newpage

{{\begin {small}\par
{ \center  Київський національний університет імені Тараса Шевченка \par}
{\parbox {5 cm}{Освітня програма \hfill}{\parbox {7 cm}{Інформатика \hfill}}\parbox {6 cm}{\hfill Семестр 6}}\par
{\parbox {5 cm} {Навчальний предмет \hfill}\parbox {7 cm}{Розробка програмного забезпечення \hfill}\parbox {6cm}{\hfill}\par}\vspace{-7pt} \end{small}}{\center Екзаменаційний білет \No \textbf
{ 70 }\par}}
{
\begin {large}
{\begin {center}\textbf{Завдання 1}\end {center}}\par
Так звані трикутники балансування зв'язують між собою:\par
\vspace{2mm}



(\,1\,) 4 або 3 параметри, при цьому може бути відсутній обсяг %good



(\,2\,) Завжди рівно 4 параметри%bad



(\,3\,) Завжди рівно 3 параметри%bad



(\,4\,) 4 або 3 параметри, при цьому може бути відсутня якість%bad


{\begin {center}\textbf{Завдання 2}\end {center}}\par
Який з наведених фрагментів коду явно потребує поліпшення (форматування рядків та типи до уваги не брати)?\par
\vspace{2mm}



(\,1\,) \{double avg\_temperature;…\} %good



(\,2\,) \{int a\_, a\_\_; …\}%bad



(\,3\,) void load\_data(…);%good



(\,4\,) void help();%good


{\begin {center}\textbf{Завдання 3}\end {center}}\par
Що є характерним для структурного програмування?\par
\vspace{2mm}



(\,1\,) Використання оператора безумовного переходу. %bad



(\,2\,) Програма зображується ієрархічною структурою блоків.%good



(\,3\,) Використання механізму винятків. %bad



(\,4\,) У коді мають використовуватись типи даних, що є класами. %bad


{\begin {center}\textbf{Завдання 4}\end {center}}\par
Вибрати всі істинні твердження (якщо існують).\par
\vspace{2mm}



(\,1\,) Неперервність проектування полягає в тім, що процес розробки не можна переривати більше ніж на вихідні. %bad



(\,2\,) Адаптивні методології має сенс використовувати для проектів, де можлива зміна вимог. %good



(\,3\,) Прогнозні та адаптивні методології використовують одні ті самі діяльності з розробки ПЗ: визначення та аналіз вимог, проектування, реалізація, тестування, впровадження, супроводження. %good



(\,4\,) У прогнозній моделі розробки розробка може виконуватись ітеративно. %good


{\begin {center}\textbf{Завдання 5}\end {center}}\par
Вибрати всі істинні твердження (якщо існують).\par
\vspace{2mm}



(\,1\,) Зв'язки між окремими модулями системи мають бути якомога більш слабкими. %good



(\,2\,) Довжина функцій не впливає на внутрішню якість коду. %bad



(\,3\,) Код з довгими функціями має меншу кількість функцій, а тому його швидше та простіше тестувати. %bad



(\,4\,) Чим вища кваліфікація програміста, тим більш складні для розуміння та довгі умовні вирази він використовує. %bad

\newpage

{\begin {center}\textbf{Завдання 6}\end {center}}\par
Вибрати всі істинні твердження (якщо існують).\par
\vspace{2mm}



(\,1\,) Ітеративна методологія розробки у загальному випадку не гарантує, що не буде надлишкового проектування. %good



(\,2\,) RUP відрізняється від Scrum тим, що ні в що не ставить персонал.%bad



(\,3\,) Розповіді користувача (user stories) – це прогресивний інструмент гнучких методологій, який нічого спільного не має з прецедентами (use case). %bad



(\,4\,) Методологія розробки Big Design Up Front погана тим, що може призводити до надлишкового проектування. %good


{\begin {center}\textbf{Завдання 7}\end {center}}\par
Вибрати всі істинні твердження (якщо існують).\par
\vspace{2mm}



(\,1\,) Тільки RUP використовує UML.%bad



(\,2\,) За використання гнучких методологій розробки припиняти розробку не рекомендується. %good



(\,3\,) Розробка за гнучкими методологіями цінить людей та взаємодію між ними. %good



(\,4\,) Програмування в парах (pair programming) підвищує якість коду та дозволяє економити на виправленні помилок. %good


{\begin {center}\textbf{Завдання 8}\end {center}}\par
Що з наступного не має відношення до Scrum або не виконується/не використовується за використання оригінального Scrum?\par
\vspace{2mm}



(\,1\,) порядок виконання запланованої на спрінт роботи визначають розробники%bad



(\,2\,) розповіді користувача%bad



(\,3\,) під час спрінту зміни, що впливають на досягнення мети спрінту, зазвичай, не допускаються%bad



(\,4\,) блочні тести%bad


{\begin {center}\textbf{Завдання 9}\end {center}}\par
Що з наведеного нижче не розглядалось та навіть не згадувалось у курсі?\par
\vspace{2mm}



(\,1\,) Cleanroom%bad



(\,2\,) sashimi%bad



(\,3\,) LeSS%bad



(\,4\,) V-модель%bad


{\begin {center}\textbf{Завдання 10}\end {center}}\par
Вибрати всі істинні твердження (якщо існують).\par
\vspace{2mm}



(\,1\,) Тестовий приклад складається з набору вхідних даних, умов виконання та очікуваних результатів.%good



(\,2\,) Відсутність знайдених дефектів на початкових стадіях розробки може бути свідченням поганого покриття коду тестами.%good



(\,3\,) За провалених димових тестів решта тестів не виконується і код повертається на доробку.%good



(\,4\,) Атрибути якості можна використовувати при уточненні вимог.%good



\end {large}}
{\vspace{5mm}\smallskip \begin {small} Затверджено на засіданні кафедри теоретичної кібернетики, протокол \No 6 від   29.01.2021 р.\par\smallskip\smallskip
{\parbox {6.5 cm}{Зав. кафедри \dotfill Ю.В.~Крак}}\hspace {0.7cm}{\hbox to 6.5 cm{ Екзаменатор \dotfill Т.О.~Карнаух}} \end{small}}
\newpage

{{\begin {small}\par
{ \center  Київський національний університет імені Тараса Шевченка \par}
{\parbox {5 cm}{Освітня програма \hfill}{\parbox {7 cm}{Інформатика \hfill}}\parbox {6 cm}{\hfill Семестр 6}}\par
{\parbox {5 cm} {Навчальний предмет \hfill}\parbox {7 cm}{Розробка програмного забезпечення \hfill}\parbox {6cm}{\hfill}\par}\vspace{-7pt} \end{small}}{\center Екзаменаційний білет \No \textbf
{ 71 }\par}}
{
\begin {large}
{\begin {center}\textbf{Завдання 1}\end {center}}\par
Так звані трикутники балансування зв'язують між собою:\par
\vspace{2mm}



(\,1\,) 4 або 3 параметри, при цьому може бути відсутній обсяг %good



(\,2\,) Завжди рівно 3 параметри%bad



(\,3\,) 4 або 3 параметри, при цьому може бути відсутній час%bad



(\,4\,) Завжди рівно 4 параметри%bad


{\begin {center}\textbf{Завдання 2}\end {center}}\par
Який з наведених фрагментів коду явно потребує поліпшення (форматування рядків та типи до уваги не брати)?\par
\vspace{2mm}



(\,1\,) \{int input\_file, file\_to\_output; …\}%bad



(\,2\,) void f13(…);%bad



(\,3\,) \{int x; int y; …\}%good



(\,4\,) \{double max\_power, min\_power;…\} %good


{\begin {center}\textbf{Завдання 3}\end {center}}\par
Що є характерним для структурного програмування?\par
\vspace{2mm}



(\,1\,) Заборона використання скалярних (простих) типів даних. %bad



(\,2\,) Переходи з середини одного блоку в середину іншого неможливі. %good



(\,3\,) Використання механізму винятків. %bad



(\,4\,) Використання оператора безумовного переходу. %bad


{\begin {center}\textbf{Завдання 4}\end {center}}\par
Вибрати всі істинні твердження (якщо існують).\par
\vspace{2mm}



(\,1\,) Для адаптивних моделей розробки оцінки часу та бюджету не виконують, бо вимоги є змінними. %bad



(\,2\,) Адаптивні моделі розробки не допускають зміну вимог до продукту, але продукт може розроблятись поетапно. %bad



(\,3\,) Прогнозна модель розробки не може одночасно бути інкрементною. %bad



(\,4\,) Для адаптивних методологій розробки, на відміну від прогнозних, попередні оцінки бюджету та/або часу не виконують, бо вони все одно зміняться. %bad


{\begin {center}\textbf{Завдання 5}\end {center}}\par
Вибрати всі істинні твердження (якщо існують).\par
\vspace{2mm}



(\,1\,) Код з довгими функціями має меншу кількість функцій, а тому для нього треба менше тестів. %bad



(\,2\,) Одним з прийомів спрощення складної умовної логіки є використання шаблону State. %good



(\,3\,) Складні умови виникають в реальних задачах доволі часто, а тому є доволі природніми і не можуть вважатись недоліком коду. %bad



(\,4\,) Залежності між різними модулями системи мають бути якомога більш сильними, бо інакше вона розпадеться. %bad

\newpage

{\begin {center}\textbf{Завдання 6}\end {center}}\par
Вибрати всі істинні твердження (якщо існують).\par
\vspace{2mm}



(\,1\,) Якщо виникають проблеми з модифікацією коду, то слід відразу змінити команду розробників на більш кваліфіковану.%bad



(\,2\,) Простота – мистецтво максимізації незробленої роботи. %good



(\,3\,) Сьогодні водоспадна методологія не має права на існування. %bad



(\,4\,) Водоспадна методологія розробки може призводити до надлишкового проектування. %good


{\begin {center}\textbf{Завдання 7}\end {center}}\par
Вибрати всі істинні твердження (якщо існують).\par
\vspace{2mm}



(\,1\,) Гнучкі методології розробки намагаються якомога швидше визначити архітектуру системи, для чого на початку завжди виконується початкове проектування системи в цілому, а тільки потім переходять до ітеративного уточнення деталей та кодування. %bad



(\,2\,) Гнучкі методи розробки нехтують проміжними технічними артефактами та документацією виключно через їх непотрібність. %bad



(\,3\,) Водоспадна методологія може мати тенденцію до надлишкового проектування. %good



(\,4\,) Гнучкі методології базуються на комунікації розробників, а тому передбачають виключно колективне володіння кодом. %bad


{\begin {center}\textbf{Завдання 8}\end {center}}\par
Що з наступного не має відношення до Scrum або не виконується/не використовується за використання оригінального Scrum?\par
\vspace{2mm}



(\,1\,) порядок виконання запланованої на спрінт роботи визначає власник продукту%good



(\,2\,) скрам-мастер визначає, які задачі команда буде виконувати на наступному спрінті%good



(\,3\,) під час спрінту дозволяється внесення змін до вимог, що вже реалізуються на поточному спрінті%good



(\,4\,) виділення одного з учасників команди як системного архітектора%good


{\begin {center}\textbf{Завдання 9}\end {center}}\par
Що з наведеного нижче не розглядалось та навіть не згадувалось у курсі?\par
\vspace{2mm}



(\,1\,) LOAN%good



(\,2\,) DSDM%bad



(\,3\,) Scrum%bad



(\,4\,) Lean%bad


{\begin {center}\textbf{Завдання 10}\end {center}}\par
Вибрати всі істинні твердження (якщо існують).\par
\vspace{2mm}



(\,1\,) Модель якості при використанні є ієрархічною.%good



(\,2\,) Забезпечення якості це не тільки тестування.%good



(\,3\,) Модель якості продукту є ієрархічною.%good



(\,4\,) Користувачі бувають вторинними.%good



\end {large}}
{\vspace{5mm}\smallskip \begin {small} Затверджено на засіданні кафедри теоретичної кібернетики, протокол \No 6 від   29.01.2021 р.\par\smallskip\smallskip
{\parbox {6.5 cm}{Зав. кафедри \dotfill Ю.В.~Крак}}\hspace {0.7cm}{\hbox to 6.5 cm{ Екзаменатор \dotfill Т.О.~Карнаух}} \end{small}}
\newpage

{{\begin {small}\par
{ \center  Київський національний університет імені Тараса Шевченка \par}
{\parbox {5 cm}{Освітня програма \hfill}{\parbox {7 cm}{Інформатика \hfill}}\parbox {6 cm}{\hfill Семестр 6}}\par
{\parbox {5 cm} {Навчальний предмет \hfill}\parbox {7 cm}{Розробка програмного забезпечення \hfill}\parbox {6cm}{\hfill}\par}\vspace{-7pt} \end{small}}{\center Екзаменаційний білет \No \textbf
{ 72 }\par}}
{
\begin {large}
{\begin {center}\textbf{Завдання 1}\end {center}}\par
Так звані трикутники балансування зв'язують між собою:\par
\vspace{2mm}



(\,1\,) 4 або 3 параметри, при цьому може бути відсутній обсяг %good



(\,2\,) Завжди рівно 3 параметри%bad



(\,3\,) 4 або 3 параметри, при цьому може бути відсутня якість%bad



(\,4\,) Завжди рівно 4 параметри%bad


{\begin {center}\textbf{Завдання 2}\end {center}}\par
Який з наведених фрагментів коду явно потребує поліпшення (форматування рядків та типи до уваги не брати)?\par
\vspace{2mm}



(\,1\,) \{int xtrmprg; …\}%bad



(\,2\,) \{fstream input\_file, output\_file;…\} %good



(\,3\,) void ff(…);%bad



(\,4\,) \{fstream inputFile, outputFile;…\} %good


{\begin {center}\textbf{Завдання 3}\end {center}}\par
Що є характерним для структурного програмування?\par
\vspace{2mm}



(\,1\,) У коді мають використовуватись структурні типи даних. %bad



(\,2\,) Кожен блок коду повинен мати єдиний вхід та єдиний вихід.%good



(\,3\,) У коді мають використовуватись типи структур. %bad



(\,4\,) У коді мають використовуватись типи даних, що є класами. %bad


{\begin {center}\textbf{Завдання 4}\end {center}}\par
Вибрати всі істинні твердження (якщо існують).\par
\vspace{2mm}



(\,1\,) Неперервне проектування створює специфікацію проекту перед початком розробки/ітерації, а під час ітерації його неперервно модифікує.%bad



(\,2\,) Для адаптивних методологій прогнозувати результат чи терміни його отримання абсолютно неможливо. %bad



(\,3\,) Неперервне проектування має бути ітеративним та інкрементним одночасно. %good



(\,4\,) Адаптивні методології не передбачають визначення вимог. %bad


{\begin {center}\textbf{Завдання 5}\end {center}}\par
Вибрати всі істинні твердження (якщо існують).\par
\vspace{2mm}



(\,1\,) Чим менше модулів у проекті, тим краще, бо меншу кількість модулів треба інтегрувати в єдине ціле. %bad



(\,2\,) Чим менше класів у коді, тим менше класів треба налагоджувати, а тому класи за можливості слід об'єднувати. %bad



(\,3\,) Правильно спроектована функція повинна мати рівно один обов'язок. %good



(\,4\,) Повторюваний код дозволяє програмісту краще зрозуміти його логіку та вибрати з кількох варіантів найкращий. %bad

\newpage

{\begin {center}\textbf{Завдання 6}\end {center}}\par
Вибрати всі істинні твердження (якщо існують).\par
\vspace{2mm}



(\,1\,) Гнучкі методології передбачають часті випуски продукту. %good



(\,2\,) RUP є прогнозною методологією. %bad



(\,3\,) Гнучкі методології розробки є стійкими до припинення розробки на певний час. %bad



(\,4\,) RUP є ітеративним та інкрементним.%good


{\begin {center}\textbf{Завдання 7}\end {center}}\par
Вибрати всі істинні твердження (якщо існують).\par
\vspace{2mm}



(\,1\,) RUP є архітектурно-орієнтованим. %good



(\,2\,) За використання гнучких методологій кожен розробник робить, що хоче. %bad



(\,3\,) За застосування адаптивних методологій розробки, так само як і за прогнозних, намагаються якомога швидше визначити архітектуру системи. Відмінність полягає в тім, у який спосіб це виконується.%good



(\,4\,) На відміну від Scrum, RUP є компонентно-орієнтованим. %bad


{\begin {center}\textbf{Завдання 8}\end {center}}\par
Що з наступного не має відношення до Scrum або не виконується/не використовується за використання оригінального Scrum?\par
\vspace{2mm}



(\,1\,) виділення в команді одного чи кількох проектувальників%good



(\,2\,) правила визначення готовності (Definition of Done) %bad



(\,3\,) канбан%bad



(\,4\,) скрам-мастер визначає пріоритети задач%good


{\begin {center}\textbf{Завдання 9}\end {center}}\par
Що з наведеного нижче не розглядалось та навіть не згадувалось у курсі?\par
\vspace{2mm}



(\,1\,) LAZY%good



(\,2\,) модель хаосу%bad



(\,3\,) XP%bad



(\,4\,) SAFe%bad


{\begin {center}\textbf{Завдання 10}\end {center}}\par
Вибрати всі істинні твердження (якщо існують).\par
\vspace{2mm}



(\,1\,) Для тестування код треба ізолювати. Для того, щоб ізолювати, треба мати тести.%good



(\,2\,) Тестування вимог може породжувати їх уточнення.%good



(\,3\,) Тестування буває позитивне та негативне.%good



(\,4\,) Якість буває зовнішньою.%good



\end {large}}
{\vspace{5mm}\smallskip \begin {small} Затверджено на засіданні кафедри теоретичної кібернетики, протокол \No 6 від   29.01.2021 р.\par\smallskip\smallskip
{\parbox {6.5 cm}{Зав. кафедри \dotfill Ю.В.~Крак}}\hspace {0.7cm}{\hbox to 6.5 cm{ Екзаменатор \dotfill Т.О.~Карнаух}} \end{small}}
\newpage

{{\begin {small}\par
{ \center  Київський національний університет імені Тараса Шевченка \par}
{\parbox {5 cm}{Освітня програма \hfill}{\parbox {7 cm}{Інформатика \hfill}}\parbox {6 cm}{\hfill Семестр 6}}\par
{\parbox {5 cm} {Навчальний предмет \hfill}\parbox {7 cm}{Розробка програмного забезпечення \hfill}\parbox {6cm}{\hfill}\par}\vspace{-7pt} \end{small}}{\center Екзаменаційний білет \No \textbf
{ 73 }\par}}
{
\begin {large}
{\begin {center}\textbf{Завдання 1}\end {center}}\par
Що з наведеного найкраще характеризує рефакторінг?\par
\vspace{2mm}



(\,1\,) Докорінна переробка всього коду з нуля,  що не змінює його видиму поведінку%bad



(\,2\,) Зміна внутрішньої структури коду, що не змінює його видиму поведінку%good



(\,3\,) Довільна переробка коду, що виникає у процесі розробки%bad



(\,4\,) Зміна внутрішньої структури коду під час усунення помилок%bad


{\begin {center}\textbf{Завдання 2}\end {center}}\par
Який з наведених фрагментів коду явно потребує поліпшення (форматування рядків та типи до уваги не брати)?\par
\vspace{2mm}



(\,1\,) void help();%good



(\,2\,) \{int xtrmprg; …\}%bad



(\,3\,) \{fstream inputFile, outputFile;…\} %good



(\,4\,) \{int input\_file, file\_to\_output; …\}%bad


{\begin {center}\textbf{Завдання 3}\end {center}}\par
Що є характерним для структурного програмування?\par
\vspace{2mm}



(\,1\,) У коді мають використовуватись типи структур, класи використовувати заборонено. %bad



(\,2\,) У коді мають використовуватись типи структур. %bad



(\,3\,) Використання оператора безумовного переходу. %bad



(\,4\,) Програма зображується ієрархічною структурою блоків.%good


{\begin {center}\textbf{Завдання 4}\end {center}}\par
Вибрати всі істинні твердження (якщо існують).\par
\vspace{2mm}



(\,1\,) І прогнозні, і адаптивні методології передбачають планування виконуваних робіт. %good



(\,2\,) Прогнозна модель розробки не може одночасно бути інкрементною. %bad



(\,3\,) Адаптивні методології не передбачають визначення вимог. %bad



(\,4\,) Для адаптивних методологій прогнозувати результат чи терміни його отримання абсолютно неможливо. %bad


{\begin {center}\textbf{Завдання 5}\end {center}}\par
Вибрати всі істинні твердження (якщо існують).\par
\vspace{2mm}



(\,1\,) Код, що виконує однакову послідовність логічних кроків, у більшості випадків є повторюваним.%good



(\,2\,) Код програми пишеться для комп'ютера, тому він має бути синтаксично та семантично коректним. Інше значення не має. %bad



(\,3\,) Код з довгими функціями має меншу кількість функцій, а тому для нього треба менше тестів. %bad



(\,4\,) Повторюваний код дозволяє програмісту краще зрозуміти його логіку та вибрати з кількох варіантів найкращий. %bad

\newpage

{\begin {center}\textbf{Завдання 6}\end {center}}\par
Вибрати всі істинні твердження (якщо існують).\par
\vspace{2mm}



(\,1\,) Програмування в парах (pair programming) зводиться до того, що кожним фрагментом коду володіють два програмісти. %bad



(\,2\,) Адаптивні методології розробки дозволяють уникнути надлишкового проектування. %good



(\,3\,) RUP є керованим прецедентами. %good



(\,4\,) Використання гнучких методологій вимагає високої кваліфікації персоналу. %good


{\begin {center}\textbf{Завдання 7}\end {center}}\par
Вибрати всі істинні твердження (якщо існують).\par
\vspace{2mm}



(\,1\,) Гнучкі методології розробки не допускають використання UML. %bad



(\,2\,) Водоспадна методологія розробки може призводити до недостатнього проектування. %bad



(\,3\,) Гнучкі методології розробки допускають суттєві зміни вимог, проте все одно намагаються якомога швидше зафіксувати архітектуру. %good



(\,4\,) Поступове уточнення вимог замовника полегшує адаптацію до змін. %good


{\begin {center}\textbf{Завдання 8}\end {center}}\par
Що з наступного не має відношення до Scrum або не виконується/не використовується за використання оригінального Scrum?\par
\vspace{2mm}



(\,1\,) власник продукту остаточно затверджує проект системи%good



(\,2\,) журнал невиконаних задач (backlog) спрінту%bad



(\,3\,) скрам-мастер визначає час, необхідний команді для реалізації конкретної задачі%good



(\,4\,) щоденний скрам%bad


{\begin {center}\textbf{Завдання 9}\end {center}}\par
Що з наведеного нижче не розглядалось та навіть не згадувалось у курсі?\par
\vspace{2mm}



(\,1\,) TDD%bad



(\,2\,) Scrum of scrums%bad



(\,3\,) X-модель%good



(\,4\,) poker%bad


{\begin {center}\textbf{Завдання 10}\end {center}}\par
Вибрати всі істинні твердження (якщо існують).\par
\vspace{2mm}



(\,1\,) Якість буває внутрішньою.%good



(\,2\,) Тестовий приклад складається з набору вхідних даних, умов виконання та очікуваних результатів.%good



(\,3\,) Модель якості при використанні є ієрархічною.%good



(\,4\,) Атрибути якості можна використовувати при уточненні вимог.%good



\end {large}}
{\vspace{5mm}\smallskip \begin {small} Затверджено на засіданні кафедри теоретичної кібернетики, протокол \No 6 від   29.01.2021 р.\par\smallskip\smallskip
{\parbox {6.5 cm}{Зав. кафедри \dotfill Ю.В.~Крак}}\hspace {0.7cm}{\hbox to 6.5 cm{ Екзаменатор \dotfill Т.О.~Карнаух}} \end{small}}
\newpage

{{\begin {small}\par
{ \center  Київський національний університет імені Тараса Шевченка \par}
{\parbox {5 cm}{Освітня програма \hfill}{\parbox {7 cm}{Інформатика \hfill}}\parbox {6 cm}{\hfill Семестр 6}}\par
{\parbox {5 cm} {Навчальний предмет \hfill}\parbox {7 cm}{Розробка програмного забезпечення \hfill}\parbox {6cm}{\hfill}\par}\vspace{-7pt} \end{small}}{\center Екзаменаційний білет \No \textbf
{ 74 }\par}}
{
\begin {large}
{\begin {center}\textbf{Завдання 1}\end {center}}\par
Так звані трикутники балансування зв'язують між собою:\par
\vspace{2mm}



(\,1\,) Завжди рівно 4 параметри%bad



(\,2\,) 4 або 3 параметри, при цьому може бути відсутній час%bad



(\,3\,) 4 або 3 параметри, при цьому може бути відсутній обсяг %good



(\,4\,) 4 або 3 параметри, при цьому може бути відсутня якість%bad


{\begin {center}\textbf{Завдання 2}\end {center}}\par
Який з наведених фрагментів коду явно потребує поліпшення (форматування рядків та типи до уваги не брати)?\par
\vspace{2mm}



(\,1\,) \{int x; int y; …\}%good



(\,2\,) \{double the\_best\_value\_anyone\_ever\_have;\} %bad



(\,3\,) void ff(…);%bad



(\,4\,) \{fstream input\_file, outputFile; …\}%bad


{\begin {center}\textbf{Завдання 3}\end {center}}\par
Що є характерним для структурного програмування?\par
\vspace{2mm}



(\,1\,) Заборона використання скалярних (простих) типів даних. %bad



(\,2\,) Кожен блок коду повинен мати єдиний вхід та єдиний вихід.%good



(\,3\,) У коді мають використовуватись структурні типи даних. %bad



(\,4\,) У коді мають використовуватись типи даних, що є класами. %bad


{\begin {center}\textbf{Завдання 4}\end {center}}\par
Вибрати всі істинні твердження (якщо існують).\par
\vspace{2mm}



(\,1\,) Адаптивні моделі розробки допускають зміну вимог до продукту. %good



(\,2\,) Прогнозні та адаптивні методології використовують одні ті самі діяльності з розробки ПЗ: визначення та аналіз вимог, проектування, реалізація, тестування, впровадження, супроводження. %good



(\,3\,) У прогнозній моделі розробки розробка не може виконуватись ітеративно. %bad



(\,4\,) Неперервність проектування полягає в тім, що процес розробки не можна переривати більше ніж на вихідні. %bad


{\begin {center}\textbf{Завдання 5}\end {center}}\par
Вибрати всі істинні твердження (якщо існують).\par
\vspace{2mm}



(\,1\,) Одним з засобів позбутись повторень у коді є використання успадкування. %good



(\,2\,) Складні комбінації умов починають керувати логікою виконання, тому що такими є вимоги до функціонування ПЗ. Оскільки ПЗ має якнайкраще забезпечувати потреби замовника, то складні комбінації умов у коді вважатись недоліком коду не можуть. %bad



(\,3\,) Гарний код має читатись як розмовна мова. %good



(\,4\,) Якщо кожна функція бібліотеки буде здатна в залежності від параметрів виконувати кілька різних дій, то загальна кількість функцій у ній буде менша і її буде простіше використовувати. %bad

\newpage

{\begin {center}\textbf{Завдання 6}\end {center}}\par
Вибрати всі істинні твердження (якщо існують).\par
\vspace{2mm}



(\,1\,) Розробка за гнучкими методологіями цінить співробітництво з замовником та піддатлівість змінам. %good



(\,2\,) Багато хто з тих, що використовує Scrum, складає архітектурну модель розроблюваної системи. %good



(\,3\,) Використання RUP не заперечує часті випуски продукту, але й не наполягає на них. %good



(\,4\,) Якщо виникають проблеми з модифікацією коду, то слід розглянути внесення змін у проект. %good


{\begin {center}\textbf{Завдання 7}\end {center}}\par
Вибрати всі істинні твердження (якщо існують).\par
\vspace{2mm}



(\,1\,) Гнучкі методології повністю нехтують документацією. %bad



(\,2\,) За програмування в парах (pair programming) поки один колега набирає код на комп'ютері інший може піти пограти в настільний теніс. %bad



(\,3\,) Програмування в парах (pair programming) спрощує введення в команду нових розробників. %good



(\,4\,) Scrum передбачає колективне володіння кодом. %good


{\begin {center}\textbf{Завдання 8}\end {center}}\par
Що з наступного не має відношення до Scrum або не виконується/не використовується за використання оригінального Scrum?\par
\vspace{2mm}



(\,1\,) розповіді користувача%bad



(\,2\,) планування спрінту%bad



(\,3\,) покер планування (Planning Poker)%bad



(\,4\,) скрам-мастер остаточно затверджує проект системи%good


{\begin {center}\textbf{Завдання 9}\end {center}}\par
Що з наведеного нижче не розглядалось та навіть не згадувалось у курсі?\par
\vspace{2mm}



(\,1\,) водоспадна модель%bad



(\,2\,) sashimi%bad



(\,3\,) BDUF%bad



(\,4\,) PRINCE2%good


{\begin {center}\textbf{Завдання 10}\end {center}}\par
Вибрати всі істинні твердження (якщо існують).\par
\vspace{2mm}



(\,1\,) Якість можна визначати через атрибути якості.%good



(\,2\,) Тестування вимог може породжувати їх уточнення.%good



(\,3\,) Модель якості даних є ієрархічною.%good



(\,4\,) За провалених димових тестів решта тестів не виконується і код повертається на доробку.%good



\end {large}}
{\vspace{5mm}\smallskip \begin {small} Затверджено на засіданні кафедри теоретичної кібернетики, протокол \No 6 від   29.01.2021 р.\par\smallskip\smallskip
{\parbox {6.5 cm}{Зав. кафедри \dotfill Ю.В.~Крак}}\hspace {0.7cm}{\hbox to 6.5 cm{ Екзаменатор \dotfill Т.О.~Карнаух}} \end{small}}
\newpage

{{\begin {small}\par
{ \center  Київський національний університет імені Тараса Шевченка \par}
{\parbox {5 cm}{Освітня програма \hfill}{\parbox {7 cm}{Інформатика \hfill}}\parbox {6 cm}{\hfill Семестр 6}}\par
{\parbox {5 cm} {Навчальний предмет \hfill}\parbox {7 cm}{Розробка програмного забезпечення \hfill}\parbox {6cm}{\hfill}\par}\vspace{-7pt} \end{small}}{\center Екзаменаційний білет \No \textbf
{ 75 }\par}}
{
\begin {large}
{\begin {center}\textbf{Завдання 1}\end {center}}\par
Шаблон проектування:\par
\vspace{2mm}



(\,1\,) Формулює ідею розв'язання проблеми та може пропонувати одну чи кілька її реалізацій%good



(\,2\,) Є шаблоном, в який треба підставити власні ідентифікатори%bad



(\,3\,) Задає розв'язання типової проблеми до найдрібніших подробиць%bad



(\,4\,) Тільки формулює ідею розв'язання проблеми%bad


{\begin {center}\textbf{Завдання 2}\end {center}}\par
Який з наведених фрагментів коду явно потребує поліпшення (форматування рядків та типи до уваги не брати)?\par
\vspace{2mm}



(\,1\,) \{fstream input\_file, output\_file;…\} %good



(\,2\,) void load\_data(T \& arg1, double arg2, double arg3, double arg4); %bad



(\,3\,) \{double max\_power, min\_power;…\} %good



(\,4\,) \{double avg\_temperature;…\} %good


{\begin {center}\textbf{Завдання 3}\end {center}}\par
Що є характерним для структурного програмування?\par
\vspace{2mm}



(\,1\,) Використання механізму винятків. %bad



(\,2\,) У коді мають використовуватись типи структур, класи використовувати заборонено. %bad



(\,3\,) Переходи з середини одного блоку в середину іншого неможливі. %good



(\,4\,) У коді мають використовуватись типи даних, що є класами. %bad


{\begin {center}\textbf{Завдання 4}\end {center}}\par
Вибрати всі істинні твердження (якщо існують).\par
\vspace{2mm}



(\,1\,) В адаптивних методологіях планування робіт не виконується. %bad



(\,2\,) Гнучкі методології не передбачають планування робіт. %bad



(\,3\,) Адаптивні моделі розробки не допускають зміну вимог до продукту, але продукт може розроблятись поетапно. %bad



(\,4\,) Прогнозні моделі розробки з самого початку фіксують вимоги до продукту. %good


{\begin {center}\textbf{Завдання 5}\end {center}}\par
Вибрати всі істинні твердження (якщо існують).\par
\vspace{2mm}



(\,1\,) Одним з прийомів спрощення складної умовної логіки є використання шаблону State. %good



(\,2\,) Одним з прийомів спрощення складних умов є їх декомпозиція. %good



(\,3\,) Довгі функції доволі часто приховують складну логіку обробки або повторюваний код. %good



(\,4\,) Чим менше модулів у проекті, тим краще, бо меншу кількість модулів треба інтегрувати в єдине ціле. %bad

\newpage

{\begin {center}\textbf{Завдання 6}\end {center}}\par
Вибрати всі істинні твердження (якщо існують).\par
\vspace{2mm}



(\,1\,) Розробка за гнучкими методологіями ніколи не слідує плану. %bad



(\,2\,) Висока внутрішня якість коду є марною тратою часу на те, що замовник не побачить та/або не зможе зацінити. %bad



(\,3\,) Водоспадна методологія не призначена для внесення змін у вимоги до системи. %good



(\,4\,) Висока внутрішня якість коду здешевшує розробку, бо дозволяє економити на вартості та часі майбутніх змін. %good


{\begin {center}\textbf{Завдання 7}\end {center}}\par
Вибрати всі істинні твердження (якщо існують).\par
\vspace{2mm}



(\,1\,) Одна з принципових відмінностей RUP від Scrum полягає у використанні численних проміжних технічних артефактів. %good



(\,2\,) Розповіді користувача (user stories) це менш формалізований порівняно з прецедентами спосіб опису вимог до системи. %good



(\,3\,) Основна відмінність V-моделі від водоспадної полягає в тім, що діяльності з тестування виконуються паралельно з усіма іншими діяльностями з розробки. %good



(\,4\,) Без рефакторінгу неперервне проектування неможливе. %good


{\begin {center}\textbf{Завдання 8}\end {center}}\par
Що з наступного не має відношення до Scrum або не виконується/не використовується за використання оригінального Scrum?\par
\vspace{2mm}



(\,1\,) приведення у належний вигляд журналу невиконаних задач (backlog grooming) %bad



(\,2\,) залучення зовнішніх фахівців для побудови архітектури системи%good



(\,3\,) канбан%bad



(\,4\,) під час спрінту дозволяється внесення змін до вимог, що вже реалізуються на поточному спрінті%good


{\begin {center}\textbf{Завдання 9}\end {center}}\par
Що з наведеного нижче не розглядалось та навіть не згадувалось у курсі?\par
\vspace{2mm}



(\,1\,) LeSS%bad



(\,2\,) RUP%bad



(\,3\,) RAD%bad



(\,4\,) CRAZY%good


{\begin {center}\textbf{Завдання 10}\end {center}}\par
Вибрати всі істинні твердження (якщо існують).\par
\vspace{2mm}



(\,1\,) Модель якості продукту є ієрархічною.%good



(\,2\,) Розглядають модель якості продукту.%good



(\,3\,) Якість можна визначати через дефекти.%good



(\,4\,) Якість буває зовнішньою.%good



\end {large}}
{\vspace{5mm}\smallskip \begin {small} Затверджено на засіданні кафедри теоретичної кібернетики, протокол \No 6 від   29.01.2021 р.\par\smallskip\smallskip
{\parbox {6.5 cm}{Зав. кафедри \dotfill Ю.В.~Крак}}\hspace {0.7cm}{\hbox to 6.5 cm{ Екзаменатор \dotfill Т.О.~Карнаух}} \end{small}}
\newpage

{{\begin {small}\par
{ \center  Київський національний університет імені Тараса Шевченка \par}
{\parbox {5 cm}{Освітня програма \hfill}{\parbox {7 cm}{Інформатика \hfill}}\parbox {6 cm}{\hfill Семестр 6}}\par
{\parbox {5 cm} {Навчальний предмет \hfill}\parbox {7 cm}{Розробка програмного забезпечення \hfill}\parbox {6cm}{\hfill}\par}\vspace{-7pt} \end{small}}{\center Екзаменаційний білет \No \textbf
{ 76 }\par}}
{
\begin {large}
{\begin {center}\textbf{Завдання 1}\end {center}}\par
Що є метою рефакторингу?\par
\vspace{2mm}



(\,1\,) Зменшення кількості рядків коду%bad



(\,2\,) Поліпшення внутрішньої структури коду%good



(\,3\,) Усунення помилок у коді%bad



(\,4\,) Додавання нової функціональності%bad


{\begin {center}\textbf{Завдання 2}\end {center}}\par
Який з наведених фрагментів коду явно потребує поліпшення (форматування рядків та типи до уваги не брати)?\par
\vspace{2mm}



(\,1\,) void f13(…);%bad



(\,2\,) void load\_data(…);%good



(\,3\,) \{int a\_, a\_\_; …\}%bad



(\,4\,) \{double maxPower, minPower;…\} %good


{\begin {center}\textbf{Завдання 3}\end {center}}\par
Що є характерним для структурного програмування?\par
\vspace{2mm}



(\,1\,) У коді мають використовуватись типи структур. %bad



(\,2\,) У коді мають використовуватись структурні типи даних. %bad



(\,3\,) Заборона використання скалярних (простих) типів даних. %bad



(\,4\,) Переходи з середини одного блоку в середину іншого неможливі. %good


{\begin {center}\textbf{Завдання 4}\end {center}}\par
Вибрати всі істинні твердження (якщо існують).\par
\vspace{2mm}



(\,1\,) Ітеративні та інкрементні методи розробки є синонімом неперервного проектування. %bad



(\,2\,) Для адаптивних методологій розробки, на відміну від прогнозних, попередні оцінки бюджету та/або часу не виконують, бо вони все одно зміняться. %bad



(\,3\,) У випадку, коли вимоги до продукту відомі заздалегідь та є фіксованими контрактом, адаптивні методології розробки використовувати не можна. %bad



(\,4\,) Адаптивна модель розробки адаптує вимоги до продукту до наявної команди розробників. %bad


{\begin {center}\textbf{Завдання 5}\end {center}}\par
Вибрати всі істинні твердження (якщо існують).\par
\vspace{2mm}



(\,1\,) Код, що виконує однакову послідовність логічних кроків, у більшості випадків є повторюваним.%good



(\,2\,) Гарний код не повинен містити різних фрагментів, що складаються з кількох інструкцій та виконують однакове перетворення даних. %good



(\,3\,) Код з довгими функціями має меншу кількість функцій, а тому його швидше та простіше тестувати. %bad



(\,4\,) Код програми пишеться для комп'ютера, тому він має бути синтаксично та семантично коректним. Інше значення не має. %bad

\newpage

{\begin {center}\textbf{Завдання 6}\end {center}}\par
Вибрати всі істинні твердження (якщо існують).\par
\vspace{2mm}



(\,1\,) Ітеративна та інкрементна розробка передбачає наявність у тому чи іншому вигляді етапу ініціалізації, ітерацій та керуючої таблиці проекту. %good



(\,2\,) Перша версія коду не завжди може бути якісною -- і це нормально. %good



(\,3\,) RUP принципово відрізняється від Scrum тим, що є керованим прецедентами. У той же час у Scrum при плануванні розробки прецеденти (чи їх аналоги) до уваги не беруться. %bad



(\,4\,) Модель хаосу є методологією розробки програмного забезпечення. %bad


{\begin {center}\textbf{Завдання 7}\end {center}}\par
Вибрати всі істинні твердження (якщо існують).\par
\vspace{2mm}



(\,1\,) Використання гнучких методологій при розробці надвеликих проектів не несе в собі жодних ризиків. %bad



(\,2\,) Гнучкі методи розробки нехтують проміжними технічними артефактами та документацією через те, що за зміни вимог та/або проекту треба витрачати додатковий час на підтримання їх у відповідному стані. %good



(\,3\,) Код, що себе сам документує, -- це код, який виводить користувачу документацію на себе. %bad



(\,4\,) RUP є компонентно-орієнтованим. %good


{\begin {center}\textbf{Завдання 8}\end {center}}\par
Що з наступного не має відношення до Scrum або не виконується/не використовується за використання оригінального Scrum?\par
\vspace{2mm}



(\,1\,) покер планування (Planning Poker)%bad



(\,2\,) скрам-мастер призначає відповідальних за конкретні задачі%good



(\,3\,) скрам-мастер визначає час, необхідний команді для реалізації конкретної задачі%good



(\,4\,) порядок виконання запланованої на спрінт роботи визначають розробники%bad


{\begin {center}\textbf{Завдання 9}\end {center}}\par
Що з наведеного нижче не розглядалось та навіть не згадувалось у курсі?\par
\vspace{2mm}



(\,1\,) X-модель%good



(\,2\,) Scrum%bad



(\,3\,) BDD%bad



(\,4\,) LOAN%good


{\begin {center}\textbf{Завдання 10}\end {center}}\par
Вибрати всі істинні твердження (якщо існують).\par
\vspace{2mm}



(\,1\,) Користувачі бувають основними.%good



(\,2\,) Можна розглядати якість під час використання.%good



(\,3\,) Користувачі бувають вторинними.%good



(\,4\,) Тестування буває позитивне та негативне.%good



\end {large}}
{\vspace{5mm}\smallskip \begin {small} Затверджено на засіданні кафедри теоретичної кібернетики, протокол \No 6 від   29.01.2021 р.\par\smallskip\smallskip
{\parbox {6.5 cm}{Зав. кафедри \dotfill Ю.В.~Крак}}\hspace {0.7cm}{\hbox to 6.5 cm{ Екзаменатор \dotfill Т.О.~Карнаух}} \end{small}}
\newpage

{{\begin {small}\par
{ \center  Київський національний університет імені Тараса Шевченка \par}
{\parbox {5 cm}{Освітня програма \hfill}{\parbox {7 cm}{Інформатика \hfill}}\parbox {6 cm}{\hfill Семестр 6}}\par
{\parbox {5 cm} {Навчальний предмет \hfill}\parbox {7 cm}{Розробка програмного забезпечення \hfill}\parbox {6cm}{\hfill}\par}\vspace{-7pt} \end{small}}{\center Екзаменаційний білет \No \textbf
{ 77 }\par}}
{
\begin {large}
{\begin {center}\textbf{Завдання 1}\end {center}}\par
Так звані трикутники балансування зв'язують між собою:\par
\vspace{2mm}



(\,1\,) 4 або 3 параметри, при цьому може бути відсутній обсяг %good



(\,2\,) Завжди рівно 3 параметри%bad



(\,3\,) 4 або 3 параметри, при цьому може бути відсутній час%bad



(\,4\,) 4 або 3 параметри, при цьому може бути відсутня якість%bad


{\begin {center}\textbf{Завдання 2}\end {center}}\par
Який з наведених фрагментів коду явно потребує поліпшення (форматування рядків та типи до уваги не брати)?\par
\vspace{2mm}



(\,1\,) void load\_data(T \& arg1, double arg2, double arg3, double arg4); %bad



(\,2\,) \{int x; int y; …\}%good



(\,3\,) \{fstream input\_file, output\_file;…\} %good



(\,4\,) \{double avg\_temperature;…\} %good


{\begin {center}\textbf{Завдання 3}\end {center}}\par
Що є характерним для структурного програмування?\par
\vspace{2mm}



(\,1\,) Кожен блок коду повинен мати єдиний вхід та єдиний вихід.%good



(\,2\,) Використання механізму винятків. %bad



(\,3\,) Використання оператора безумовного переходу. %bad



(\,4\,) У коді мають використовуватись типи структур, класи використовувати заборонено. %bad


{\begin {center}\textbf{Завдання 4}\end {center}}\par
Вибрати всі істинні твердження (якщо існують).\par
\vspace{2mm}



(\,1\,) Неперервне проектування здійснює створення та модифікацію проекту системи в міру розробки. %good



(\,2\,) Адаптивність методології проявляється в тому, що за зміни вимог розробка починається заново. %bad



(\,3\,) Адаптивні методології має сенс використовувати для проектів, де можлива зміна вимог. %good



(\,4\,) Неперервне проектування має бути ітеративним та інкрементним одночасно. %good


{\begin {center}\textbf{Завдання 5}\end {center}}\par
Вибрати всі істинні твердження (якщо існують).\par
\vspace{2mm}



(\,1\,) Одним з прийомів спрощення умовної логіки є використання поліморфізму. %good



(\,2\,) Втілення в коді різних підходів до одного того самого перетворення даних дозволяє програмісту продемонструвати свою професійну підготовку з найкращого боку. %bad



(\,3\,) Складні умови виникають в реальних задачах доволі часто, а тому є доволі природніми і не можуть вважатись недоліком коду. %bad



(\,4\,) Довжина функцій не впливає на внутрішню якість коду. %bad

\newpage

{\begin {center}\textbf{Завдання 6}\end {center}}\par
Вибрати всі істинні твердження (якщо існують).\par
\vspace{2mm}



(\,1\,) За відмови від документації одним з методом збереження знань про систему є постійне пересування персоналу між її частинами. %good



(\,2\,) Розробка за гнучкими методологіями не передбачає узгодження умов контракту. %bad



(\,3\,) Тільки RUP використовує UML.%bad



(\,4\,) Адаптивні методології розробки дозволяють уникнути надлишкового проектування. %good


{\begin {center}\textbf{Завдання 7}\end {center}}\par
Вибрати всі істинні твердження (якщо існують).\par
\vspace{2mm}



(\,1\,) Модель хаосу є методологією розробки програмного забезпечення. %bad



(\,2\,) Гнучкі методології повністю нехтують документацією. %bad



(\,3\,) Програмування в парах (pair programming) зводиться до того, що кожним фрагментом коду володіють два програмісти. %bad



(\,4\,) Без рефакторінгу неперервне проектування неможливе. %good


{\begin {center}\textbf{Завдання 8}\end {center}}\par
Що з наступного не має відношення до Scrum або не виконується/не використовується за використання оригінального Scrum?\par
\vspace{2mm}



(\,1\,) ретроспектива спрінту%bad



(\,2\,) бачення продукту (vision) %bad



(\,3\,) під час спрінту зміни, що впливають на досягнення мети спрінту, зазвичай, не допускаються%bad



(\,4\,) скрам-мастер визначає пріоритети задач%good


{\begin {center}\textbf{Завдання 9}\end {center}}\par
Що з наведеного нижче не розглядалось та навіть не згадувалось у курсі?\par
\vspace{2mm}



(\,1\,) Scrum of scrums%bad



(\,2\,) LAZY%good



(\,3\,) водоспадна модель%bad



(\,4\,) PRINCE2%good


{\begin {center}\textbf{Завдання 10}\end {center}}\par
Вибрати всі істинні твердження (якщо існують).\par
\vspace{2mm}



(\,1\,) Розглядають модель якості при використанні.%good



(\,2\,) Відсутність знайдених дефектів на початкових стадіях розробки може бути свідченням поганого покриття коду тестами.%good



(\,3\,) Коли нема можливості перевірити всі можливі комбінації, тестують їх попарні (pairwise) комбінації. Це може значно зменшити кількість тестів.%good



(\,4\,) І верифікація, і валідація можуть використовувати тестування.%good



\end {large}}
{\vspace{5mm}\smallskip \begin {small} Затверджено на засіданні кафедри теоретичної кібернетики, протокол \No 6 від   29.01.2021 р.\par\smallskip\smallskip
{\parbox {6.5 cm}{Зав. кафедри \dotfill Ю.В.~Крак}}\hspace {0.7cm}{\hbox to 6.5 cm{ Екзаменатор \dotfill Т.О.~Карнаух}} \end{small}}
\newpage

{{\begin {small}\par
{ \center  Київський національний університет імені Тараса Шевченка \par}
{\parbox {5 cm}{Освітня програма \hfill}{\parbox {7 cm}{Інформатика \hfill}}\parbox {6 cm}{\hfill Семестр 6}}\par
{\parbox {5 cm} {Навчальний предмет \hfill}\parbox {7 cm}{Розробка програмного забезпечення \hfill}\parbox {6cm}{\hfill}\par}\vspace{-7pt} \end{small}}{\center Екзаменаційний білет \No \textbf
{ 78 }\par}}
{
\begin {large}
{\begin {center}\textbf{Завдання 1}\end {center}}\par
Так звані трикутники балансування зв'язують між собою:\par
\vspace{2mm}



(\,1\,) Завжди рівно 3 параметри%bad



(\,2\,) 4 або 3 параметри, при цьому може бути відсутня якість%bad



(\,3\,) 4 або 3 параметри, при цьому може бути відсутній обсяг %good



(\,4\,) 4 або 3 параметри, при цьому може бути відсутній час%bad


{\begin {center}\textbf{Завдання 2}\end {center}}\par
Який з наведених фрагментів коду явно потребує поліпшення (форматування рядків та типи до уваги не брати)?\par
\vspace{2mm}



(\,1\,) \{double maxPower, minPower;…\} %good



(\,2\,) void help();%good



(\,3\,) \{fstream input\_file, outputFile; …\}%bad



(\,4\,) void ff(…);%bad


{\begin {center}\textbf{Завдання 3}\end {center}}\par
Що є характерним для структурного програмування?\par
\vspace{2mm}



(\,1\,) Програма зображується ієрархічною структурою блоків.%good



(\,2\,) Використання оператора безумовного переходу. %bad



(\,3\,) У коді мають використовуватись типи структур, класи використовувати заборонено. %bad



(\,4\,) У коді мають використовуватись структурні типи даних. %bad


{\begin {center}\textbf{Завдання 4}\end {center}}\par
Вибрати всі істинні твердження (якщо існують).\par
\vspace{2mm}



(\,1\,) Неперервне проектування створює специфікацію проекту перед початком розробки/ітерації, а під час ітерації його неперервно модифікує.%bad



(\,2\,) Прогнозна модель розробки може одночасно бути інкрементною. %good



(\,3\,) У прогнозній моделі розробки розробка може виконуватись ітеративно. %good



(\,4\,) Для адаптивних моделей розробки оцінки часу та бюджету не виконують, бо вимоги є змінними. %bad


{\begin {center}\textbf{Завдання 5}\end {center}}\par
Вибрати всі істинні твердження (якщо існують).\par
\vspace{2mm}



(\,1\,) Зв'язки між окремими модулями системи мають бути якомога більш слабкими. %good



(\,2\,) Чим менше класів у коді, тим менше класів треба налагоджувати, а тому класи за можливості слід об'єднувати. %bad



(\,3\,) Чим вища кваліфікація програміста, тим більш складні для розуміння та довгі умовні вирази він використовує. %bad



(\,4\,) Повторюваний код має повторюватись у точності до символу, інакше він повторюваним вважатись не може. %bad

\newpage

{\begin {center}\textbf{Завдання 6}\end {center}}\par
Вибрати всі істинні твердження (якщо існують).\par
\vspace{2mm}



(\,1\,) Гнучкі методології базуються на комунікації розробників, а тому передбачають виключно колективне володіння кодом. %bad



(\,2\,) Якщо виникають проблеми з модифікацією коду, то слід розглянути внесення змін у проект. %good



(\,3\,) Гнучкі методології розробки є стійкими до припинення розробки на певний час. %bad



(\,4\,) Розповіді користувача (user stories) це менш формалізований порівняно з прецедентами спосіб опису вимог до системи. %good


{\begin {center}\textbf{Завдання 7}\end {center}}\par
Вибрати всі істинні твердження (якщо існують).\par
\vspace{2mm}



(\,1\,) Розробка за гнучкими методологіями ніколи не слідує плану. %bad



(\,2\,) Простота – мистецтво максимізації незробленої роботи. %good



(\,3\,) RUP є архітектурно-орієнтованим. %good



(\,4\,) Методологія розробки Big Design Up Front погана тим, що може призводити до недостатнього проектування.%bad


{\begin {center}\textbf{Завдання 8}\end {center}}\par
Що з наступного не має відношення до Scrum або не виконується/не використовується за використання оригінального Scrum?\par
\vspace{2mm}



(\,1\,) порядок виконання запланованої на спрінт роботи визначає власник продукту%good



(\,2\,) виділення в команді одного чи кількох проектувальників%good



(\,3\,) команда визначає, які задачі вона буде виконувати на наступному спрінті%good



(\,4\,) щоденний скрам%bad


{\begin {center}\textbf{Завдання 9}\end {center}}\par
Що з наведеного нижче не розглядалось та навіть не згадувалось у курсі?\par
\vspace{2mm}



(\,1\,) poker%bad



(\,2\,) модель хаосу%bad



(\,3\,) RAD%bad



(\,4\,) Lean%bad


{\begin {center}\textbf{Завдання 10}\end {center}}\par
Вибрати всі істинні твердження (якщо існують).\par
\vspace{2mm}



(\,1\,) Забезпечення якості це не тільки тестування.%good



(\,2\,) Користувачі бувають непрямими.%good



(\,3\,) Розглядають модель якості даних.%good



(\,4\,) Для тестування код треба ізолювати. Для того, щоб ізолювати, треба мати тести.%good



\end {large}}
{\vspace{5mm}\smallskip \begin {small} Затверджено на засіданні кафедри теоретичної кібернетики, протокол \No 6 від   29.01.2021 р.\par\smallskip\smallskip
{\parbox {6.5 cm}{Зав. кафедри \dotfill Ю.В.~Крак}}\hspace {0.7cm}{\hbox to 6.5 cm{ Екзаменатор \dotfill Т.О.~Карнаух}} \end{small}}
\newpage

{{\begin {small}\par
{ \center  Київський національний університет імені Тараса Шевченка \par}
{\parbox {5 cm}{Освітня програма \hfill}{\parbox {7 cm}{Інформатика \hfill}}\parbox {6 cm}{\hfill Семестр 6}}\par
{\parbox {5 cm} {Навчальний предмет \hfill}\parbox {7 cm}{Розробка програмного забезпечення \hfill}\parbox {6cm}{\hfill}\par}\vspace{-7pt} \end{small}}{\center Екзаменаційний білет \No \textbf
{ 79 }\par}}
{
\begin {large}
{\begin {center}\textbf{Завдання 1}\end {center}}\par
Так звані трикутники балансування зв'язують між собою:\par
\vspace{2mm}



(\,1\,) Завжди рівно 4 параметри%bad



(\,2\,) Завжди рівно 3 параметри%bad



(\,3\,) 4 або 3 параметри, при цьому може бути відсутня якість%bad



(\,4\,) 4 або 3 параметри, при цьому може бути відсутній обсяг %good


{\begin {center}\textbf{Завдання 2}\end {center}}\par
Який з наведених фрагментів коду явно потребує поліпшення (форматування рядків та типи до уваги не брати)?\par
\vspace{2mm}



(\,1\,) \{int a\_, a\_\_; …\}%bad



(\,2\,) \{double max\_power, min\_power;…\} %good



(\,3\,) \{double the\_best\_value\_anyone\_ever\_have;\} %bad



(\,4\,) \{int xtrmprg; …\}%bad


{\begin {center}\textbf{Завдання 3}\end {center}}\par
Що є характерним для структурного програмування?\par
\vspace{2mm}



(\,1\,) Програма зображується ієрархічною структурою блоків.%good



(\,2\,) У коді мають використовуватись типи структур. %bad



(\,3\,) Використання механізму винятків. %bad



(\,4\,) Заборона використання скалярних (простих) типів даних. %bad


{\begin {center}\textbf{Завдання 4}\end {center}}\par
Вибрати всі істинні твердження (якщо існують).\par
\vspace{2mm}



(\,1\,) У випадку, коли вимоги до продукту відомі заздалегідь та є фіксованими контрактом, адаптивні методології розробки використовувати не можна. %bad



(\,2\,) Гнучкі методології не передбачають планування робіт. %bad



(\,3\,) У прогнозній моделі розробки розробка не може виконуватись ітеративно. %bad



(\,4\,) Неперервне проектування створює специфікацію проекту перед початком розробки/ітерації, а під час ітерації його неперервно модифікує.%bad


{\begin {center}\textbf{Завдання 5}\end {center}}\par
Вибрати всі істинні твердження (якщо існують).\par
\vspace{2mm}



(\,1\,) Від довгих функцій можна позбутися, розклавши їх у послідовність викликів дрібніших функцій. %good



(\,2\,) Одним з засобів позбутись повторень у коді є використання шаблону Template Method (але цей шаблон здатен боротися і з іншими недоліками коду). %good



(\,3\,) Одним з засобів позбутись повторень у коді є використання функцій. %good



(\,4\,) Залежності всередині модуля можуть бути сильними. %good

\newpage

{\begin {center}\textbf{Завдання 6}\end {center}}\par
Вибрати всі істинні твердження (якщо існують).\par
\vspace{2mm}



(\,1\,) Код, що себе сам документує, -- це код, який виводить користувачу документацію на себе. %bad



(\,2\,) Розробка за гнучкими методологіями не передбачає узгодження умов контракту. %bad



(\,3\,) Гнучкі методології можуть передбачати як спільне, так і індивідуальне володіння кодом. %good



(\,4\,) Розробка за гнучкими методологіями цінить співробітництво з замовником та піддатлівість змінам. %good


{\begin {center}\textbf{Завдання 7}\end {center}}\par
Вибрати всі істинні твердження (якщо існують).\par
\vspace{2mm}



(\,1\,) Висока внутрішня якість коду є марною тратою часу на те, що замовник не побачить та/або не зможе зацінити. %bad



(\,2\,) Scrum передбачає колективне володіння кодом. %good



(\,3\,) Розробка за гнучкими методологіями цінить людей та взаємодію між ними. %good



(\,4\,) RUP є ітеративним та інкрементним.%good


{\begin {center}\textbf{Завдання 8}\end {center}}\par
Що з наступного не має відношення до Scrum або не виконується/не використовується за використання оригінального Scrum?\par
\vspace{2mm}



(\,1\,) обернені вимоги%bad



(\,2\,) розповіді користувача%bad



(\,3\,) журнал невиконаних задач (backlog) спрінту%bad



(\,4\,) власник продукту призначає час, необхідний команді для реалізації конкретної задачі%good


{\begin {center}\textbf{Завдання 9}\end {center}}\par
Що з наведеного нижче не розглядалось та навіть не згадувалось у курсі?\par
\vspace{2mm}



(\,1\,) XP%bad



(\,2\,) LeSS%bad



(\,3\,) TDD%bad



(\,4\,) sashimi%bad


{\begin {center}\textbf{Завдання 10}\end {center}}\par
Вибрати всі істинні твердження (якщо існують).\par
\vspace{2mm}



(\,1\,) Розглядають модель якості даних.%good



(\,2\,) Розглядають модель якості при використанні.%good



(\,3\,) Коли нема можливості перевірити всі можливі комбінації, тестують їх попарні (pairwise) комбінації. Це може значно зменшити кількість тестів.%good



(\,4\,) Модель якості продукту є ієрархічною.%good



\end {large}}
{\vspace{5mm}\smallskip \begin {small} Затверджено на засіданні кафедри теоретичної кібернетики, протокол \No 6 від   29.01.2021 р.\par\smallskip\smallskip
{\parbox {6.5 cm}{Зав. кафедри \dotfill Ю.В.~Крак}}\hspace {0.7cm}{\hbox to 6.5 cm{ Екзаменатор \dotfill Т.О.~Карнаух}} \end{small}}
\newpage

{{\begin {small}\par
{ \center  Київський національний університет імені Тараса Шевченка \par}
{\parbox {5 cm}{Освітня програма \hfill}{\parbox {7 cm}{Інформатика \hfill}}\parbox {6 cm}{\hfill Семестр 6}}\par
{\parbox {5 cm} {Навчальний предмет \hfill}\parbox {7 cm}{Розробка програмного забезпечення \hfill}\parbox {6cm}{\hfill}\par}\vspace{-7pt} \end{small}}{\center Екзаменаційний білет \No \textbf
{ 80 }\par}}
{
\begin {large}
{\begin {center}\textbf{Завдання 1}\end {center}}\par
Шаблон проектування:\par
\vspace{2mm}



(\,1\,) Є шаблоном, в який треба підставити власні ідентифікатори%bad



(\,2\,) Задає розв'язання типової проблеми до найдрібніших подробиць%bad



(\,3\,) Формулює ідею розв'язання проблеми та може пропонувати одну чи кілька її реалізацій%good



(\,4\,) Тільки формулює ідею розв'язання проблеми%bad


{\begin {center}\textbf{Завдання 2}\end {center}}\par
Який з наведених фрагментів коду явно потребує поліпшення (форматування рядків та типи до уваги не брати)?\par
\vspace{2mm}



(\,1\,) \{fstream inputFile, outputFile;…\} %good



(\,2\,) void f13(…);%bad



(\,3\,) \{int input\_file, file\_to\_output; …\}%bad



(\,4\,) void load\_data(…);%good


{\begin {center}\textbf{Завдання 3}\end {center}}\par
Що є характерним для структурного програмування?\par
\vspace{2mm}



(\,1\,) Переходи з середини одного блоку в середину іншого неможливі. %good



(\,2\,) У коді мають використовуватись типи даних, що є класами. %bad



(\,3\,) У коді мають використовуватись типи структур. %bad



(\,4\,) Використання механізму винятків. %bad


{\begin {center}\textbf{Завдання 4}\end {center}}\par
Вибрати всі істинні твердження (якщо існують).\par
\vspace{2mm}



(\,1\,) Прогнозні моделі розробки з самого початку фіксують вимоги до продукту. %good



(\,2\,) Адаптивні методології не передбачають визначення вимог. %bad



(\,3\,) І прогнозні, і адаптивні методології передбачають планування виконуваних робіт. %good



(\,4\,) Адаптивна модель розробки адаптує вимоги до продукту до наявної команди розробників. %bad


{\begin {center}\textbf{Завдання 5}\end {center}}\par
Вибрати всі істинні твердження (якщо існують).\par
\vspace{2mm}



(\,1\,) Правильно спроектована функція повинна мати рівно один обов'язок. %good



(\,2\,) Залежності між різними модулями системи мають бути якомога більш сильними, бо інакше вона розпадеться. %bad



(\,3\,) Одним з прийомів спрощення складної умовної логіки є використання шаблону Strategy. %good



(\,4\,) Складні комбінації умов починають керувати логікою виконання, тому що такими є вимоги до функціонування ПЗ. Оскільки ПЗ має якнайкраще забезпечувати потреби замовника, то складні комбінації умов у коді вважатись недоліком коду не можуть. %bad

\newpage

{\begin {center}\textbf{Завдання 6}\end {center}}\par
Вибрати всі істинні твердження (якщо існують).\par
\vspace{2mm}



(\,1\,) Використання RUP не заперечує часті випуски продукту, але й не наполягає на них. %good



(\,2\,) Висока внутрішня якість коду здешевшує розробку, бо дозволяє економити на вартості та часі майбутніх змін. %good



(\,3\,) Одна з принципових відмінностей RUP від Scrum полягає у використанні численних проміжних технічних артефактів. %good



(\,4\,) Методологія розробки Big Design Up Front погана тим, що може призводити до надлишкового проектування. %good


{\begin {center}\textbf{Завдання 7}\end {center}}\par
Вибрати всі істинні твердження (якщо існують).\par
\vspace{2mm}



(\,1\,) Сьогодні водоспадна методологія не має права на існування. %bad



(\,2\,) Використання гнучких методологій при розробці надвеликих проектів не несе в собі жодних ризиків. %bad



(\,3\,) Перша версія коду не завжди може бути якісною -- і це нормально. %good



(\,4\,) Гнучкі методології передбачають часті випуски продукту. %good


{\begin {center}\textbf{Завдання 8}\end {center}}\par
Що з наступного не має відношення до Scrum або не виконується/не використовується за використання оригінального Scrum?\par
\vspace{2mm}



(\,1\,) виділення в команді одного чи кількох тестувальників%good



(\,2\,) порядок виконання запланованої на спрінт роботи визначає скрам-мастер%good



(\,3\,) скрам-мастер остаточно затверджує проект системи%good



(\,4\,) скрам-мастер визначає, які задачі команда буде виконувати на наступному спрінті%good


{\begin {center}\textbf{Завдання 9}\end {center}}\par
Що з наведеного нижче не розглядалось та навіть не згадувалось у курсі?\par
\vspace{2mm}



(\,1\,) Cleanroom%bad



(\,2\,) DSDM%bad



(\,3\,) CRAZY%good



(\,4\,) SAFe%bad


{\begin {center}\textbf{Завдання 10}\end {center}}\par
Вибрати всі істинні твердження (якщо існують).\par
\vspace{2mm}



(\,1\,) Можна розглядати якість під час використання.%good



(\,2\,) Для тестування код треба ізолювати. Для того, щоб ізолювати, треба мати тести.%good



(\,3\,) Користувачі бувають непрямими.%good



(\,4\,) Забезпечення якості це не тільки тестування.%good



\end {large}}
{\vspace{5mm}\smallskip \begin {small} Затверджено на засіданні кафедри теоретичної кібернетики, протокол \No 6 від   29.01.2021 р.\par\smallskip\smallskip
{\parbox {6.5 cm}{Зав. кафедри \dotfill Ю.В.~Крак}}\hspace {0.7cm}{\hbox to 6.5 cm{ Екзаменатор \dotfill Т.О.~Карнаух}} \end{small}}
\newpage

{{\begin {small}\par
{ \center  Київський національний університет імені Тараса Шевченка \par}
{\parbox {5 cm}{Освітня програма \hfill}{\parbox {7 cm}{Інформатика \hfill}}\parbox {6 cm}{\hfill Семестр 6}}\par
{\parbox {5 cm} {Навчальний предмет \hfill}\parbox {7 cm}{Розробка програмного забезпечення \hfill}\parbox {6cm}{\hfill}\par}\vspace{-7pt} \end{small}}{\center Екзаменаційний білет \No \textbf
{ 81 }\par}}
{
\begin {large}
{\begin {center}\textbf{Завдання 1}\end {center}}\par
Так звані трикутники балансування зв'язують між собою:\par
\vspace{2mm}



(\,1\,) Завжди рівно 4 параметри%bad



(\,2\,) 4 або 3 параметри, при цьому може бути відсутня якість%bad



(\,3\,) 4 або 3 параметри, при цьому може бути відсутній час%bad



(\,4\,) 4 або 3 параметри, при цьому може бути відсутній обсяг %good


{\begin {center}\textbf{Завдання 2}\end {center}}\par
Який з наведених фрагментів коду явно потребує поліпшення (форматування рядків та типи до уваги не брати)?\par
\vspace{2mm}



(\,1\,) \{double max\_power, min\_power;…\} %good



(\,2\,) \{double the\_best\_value\_anyone\_ever\_have;\} %bad



(\,3\,) void ff(…);%bad



(\,4\,) void load\_data(T \& arg1, double arg2, double arg3, double arg4); %bad


{\begin {center}\textbf{Завдання 3}\end {center}}\par
Що є характерним для структурного програмування?\par
\vspace{2mm}



(\,1\,) Кожен блок коду повинен мати єдиний вхід та єдиний вихід.%good



(\,2\,) У коді мають використовуватись типи даних, що є класами. %bad



(\,3\,) У коді мають використовуватись структурні типи даних. %bad



(\,4\,) У коді мають використовуватись типи структур, класи використовувати заборонено. %bad


{\begin {center}\textbf{Завдання 4}\end {center}}\par
Вибрати всі істинні твердження (якщо існують).\par
\vspace{2mm}



(\,1\,) У прогнозній моделі розробки розробка може виконуватись ітеративно. %good



(\,2\,) Адаптивність методології проявляється в тому, що за зміни вимог розробка починається заново. %bad



(\,3\,) Адаптивні методології має сенс використовувати для проектів, де можлива зміна вимог. %good



(\,4\,) Неперервне проектування має бути ітеративним та інкрементним одночасно. %good


{\begin {center}\textbf{Завдання 5}\end {center}}\par
Вибрати всі істинні твердження (якщо існують).\par
\vspace{2mm}



(\,1\,) Код з довгими функціями має меншу кількість функцій, а тому для нього треба менше тестів. %bad



(\,2\,) Повторюваний код дозволяє програмісту краще зрозуміти його логіку та вибрати з кількох варіантів найкращий. %bad



(\,3\,) Втілення в коді різних підходів до одного того самого перетворення даних дозволяє програмісту продемонструвати свою професійну підготовку з найкращого боку. %bad



(\,4\,) Одним з прийомів спрощення умовної логіки є використання поліморфізму. %good

\newpage

{\begin {center}\textbf{Завдання 6}\end {center}}\par
Вибрати всі істинні твердження (якщо існують).\par
\vspace{2mm}



(\,1\,) Гнучкі методології розробки намагаються якомога швидше визначити архітектуру системи, для чого на початку завжди виконується початкове проектування системи в цілому, а тільки потім переходять до ітеративного уточнення деталей та кодування. %bad



(\,2\,) Ітеративна методологія розробки у загальному випадку не гарантує, що не буде надлишкового проектування. %good



(\,3\,) Гнучкі методи розробки нехтують проміжними технічними артефактами та документацією виключно через їх непотрібність. %bad



(\,4\,) RUP є компонентно-орієнтованим. %good


{\begin {center}\textbf{Завдання 7}\end {center}}\par
Вибрати всі істинні твердження (якщо існують).\par
\vspace{2mm}



(\,1\,) Гнучкі методології розробки допускають суттєві зміни вимог, проте все одно намагаються якомога швидше зафіксувати архітектуру. %good



(\,2\,) Гнучкі методи розробки нехтують проміжними технічними артефактами та документацією через те, що за зміни вимог та/або проекту треба витрачати додатковий час на підтримання їх у відповідному стані. %good



(\,3\,) Водоспадна методологія розробки може призводити до надлишкового проектування. %good



(\,4\,) RUP є прогнозною методологією. %bad


{\begin {center}\textbf{Завдання 8}\end {center}}\par
Що з наступного не має відношення до Scrum або не виконується/не використовується за використання оригінального Scrum?\par
\vspace{2mm}



(\,1\,) блочні тести%bad



(\,2\,) власник продукту призначає відповідальних за конкретні задачі%good



(\,3\,) планування спрінту%bad



(\,4\,) журнал невиконаних задач (backlog) продукту%bad


{\begin {center}\textbf{Завдання 9}\end {center}}\par
Що з наведеного нижче не розглядалось та навіть не згадувалось у курсі?\par
\vspace{2mm}



(\,1\,) спіральна модель%bad



(\,2\,) CRAZY%good



(\,3\,) PRINCE2%good



(\,4\,) FDD%bad


{\begin {center}\textbf{Завдання 10}\end {center}}\par
Вибрати всі істинні твердження (якщо існують).\par
\vspace{2mm}



(\,1\,) Користувачі бувають основними.%good



(\,2\,) Відсутність знайдених дефектів на початкових стадіях розробки може бути свідченням поганого покриття коду тестами.%good



(\,3\,) За провалених димових тестів решта тестів не виконується і код повертається на доробку.%good



(\,4\,) Якість буває зовнішньою.%good



\end {large}}
{\vspace{5mm}\smallskip \begin {small} Затверджено на засіданні кафедри теоретичної кібернетики, протокол \No 6 від   29.01.2021 р.\par\smallskip\smallskip
{\parbox {6.5 cm}{Зав. кафедри \dotfill Ю.В.~Крак}}\hspace {0.7cm}{\hbox to 6.5 cm{ Екзаменатор \dotfill Т.О.~Карнаух}} \end{small}}
\newpage

{{\begin {small}\par
{ \center  Київський національний університет імені Тараса Шевченка \par}
{\parbox {5 cm}{Освітня програма \hfill}{\parbox {7 cm}{Інформатика \hfill}}\parbox {6 cm}{\hfill Семестр 6}}\par
{\parbox {5 cm} {Навчальний предмет \hfill}\parbox {7 cm}{Розробка програмного забезпечення \hfill}\parbox {6cm}{\hfill}\par}\vspace{-7pt} \end{small}}{\center Екзаменаційний білет \No \textbf
{ 82 }\par}}
{
\begin {large}
{\begin {center}\textbf{Завдання 1}\end {center}}\par
Що є метою рефакторингу?\par
\vspace{2mm}



(\,1\,) Поліпшення внутрішньої структури коду%good



(\,2\,) Зменшення кількості рядків коду%bad



(\,3\,) Усунення помилок у коді%bad



(\,4\,) Додавання нової функціональності%bad


{\begin {center}\textbf{Завдання 2}\end {center}}\par
Який з наведених фрагментів коду явно потребує поліпшення (форматування рядків та типи до уваги не брати)?\par
\vspace{2mm}



(\,1\,) \{fstream inputFile, outputFile;…\} %good



(\,2\,) \{int a\_, a\_\_; …\}%bad



(\,3\,) void help();%good



(\,4\,) \{fstream input\_file, outputFile; …\}%bad


{\begin {center}\textbf{Завдання 3}\end {center}}\par
Що є характерним для структурного програмування?\par
\vspace{2mm}



(\,1\,) Заборона використання скалярних (простих) типів даних. %bad



(\,2\,) Кожен блок коду повинен мати єдиний вхід та єдиний вихід.%good



(\,3\,) Використання оператора безумовного переходу. %bad



(\,4\,) Заборона використання скалярних (простих) типів даних. %bad


{\begin {center}\textbf{Завдання 4}\end {center}}\par
Вибрати всі істинні твердження (якщо існують).\par
\vspace{2mm}



(\,1\,) Прогнозна модель розробки не може одночасно бути інкрементною. %bad



(\,2\,) Адаптивні моделі розробки не допускають зміну вимог до продукту, але продукт може розроблятись поетапно. %bad



(\,3\,) В адаптивних методологіях планування робіт не виконується. %bad



(\,4\,) Ітеративні та інкрементні методи розробки є синонімом неперервного проектування. %bad


{\begin {center}\textbf{Завдання 5}\end {center}}\par
Вибрати всі істинні твердження (якщо існують).\par
\vspace{2mm}



(\,1\,) Гарний код має читатись як розмовна мова. %good



(\,2\,) Від довгих функцій можна позбутися, розклавши їх у послідовність викликів дрібніших функцій. %good



(\,3\,) Одним з засобів позбутись повторень у коді є використання функцій. %good



(\,4\,) Повторюваний код має повторюватись у точності до символу, інакше він повторюваним вважатись не може. %bad

\newpage

{\begin {center}\textbf{Завдання 6}\end {center}}\par
Вибрати всі істинні твердження (якщо існують).\par
\vspace{2mm}



(\,1\,) Програмування в парах (pair programming) підвищує якість коду та дозволяє економити на виправленні помилок. %good



(\,2\,) За невизначених або піддатливих змінам вимог краще застосовувати гнучкі методології розробки. %good



(\,3\,) За відмови від документації одним з методом збереження знань про систему є постійне пересування персоналу між її частинами. %good



(\,4\,) RUP є керованим прецедентами. %good


{\begin {center}\textbf{Завдання 7}\end {center}}\par
Вибрати всі істинні твердження (якщо існують).\par
\vspace{2mm}



(\,1\,) Багато хто з тих, що використовує Scrum, складає архітектурну модель розроблюваної системи. %good



(\,2\,) Гнучкі методології розробки не допускають використання UML. %bad



(\,3\,) На відміну від Scrum, RUP є компонентно-орієнтованим. %bad



(\,4\,) Програмування в парах (pair programming) спрощує введення в команду нових розробників. %good


{\begin {center}\textbf{Завдання 8}\end {center}}\par
Що з наступного не має відношення до Scrum або не виконується/не використовується за використання оригінального Scrum?\par
\vspace{2mm}



(\,1\,) правила визначення готовності (Definition of Done) %bad



(\,2\,) виділення одного з учасників команди як системного архітектора%good



(\,3\,) власник продукту остаточно затверджує проект системи%good



(\,4\,) розповіді користувача%bad


{\begin {center}\textbf{Завдання 9}\end {center}}\par
Що з наведеного нижче не розглядалось та навіть не згадувалось у курсі?\par
\vspace{2mm}



(\,1\,) V-модель%bad



(\,2\,) LAZY%good



(\,3\,) BDUF%bad



(\,4\,) poker%bad


{\begin {center}\textbf{Завдання 10}\end {center}}\par
Вибрати всі істинні твердження (якщо існують).\par
\vspace{2mm}



(\,1\,) Користувачі бувають вторинними.%good



(\,2\,) Розглядають модель якості продукту.%good



(\,3\,) Тестовий приклад складається з набору вхідних даних, умов виконання та очікуваних результатів.%good



(\,4\,) І верифікація, і валідація можуть використовувати тестування.%good



\end {large}}
{\vspace{5mm}\smallskip \begin {small} Затверджено на засіданні кафедри теоретичної кібернетики, протокол \No 6 від   29.01.2021 р.\par\smallskip\smallskip
{\parbox {6.5 cm}{Зав. кафедри \dotfill Ю.В.~Крак}}\hspace {0.7cm}{\hbox to 6.5 cm{ Екзаменатор \dotfill Т.О.~Карнаух}} \end{small}}
\newpage

{{\begin {small}\par
{ \center  Київський національний університет імені Тараса Шевченка \par}
{\parbox {5 cm}{Освітня програма \hfill}{\parbox {7 cm}{Інформатика \hfill}}\parbox {6 cm}{\hfill Семестр 6}}\par
{\parbox {5 cm} {Навчальний предмет \hfill}\parbox {7 cm}{Розробка програмного забезпечення \hfill}\parbox {6cm}{\hfill}\par}\vspace{-7pt} \end{small}}{\center Екзаменаційний білет \No \textbf
{ 83 }\par}}
{
\begin {large}
{\begin {center}\textbf{Завдання 1}\end {center}}\par
Так звані трикутники балансування зв'язують між собою:\par
\vspace{2mm}



(\,1\,) 4 або 3 параметри, при цьому може бути відсутній обсяг %good



(\,2\,) Завжди рівно 4 параметри%bad



(\,3\,) 4 або 3 параметри, при цьому може бути відсутній час%bad



(\,4\,) Завжди рівно 3 параметри%bad


{\begin {center}\textbf{Завдання 2}\end {center}}\par
Який з наведених фрагментів коду явно потребує поліпшення (форматування рядків та типи до уваги не брати)?\par
\vspace{2mm}



(\,1\,) \{int xtrmprg; …\}%bad



(\,2\,) void f13(…);%bad



(\,3\,) \{int input\_file, file\_to\_output; …\}%bad



(\,4\,) \{fstream input\_file, output\_file;…\} %good


{\begin {center}\textbf{Завдання 3}\end {center}}\par
Що є характерним для структурного програмування?\par
\vspace{2mm}



(\,1\,) Використання оператора безумовного переходу. %bad



(\,2\,) Переходи з середини одного блоку в середину іншого неможливі. %good



(\,3\,) У коді мають використовуватись типи структур. %bad



(\,4\,) У коді мають використовуватись типи даних, що є класами. %bad


{\begin {center}\textbf{Завдання 4}\end {center}}\par
Вибрати всі істинні твердження (якщо існують).\par
\vspace{2mm}



(\,1\,) Для адаптивних моделей розробки оцінки часу та бюджету не виконують, бо вимоги є змінними. %bad



(\,2\,) Неперервність проектування полягає в тім, що процес розробки не можна переривати більше ніж на вихідні. %bad



(\,3\,) Адаптивні моделі розробки допускають зміну вимог до продукту. %good



(\,4\,) Для адаптивних методологій розробки, на відміну від прогнозних, попередні оцінки бюджету та/або часу не виконують, бо вони все одно зміняться. %bad


{\begin {center}\textbf{Завдання 5}\end {center}}\par
Вибрати всі істинні твердження (якщо існують).\par
\vspace{2mm}



(\,1\,) Складні умови виникають в реальних задачах доволі часто, а тому є доволі природніми і не можуть вважатись недоліком коду. %bad



(\,2\,) Код програми пишеться для комп'ютера, тому він має бути синтаксично та семантично коректним. Інше значення не має. %bad



(\,3\,) Зв'язки між окремими модулями системи мають бути якомога більш слабкими. %good



(\,4\,) Одним з прийомів спрощення складних умов є їх декомпозиція. %good

\newpage

{\begin {center}\textbf{Завдання 6}\end {center}}\par
Вибрати всі істинні твердження (якщо існують).\par
\vspace{2mm}



(\,1\,) Водоспадна методологія розробки може призводити до недостатнього проектування. %bad



(\,2\,) За програмування в парах (pair programming) поки один колега набирає код на комп'ютері інший може піти пограти в настільний теніс. %bad



(\,3\,) Ітеративна та інкрементна розробка передбачає наявність у тому чи іншому вигляді етапу ініціалізації, ітерацій та керуючої таблиці проекту. %good



(\,4\,) RUP принципово відрізняється від Scrum тим, що є керованим прецедентами. У той же час у Scrum при плануванні розробки прецеденти (чи їх аналоги) до уваги не беруться. %bad


{\begin {center}\textbf{Завдання 7}\end {center}}\par
Вибрати всі істинні твердження (якщо існують).\par
\vspace{2mm}



(\,1\,) Поступове уточнення вимог замовника полегшує адаптацію до змін. %good



(\,2\,) Основна відмінність V-моделі від водоспадної полягає в тім, що діяльності з тестування виконуються паралельно з усіма іншими діяльностями з розробки. %good



(\,3\,) Розповіді користувача (user stories) – це прогресивний інструмент гнучких методологій, який нічого спільного не має з прецедентами (use case). %bad



(\,4\,) Водоспадна методологія не призначена для внесення змін у вимоги до системи. %good


{\begin {center}\textbf{Завдання 8}\end {center}}\par
Що з наступного не має відношення до Scrum або не виконується/не використовується за використання оригінального Scrum?\par
\vspace{2mm}



(\,1\,) приведення у належний вигляд журналу невиконаних задач (backlog grooming) %bad



(\,2\,) побудова різнопланових моделей системи%good



(\,3\,) журнал невиконаних задач (backlog) продукту%bad



(\,4\,) виділення одного з учасників команди як системного архітектора%good


{\begin {center}\textbf{Завдання 9}\end {center}}\par
Що з наведеного нижче не розглядалось та навіть не згадувалось у курсі?\par
\vspace{2mm}



(\,1\,) RUP%bad



(\,2\,) V-модель%bad



(\,3\,) LeSS%bad



(\,4\,) X-модель%good


{\begin {center}\textbf{Завдання 10}\end {center}}\par
Вибрати всі істинні твердження (якщо існують).\par
\vspace{2mm}



(\,1\,) Тестування буває позитивне та негативне.%good



(\,2\,) Атрибути якості можна використовувати при уточненні вимог.%good



(\,3\,) Модель якості даних є ієрархічною.%good



(\,4\,) Якість можна визначати через атрибути якості.%good



\end {large}}
{\vspace{5mm}\smallskip \begin {small} Затверджено на засіданні кафедри теоретичної кібернетики, протокол \No 6 від   29.01.2021 р.\par\smallskip\smallskip
{\parbox {6.5 cm}{Зав. кафедри \dotfill Ю.В.~Крак}}\hspace {0.7cm}{\hbox to 6.5 cm{ Екзаменатор \dotfill Т.О.~Карнаух}} \end{small}}
\newpage

{{\begin {small}\par
{ \center  Київський національний університет імені Тараса Шевченка \par}
{\parbox {5 cm}{Освітня програма \hfill}{\parbox {7 cm}{Інформатика \hfill}}\parbox {6 cm}{\hfill Семестр 6}}\par
{\parbox {5 cm} {Навчальний предмет \hfill}\parbox {7 cm}{Розробка програмного забезпечення \hfill}\parbox {6cm}{\hfill}\par}\vspace{-7pt} \end{small}}{\center Екзаменаційний білет \No \textbf
{ 84 }\par}}
{
\begin {large}
{\begin {center}\textbf{Завдання 1}\end {center}}\par
Що з наведеного найкраще характеризує рефакторінг?\par
\vspace{2mm}



(\,1\,) Докорінна переробка всього коду з нуля,  що не змінює його видиму поведінку%bad



(\,2\,) Довільна переробка коду, що виникає у процесі розробки%bad



(\,3\,) Зміна внутрішньої структури коду під час усунення помилок%bad



(\,4\,) Зміна внутрішньої структури коду, що не змінює його видиму поведінку%good


{\begin {center}\textbf{Завдання 2}\end {center}}\par
Який з наведених фрагментів коду явно потребує поліпшення (форматування рядків та типи до уваги не брати)?\par
\vspace{2mm}



(\,1\,) \{double maxPower, minPower;…\} %good



(\,2\,) \{int x; int y; …\}%good



(\,3\,) \{double avg\_temperature;…\} %good



(\,4\,) void load\_data(…);%good


{\begin {center}\textbf{Завдання 3}\end {center}}\par
Що є характерним для структурного програмування?\par
\vspace{2mm}



(\,1\,) Використання механізму винятків. %bad



(\,2\,) У коді мають використовуватись типи структур, класи використовувати заборонено. %bad



(\,3\,) Програма зображується ієрархічною структурою блоків.%good



(\,4\,) У коді мають використовуватись структурні типи даних. %bad


{\begin {center}\textbf{Завдання 4}\end {center}}\par
Вибрати всі істинні твердження (якщо існують).\par
\vspace{2mm}



(\,1\,) Прогнозна модель розробки може одночасно бути інкрементною. %good



(\,2\,) Для адаптивних методологій прогнозувати результат чи терміни його отримання абсолютно неможливо. %bad



(\,3\,) Неперервне проектування здійснює створення та модифікацію проекту системи в міру розробки. %good



(\,4\,) Прогнозні та адаптивні методології використовують одні ті самі діяльності з розробки ПЗ: визначення та аналіз вимог, проектування, реалізація, тестування, впровадження, супроводження. %good


{\begin {center}\textbf{Завдання 5}\end {center}}\par
Вибрати всі істинні твердження (якщо існують).\par
\vspace{2mm}



(\,1\,) Чим менше модулів у проекті, тим краще, бо меншу кількість модулів треба інтегрувати в єдине ціле. %bad



(\,2\,) Одним з засобів позбутись повторень у коді є використання успадкування. %good



(\,3\,) Якщо кожна функція бібліотеки буде здатна в залежності від параметрів виконувати кілька різних дій, то загальна кількість функцій у ній буде менша і її буде простіше використовувати. %bad



(\,4\,) Чим менше класів у коді, тим менше класів треба налагоджувати, а тому класи за можливості слід об'єднувати. %bad

\newpage

{\begin {center}\textbf{Завдання 6}\end {center}}\par
Вибрати всі істинні твердження (якщо існують).\par
\vspace{2mm}



(\,1\,) RUP відрізняється від Scrum тим, що ні в що не ставить персонал.%bad



(\,2\,) Якщо виникають проблеми з модифікацією коду, то слід відразу змінити команду розробників на більш кваліфіковану.%bad



(\,3\,) Водоспадна методологія може мати тенденцію до надлишкового проектування. %good



(\,4\,) Використання гнучких методологій вимагає високої кваліфікації персоналу. %good


{\begin {center}\textbf{Завдання 7}\end {center}}\par
Вибрати всі істинні твердження (якщо існують).\par
\vspace{2mm}



(\,1\,) За використання гнучких методологій розробки припиняти розробку не рекомендується. %good



(\,2\,) За застосування адаптивних методологій розробки, так само як і за прогнозних, намагаються якомога швидше визначити архітектуру системи. Відмінність полягає в тім, у який спосіб це виконується.%good



(\,3\,) За використання гнучких методологій кожен розробник робить, що хоче. %bad



(\,4\,) Програмування в парах (pair programming) підвищує якість коду та дозволяє економити на виправленні помилок. %good


{\begin {center}\textbf{Завдання 8}\end {center}}\par
Що з наступного не має відношення до Scrum або не виконується/не використовується за використання оригінального Scrum?\par
\vspace{2mm}



(\,1\,) розповіді користувача%bad



(\,2\,) порядок виконання запланованої на спрінт роботи визначає скрам-мастер%good



(\,3\,) обернені вимоги%bad



(\,4\,) виділення в команді одного чи кількох проектувальників%good


{\begin {center}\textbf{Завдання 9}\end {center}}\par
Що з наведеного нижче не розглядалось та навіть не згадувалось у курсі?\par
\vspace{2mm}



(\,1\,) RAD%bad



(\,2\,) Lean%bad



(\,3\,) RUP%bad



(\,4\,) SAFe%bad


{\begin {center}\textbf{Завдання 10}\end {center}}\par
Вибрати всі істинні твердження (якщо існують).\par
\vspace{2mm}



(\,1\,) Якість можна визначати через дефекти.%good



(\,2\,) Якість буває внутрішньою.%good



(\,3\,) Тестування вимог може породжувати їх уточнення.%good



(\,4\,) Модель якості при використанні є ієрархічною.%good



\end {large}}
{\vspace{5mm}\smallskip \begin {small} Затверджено на засіданні кафедри теоретичної кібернетики, протокол \No 6 від   29.01.2021 р.\par\smallskip\smallskip
{\parbox {6.5 cm}{Зав. кафедри \dotfill Ю.В.~Крак}}\hspace {0.7cm}{\hbox to 6.5 cm{ Екзаменатор \dotfill Т.О.~Карнаух}} \end{small}}
\newpage

{{\begin {small}\par
{ \center  Київський національний університет імені Тараса Шевченка \par}
{\parbox {5 cm}{Освітня програма \hfill}{\parbox {7 cm}{Інформатика \hfill}}\parbox {6 cm}{\hfill Семестр 6}}\par
{\parbox {5 cm} {Навчальний предмет \hfill}\parbox {7 cm}{Розробка програмного забезпечення \hfill}\parbox {6cm}{\hfill}\par}\vspace{-7pt} \end{small}}{\center Екзаменаційний білет \No \textbf
{ 85 }\par}}
{
\begin {large}
{\begin {center}\textbf{Завдання 1}\end {center}}\par
Так звані трикутники балансування зв'язують між собою:\par
\vspace{2mm}



(\,1\,) Завжди рівно 4 параметри%bad



(\,2\,) 4 або 3 параметри, при цьому може бути відсутня якість%bad



(\,3\,) 4 або 3 параметри, при цьому може бути відсутній обсяг %good



(\,4\,) Завжди рівно 3 параметри%bad


{\begin {center}\textbf{Завдання 2}\end {center}}\par
Який з наведених фрагментів коду явно потребує поліпшення (форматування рядків та типи до уваги не брати)?\par
\vspace{2mm}



(\,1\,) void load\_data(T \& arg1, double arg2, double arg3, double arg4); %bad



(\,2\,) \{int xtrmprg; …\}%bad



(\,3\,) \{int input\_file, file\_to\_output; …\}%bad



(\,4\,) \{double avg\_temperature;…\} %good


{\begin {center}\textbf{Завдання 3}\end {center}}\par
Що є характерним для структурного програмування?\par
\vspace{2mm}



(\,1\,) Використання оператора безумовного переходу. %bad



(\,2\,) У коді мають використовуватись структурні типи даних. %bad



(\,3\,) Кожен блок коду повинен мати єдиний вхід та єдиний вихід.%good



(\,4\,) У коді мають використовуватись типи структур. %bad


{\begin {center}\textbf{Завдання 4}\end {center}}\par
Вибрати всі істинні твердження (якщо існують).\par
\vspace{2mm}



(\,1\,) Гнучкі методології не передбачають планування робіт. %bad



(\,2\,) Адаптивні методології має сенс використовувати для проектів, де можлива зміна вимог. %good



(\,3\,) Неперервність проектування полягає в тім, що процес розробки не можна переривати більше ніж на вихідні. %bad



(\,4\,) Для адаптивних методологій прогнозувати результат чи терміни його отримання абсолютно неможливо. %bad


{\begin {center}\textbf{Завдання 5}\end {center}}\par
Вибрати всі істинні твердження (якщо існують).\par
\vspace{2mm}



(\,1\,) Код з довгими функціями має меншу кількість функцій, а тому його швидше та простіше тестувати. %bad



(\,2\,) Залежності всередині модуля можуть бути сильними. %good



(\,3\,) Код, що виконує однакову послідовність логічних кроків, у більшості випадків є повторюваним.%good



(\,4\,) Одним з прийомів спрощення складної умовної логіки є використання шаблону State. %good

\newpage

{\begin {center}\textbf{Завдання 6}\end {center}}\par
Вибрати всі істинні твердження (якщо існують).\par
\vspace{2mm}



(\,1\,) Гнучкі методології повністю нехтують документацією. %bad



(\,2\,) Ітеративна методологія розробки у загальному випадку не гарантує, що не буде надлишкового проектування. %good



(\,3\,) RUP є прогнозною методологією. %bad



(\,4\,) За невизначених або піддатливих змінам вимог краще застосовувати гнучкі методології розробки. %good


{\begin {center}\textbf{Завдання 7}\end {center}}\par
Вибрати всі істинні твердження (якщо існують).\par
\vspace{2mm}



(\,1\,) Водоспадна методологія не призначена для внесення змін у вимоги до системи. %good



(\,2\,) Гнучкі методології передбачають часті випуски продукту. %good



(\,3\,) Водоспадна методологія може мати тенденцію до надлишкового проектування. %good



(\,4\,) RUP відрізняється від Scrum тим, що ні в що не ставить персонал.%bad


{\begin {center}\textbf{Завдання 8}\end {center}}\par
Що з наступного не має відношення до Scrum або не виконується/не використовується за використання оригінального Scrum?\par
\vspace{2mm}



(\,1\,) порядок виконання запланованої на спрінт роботи визначають розробники%bad



(\,2\,) побудова різнопланових моделей системи%good



(\,3\,) порядок виконання запланованої на спрінт роботи визначає власник продукту%good



(\,4\,) під час спрінту дозволяється внесення змін до вимог, що вже реалізуються на поточному спрінті%good


{\begin {center}\textbf{Завдання 9}\end {center}}\par
Що з наведеного нижче не розглядалось та навіть не згадувалось у курсі?\par
\vspace{2mm}



(\,1\,) LOAN%good



(\,2\,) BDD%bad



(\,3\,) DSDM%bad



(\,4\,) sashimi%bad


{\begin {center}\textbf{Завдання 10}\end {center}}\par
Вибрати всі істинні твердження (якщо існують).\par
\vspace{2mm}



(\,1\,) Якість буває внутрішньою.%good



(\,2\,) Розглядають модель якості даних.%good



(\,3\,) Користувачі бувають вторинними.%good



(\,4\,) Розглядають модель якості продукту.%good



\end {large}}
{\vspace{5mm}\smallskip \begin {small} Затверджено на засіданні кафедри теоретичної кібернетики, протокол \No 6 від   29.01.2021 р.\par\smallskip\smallskip
{\parbox {6.5 cm}{Зав. кафедри \dotfill Ю.В.~Крак}}\hspace {0.7cm}{\hbox to 6.5 cm{ Екзаменатор \dotfill Т.О.~Карнаух}} \end{small}}
\newpage

{{\begin {small}\par
{ \center  Київський національний університет імені Тараса Шевченка \par}
{\parbox {5 cm}{Освітня програма \hfill}{\parbox {7 cm}{Інформатика \hfill}}\parbox {6 cm}{\hfill Семестр 6}}\par
{\parbox {5 cm} {Навчальний предмет \hfill}\parbox {7 cm}{Розробка програмного забезпечення \hfill}\parbox {6cm}{\hfill}\par}\vspace{-7pt} \end{small}}{\center Екзаменаційний білет \No \textbf
{ 86 }\par}}
{
\begin {large}
{\begin {center}\textbf{Завдання 1}\end {center}}\par
Так звані трикутники балансування зв'язують між собою:\par
\vspace{2mm}



(\,1\,) Завжди рівно 3 параметри%bad



(\,2\,) 4 або 3 параметри, при цьому може бути відсутній час%bad



(\,3\,) 4 або 3 параметри, при цьому може бути відсутня якість%bad



(\,4\,) 4 або 3 параметри, при цьому може бути відсутній обсяг %good


{\begin {center}\textbf{Завдання 2}\end {center}}\par
Який з наведених фрагментів коду явно потребує поліпшення (форматування рядків та типи до уваги не брати)?\par
\vspace{2mm}



(\,1\,) \{double max\_power, min\_power;…\} %good



(\,2\,) \{int x; int y; …\}%good



(\,3\,) \{fstream inputFile, outputFile;…\} %good



(\,4\,) \{int a\_, a\_\_; …\}%bad


{\begin {center}\textbf{Завдання 3}\end {center}}\par
Що є характерним для структурного програмування?\par
\vspace{2mm}



(\,1\,) Заборона використання скалярних (простих) типів даних. %bad



(\,2\,) Використання механізму винятків. %bad



(\,3\,) У коді мають використовуватись типи структур, класи використовувати заборонено. %bad



(\,4\,) Програма зображується ієрархічною структурою блоків.%good


{\begin {center}\textbf{Завдання 4}\end {center}}\par
Вибрати всі істинні твердження (якщо існують).\par
\vspace{2mm}



(\,1\,) Адаптивність методології проявляється в тому, що за зміни вимог розробка починається заново. %bad



(\,2\,) Неперервне проектування здійснює створення та модифікацію проекту системи в міру розробки. %good



(\,3\,) Для адаптивних методологій розробки, на відміну від прогнозних, попередні оцінки бюджету та/або часу не виконують, бо вони все одно зміняться. %bad



(\,4\,) Для адаптивних моделей розробки оцінки часу та бюджету не виконують, бо вимоги є змінними. %bad


{\begin {center}\textbf{Завдання 5}\end {center}}\par
Вибрати всі істинні твердження (якщо існують).\par
\vspace{2mm}



(\,1\,) Довгі функції доволі часто приховують складну логіку обробки або повторюваний код. %good



(\,2\,) Довжина функцій не впливає на внутрішню якість коду. %bad



(\,3\,) Одним з засобів позбутись повторень у коді є використання шаблону Template Method (але цей шаблон здатен боротися і з іншими недоліками коду). %good



(\,4\,) Правильно спроектована функція повинна мати рівно один обов'язок. %good

\newpage

{\begin {center}\textbf{Завдання 6}\end {center}}\par
Вибрати всі істинні твердження (якщо існують).\par
\vspace{2mm}



(\,1\,) RUP є керованим прецедентами. %good



(\,2\,) Розробка за гнучкими методологіями цінить людей та взаємодію між ними. %good



(\,3\,) За використання гнучких методологій розробки припиняти розробку не рекомендується. %good



(\,4\,) Розповіді користувача (user stories) це менш формалізований порівняно з прецедентами спосіб опису вимог до системи. %good


{\begin {center}\textbf{Завдання 7}\end {center}}\par
Вибрати всі істинні твердження (якщо існують).\par
\vspace{2mm}



(\,1\,) Код, що себе сам документує, -- це код, який виводить користувачу документацію на себе. %bad



(\,2\,) Гнучкі методології базуються на комунікації розробників, а тому передбачають виключно колективне володіння кодом. %bad



(\,3\,) Якщо виникають проблеми з модифікацією коду, то слід відразу змінити команду розробників на більш кваліфіковану.%bad



(\,4\,) Гнучкі методології розробки намагаються якомога швидше визначити архітектуру системи, для чого на початку завжди виконується початкове проектування системи в цілому, а тільки потім переходять до ітеративного уточнення деталей та кодування. %bad


{\begin {center}\textbf{Завдання 8}\end {center}}\par
Що з наступного не має відношення до Scrum або не виконується/не використовується за використання оригінального Scrum?\par
\vspace{2mm}



(\,1\,) скрам-мастер визначає, які задачі команда буде виконувати на наступному спрінті%good



(\,2\,) розповіді користувача%bad



(\,3\,) планування спрінту%bad



(\,4\,) бачення продукту (vision) %bad


{\begin {center}\textbf{Завдання 9}\end {center}}\par
Що з наведеного нижче не розглядалось та навіть не згадувалось у курсі?\par
\vspace{2mm}



(\,1\,) SAFe%bad



(\,2\,) Scrum of scrums%bad



(\,3\,) LeSS%bad



(\,4\,) Cleanroom%bad


{\begin {center}\textbf{Завдання 10}\end {center}}\par
Вибрати всі істинні твердження (якщо існують).\par
\vspace{2mm}



(\,1\,) Можна розглядати якість під час використання.%good



(\,2\,) Користувачі бувають непрямими.%good



(\,3\,) Якість можна визначати через дефекти.%good



(\,4\,) Розглядають модель якості при використанні.%good



\end {large}}
{\vspace{5mm}\smallskip \begin {small} Затверджено на засіданні кафедри теоретичної кібернетики, протокол \No 6 від   29.01.2021 р.\par\smallskip\smallskip
{\parbox {6.5 cm}{Зав. кафедри \dotfill Ю.В.~Крак}}\hspace {0.7cm}{\hbox to 6.5 cm{ Екзаменатор \dotfill Т.О.~Карнаух}} \end{small}}
\newpage

{{\begin {small}\par
{ \center  Київський національний університет імені Тараса Шевченка \par}
{\parbox {5 cm}{Освітня програма \hfill}{\parbox {7 cm}{Інформатика \hfill}}\parbox {6 cm}{\hfill Семестр 6}}\par
{\parbox {5 cm} {Навчальний предмет \hfill}\parbox {7 cm}{Розробка програмного забезпечення \hfill}\parbox {6cm}{\hfill}\par}\vspace{-7pt} \end{small}}{\center Екзаменаційний білет \No \textbf
{ 87 }\par}}
{
\begin {large}
{\begin {center}\textbf{Завдання 1}\end {center}}\par
Так звані трикутники балансування зв'язують між собою:\par
\vspace{2mm}



(\,1\,) 4 або 3 параметри, при цьому може бути відсутній обсяг %good



(\,2\,) Завжди рівно 4 параметри%bad



(\,3\,) Завжди рівно 3 параметри%bad



(\,4\,) 4 або 3 параметри, при цьому може бути відсутній час%bad


{\begin {center}\textbf{Завдання 2}\end {center}}\par
Який з наведених фрагментів коду явно потребує поліпшення (форматування рядків та типи до уваги не брати)?\par
\vspace{2mm}



(\,1\,) \{double the\_best\_value\_anyone\_ever\_have;\} %bad



(\,2\,) void f13(…);%bad



(\,3\,) \{double maxPower, minPower;…\} %good



(\,4\,) void ff(…);%bad


{\begin {center}\textbf{Завдання 3}\end {center}}\par
Що є характерним для структурного програмування?\par
\vspace{2mm}



(\,1\,) Переходи з середини одного блоку в середину іншого неможливі. %good



(\,2\,) У коді мають використовуватись типи даних, що є класами. %bad



(\,3\,) У коді мають використовуватись типи структур, класи використовувати заборонено. %bad



(\,4\,) Заборона використання скалярних (простих) типів даних. %bad


{\begin {center}\textbf{Завдання 4}\end {center}}\par
Вибрати всі істинні твердження (якщо існують).\par
\vspace{2mm}



(\,1\,) У прогнозній моделі розробки розробка може виконуватись ітеративно. %good



(\,2\,) Ітеративні та інкрементні методи розробки є синонімом неперервного проектування. %bad



(\,3\,) У прогнозній моделі розробки розробка не може виконуватись ітеративно. %bad



(\,4\,) Прогнозна модель розробки може одночасно бути інкрементною. %good


{\begin {center}\textbf{Завдання 5}\end {center}}\par
Вибрати всі істинні твердження (якщо існують).\par
\vspace{2mm}



(\,1\,) Чим вища кваліфікація програміста, тим більш складні для розуміння та довгі умовні вирази він використовує. %bad



(\,2\,) Залежності між різними модулями системи мають бути якомога більш сильними, бо інакше вона розпадеться. %bad



(\,3\,) Гарний код не повинен містити різних фрагментів, що складаються з кількох інструкцій та виконують однакове перетворення даних. %good



(\,4\,) Одним з прийомів спрощення складної умовної логіки є використання шаблону Strategy. %good

\newpage

{\begin {center}\textbf{Завдання 6}\end {center}}\par
Вибрати всі істинні твердження (якщо існують).\par
\vspace{2mm}



(\,1\,) Розробка за гнучкими методологіями ніколи не слідує плану. %bad



(\,2\,) На відміну від Scrum, RUP є компонентно-орієнтованим. %bad



(\,3\,) Сьогодні водоспадна методологія не має права на існування. %bad



(\,4\,) Адаптивні методології розробки дозволяють уникнути надлишкового проектування. %good


{\begin {center}\textbf{Завдання 7}\end {center}}\par
Вибрати всі істинні твердження (якщо існують).\par
\vspace{2mm}



(\,1\,) Основна відмінність V-моделі від водоспадної полягає в тім, що діяльності з тестування виконуються паралельно з усіма іншими діяльностями з розробки. %good



(\,2\,) Тільки RUP використовує UML.%bad



(\,3\,) За відмови від документації одним з методом збереження знань про систему є постійне пересування персоналу між її частинами. %good



(\,4\,) Використання RUP не заперечує часті випуски продукту, але й не наполягає на них. %good


{\begin {center}\textbf{Завдання 8}\end {center}}\par
Що з наступного не має відношення до Scrum або не виконується/не використовується за використання оригінального Scrum?\par
\vspace{2mm}



(\,1\,) скрам-мастер остаточно затверджує проект системи%good



(\,2\,) виділення в команді одного чи кількох тестувальників%good



(\,3\,) правила визначення готовності (Definition of Done) %bad



(\,4\,) канбан%bad


{\begin {center}\textbf{Завдання 9}\end {center}}\par
Що з наведеного нижче не розглядалось та навіть не згадувалось у курсі?\par
\vspace{2mm}



(\,1\,) poker%bad



(\,2\,) XP%bad



(\,3\,) PRINCE2%good



(\,4\,) FDD%bad


{\begin {center}\textbf{Завдання 10}\end {center}}\par
Вибрати всі істинні твердження (якщо існують).\par
\vspace{2mm}



(\,1\,) Тестування вимог може породжувати їх уточнення.%good



(\,2\,) Тестовий приклад складається з набору вхідних даних, умов виконання та очікуваних результатів.%good



(\,3\,) Якість буває зовнішньою.%good



(\,4\,) Користувачі бувають основними.%good



\end {large}}
{\vspace{5mm}\smallskip \begin {small} Затверджено на засіданні кафедри теоретичної кібернетики, протокол \No 6 від   29.01.2021 р.\par\smallskip\smallskip
{\parbox {6.5 cm}{Зав. кафедри \dotfill Ю.В.~Крак}}\hspace {0.7cm}{\hbox to 6.5 cm{ Екзаменатор \dotfill Т.О.~Карнаух}} \end{small}}
\newpage

{{\begin {small}\par
{ \center  Київський національний університет імені Тараса Шевченка \par}
{\parbox {5 cm}{Освітня програма \hfill}{\parbox {7 cm}{Інформатика \hfill}}\parbox {6 cm}{\hfill Семестр 6}}\par
{\parbox {5 cm} {Навчальний предмет \hfill}\parbox {7 cm}{Розробка програмного забезпечення \hfill}\parbox {6cm}{\hfill}\par}\vspace{-7pt} \end{small}}{\center Екзаменаційний білет \No \textbf
{ 88 }\par}}
{
\begin {large}
{\begin {center}\textbf{Завдання 1}\end {center}}\par
Що з наведеного найкраще характеризує рефакторінг?\par
\vspace{2mm}



(\,1\,) Зміна внутрішньої структури коду під час усунення помилок%bad



(\,2\,) Докорінна переробка всього коду з нуля,  що не змінює його видиму поведінку%bad



(\,3\,) Довільна переробка коду, що виникає у процесі розробки%bad



(\,4\,) Зміна внутрішньої структури коду, що не змінює його видиму поведінку%good


{\begin {center}\textbf{Завдання 2}\end {center}}\par
Який з наведених фрагментів коду явно потребує поліпшення (форматування рядків та типи до уваги не брати)?\par
\vspace{2mm}



(\,1\,) \{fstream input\_file, outputFile; …\}%bad



(\,2\,) \{fstream input\_file, output\_file;…\} %good



(\,3\,) void help();%good



(\,4\,) void load\_data(…);%good


{\begin {center}\textbf{Завдання 3}\end {center}}\par
Що є характерним для структурного програмування?\par
\vspace{2mm}



(\,1\,) У коді мають використовуватись структурні типи даних. %bad



(\,2\,) У коді мають використовуватись типи даних, що є класами. %bad



(\,3\,) Переходи з середини одного блоку в середину іншого неможливі. %good



(\,4\,) Використання оператора безумовного переходу. %bad


{\begin {center}\textbf{Завдання 4}\end {center}}\par
Вибрати всі істинні твердження (якщо існують).\par
\vspace{2mm}



(\,1\,) У випадку, коли вимоги до продукту відомі заздалегідь та є фіксованими контрактом, адаптивні методології розробки використовувати не можна. %bad



(\,2\,) Адаптивні методології не передбачають визначення вимог. %bad



(\,3\,) Неперервне проектування створює специфікацію проекту перед початком розробки/ітерації, а під час ітерації його неперервно модифікує.%bad



(\,4\,) І прогнозні, і адаптивні методології передбачають планування виконуваних робіт. %good


{\begin {center}\textbf{Завдання 5}\end {center}}\par
Вибрати всі істинні твердження (якщо існують).\par
\vspace{2mm}



(\,1\,) Одним з прийомів спрощення складної умовної логіки є використання шаблону State. %good



(\,2\,) Код, що виконує однакову послідовність логічних кроків, у більшості випадків є повторюваним.%good



(\,3\,) Якщо кожна функція бібліотеки буде здатна в залежності від параметрів виконувати кілька різних дій, то загальна кількість функцій у ній буде менша і її буде простіше використовувати. %bad



(\,4\,) Код з довгими функціями має меншу кількість функцій, а тому його швидше та простіше тестувати. %bad

\newpage

{\begin {center}\textbf{Завдання 6}\end {center}}\par
Вибрати всі істинні твердження (якщо існують).\par
\vspace{2mm}



(\,1\,) Гнучкі методи розробки нехтують проміжними технічними артефактами та документацією виключно через їх непотрібність. %bad



(\,2\,) Багато хто з тих, що використовує Scrum, складає архітектурну модель розроблюваної системи. %good



(\,3\,) Методологія розробки Big Design Up Front погана тим, що може призводити до недостатнього проектування.%bad



(\,4\,) Гнучкі методології розробки не допускають використання UML. %bad


{\begin {center}\textbf{Завдання 7}\end {center}}\par
Вибрати всі істинні твердження (якщо існують).\par
\vspace{2mm}



(\,1\,) Перша версія коду не завжди може бути якісною -- і це нормально. %good



(\,2\,) Методологія розробки Big Design Up Front погана тим, що може призводити до надлишкового проектування. %good



(\,3\,) Розповіді користувача (user stories) – це прогресивний інструмент гнучких методологій, який нічого спільного не має з прецедентами (use case). %bad



(\,4\,) Водоспадна методологія розробки може призводити до надлишкового проектування. %good


{\begin {center}\textbf{Завдання 8}\end {center}}\par
Що з наступного не має відношення до Scrum або не виконується/не використовується за використання оригінального Scrum?\par
\vspace{2mm}



(\,1\,) приведення у належний вигляд журналу невиконаних задач (backlog grooming) %bad



(\,2\,) скрам-мастер призначає відповідальних за конкретні задачі%good



(\,3\,) власник продукту призначає час, необхідний команді для реалізації конкретної задачі%good



(\,4\,) скрам-мастер визначає час, необхідний команді для реалізації конкретної задачі%good


{\begin {center}\textbf{Завдання 9}\end {center}}\par
Що з наведеного нижче не розглядалось та навіть не згадувалось у курсі?\par
\vspace{2mm}



(\,1\,) X-модель%good



(\,2\,) Lean%bad



(\,3\,) Scrum%bad



(\,4\,) BDUF%bad


{\begin {center}\textbf{Завдання 10}\end {center}}\par
Вибрати всі істинні твердження (якщо існують).\par
\vspace{2mm}



(\,1\,) Для тестування код треба ізолювати. Для того, щоб ізолювати, треба мати тести.%good



(\,2\,) І верифікація, і валідація можуть використовувати тестування.%good



(\,3\,) Забезпечення якості це не тільки тестування.%good



(\,4\,) За провалених димових тестів решта тестів не виконується і код повертається на доробку.%good



\end {large}}
{\vspace{5mm}\smallskip \begin {small} Затверджено на засіданні кафедри теоретичної кібернетики, протокол \No 6 від   29.01.2021 р.\par\smallskip\smallskip
{\parbox {6.5 cm}{Зав. кафедри \dotfill Ю.В.~Крак}}\hspace {0.7cm}{\hbox to 6.5 cm{ Екзаменатор \dotfill Т.О.~Карнаух}} \end{small}}
\newpage

{{\begin {small}\par
{ \center  Київський національний університет імені Тараса Шевченка \par}
{\parbox {5 cm}{Освітня програма \hfill}{\parbox {7 cm}{Інформатика \hfill}}\parbox {6 cm}{\hfill Семестр 6}}\par
{\parbox {5 cm} {Навчальний предмет \hfill}\parbox {7 cm}{Розробка програмного забезпечення \hfill}\parbox {6cm}{\hfill}\par}\vspace{-7pt} \end{small}}{\center Екзаменаційний білет \No \textbf
{ 89 }\par}}
{
\begin {large}
{\begin {center}\textbf{Завдання 1}\end {center}}\par
Так звані трикутники балансування зв'язують між собою:\par
\vspace{2mm}



(\,1\,) 4 або 3 параметри, при цьому може бути відсутній обсяг %good



(\,2\,) Завжди рівно 4 параметри%bad



(\,3\,) 4 або 3 параметри, при цьому може бути відсутній час%bad



(\,4\,) 4 або 3 параметри, при цьому може бути відсутня якість%bad


{\begin {center}\textbf{Завдання 2}\end {center}}\par
Який з наведених фрагментів коду явно потребує поліпшення (форматування рядків та типи до уваги не брати)?\par
\vspace{2mm}



(\,1\,) \{int xtrmprg; …\}%bad



(\,2\,) \{int x; int y; …\}%good



(\,3\,) \{double maxPower, minPower;…\} %good



(\,4\,) \{double avg\_temperature;…\} %good


{\begin {center}\textbf{Завдання 3}\end {center}}\par
Що є характерним для структурного програмування?\par
\vspace{2mm}



(\,1\,) Використання механізму винятків. %bad



(\,2\,) У коді мають використовуватись типи структур. %bad



(\,3\,) Кожен блок коду повинен мати єдиний вхід та єдиний вихід.%good



(\,4\,) У коді мають використовуватись структурні типи даних. %bad


{\begin {center}\textbf{Завдання 4}\end {center}}\par
Вибрати всі істинні твердження (якщо існують).\par
\vspace{2mm}



(\,1\,) Прогнозна модель розробки не може одночасно бути інкрементною. %bad



(\,2\,) Неперервне проектування має бути ітеративним та інкрементним одночасно. %good



(\,3\,) Адаптивні моделі розробки не допускають зміну вимог до продукту, але продукт може розроблятись поетапно. %bad



(\,4\,) Адаптивна модель розробки адаптує вимоги до продукту до наявної команди розробників. %bad


{\begin {center}\textbf{Завдання 5}\end {center}}\par
Вибрати всі істинні твердження (якщо існують).\par
\vspace{2mm}



(\,1\,) Чим менше класів у коді, тим менше класів треба налагоджувати, а тому класи за можливості слід об'єднувати. %bad



(\,2\,) Складні комбінації умов починають керувати логікою виконання, тому що такими є вимоги до функціонування ПЗ. Оскільки ПЗ має якнайкраще забезпечувати потреби замовника, то складні комбінації умов у коді вважатись недоліком коду не можуть. %bad



(\,3\,) Одним з засобів позбутись повторень у коді є використання шаблону Template Method (але цей шаблон здатен боротися і з іншими недоліками коду). %good



(\,4\,) Залежності всередині модуля можуть бути сильними. %good

\newpage

{\begin {center}\textbf{Завдання 6}\end {center}}\par
Вибрати всі істинні твердження (якщо існують).\par
\vspace{2mm}



(\,1\,) Програмування в парах (pair programming) зводиться до того, що кожним фрагментом коду володіють два програмісти. %bad



(\,2\,) Без рефакторінгу неперервне проектування неможливе. %good



(\,3\,) За використання гнучких методологій кожен розробник робить, що хоче. %bad



(\,4\,) Програмування в парах (pair programming) спрощує введення в команду нових розробників. %good


{\begin {center}\textbf{Завдання 7}\end {center}}\par
Вибрати всі істинні твердження (якщо існують).\par
\vspace{2mm}



(\,1\,) Розробка за гнучкими методологіями не передбачає узгодження умов контракту. %bad



(\,2\,) Поступове уточнення вимог замовника полегшує адаптацію до змін. %good



(\,3\,) Використання гнучких методологій вимагає високої кваліфікації персоналу. %good



(\,4\,) Ітеративна та інкрементна розробка передбачає наявність у тому чи іншому вигляді етапу ініціалізації, ітерацій та керуючої таблиці проекту. %good


{\begin {center}\textbf{Завдання 8}\end {center}}\par
Що з наступного не має відношення до Scrum або не виконується/не використовується за використання оригінального Scrum?\par
\vspace{2mm}



(\,1\,) власник продукту остаточно затверджує проект системи%good



(\,2\,) власник продукту призначає відповідальних за конкретні задачі%good



(\,3\,) ретроспектива спрінту%bad



(\,4\,) залучення зовнішніх фахівців для побудови архітектури системи%good


{\begin {center}\textbf{Завдання 9}\end {center}}\par
Що з наведеного нижче не розглядалось та навіть не згадувалось у курсі?\par
\vspace{2mm}



(\,1\,) CRAZY%good



(\,2\,) модель хаосу%bad



(\,3\,) спіральна модель%bad



(\,4\,) sashimi%bad


{\begin {center}\textbf{Завдання 10}\end {center}}\par
Вибрати всі істинні твердження (якщо існують).\par
\vspace{2mm}



(\,1\,) Модель якості при використанні є ієрархічною.%good



(\,2\,) Якість можна визначати через атрибути якості.%good



(\,3\,) Відсутність знайдених дефектів на початкових стадіях розробки може бути свідченням поганого покриття коду тестами.%good



(\,4\,) Модель якості даних є ієрархічною.%good



\end {large}}
{\vspace{5mm}\smallskip \begin {small} Затверджено на засіданні кафедри теоретичної кібернетики, протокол \No 6 від   29.01.2021 р.\par\smallskip\smallskip
{\parbox {6.5 cm}{Зав. кафедри \dotfill Ю.В.~Крак}}\hspace {0.7cm}{\hbox to 6.5 cm{ Екзаменатор \dotfill Т.О.~Карнаух}} \end{small}}
\newpage

{{\begin {small}\par
{ \center  Київський національний університет імені Тараса Шевченка \par}
{\parbox {5 cm}{Освітня програма \hfill}{\parbox {7 cm}{Інформатика \hfill}}\parbox {6 cm}{\hfill Семестр 6}}\par
{\parbox {5 cm} {Навчальний предмет \hfill}\parbox {7 cm}{Розробка програмного забезпечення \hfill}\parbox {6cm}{\hfill}\par}\vspace{-7pt} \end{small}}{\center Екзаменаційний білет \No \textbf
{ 90 }\par}}
{
\begin {large}
{\begin {center}\textbf{Завдання 1}\end {center}}\par
Шаблон проектування:\par
\vspace{2mm}



(\,1\,) Є шаблоном, в який треба підставити власні ідентифікатори%bad



(\,2\,) Формулює ідею розв'язання проблеми та може пропонувати одну чи кілька її реалізацій%good



(\,3\,) Тільки формулює ідею розв'язання проблеми%bad



(\,4\,) Задає розв'язання типової проблеми до найдрібніших подробиць%bad


{\begin {center}\textbf{Завдання 2}\end {center}}\par
Який з наведених фрагментів коду явно потребує поліпшення (форматування рядків та типи до уваги не брати)?\par
\vspace{2mm}



(\,1\,) \{double max\_power, min\_power;…\} %good



(\,2\,) \{fstream inputFile, outputFile;…\} %good



(\,3\,) \{int a\_, a\_\_; …\}%bad



(\,4\,) void help();%good


{\begin {center}\textbf{Завдання 3}\end {center}}\par
Що є характерним для структурного програмування?\par
\vspace{2mm}



(\,1\,) У коді мають використовуватись типи структур, класи використовувати заборонено. %bad



(\,2\,) Програма зображується ієрархічною структурою блоків.%good



(\,3\,) У коді мають використовуватись типи даних, що є класами. %bad



(\,4\,) У коді мають використовуватись типи структур. %bad


{\begin {center}\textbf{Завдання 4}\end {center}}\par
Вибрати всі істинні твердження (якщо існують).\par
\vspace{2mm}



(\,1\,) В адаптивних методологіях планування робіт не виконується. %bad



(\,2\,) Прогнозні та адаптивні методології використовують одні ті самі діяльності з розробки ПЗ: визначення та аналіз вимог, проектування, реалізація, тестування, впровадження, супроводження. %good



(\,3\,) Адаптивні моделі розробки допускають зміну вимог до продукту. %good



(\,4\,) Прогнозні моделі розробки з самого початку фіксують вимоги до продукту. %good


{\begin {center}\textbf{Завдання 5}\end {center}}\par
Вибрати всі істинні твердження (якщо існують).\par
\vspace{2mm}



(\,1\,) Довгі функції доволі часто приховують складну логіку обробки або повторюваний код. %good



(\,2\,) Повторюваний код має повторюватись у точності до символу, інакше він повторюваним вважатись не може. %bad



(\,3\,) Одним з засобів позбутись повторень у коді є використання функцій. %good



(\,4\,) Гарний код має читатись як розмовна мова. %good

\newpage

{\begin {center}\textbf{Завдання 6}\end {center}}\par
Вибрати всі істинні твердження (якщо існують).\par
\vspace{2mm}



(\,1\,) Гнучкі методології можуть передбачати як спільне, так і індивідуальне володіння кодом. %good



(\,2\,) За застосування адаптивних методологій розробки, так само як і за прогнозних, намагаються якомога швидше визначити архітектуру системи. Відмінність полягає в тім, у який спосіб це виконується.%good



(\,3\,) Розробка за гнучкими методологіями цінить співробітництво з замовником та піддатлівість змінам. %good



(\,4\,) Одна з принципових відмінностей RUP від Scrum полягає у використанні численних проміжних технічних артефактів. %good


{\begin {center}\textbf{Завдання 7}\end {center}}\par
Вибрати всі істинні твердження (якщо існують).\par
\vspace{2mm}



(\,1\,) Водоспадна методологія розробки може призводити до недостатнього проектування. %bad



(\,2\,) Простота – мистецтво максимізації незробленої роботи. %good



(\,3\,) Висока внутрішня якість коду є марною тратою часу на те, що замовник не побачить та/або не зможе зацінити. %bad



(\,4\,) RUP є компонентно-орієнтованим. %good


{\begin {center}\textbf{Завдання 8}\end {center}}\par
Що з наступного не має відношення до Scrum або не виконується/не використовується за використання оригінального Scrum?\par
\vspace{2mm}



(\,1\,) блочні тести%bad



(\,2\,) команда визначає, які задачі вона буде виконувати на наступному спрінті%good



(\,3\,) щоденний скрам%bad



(\,4\,) скрам-мастер визначає пріоритети задач%good


{\begin {center}\textbf{Завдання 9}\end {center}}\par
Що з наведеного нижче не розглядалось та навіть не згадувалось у курсі?\par
\vspace{2mm}



(\,1\,) LOAN%good



(\,2\,) TDD%bad



(\,3\,) LAZY%good



(\,4\,) водоспадна модель%bad


{\begin {center}\textbf{Завдання 10}\end {center}}\par
Вибрати всі істинні твердження (якщо існують).\par
\vspace{2mm}



(\,1\,) Модель якості продукту є ієрархічною.%good



(\,2\,) Коли нема можливості перевірити всі можливі комбінації, тестують їх попарні (pairwise) комбінації. Це може значно зменшити кількість тестів.%good



(\,3\,) Тестування буває позитивне та негативне.%good



(\,4\,) Атрибути якості можна використовувати при уточненні вимог.%good



\end {large}}
{\vspace{5mm}\smallskip \begin {small} Затверджено на засіданні кафедри теоретичної кібернетики, протокол \No 6 від   29.01.2021 р.\par\smallskip\smallskip
{\parbox {6.5 cm}{Зав. кафедри \dotfill Ю.В.~Крак}}\hspace {0.7cm}{\hbox to 6.5 cm{ Екзаменатор \dotfill Т.О.~Карнаух}} \end{small}}
\newpage

{{\begin {small}\par
{ \center  Київський національний університет імені Тараса Шевченка \par}
{\parbox {5 cm}{Освітня програма \hfill}{\parbox {7 cm}{Інформатика \hfill}}\parbox {6 cm}{\hfill Семестр 6}}\par
{\parbox {5 cm} {Навчальний предмет \hfill}\parbox {7 cm}{Розробка програмного забезпечення \hfill}\parbox {6cm}{\hfill}\par}\vspace{-7pt} \end{small}}{\center Екзаменаційний білет \No \textbf
{ 91 }\par}}
{
\begin {large}
{\begin {center}\textbf{Завдання 1}\end {center}}\par
Що є метою рефакторингу?\par
\vspace{2mm}



(\,1\,) Усунення помилок у коді%bad



(\,2\,) Додавання нової функціональності%bad



(\,3\,) Зменшення кількості рядків коду%bad



(\,4\,) Поліпшення внутрішньої структури коду%good


{\begin {center}\textbf{Завдання 2}\end {center}}\par
Який з наведених фрагментів коду явно потребує поліпшення (форматування рядків та типи до уваги не брати)?\par
\vspace{2mm}



(\,1\,) void ff(…);%bad



(\,2\,) \{int input\_file, file\_to\_output; …\}%bad



(\,3\,) \{fstream input\_file, outputFile; …\}%bad



(\,4\,) void load\_data(T \& arg1, double arg2, double arg3, double arg4); %bad


{\begin {center}\textbf{Завдання 3}\end {center}}\par
Що є характерним для структурного програмування?\par
\vspace{2mm}



(\,1\,) Використання механізму винятків. %bad



(\,2\,) Кожен блок коду повинен мати єдиний вхід та єдиний вихід.%good



(\,3\,) Використання оператора безумовного переходу. %bad



(\,4\,) Заборона використання скалярних (простих) типів даних. %bad


{\begin {center}\textbf{Завдання 4}\end {center}}\par
Вибрати всі істинні твердження (якщо існують).\par
\vspace{2mm}



(\,1\,) Прогнозні моделі розробки з самого початку фіксують вимоги до продукту. %good



(\,2\,) Для адаптивних моделей розробки оцінки часу та бюджету не виконують, бо вимоги є змінними. %bad



(\,3\,) Неперервне проектування має бути ітеративним та інкрементним одночасно. %good



(\,4\,) Неперервне проектування здійснює створення та модифікацію проекту системи в міру розробки. %good


{\begin {center}\textbf{Завдання 5}\end {center}}\par
Вибрати всі істинні твердження (якщо існують).\par
\vspace{2mm}



(\,1\,) Код з довгими функціями має меншу кількість функцій, а тому для нього треба менше тестів. %bad



(\,2\,) Зв'язки між окремими модулями системи мають бути якомога більш слабкими. %good



(\,3\,) Гарний код не повинен містити різних фрагментів, що складаються з кількох інструкцій та виконують однакове перетворення даних. %good



(\,4\,) Одним з прийомів спрощення складних умов є їх декомпозиція. %good

\newpage

{\begin {center}\textbf{Завдання 6}\end {center}}\par
Вибрати всі істинні твердження (якщо існують).\par
\vspace{2mm}



(\,1\,) RUP принципово відрізняється від Scrum тим, що є керованим прецедентами. У той же час у Scrum при плануванні розробки прецеденти (чи їх аналоги) до уваги не беруться. %bad



(\,2\,) RUP є ітеративним та інкрементним.%good



(\,3\,) За програмування в парах (pair programming) поки один колега набирає код на комп'ютері інший може піти пограти в настільний теніс. %bad



(\,4\,) Гнучкі методології розробки допускають суттєві зміни вимог, проте все одно намагаються якомога швидше зафіксувати архітектуру. %good


{\begin {center}\textbf{Завдання 7}\end {center}}\par
Вибрати всі істинні твердження (якщо існують).\par
\vspace{2mm}



(\,1\,) Scrum передбачає колективне володіння кодом. %good



(\,2\,) Гнучкі методи розробки нехтують проміжними технічними артефактами та документацією через те, що за зміни вимог та/або проекту треба витрачати додатковий час на підтримання їх у відповідному стані. %good



(\,3\,) Якщо виникають проблеми з модифікацією коду, то слід розглянути внесення змін у проект. %good



(\,4\,) Гнучкі методології розробки є стійкими до припинення розробки на певний час. %bad


{\begin {center}\textbf{Завдання 8}\end {center}}\par
Що з наступного не має відношення до Scrum або не виконується/не використовується за використання оригінального Scrum?\par
\vspace{2mm}



(\,1\,) покер планування (Planning Poker)%bad



(\,2\,) журнал невиконаних задач (backlog) спрінту%bad



(\,3\,) під час спрінту зміни, що впливають на досягнення мети спрінту, зазвичай, не допускаються%bad



(\,4\,) під час спрінту зміни, що впливають на досягнення мети спрінту, зазвичай, не допускаються%bad


{\begin {center}\textbf{Завдання 9}\end {center}}\par
Що з наведеного нижче не розглядалось та навіть не згадувалось у курсі?\par
\vspace{2mm}



(\,1\,) Scrum of scrums%bad



(\,2\,) Lean%bad



(\,3\,) BDUF%bad



(\,4\,) SAFe%bad


{\begin {center}\textbf{Завдання 10}\end {center}}\par
Вибрати всі істинні твердження (якщо існують).\par
\vspace{2mm}



(\,1\,) Коли нема можливості перевірити всі можливі комбінації, тестують їх попарні (pairwise) комбінації. Це може значно зменшити кількість тестів.%good



(\,2\,) Якість можна визначати через дефекти.%good



(\,3\,) Розглядають модель якості при використанні.%good



(\,4\,) Якість можна визначати через атрибути якості.%good



\end {large}}
{\vspace{5mm}\smallskip \begin {small} Затверджено на засіданні кафедри теоретичної кібернетики, протокол \No 6 від   29.01.2021 р.\par\smallskip\smallskip
{\parbox {6.5 cm}{Зав. кафедри \dotfill Ю.В.~Крак}}\hspace {0.7cm}{\hbox to 6.5 cm{ Екзаменатор \dotfill Т.О.~Карнаух}} \end{small}}
\newpage

{{\begin {small}\par
{ \center  Київський національний університет імені Тараса Шевченка \par}
{\parbox {5 cm}{Освітня програма \hfill}{\parbox {7 cm}{Інформатика \hfill}}\parbox {6 cm}{\hfill Семестр 6}}\par
{\parbox {5 cm} {Навчальний предмет \hfill}\parbox {7 cm}{Розробка програмного забезпечення \hfill}\parbox {6cm}{\hfill}\par}\vspace{-7pt} \end{small}}{\center Екзаменаційний білет \No \textbf
{ 92 }\par}}
{
\begin {large}
{\begin {center}\textbf{Завдання 1}\end {center}}\par
Так звані трикутники балансування зв'язують між собою:\par
\vspace{2mm}



(\,1\,) Завжди рівно 3 параметри%bad



(\,2\,) 4 або 3 параметри, при цьому може бути відсутній час%bad



(\,3\,) 4 або 3 параметри, при цьому може бути відсутній обсяг %good



(\,4\,) 4 або 3 параметри, при цьому може бути відсутня якість%bad


{\begin {center}\textbf{Завдання 2}\end {center}}\par
Який з наведених фрагментів коду явно потребує поліпшення (форматування рядків та типи до уваги не брати)?\par
\vspace{2mm}



(\,1\,) void f13(…);%bad



(\,2\,) \{fstream input\_file, output\_file;…\} %good



(\,3\,) \{double the\_best\_value\_anyone\_ever\_have;\} %bad



(\,4\,) void load\_data(…);%good


{\begin {center}\textbf{Завдання 3}\end {center}}\par
Що є характерним для структурного програмування?\par
\vspace{2mm}



(\,1\,) Переходи з середини одного блоку в середину іншого неможливі. %good



(\,2\,) У коді мають використовуватись структурні типи даних. %bad



(\,3\,) Використання механізму винятків. %bad



(\,4\,) У коді мають використовуватись типи структур, класи використовувати заборонено. %bad


{\begin {center}\textbf{Завдання 4}\end {center}}\par
Вибрати всі істинні твердження (якщо існують).\par
\vspace{2mm}



(\,1\,) Адаптивні моделі розробки не допускають зміну вимог до продукту, але продукт може розроблятись поетапно. %bad



(\,2\,) Ітеративні та інкрементні методи розробки є синонімом неперервного проектування. %bad



(\,3\,) В адаптивних методологіях планування робіт не виконується. %bad



(\,4\,) Адаптивні методології має сенс використовувати для проектів, де можлива зміна вимог. %good


{\begin {center}\textbf{Завдання 5}\end {center}}\par
Вибрати всі істинні твердження (якщо існують).\par
\vspace{2mm}



(\,1\,) Чим вища кваліфікація програміста, тим більш складні для розуміння та довгі умовні вирази він використовує. %bad



(\,2\,) Втілення в коді різних підходів до одного того самого перетворення даних дозволяє програмісту продемонструвати свою професійну підготовку з найкращого боку. %bad



(\,3\,) Залежності між різними модулями системи мають бути якомога більш сильними, бо інакше вона розпадеться. %bad



(\,4\,) Від довгих функцій можна позбутися, розклавши їх у послідовність викликів дрібніших функцій. %good

\newpage

{\begin {center}\textbf{Завдання 6}\end {center}}\par
Вибрати всі істинні твердження (якщо існують).\par
\vspace{2mm}



(\,1\,) RUP є архітектурно-орієнтованим. %good



(\,2\,) Використання гнучких методологій при розробці надвеликих проектів не несе в собі жодних ризиків. %bad



(\,3\,) Висока внутрішня якість коду здешевшує розробку, бо дозволяє економити на вартості та часі майбутніх змін. %good



(\,4\,) Модель хаосу є методологією розробки програмного забезпечення. %bad


{\begin {center}\textbf{Завдання 7}\end {center}}\par
Вибрати всі істинні твердження (якщо існують).\par
\vspace{2mm}



(\,1\,) RUP є компонентно-орієнтованим. %good



(\,2\,) Код, що себе сам документує, -- це код, який виводить користувачу документацію на себе. %bad



(\,3\,) RUP є ітеративним та інкрементним.%good



(\,4\,) RUP відрізняється від Scrum тим, що ні в що не ставить персонал.%bad


{\begin {center}\textbf{Завдання 8}\end {center}}\par
Що з наступного не має відношення до Scrum або не виконується/не використовується за використання оригінального Scrum?\par
\vspace{2mm}



(\,1\,) ретроспектива спрінту%bad



(\,2\,) порядок виконання запланованої на спрінт роботи визначає скрам-мастер%good



(\,3\,) власник продукту призначає час, необхідний команді для реалізації конкретної задачі%good



(\,4\,) порядок виконання запланованої на спрінт роботи визначає власник продукту%good


{\begin {center}\textbf{Завдання 9}\end {center}}\par
Що з наведеного нижче не розглядалось та навіть не згадувалось у курсі?\par
\vspace{2mm}



(\,1\,) RUP%bad



(\,2\,) RAD%bad



(\,3\,) LOAN%good



(\,4\,) LeSS%bad


{\begin {center}\textbf{Завдання 10}\end {center}}\par
Вибрати всі істинні твердження (якщо існують).\par
\vspace{2mm}



(\,1\,) Модель якості продукту є ієрархічною.%good



(\,2\,) Відсутність знайдених дефектів на початкових стадіях розробки може бути свідченням поганого покриття коду тестами.%good



(\,3\,) Розглядають модель якості даних.%good



(\,4\,) Для тестування код треба ізолювати. Для того, щоб ізолювати, треба мати тести.%good



\end {large}}
{\vspace{5mm}\smallskip \begin {small} Затверджено на засіданні кафедри теоретичної кібернетики, протокол \No 6 від   29.01.2021 р.\par\smallskip\smallskip
{\parbox {6.5 cm}{Зав. кафедри \dotfill Ю.В.~Крак}}\hspace {0.7cm}{\hbox to 6.5 cm{ Екзаменатор \dotfill Т.О.~Карнаух}} \end{small}}
\newpage

{{\begin {small}\par
{ \center  Київський національний університет імені Тараса Шевченка \par}
{\parbox {5 cm}{Освітня програма \hfill}{\parbox {7 cm}{Інформатика \hfill}}\parbox {6 cm}{\hfill Семестр 6}}\par
{\parbox {5 cm} {Навчальний предмет \hfill}\parbox {7 cm}{Розробка програмного забезпечення \hfill}\parbox {6cm}{\hfill}\par}\vspace{-7pt} \end{small}}{\center Екзаменаційний білет \No \textbf
{ 93 }\par}}
{
\begin {large}
{\begin {center}\textbf{Завдання 1}\end {center}}\par
Так звані трикутники балансування зв'язують між собою:\par
\vspace{2mm}



(\,1\,) 4 або 3 параметри, при цьому може бути відсутній обсяг %good



(\,2\,) Завжди рівно 4 параметри%bad



(\,3\,) 4 або 3 параметри, при цьому може бути відсутня якість%bad



(\,4\,) Завжди рівно 3 параметри%bad


{\begin {center}\textbf{Завдання 2}\end {center}}\par
Який з наведених фрагментів коду явно потребує поліпшення (форматування рядків та типи до уваги не брати)?\par
\vspace{2mm}



(\,1\,) \{double maxPower, minPower;…\} %good



(\,2\,) \{int a\_, a\_\_; …\}%bad



(\,3\,) void help();%good



(\,4\,) \{int input\_file, file\_to\_output; …\}%bad


{\begin {center}\textbf{Завдання 3}\end {center}}\par
Що є характерним для структурного програмування?\par
\vspace{2mm}



(\,1\,) У коді мають використовуватись типи даних, що є класами. %bad



(\,2\,) Використання оператора безумовного переходу. %bad



(\,3\,) Програма зображується ієрархічною структурою блоків.%good



(\,4\,) У коді мають використовуватись типи структур. %bad


{\begin {center}\textbf{Завдання 4}\end {center}}\par
Вибрати всі істинні твердження (якщо існують).\par
\vspace{2mm}



(\,1\,) У прогнозній моделі розробки розробка не може виконуватись ітеративно. %bad



(\,2\,) Адаптивна модель розробки адаптує вимоги до продукту до наявної команди розробників. %bad



(\,3\,) І прогнозні, і адаптивні методології передбачають планування виконуваних робіт. %good



(\,4\,) Неперервність проектування полягає в тім, що процес розробки не можна переривати більше ніж на вихідні. %bad


{\begin {center}\textbf{Завдання 5}\end {center}}\par
Вибрати всі істинні твердження (якщо існують).\par
\vspace{2mm}



(\,1\,) Довжина функцій не впливає на внутрішню якість коду. %bad



(\,2\,) Правильно спроектована функція повинна мати рівно один обов'язок. %good



(\,3\,) Повторюваний код дозволяє програмісту краще зрозуміти його логіку та вибрати з кількох варіантів найкращий. %bad



(\,4\,) Одним з прийомів спрощення складної умовної логіки є використання шаблону Strategy. %good

\newpage

{\begin {center}\textbf{Завдання 6}\end {center}}\par
Вибрати всі істинні твердження (якщо існують).\par
\vspace{2mm}



(\,1\,) Без рефакторінгу неперервне проектування неможливе. %good



(\,2\,) Використання гнучких методологій при розробці надвеликих проектів не несе в собі жодних ризиків. %bad



(\,3\,) Багато хто з тих, що використовує Scrum, складає архітектурну модель розроблюваної системи. %good



(\,4\,) Якщо виникають проблеми з модифікацією коду, то слід розглянути внесення змін у проект. %good


{\begin {center}\textbf{Завдання 7}\end {center}}\par
Вибрати всі істинні твердження (якщо існують).\par
\vspace{2mm}



(\,1\,) Водоспадна методологія може мати тенденцію до надлишкового проектування. %good



(\,2\,) Висока внутрішня якість коду є марною тратою часу на те, що замовник не побачить та/або не зможе зацінити. %bad



(\,3\,) За відмови від документації одним з методом збереження знань про систему є постійне пересування персоналу між її частинами. %good



(\,4\,) Одна з принципових відмінностей RUP від Scrum полягає у використанні численних проміжних технічних артефактів. %good


{\begin {center}\textbf{Завдання 8}\end {center}}\par
Що з наступного не має відношення до Scrum або не виконується/не використовується за використання оригінального Scrum?\par
\vspace{2mm}



(\,1\,) скрам-мастер остаточно затверджує проект системи%good



(\,2\,) обернені вимоги%bad



(\,3\,) команда визначає, які задачі вона буде виконувати на наступному спрінті%good



(\,4\,) власник продукту призначає відповідальних за конкретні задачі%good


{\begin {center}\textbf{Завдання 9}\end {center}}\par
Що з наведеного нижче не розглядалось та навіть не згадувалось у курсі?\par
\vspace{2mm}



(\,1\,) CRAZY%good



(\,2\,) LAZY%good



(\,3\,) Cleanroom%bad



(\,4\,) спіральна модель%bad


{\begin {center}\textbf{Завдання 10}\end {center}}\par
Вибрати всі істинні твердження (якщо існують).\par
\vspace{2mm}



(\,1\,) Забезпечення якості це не тільки тестування.%good



(\,2\,) Можна розглядати якість під час використання.%good



(\,3\,) Розглядають модель якості продукту.%good



(\,4\,) Тестування вимог може породжувати їх уточнення.%good



\end {large}}
{\vspace{5mm}\smallskip \begin {small} Затверджено на засіданні кафедри теоретичної кібернетики, протокол \No 6 від   29.01.2021 р.\par\smallskip\smallskip
{\parbox {6.5 cm}{Зав. кафедри \dotfill Ю.В.~Крак}}\hspace {0.7cm}{\hbox to 6.5 cm{ Екзаменатор \dotfill Т.О.~Карнаух}} \end{small}}
\newpage

{{\begin {small}\par
{ \center  Київський національний університет імені Тараса Шевченка \par}
{\parbox {5 cm}{Освітня програма \hfill}{\parbox {7 cm}{Інформатика \hfill}}\parbox {6 cm}{\hfill Семестр 6}}\par
{\parbox {5 cm} {Навчальний предмет \hfill}\parbox {7 cm}{Розробка програмного забезпечення \hfill}\parbox {6cm}{\hfill}\par}\vspace{-7pt} \end{small}}{\center Екзаменаційний білет \No \textbf
{ 94 }\par}}
{
\begin {large}
{\begin {center}\textbf{Завдання 1}\end {center}}\par
Що з наведеного найкраще характеризує рефакторінг?\par
\vspace{2mm}



(\,1\,) Довільна переробка коду, що виникає у процесі розробки%bad



(\,2\,) Зміна внутрішньої структури коду, що не змінює його видиму поведінку%good



(\,3\,) Зміна внутрішньої структури коду під час усунення помилок%bad



(\,4\,) Докорінна переробка всього коду з нуля,  що не змінює його видиму поведінку%bad


{\begin {center}\textbf{Завдання 2}\end {center}}\par
Який з наведених фрагментів коду явно потребує поліпшення (форматування рядків та типи до уваги не брати)?\par
\vspace{2mm}



(\,1\,) void f13(…);%bad



(\,2\,) \{double max\_power, min\_power;…\} %good



(\,3\,) \{double avg\_temperature;…\} %good



(\,4\,) \{int xtrmprg; …\}%bad


{\begin {center}\textbf{Завдання 3}\end {center}}\par
Що є характерним для структурного програмування?\par
\vspace{2mm}



(\,1\,) Заборона використання скалярних (простих) типів даних. %bad



(\,2\,) У коді мають використовуватись типи даних, що є класами. %bad



(\,3\,) У коді мають використовуватись структурні типи даних. %bad



(\,4\,) Програма зображується ієрархічною структурою блоків.%good


{\begin {center}\textbf{Завдання 4}\end {center}}\par
Вибрати всі істинні твердження (якщо існують).\par
\vspace{2mm}



(\,1\,) Адаптивні моделі розробки допускають зміну вимог до продукту. %good



(\,2\,) Прогнозні та адаптивні методології використовують одні ті самі діяльності з розробки ПЗ: визначення та аналіз вимог, проектування, реалізація, тестування, впровадження, супроводження. %good



(\,3\,) Прогнозна модель розробки не може одночасно бути інкрементною. %bad



(\,4\,) Адаптивність методології проявляється в тому, що за зміни вимог розробка починається заново. %bad


{\begin {center}\textbf{Завдання 5}\end {center}}\par
Вибрати всі істинні твердження (якщо існують).\par
\vspace{2mm}



(\,1\,) Складні умови виникають в реальних задачах доволі часто, а тому є доволі природніми і не можуть вважатись недоліком коду. %bad



(\,2\,) Чим менше модулів у проекті, тим краще, бо меншу кількість модулів треба інтегрувати в єдине ціле. %bad



(\,3\,) Одним з прийомів спрощення умовної логіки є використання поліморфізму. %good



(\,4\,) Одним з засобів позбутись повторень у коді є використання успадкування. %good

\newpage

{\begin {center}\textbf{Завдання 6}\end {center}}\par
Вибрати всі істинні твердження (якщо існують).\par
\vspace{2mm}



(\,1\,) Перша версія коду не завжди може бути якісною -- і це нормально. %good



(\,2\,) Програмування в парах (pair programming) підвищує якість коду та дозволяє економити на виправленні помилок. %good



(\,3\,) Гнучкі методології передбачають часті випуски продукту. %good



(\,4\,) Гнучкі методології базуються на комунікації розробників, а тому передбачають виключно колективне володіння кодом. %bad


{\begin {center}\textbf{Завдання 7}\end {center}}\par
Вибрати всі істинні твердження (якщо існують).\par
\vspace{2mm}



(\,1\,) Розповіді користувача (user stories) це менш формалізований порівняно з прецедентами спосіб опису вимог до системи. %good



(\,2\,) Поступове уточнення вимог замовника полегшує адаптацію до змін. %good



(\,3\,) Водоспадна методологія розробки може призводити до недостатнього проектування. %bad



(\,4\,) Гнучкі методи розробки нехтують проміжними технічними артефактами та документацією через те, що за зміни вимог та/або проекту треба витрачати додатковий час на підтримання їх у відповідному стані. %good


{\begin {center}\textbf{Завдання 8}\end {center}}\par
Що з наступного не має відношення до Scrum або не виконується/не використовується за використання оригінального Scrum?\par
\vspace{2mm}



(\,1\,) виділення одного з учасників команди як системного архітектора%good



(\,2\,) бачення продукту (vision) %bad



(\,3\,) побудова різнопланових моделей системи%good



(\,4\,) скрам-мастер призначає відповідальних за конкретні задачі%good


{\begin {center}\textbf{Завдання 9}\end {center}}\par
Що з наведеного нижче не розглядалось та навіть не згадувалось у курсі?\par
\vspace{2mm}



(\,1\,) модель хаосу%bad



(\,2\,) sashimi%bad



(\,3\,) poker%bad



(\,4\,) V-модель%bad


{\begin {center}\textbf{Завдання 10}\end {center}}\par
Вибрати всі істинні твердження (якщо існують).\par
\vspace{2mm}



(\,1\,) Користувачі бувають основними.%good



(\,2\,) Модель якості при використанні є ієрархічною.%good



(\,3\,) Якість буває внутрішньою.%good



(\,4\,) Користувачі бувають вторинними.%good



\end {large}}
{\vspace{5mm}\smallskip \begin {small} Затверджено на засіданні кафедри теоретичної кібернетики, протокол \No 6 від   29.01.2021 р.\par\smallskip\smallskip
{\parbox {6.5 cm}{Зав. кафедри \dotfill Ю.В.~Крак}}\hspace {0.7cm}{\hbox to 6.5 cm{ Екзаменатор \dotfill Т.О.~Карнаух}} \end{small}}
\newpage

{{\begin {small}\par
{ \center  Київський національний університет імені Тараса Шевченка \par}
{\parbox {5 cm}{Освітня програма \hfill}{\parbox {7 cm}{Інформатика \hfill}}\parbox {6 cm}{\hfill Семестр 6}}\par
{\parbox {5 cm} {Навчальний предмет \hfill}\parbox {7 cm}{Розробка програмного забезпечення \hfill}\parbox {6cm}{\hfill}\par}\vspace{-7pt} \end{small}}{\center Екзаменаційний білет \No \textbf
{ 95 }\par}}
{
\begin {large}
{\begin {center}\textbf{Завдання 1}\end {center}}\par
Так звані трикутники балансування зв'язують між собою:\par
\vspace{2mm}



(\,1\,) 4 або 3 параметри, при цьому може бути відсутній обсяг %good



(\,2\,) Завжди рівно 4 параметри%bad



(\,3\,) 4 або 3 параметри, при цьому може бути відсутня якість%bad



(\,4\,) 4 або 3 параметри, при цьому може бути відсутній час%bad


{\begin {center}\textbf{Завдання 2}\end {center}}\par
Який з наведених фрагментів коду явно потребує поліпшення (форматування рядків та типи до уваги не брати)?\par
\vspace{2mm}



(\,1\,) void load\_data(T \& arg1, double arg2, double arg3, double arg4); %bad



(\,2\,) \{int x; int y; …\}%good



(\,3\,) void load\_data(…);%good



(\,4\,) \{double the\_best\_value\_anyone\_ever\_have;\} %bad


{\begin {center}\textbf{Завдання 3}\end {center}}\par
Що є характерним для структурного програмування?\par
\vspace{2mm}



(\,1\,) Використання механізму винятків. %bad



(\,2\,) Кожен блок коду повинен мати єдиний вхід та єдиний вихід.%good



(\,3\,) У коді мають використовуватись типи структур. %bad



(\,4\,) Використання оператора безумовного переходу. %bad


{\begin {center}\textbf{Завдання 4}\end {center}}\par
Вибрати всі істинні твердження (якщо існують).\par
\vspace{2mm}



(\,1\,) У прогнозній моделі розробки розробка може виконуватись ітеративно. %good



(\,2\,) Адаптивні методології не передбачають визначення вимог. %bad



(\,3\,) Неперервне проектування створює специфікацію проекту перед початком розробки/ітерації, а під час ітерації його неперервно модифікує.%bad



(\,4\,) Для адаптивних методологій прогнозувати результат чи терміни його отримання абсолютно неможливо. %bad


{\begin {center}\textbf{Завдання 5}\end {center}}\par
Вибрати всі істинні твердження (якщо існують).\par
\vspace{2mm}



(\,1\,) Код програми пишеться для комп'ютера, тому він має бути синтаксично та семантично коректним. Інше значення не має. %bad



(\,2\,) Код, що виконує однакову послідовність логічних кроків, у більшості випадків є повторюваним.%good



(\,3\,) Код з довгими функціями має меншу кількість функцій, а тому його швидше та простіше тестувати. %bad



(\,4\,) Код програми пишеться для комп'ютера, тому він має бути синтаксично та семантично коректним. Інше значення не має. %bad

\newpage

{\begin {center}\textbf{Завдання 6}\end {center}}\par
Вибрати всі істинні твердження (якщо існують).\par
\vspace{2mm}



(\,1\,) Використання RUP не заперечує часті випуски продукту, але й не наполягає на них. %good



(\,2\,) Ітеративна методологія розробки у загальному випадку не гарантує, що не буде надлишкового проектування. %good



(\,3\,) Методологія розробки Big Design Up Front погана тим, що може призводити до надлишкового проектування. %good



(\,4\,) RUP є керованим прецедентами. %good


{\begin {center}\textbf{Завдання 7}\end {center}}\par
Вибрати всі істинні твердження (якщо існують).\par
\vspace{2mm}



(\,1\,) Розповіді користувача (user stories) – це прогресивний інструмент гнучких методологій, який нічого спільного не має з прецедентами (use case). %bad



(\,2\,) За використання гнучких методологій розробки припиняти розробку не рекомендується. %good



(\,3\,) RUP принципово відрізняється від Scrum тим, що є керованим прецедентами. У той же час у Scrum при плануванні розробки прецеденти (чи їх аналоги) до уваги не беруться. %bad



(\,4\,) Методологія розробки Big Design Up Front погана тим, що може призводити до недостатнього проектування.%bad


{\begin {center}\textbf{Завдання 8}\end {center}}\par
Що з наступного не має відношення до Scrum або не виконується/не використовується за використання оригінального Scrum?\par
\vspace{2mm}



(\,1\,) скрам-мастер визначає час, необхідний команді для реалізації конкретної задачі%good



(\,2\,) розповіді користувача%bad



(\,3\,) канбан%bad



(\,4\,) залучення зовнішніх фахівців для побудови архітектури системи%good


{\begin {center}\textbf{Завдання 9}\end {center}}\par
Що з наведеного нижче не розглядалось та навіть не згадувалось у курсі?\par
\vspace{2mm}



(\,1\,) X-модель%good



(\,2\,) Scrum%bad



(\,3\,) PRINCE2%good



(\,4\,) BDD%bad


{\begin {center}\textbf{Завдання 10}\end {center}}\par
Вибрати всі істинні твердження (якщо існують).\par
\vspace{2mm}



(\,1\,) І верифікація, і валідація можуть використовувати тестування.%good



(\,2\,) Тестування буває позитивне та негативне.%good



(\,3\,) Тестовий приклад складається з набору вхідних даних, умов виконання та очікуваних результатів.%good



(\,4\,) Модель якості даних є ієрархічною.%good



\end {large}}
{\vspace{5mm}\smallskip \begin {small} Затверджено на засіданні кафедри теоретичної кібернетики, протокол \No 6 від   29.01.2021 р.\par\smallskip\smallskip
{\parbox {6.5 cm}{Зав. кафедри \dotfill Ю.В.~Крак}}\hspace {0.7cm}{\hbox to 6.5 cm{ Екзаменатор \dotfill Т.О.~Карнаух}} \end{small}}
\newpage

{{\begin {small}\par
{ \center  Київський національний університет імені Тараса Шевченка \par}
{\parbox {5 cm}{Освітня програма \hfill}{\parbox {7 cm}{Інформатика \hfill}}\parbox {6 cm}{\hfill Семестр 6}}\par
{\parbox {5 cm} {Навчальний предмет \hfill}\parbox {7 cm}{Розробка програмного забезпечення \hfill}\parbox {6cm}{\hfill}\par}\vspace{-7pt} \end{small}}{\center Екзаменаційний білет \No \textbf
{ 96 }\par}}
{
\begin {large}
{\begin {center}\textbf{Завдання 1}\end {center}}\par
Що є метою рефакторингу?\par
\vspace{2mm}



(\,1\,) Додавання нової функціональності%bad



(\,2\,) Усунення помилок у коді%bad



(\,3\,) Зменшення кількості рядків коду%bad



(\,4\,) Поліпшення внутрішньої структури коду%good


{\begin {center}\textbf{Завдання 2}\end {center}}\par
Який з наведених фрагментів коду явно потребує поліпшення (форматування рядків та типи до уваги не брати)?\par
\vspace{2mm}



(\,1\,) \{fstream input\_file, output\_file;…\} %good



(\,2\,) \{fstream inputFile, outputFile;…\} %good



(\,3\,) void ff(…);%bad



(\,4\,) \{fstream input\_file, outputFile; …\}%bad


{\begin {center}\textbf{Завдання 3}\end {center}}\par
Що є характерним для структурного програмування?\par
\vspace{2mm}



(\,1\,) Заборона використання скалярних (простих) типів даних. %bad



(\,2\,) Переходи з середини одного блоку в середину іншого неможливі. %good



(\,3\,) У коді мають використовуватись типи структур, класи використовувати заборонено. %bad



(\,4\,) У коді мають використовуватись типи структур, класи використовувати заборонено. %bad


{\begin {center}\textbf{Завдання 4}\end {center}}\par
Вибрати всі істинні твердження (якщо існують).\par
\vspace{2mm}



(\,1\,) Прогнозна модель розробки може одночасно бути інкрементною. %good



(\,2\,) Гнучкі методології не передбачають планування робіт. %bad



(\,3\,) У випадку, коли вимоги до продукту відомі заздалегідь та є фіксованими контрактом, адаптивні методології розробки використовувати не можна. %bad



(\,4\,) Для адаптивних методологій розробки, на відміну від прогнозних, попередні оцінки бюджету та/або часу не виконують, бо вони все одно зміняться. %bad


{\begin {center}\textbf{Завдання 5}\end {center}}\par
Вибрати всі істинні твердження (якщо існують).\par
\vspace{2mm}



(\,1\,) Якщо кожна функція бібліотеки буде здатна в залежності від параметрів виконувати кілька різних дій, то загальна кількість функцій у ній буде менша і її буде простіше використовувати. %bad



(\,2\,) Код з довгими функціями має меншу кількість функцій, а тому для нього треба менше тестів. %bad



(\,3\,) Одним з прийомів спрощення складної умовної логіки є використання шаблону Strategy. %good



(\,4\,) Залежності між різними модулями системи мають бути якомога більш сильними, бо інакше вона розпадеться. %bad

\newpage

{\begin {center}\textbf{Завдання 6}\end {center}}\par
Вибрати всі істинні твердження (якщо існують).\par
\vspace{2mm}



(\,1\,) Якщо виникають проблеми з модифікацією коду, то слід відразу змінити команду розробників на більш кваліфіковану.%bad



(\,2\,) Тільки RUP використовує UML.%bad



(\,3\,) Адаптивні методології розробки дозволяють уникнути надлишкового проектування. %good



(\,4\,) Гнучкі методології повністю нехтують документацією. %bad


{\begin {center}\textbf{Завдання 7}\end {center}}\par
Вибрати всі істинні твердження (якщо існують).\par
\vspace{2mm}



(\,1\,) Простота – мистецтво максимізації незробленої роботи. %good



(\,2\,) Scrum передбачає колективне володіння кодом. %good



(\,3\,) Використання гнучких методологій вимагає високої кваліфікації персоналу. %good



(\,4\,) На відміну від Scrum, RUP є компонентно-орієнтованим. %bad


{\begin {center}\textbf{Завдання 8}\end {center}}\par
Що з наступного не має відношення до Scrum або не виконується/не використовується за використання оригінального Scrum?\par
\vspace{2mm}



(\,1\,) журнал невиконаних задач (backlog) продукту%bad



(\,2\,) журнал невиконаних задач (backlog) спрінту%bad



(\,3\,) власник продукту остаточно затверджує проект системи%good



(\,4\,) під час спрінту дозволяється внесення змін до вимог, що вже реалізуються на поточному спрінті%good


{\begin {center}\textbf{Завдання 9}\end {center}}\par
Що з наведеного нижче не розглядалось та навіть не згадувалось у курсі?\par
\vspace{2mm}



(\,1\,) poker%bad



(\,2\,) FDD%bad



(\,3\,) DSDM%bad



(\,4\,) SAFe%bad


{\begin {center}\textbf{Завдання 10}\end {center}}\par
Вибрати всі істинні твердження (якщо існують).\par
\vspace{2mm}



(\,1\,) За провалених димових тестів решта тестів не виконується і код повертається на доробку.%good



(\,2\,) Користувачі бувають непрямими.%good



(\,3\,) Якість буває зовнішньою.%good



(\,4\,) Атрибути якості можна використовувати при уточненні вимог.%good



\end {large}}
{\vspace{5mm}\smallskip \begin {small} Затверджено на засіданні кафедри теоретичної кібернетики, протокол \No 6 від   29.01.2021 р.\par\smallskip\smallskip
{\parbox {6.5 cm}{Зав. кафедри \dotfill Ю.В.~Крак}}\hspace {0.7cm}{\hbox to 6.5 cm{ Екзаменатор \dotfill Т.О.~Карнаух}} \end{small}}
\newpage

{{\begin {small}\par
{ \center  Київський національний університет імені Тараса Шевченка \par}
{\parbox {5 cm}{Освітня програма \hfill}{\parbox {7 cm}{Інформатика \hfill}}\parbox {6 cm}{\hfill Семестр 6}}\par
{\parbox {5 cm} {Навчальний предмет \hfill}\parbox {7 cm}{Розробка програмного забезпечення \hfill}\parbox {6cm}{\hfill}\par}\vspace{-7pt} \end{small}}{\center Екзаменаційний білет \No \textbf
{ 97 }\par}}
{
\begin {large}
{\begin {center}\textbf{Завдання 1}\end {center}}\par
Так звані трикутники балансування зв'язують між собою:\par
\vspace{2mm}



(\,1\,) 4 або 3 параметри, при цьому може бути відсутній обсяг %good



(\,2\,) 4 або 3 параметри, при цьому може бути відсутній час%bad



(\,3\,) Завжди рівно 3 параметри%bad



(\,4\,) Завжди рівно 4 параметри%bad


{\begin {center}\textbf{Завдання 2}\end {center}}\par
Який з наведених фрагментів коду явно потребує поліпшення (форматування рядків та типи до уваги не брати)?\par
\vspace{2mm}



(\,1\,) \{double avg\_temperature;…\} %good



(\,2\,) void f13(…);%bad



(\,3\,) \{double max\_power, min\_power;…\} %good



(\,4\,) void ff(…);%bad


{\begin {center}\textbf{Завдання 3}\end {center}}\par
Що є характерним для структурного програмування?\par
\vspace{2mm}



(\,1\,) У коді мають використовуватись типи даних, що є класами. %bad



(\,2\,) Програма зображується ієрархічною структурою блоків.%good



(\,3\,) Використання механізму винятків. %bad



(\,4\,) Використання оператора безумовного переходу. %bad


{\begin {center}\textbf{Завдання 4}\end {center}}\par
Вибрати всі істинні твердження (якщо існують).\par
\vspace{2mm}



(\,1\,) У прогнозній моделі розробки розробка не може виконуватись ітеративно. %bad



(\,2\,) Гнучкі методології не передбачають планування робіт. %bad



(\,3\,) Прогнозні моделі розробки з самого початку фіксують вимоги до продукту. %good



(\,4\,) Адаптивні методології має сенс використовувати для проектів, де можлива зміна вимог. %good


{\begin {center}\textbf{Завдання 5}\end {center}}\par
Вибрати всі істинні твердження (якщо існують).\par
\vspace{2mm}



(\,1\,) Повторюваний код дозволяє програмісту краще зрозуміти його логіку та вибрати з кількох варіантів найкращий. %bad



(\,2\,) Складні умови виникають в реальних задачах доволі часто, а тому є доволі природніми і не можуть вважатись недоліком коду. %bad



(\,3\,) Чим вища кваліфікація програміста, тим більш складні для розуміння та довгі умовні вирази він використовує. %bad



(\,4\,) Одним з засобів позбутись повторень у коді є використання шаблону Template Method (але цей шаблон здатен боротися і з іншими недоліками коду). %good

\newpage

{\begin {center}\textbf{Завдання 6}\end {center}}\par
Вибрати всі істинні твердження (якщо існують).\par
\vspace{2mm}



(\,1\,) Гнучкі методології розробки не допускають використання UML. %bad



(\,2\,) Розробка за гнучкими методологіями цінить людей та взаємодію між ними. %good



(\,3\,) Ітеративна та інкрементна розробка передбачає наявність у тому чи іншому вигляді етапу ініціалізації, ітерацій та керуючої таблиці проекту. %good



(\,4\,) Розробка за гнучкими методологіями цінить співробітництво з замовником та піддатлівість змінам. %good


{\begin {center}\textbf{Завдання 7}\end {center}}\par
Вибрати всі істинні твердження (якщо існують).\par
\vspace{2mm}



(\,1\,) Програмування в парах (pair programming) спрощує введення в команду нових розробників. %good



(\,2\,) Гнучкі методи розробки нехтують проміжними технічними артефактами та документацією виключно через їх непотрібність. %bad



(\,3\,) Модель хаосу є методологією розробки програмного забезпечення. %bad



(\,4\,) Водоспадна методологія не призначена для внесення змін у вимоги до системи. %good


{\begin {center}\textbf{Завдання 8}\end {center}}\par
Що з наступного не має відношення до Scrum або не виконується/не використовується за використання оригінального Scrum?\par
\vspace{2mm}



(\,1\,) скрам-мастер визначає пріоритети задач%good



(\,2\,) щоденний скрам%bad



(\,3\,) приведення у належний вигляд журналу невиконаних задач (backlog grooming) %bad



(\,4\,) планування спрінту%bad


{\begin {center}\textbf{Завдання 9}\end {center}}\par
Що з наведеного нижче не розглядалось та навіть не згадувалось у курсі?\par
\vspace{2mm}



(\,1\,) XP%bad



(\,2\,) X-модель%good



(\,3\,) водоспадна модель%bad



(\,4\,) LeSS%bad


{\begin {center}\textbf{Завдання 10}\end {center}}\par
Вибрати всі істинні твердження (якщо існують).\par
\vspace{2mm}



(\,1\,) Тестовий приклад складається з набору вхідних даних, умов виконання та очікуваних результатів.%good



(\,2\,) За провалених димових тестів решта тестів не виконується і код повертається на доробку.%good



(\,3\,) Атрибути якості можна використовувати при уточненні вимог.%good



(\,4\,) Модель якості продукту є ієрархічною.%good



\end {large}}
{\vspace{5mm}\smallskip \begin {small} Затверджено на засіданні кафедри теоретичної кібернетики, протокол \No 6 від   29.01.2021 р.\par\smallskip\smallskip
{\parbox {6.5 cm}{Зав. кафедри \dotfill Ю.В.~Крак}}\hspace {0.7cm}{\hbox to 6.5 cm{ Екзаменатор \dotfill Т.О.~Карнаух}} \end{small}}
\newpage

{{\begin {small}\par
{ \center  Київський національний університет імені Тараса Шевченка \par}
{\parbox {5 cm}{Освітня програма \hfill}{\parbox {7 cm}{Інформатика \hfill}}\parbox {6 cm}{\hfill Семестр 6}}\par
{\parbox {5 cm} {Навчальний предмет \hfill}\parbox {7 cm}{Розробка програмного забезпечення \hfill}\parbox {6cm}{\hfill}\par}\vspace{-7pt} \end{small}}{\center Екзаменаційний білет \No \textbf
{ 98 }\par}}
{
\begin {large}
{\begin {center}\textbf{Завдання 1}\end {center}}\par
Шаблон проектування:\par
\vspace{2mm}



(\,1\,) Формулює ідею розв'язання проблеми та може пропонувати одну чи кілька її реалізацій%good



(\,2\,) Тільки формулює ідею розв'язання проблеми%bad



(\,3\,) Є шаблоном, в який треба підставити власні ідентифікатори%bad



(\,4\,) Задає розв'язання типової проблеми до найдрібніших подробиць%bad


{\begin {center}\textbf{Завдання 2}\end {center}}\par
Який з наведених фрагментів коду явно потребує поліпшення (форматування рядків та типи до уваги не брати)?\par
\vspace{2mm}



(\,1\,) \{int a\_, a\_\_; …\}%bad



(\,2\,) \{double maxPower, minPower;…\} %good



(\,3\,) \{int xtrmprg; …\}%bad



(\,4\,) \{fstream input\_file, outputFile; …\}%bad


{\begin {center}\textbf{Завдання 3}\end {center}}\par
Що є характерним для структурного програмування?\par
\vspace{2mm}



(\,1\,) У коді мають використовуватись структурні типи даних. %bad



(\,2\,) Переходи з середини одного блоку в середину іншого неможливі. %good



(\,3\,) У коді мають використовуватись типи структур. %bad



(\,4\,) Заборона використання скалярних (простих) типів даних. %bad


{\begin {center}\textbf{Завдання 4}\end {center}}\par
Вибрати всі істинні твердження (якщо існують).\par
\vspace{2mm}



(\,1\,) У випадку, коли вимоги до продукту відомі заздалегідь та є фіксованими контрактом, адаптивні методології розробки використовувати не можна. %bad



(\,2\,) В адаптивних методологіях планування робіт не виконується. %bad



(\,3\,) Неперервність проектування полягає в тім, що процес розробки не можна переривати більше ніж на вихідні. %bad



(\,4\,) Прогнозні та адаптивні методології використовують одні ті самі діяльності з розробки ПЗ: визначення та аналіз вимог, проектування, реалізація, тестування, впровадження, супроводження. %good


{\begin {center}\textbf{Завдання 5}\end {center}}\par
Вибрати всі істинні твердження (якщо існують).\par
\vspace{2mm}



(\,1\,) Одним з прийомів спрощення складних умов є їх декомпозиція. %good



(\,2\,) Довжина функцій не впливає на внутрішню якість коду. %bad



(\,3\,) Чим менше модулів у проекті, тим краще, бо меншу кількість модулів треба інтегрувати в єдине ціле. %bad



(\,4\,) Повторюваний код має повторюватись у точності до символу, інакше він повторюваним вважатись не може. %bad

\newpage

{\begin {center}\textbf{Завдання 6}\end {center}}\par
Вибрати всі істинні твердження (якщо існують).\par
\vspace{2mm}



(\,1\,) За використання гнучких методологій кожен розробник робить, що хоче. %bad



(\,2\,) RUP є архітектурно-орієнтованим. %good



(\,3\,) RUP є прогнозною методологією. %bad



(\,4\,) Висока внутрішня якість коду здешевшує розробку, бо дозволяє економити на вартості та часі майбутніх змін. %good


{\begin {center}\textbf{Завдання 7}\end {center}}\par
Вибрати всі істинні твердження (якщо існують).\par
\vspace{2mm}



(\,1\,) Гнучкі методології розробки намагаються якомога швидше визначити архітектуру системи, для чого на початку завжди виконується початкове проектування системи в цілому, а тільки потім переходять до ітеративного уточнення деталей та кодування. %bad



(\,2\,) Гнучкі методології розробки допускають суттєві зміни вимог, проте все одно намагаються якомога швидше зафіксувати архітектуру. %good



(\,3\,) Гнучкі методології розробки є стійкими до припинення розробки на певний час. %bad



(\,4\,) Розробка за гнучкими методологіями не передбачає узгодження умов контракту. %bad


{\begin {center}\textbf{Завдання 8}\end {center}}\par
Що з наступного не має відношення до Scrum або не виконується/не використовується за використання оригінального Scrum?\par
\vspace{2mm}



(\,1\,) скрам-мастер визначає, які задачі команда буде виконувати на наступному спрінті%good



(\,2\,) блочні тести%bad



(\,3\,) покер планування (Planning Poker)%bad



(\,4\,) виділення в команді одного чи кількох проектувальників%good


{\begin {center}\textbf{Завдання 9}\end {center}}\par
Що з наведеного нижче не розглядалось та навіть не згадувалось у курсі?\par
\vspace{2mm}



(\,1\,) CRAZY%good



(\,2\,) Lean%bad



(\,3\,) TDD%bad



(\,4\,) Cleanroom%bad


{\begin {center}\textbf{Завдання 10}\end {center}}\par
Вибрати всі істинні твердження (якщо існують).\par
\vspace{2mm}



(\,1\,) Можна розглядати якість під час використання.%good



(\,2\,) Коли нема можливості перевірити всі можливі комбінації, тестують їх попарні (pairwise) комбінації. Це може значно зменшити кількість тестів.%good



(\,3\,) І верифікація, і валідація можуть використовувати тестування.%good



(\,4\,) Модель якості даних є ієрархічною.%good



\end {large}}
{\vspace{5mm}\smallskip \begin {small} Затверджено на засіданні кафедри теоретичної кібернетики, протокол \No 6 від   29.01.2021 р.\par\smallskip\smallskip
{\parbox {6.5 cm}{Зав. кафедри \dotfill Ю.В.~Крак}}\hspace {0.7cm}{\hbox to 6.5 cm{ Екзаменатор \dotfill Т.О.~Карнаух}} \end{small}}
\newpage

{{\begin {small}\par
{ \center  Київський національний університет імені Тараса Шевченка \par}
{\parbox {5 cm}{Освітня програма \hfill}{\parbox {7 cm}{Інформатика \hfill}}\parbox {6 cm}{\hfill Семестр 6}}\par
{\parbox {5 cm} {Навчальний предмет \hfill}\parbox {7 cm}{Розробка програмного забезпечення \hfill}\parbox {6cm}{\hfill}\par}\vspace{-7pt} \end{small}}{\center Екзаменаційний білет \No \textbf
{ 99 }\par}}
{
\begin {large}
{\begin {center}\textbf{Завдання 1}\end {center}}\par
Так звані трикутники балансування зв'язують між собою:\par
\vspace{2mm}



(\,1\,) 4 або 3 параметри, при цьому може бути відсутній час%bad



(\,2\,) 4 або 3 параметри, при цьому може бути відсутній обсяг %good



(\,3\,) Завжди рівно 4 параметри%bad



(\,4\,) Завжди рівно 3 параметри%bad


{\begin {center}\textbf{Завдання 2}\end {center}}\par
Який з наведених фрагментів коду явно потребує поліпшення (форматування рядків та типи до уваги не брати)?\par
\vspace{2mm}



(\,1\,) \{int input\_file, file\_to\_output; …\}%bad



(\,2\,) \{double the\_best\_value\_anyone\_ever\_have;\} %bad



(\,3\,) void load\_data(…);%good



(\,4\,) \{int x; int y; …\}%good


{\begin {center}\textbf{Завдання 3}\end {center}}\par
Що є характерним для структурного програмування?\par
\vspace{2mm}



(\,1\,) Використання оператора безумовного переходу. %bad



(\,2\,) Заборона використання скалярних (простих) типів даних. %bad



(\,3\,) Кожен блок коду повинен мати єдиний вхід та єдиний вихід.%good



(\,4\,) У коді мають використовуватись структурні типи даних. %bad


{\begin {center}\textbf{Завдання 4}\end {center}}\par
Вибрати всі істинні твердження (якщо існують).\par
\vspace{2mm}



(\,1\,) Неперервне проектування має бути ітеративним та інкрементним одночасно. %good



(\,2\,) І прогнозні, і адаптивні методології передбачають планування виконуваних робіт. %good



(\,3\,) Для адаптивних методологій розробки, на відміну від прогнозних, попередні оцінки бюджету та/або часу не виконують, бо вони все одно зміняться. %bad



(\,4\,) Адаптивна модель розробки адаптує вимоги до продукту до наявної команди розробників. %bad


{\begin {center}\textbf{Завдання 5}\end {center}}\par
Вибрати всі істинні твердження (якщо існують).\par
\vspace{2mm}



(\,1\,) Залежності всередині модуля можуть бути сильними. %good



(\,2\,) Втілення в коді різних підходів до одного того самого перетворення даних дозволяє програмісту продемонструвати свою професійну підготовку з найкращого боку. %bad



(\,3\,) Одним з засобів позбутись повторень у коді є використання функцій. %good



(\,4\,) Гарний код не повинен містити різних фрагментів, що складаються з кількох інструкцій та виконують однакове перетворення даних. %good

\newpage

{\begin {center}\textbf{Завдання 6}\end {center}}\par
Вибрати всі істинні твердження (якщо існують).\par
\vspace{2mm}



(\,1\,) Сьогодні водоспадна методологія не має права на існування. %bad



(\,2\,) Розробка за гнучкими методологіями ніколи не слідує плану. %bad



(\,3\,) Основна відмінність V-моделі від водоспадної полягає в тім, що діяльності з тестування виконуються паралельно з усіма іншими діяльностями з розробки. %good



(\,4\,) Водоспадна методологія розробки може призводити до надлишкового проектування. %good


{\begin {center}\textbf{Завдання 7}\end {center}}\par
Вибрати всі істинні твердження (якщо існують).\par
\vspace{2mm}



(\,1\,) Програмування в парах (pair programming) зводиться до того, що кожним фрагментом коду володіють два програмісти. %bad



(\,2\,) За невизначених або піддатливих змінам вимог краще застосовувати гнучкі методології розробки. %good



(\,3\,) Гнучкі методології можуть передбачати як спільне, так і індивідуальне володіння кодом. %good



(\,4\,) За програмування в парах (pair programming) поки один колега набирає код на комп'ютері інший може піти пограти в настільний теніс. %bad


{\begin {center}\textbf{Завдання 8}\end {center}}\par
Що з наступного не має відношення до Scrum або не виконується/не використовується за використання оригінального Scrum?\par
\vspace{2mm}



(\,1\,) розповіді користувача%bad



(\,2\,) виділення в команді одного чи кількох тестувальників%good



(\,3\,) порядок виконання запланованої на спрінт роботи визначають розробники%bad



(\,4\,) правила визначення готовності (Definition of Done) %bad


{\begin {center}\textbf{Завдання 9}\end {center}}\par
Що з наведеного нижче не розглядалось та навіть не згадувалось у курсі?\par
\vspace{2mm}



(\,1\,) LAZY%good



(\,2\,) модель хаосу%bad



(\,3\,) sashimi%bad



(\,4\,) Scrum%bad


{\begin {center}\textbf{Завдання 10}\end {center}}\par
Вибрати всі істинні твердження (якщо існують).\par
\vspace{2mm}



(\,1\,) Користувачі бувають вторинними.%good



(\,2\,) Якість буває зовнішньою.%good



(\,3\,) Розглядають модель якості даних.%good



(\,4\,) Для тестування код треба ізолювати. Для того, щоб ізолювати, треба мати тести.%good



\end {large}}
{\vspace{5mm}\smallskip \begin {small} Затверджено на засіданні кафедри теоретичної кібернетики, протокол \No 6 від   29.01.2021 р.\par\smallskip\smallskip
{\parbox {6.5 cm}{Зав. кафедри \dotfill Ю.В.~Крак}}\hspace {0.7cm}{\hbox to 6.5 cm{ Екзаменатор \dotfill Т.О.~Карнаух}} \end{small}}
\newpage

{{\begin {small}\par
{ \center  Київський національний університет імені Тараса Шевченка \par}
{\parbox {5 cm}{Освітня програма \hfill}{\parbox {7 cm}{Інформатика \hfill}}\parbox {6 cm}{\hfill Семестр 6}}\par
{\parbox {5 cm} {Навчальний предмет \hfill}\parbox {7 cm}{Розробка програмного забезпечення \hfill}\parbox {6cm}{\hfill}\par}\vspace{-7pt} \end{small}}{\center Екзаменаційний білет \No \textbf
{ 100 }\par}}
{
\begin {large}
{\begin {center}\textbf{Завдання 1}\end {center}}\par
Так звані трикутники балансування зв'язують між собою:\par
\vspace{2mm}



(\,1\,) 4 або 3 параметри, при цьому може бути відсутній обсяг %good



(\,2\,) 4 або 3 параметри, при цьому може бути відсутній час%bad



(\,3\,) 4 або 3 параметри, при цьому може бути відсутня якість%bad



(\,4\,) Завжди рівно 3 параметри%bad


{\begin {center}\textbf{Завдання 2}\end {center}}\par
Який з наведених фрагментів коду явно потребує поліпшення (форматування рядків та типи до уваги не брати)?\par
\vspace{2mm}



(\,1\,) \{fstream input\_file, output\_file;…\} %good



(\,2\,) void help();%good



(\,3\,) \{fstream inputFile, outputFile;…\} %good



(\,4\,) void load\_data(T \& arg1, double arg2, double arg3, double arg4); %bad


{\begin {center}\textbf{Завдання 3}\end {center}}\par
Що є характерним для структурного програмування?\par
\vspace{2mm}



(\,1\,) Використання механізму винятків. %bad



(\,2\,) Кожен блок коду повинен мати єдиний вхід та єдиний вихід.%good



(\,3\,) У коді мають використовуватись типи структур. %bad



(\,4\,) У коді мають використовуватись типи даних, що є класами. %bad


{\begin {center}\textbf{Завдання 4}\end {center}}\par
Вибрати всі істинні твердження (якщо існують).\par
\vspace{2mm}



(\,1\,) Адаптивні моделі розробки допускають зміну вимог до продукту. %good



(\,2\,) Ітеративні та інкрементні методи розробки є синонімом неперервного проектування. %bad



(\,3\,) Прогнозна модель розробки може одночасно бути інкрементною. %good



(\,4\,) Адаптивність методології проявляється в тому, що за зміни вимог розробка починається заново. %bad


{\begin {center}\textbf{Завдання 5}\end {center}}\par
Вибрати всі істинні твердження (якщо існують).\par
\vspace{2mm}



(\,1\,) Зв'язки між окремими модулями системи мають бути якомога більш слабкими. %good



(\,2\,) Складні комбінації умов починають керувати логікою виконання, тому що такими є вимоги до функціонування ПЗ. Оскільки ПЗ має якнайкраще забезпечувати потреби замовника, то складні комбінації умов у коді вважатись недоліком коду не можуть. %bad



(\,3\,) Чим менше класів у коді, тим менше класів треба налагоджувати, а тому класи за можливості слід об'єднувати. %bad



(\,4\,) Довгі функції доволі часто приховують складну логіку обробки або повторюваний код. %good

\newpage

{\begin {center}\textbf{Завдання 6}\end {center}}\par
Вибрати всі істинні твердження (якщо існують).\par
\vspace{2mm}



(\,1\,) За застосування адаптивних методологій розробки, так само як і за прогнозних, намагаються якомога швидше визначити архітектуру системи. Відмінність полягає в тім, у який спосіб це виконується.%good



(\,2\,) Розробка за гнучкими методологіями ніколи не слідує плану. %bad



(\,3\,) Водоспадна методологія не призначена для внесення змін у вимоги до системи. %good



(\,4\,) Scrum передбачає колективне володіння кодом. %good


{\begin {center}\textbf{Завдання 7}\end {center}}\par
Вибрати всі істинні твердження (якщо існують).\par
\vspace{2mm}



(\,1\,) Основна відмінність V-моделі від водоспадної полягає в тім, що діяльності з тестування виконуються паралельно з усіма іншими діяльностями з розробки. %good



(\,2\,) RUP є ітеративним та інкрементним.%good



(\,3\,) Сьогодні водоспадна методологія не має права на існування. %bad



(\,4\,) Гнучкі методології передбачають часті випуски продукту. %good


{\begin {center}\textbf{Завдання 8}\end {center}}\par
Що з наступного не має відношення до Scrum або не виконується/не використовується за використання оригінального Scrum?\par
\vspace{2mm}



(\,1\,) власник продукту остаточно затверджує проект системи%good



(\,2\,) планування спрінту%bad



(\,3\,) обернені вимоги%bad



(\,4\,) покер планування (Planning Poker)%bad


{\begin {center}\textbf{Завдання 9}\end {center}}\par
Що з наведеного нижче не розглядалось та навіть не згадувалось у курсі?\par
\vspace{2mm}



(\,1\,) PRINCE2%good



(\,2\,) LOAN%good



(\,3\,) V-модель%bad



(\,4\,) Scrum of scrums%bad


{\begin {center}\textbf{Завдання 10}\end {center}}\par
Вибрати всі істинні твердження (якщо існують).\par
\vspace{2mm}



(\,1\,) Розглядають модель якості при використанні.%good



(\,2\,) Розглядають модель якості продукту.%good



(\,3\,) Якість можна визначати через атрибути якості.%good



(\,4\,) Користувачі бувають основними.%good



\end {large}}
{\vspace{5mm}\smallskip \begin {small} Затверджено на засіданні кафедри теоретичної кібернетики, протокол \No 6 від   29.01.2021 р.\par\smallskip\smallskip
{\parbox {6.5 cm}{Зав. кафедри \dotfill Ю.В.~Крак}}\hspace {0.7cm}{\hbox to 6.5 cm{ Екзаменатор \dotfill Т.О.~Карнаух}} \end{small}}
\newpage

\end{document}
